\documentclass{article}
\usepackage{graphicx} % Required for inserting images
\usepackage[margin=1in]{geometry}
\usepackage{xurl}
% Fix paragraph indent — must set both:
\setlength{\parindent}{1.5em}        % global default
\usepackage{caption}

\usepackage{array}              % For inserting large multi-page tables
\usepackage{longtable}
\usepackage{setspace}

\usepackage[numbers]{natbib}    % For proper citations
\makeatletter
\makeatother
      % For referencing though out the document. Can change, look at user manual how to change
\usepackage{pdflscape}          % For inserting landscape-mode objects 
\usepackage{amsmath}            % For matrices
\usepackage{listings}           % For inserting programming code
\usepackage[figuresleft]{rotating}           % For inserting sideways tables and figures
\usepackage{lipsum}             % For dummy text. You can remove this once you have remove all of the example text. 
\usepackage{textcomp}

\makeatletter
\makeatother

% For using row-spanning and column-spanning in tables:
\usepackage{multirow}

% Additional packages not included in CMU default template
\usepackage{enumitem}
\usepackage[version=4]{mhchem}
\usepackage{marvosym}
\usepackage{booktabs}
\usepackage{threeparttablex}
\usepackage{makecell}
\usepackage{siunitx}

\usepackage{colortbl}
\usepackage{xcolor}
\definecolor{tgreen}{RGB}{119,181,96}
\definecolor{tyellow}{RGB}{255,220,100}

\usepackage[T1]{fontenc}
\renewcommand\_{\textunderscore\allowbreak}

\newcommand{\altfig}[1]{} % alt-text stored in source but not rendered
\newcommand{\altfigend}{}
\newcommand{\alteq}[1]{}
\newcommand{\alteqend}{}
\newcommand{\alttbl}[1]{}
\newcommand{\alttblend}{}
\newcommand{\altlst}[1]{}
\newcommand{\altlstend}{}
\newcommand{\tagmcbegin}[1]{}
\newcommand{\tagmcend}{}

\usepackage{authblk}

\title{U.S. lithium production can undercut marginal global supply from Chinese lepidolite on both costs and carbon emissions}
\author[1,2,*]{Anthony L. Cheng}
\author[2]{Erica R.H. Fuchs}
\author[2,3]{Valerie J. Karplus}
\author[2,4]{Jeremy J. Michalek}
\affil[1]{School of Industrial and Systems Engineering, Georgia Institute of Technology}
\affil[2]{Department of Engineering and Public Policy, Carnegie Mellon University}
\affil[3]{Scott Institute for Energy Innovation, Carnegie Mellon University}
\affil[4]{Department of Mechanical Engineering, Carnegie Mellon 
University}
\affil[*]{Corresponding Author}
\date{October 2026}

\begin{document}
\maketitle

\clearpage

% \section{Abstract}

% Background/movitation (can be cut or trimmed if needed for word limit)
% Rapid growth in batteries for electric vehicles and energy storage has intensified demand for critical minerals, particularly lithium. Although lithium is widely distributed geologically, supply chains are concentrated in China, raising supply security and environmental concerns. Within China, lepidolite-based production currently appears to be on the margin of the global supply curve, yet its cost and environmental performance are poorly characterized in academic literature, limiting assessment of whether alternative production pathways can compete.
% What we did and what we found (central importance for abstract)
% We develop an integrated, empirically-grounded, transparent bottom-up technoeconomic and life cycle greenhouse gas emission model for three projects representing distinct resource-pathway combinations: U.S. evaporative brine at Silver Peak, the only active domestic producer; U.S. sulfuric acid-leach clay at Thacker Pass, currently under construction; and Chinese sulfate roasting lepidolite at Jianxiawo, a large near-marginal producer. Resource-pathway pairings create structural cost and emissions differences: 
% Relative to modeled baseline brine-based production at Silver Peak, clay at Thacker Pass has roughly 60\% (22\% to 118\%) higher costs and 180\% (125\% to 225\%) higher greenhouse gas emissions, but lepidolite production at Jianxiawo has 120\% (100% to 150\%) higher costs and 400\% (300\% to 575\%) higher emissions, all per tonne of lithium carbonate produced.
%Silver Peak (\$5,500 [\$4,500---6,800]/t \ce{Li2CO3}; 6.0 [4.7---7.3] t~\ce{CO2eq}/t \ce{Li2CO3}) and Thacker Pass (\$8,700 [\$6,700---12,000]/t; 16.7 [13.5---19.4] t/t) compare favorably to Jianxiawo (\$12,100 [\$11,200---13,900]/t; 30.3 [24.5---40.4] t/t), where ranges reflect plausible variation in the dominant cost and emissions drivers. 
% Silver Peak (\$5,500/t \ce{Li2CO3}; 6.0 t~\ce{CO2eq}/t \ce{Li2CO3}) is cheaper and lower-emission than Jianxiawo (\$12,100/t; 30.3 t/t) across plausible ranges of the dominant drivers. Thacker Pass (\$8,700/t; 16.7 t/t) has strictly lower emissions across these ranges as well as a lower midpoint cost, though its upper cost range partially overlaps Jianxiawo's lower range. These estimates assume a generic conservative design for Thacker Pass and an optimistic Jianxiawo ore grade, suggesting the true gap is likely wider. Contrary to any assumption that Chinese operations benefit from lax environmental regulation, hazardous waste disposal requirements for roasting residues further disadvantage the pathway relative to U.S. sources. However, Thacker Pass could be exposed to shocks in sulfur prices: at prices observed at the height of the 2026 Strait of Hormuz closure, its cost rises to \$15,900/t, above Jianxiawo's modeled upper bound, though such spikes are unlikely to persist over its entire mine lifetime. These results suggest that growth of U.S. lithium production has potential to economically displace currently marginal lepidolite-based lithium production, while offering lower climate change externalities.

\section*{Summary}

% Anthony's version with edits
%Lithium supply chains are concentrated in China. Chinese lepidolite-based production acts as a marginal global supplier, yet its costs and emissions are not well characterized. We develop an integrated, process-based techno-economic and life cycle greenhouse gas emissions model for three resource-pathway combinations: U.S. evaporative brine (existing at Silver Peak), U.S. sulfuric acid-leach clay (under development at Thacker Pass), and Chinese sulfate-roasting lepidolite (existing at Jianxiawo). At midpoint estimates, U.S. brine (\$5,500/t \ce{Li2CO3}; 6.0 t~\ce{CO2eq}/t) and clay (\$8,700/t; 16.7 t/t) production are cheaper and lower emissions than Chinese lepidolite (\$12,100/t; 30.3 t/t), though plausible cost ranges for U.S. clay and Chinese lepidolite overlap by conservatively $\sim$\$800/t. These results suggest that growth of U.S. lithium production has the potential to economically displace currently marginal lepidolite-based lithium production, while offering lower climate change externalities.

% Jeremy's revision
We assess the potential for U.S. development of lithium production from clays to compete with marginal global lithium supply from Chinese lepidolite on cost and emissions. Specifically, we develop an open-source integrated process-based technoeconomic and life cycle greenhouse gas emissions model for three key resource-pathway combinations using a common framework: U.S. evaporative brine (existing at Silver Peak), U.S. sulfuric acid-leach clay (under development at Thacker Pass), and Chinese sulfate-roasting lepidolite (existing at Jianxiawo). At midpoint estimates we find that U.S. evaporative brine (\$5,500/t \ce{Li2CO3}; 6 t~\ce{CO2eq}/t) and clay (\$8,700/t; 17 t/t) production are cheaper and produce lower emissions than Chinese lepidolite (\$12,100/t; 30 t/t), though plausible cost ranges for U.S. clay and Chinese lepidolite overlap by, conservatively, about \$800/t. These results suggest that growth of U.S. lithium production has the potential to economically displace marginal lepidolite-based lithium production while offering lower climate change externalities.

\section*{Context and Scale}
Electric vehicles and grid batteries are driving growth in lithium demand, and most of the world's refined lithium chemicals are made in China. Part of China's domestic production uses lepidolite, a lithium-bearing mica that must be roasted with sulfate salts before lithium can be leached out. Because lepidolite operations shut down when prices fall and restart when prices rise, they appear to set the global lithium price. The United States is trying to build domestic supply, but it has been unclear whether U.S. projects can compete with this marginal supply on cost and climate impact. We build openly documented process models of an operating U.S. evaporative brine project, a U.S. clay project under construction, and a large Chinese lepidolite operation, and we compare their costs and greenhouse gas emissions on the same basis.

Under central assumptions, both U.S. projects produce lithium carbonate more cheaply than the lepidolite operation, with 45--80\% lower emissions. These differences follow primarily from the type of resource and the chemistry needed to process it, rather than from local labor, energy, or reagent prices, so they are likely to carry over to similar deposits elsewhere. The cheapest U.S. source, brine at Silver Peak, is small and limited by land and water rights; the clay project at Thacker Pass could supply 10\% of current global demand at full capacity. Its largest cost risk is sulfur, a byproduct of oil refining and gas processing whose price lithium producers cannot influence: prices seen during the 2026 Strait of Hormuz closure would push U.S. clay costs above Chinese lepidolite. Its emissions are dominated by lime used to remove magnesium, so cleaner lime and soda ash, not cleaner electricity, are the main decarbonization levers. Contrary to assumptions of lax regulation, Chinese hazardous waste rules for roasting residues raise lepidolite costs further.

\section*{Highlights}
\begin{itemize}
    \item At midpoint, U.S. evaporative brine and clay lithium undercut lepidolite on cost and emissions
    \item Lepidolite-based lithium emits 30 t \ce{CO2}eq/t \ce{Li2CO3}, five times U.S. evaporative brine
    \item Hazardous waste rules for roasting residues raise Chinese lepidolite costs
    \item 2026 Hormuz-level sulfur prices push U.S. clay lithium cost above lepidolite's
\end{itemize}


\clearpage
\section{Introduction}

Lithium is an essential component of nearly all current electric vehicle batteries and large-scale battery energy systems, and it is projected to experience an increase in demand by a factor of 20 to 40 between 2020 and 2040~\cite{internationalenergyagencyGlobalCriticalMinerals2024, xuFutureMaterialDemand2020, greimAssessmentLithiumCriticality2020}. While recovered materials from battery recycling could dampen demand for raw materials extraction, production from raw materials will still need to produce more than 80\% of future demands under most scenarios~\cite{xuFutureMaterialDemand2020}. Natural endowments of lithium are relatively widespread, but current production is geographically concentrated, while also mismatched with demand sources (Fig. \ref{fig:04-li_distribution}). Australia, Chile, and Argentina dominate primary extraction while China dominates refining and materials production, and yet most demand for lithium products, such as for electric vehicles, is in China, the U.S., and Europe \cite{chengElectricVehicleBattery2024, greitemeierChinasHoldLithiumion2025}. This concentration has raised concerns about supply chain resilience for countries that rely on lithium for key energy technologies \cite{internationalenergyagencyGlobalCriticalMinerals2024, olivettiLithiumIonBatterySupply2017, chengElectricVehicleBattery2024}. These concerns span both geopolitical risks such as export restrictions and resource nationalism \cite{bremmerRiseFallResource2009, prinaceraiGeographyControlDeep2024} and physical risks from natural disasters, labor strikes, or environmental impacts \cite{helbigSupplyRisksAssociated2018, schrijversReviewMethodsData2020}, leading to an increasing interest in ``reshoring'' and ``friend-shoring'' supply chains. 

\begin{figure}[!h]
    \centering
    \altfig{A series of horizontal stacked bars in one chart showing the geographic distribution of lithium across five supply chain stages: resources, reserves, mining and extraction, refining, EV cathode production, and end demand. Resources and reserves are concentrated in Chile, Australia, and Argentina. Australia and Chile dominate extraction. China dominates refining with a large majority share. End demand for cathodes and EVs is split among China, Europe, the United States, South Korea, and Japan, highlighting the mismatch between where lithium is found, where it is processed, and where it is consumed.}
        \includegraphics[width=\linewidth]{figures/04-li-distribution.png}
    \altfigend
    \caption{Geographic distribution of lithium (Li) across resources, production, and demand (2023).}
    \small{There is a significant mismatch between the location of lithium resources, reserves, lithium production, lithium refining, and end lithium demand. The following codes are used to represent countries or regions: USA, United States of America; AUS, Australia; CHL, Chile; CAN, Canada; CHN, China; EUR, Europe; ARG, Argentina; BRA, Brazil; ZIM, Zimbabwe; JPN, Japan; KOR, South Korea; Other any other country not explicitly listed here. Data for the year 2023, with resources, reserves, extraction, and the quantity of refined lithium from the USGS Mineral Commodity Survey 2025 \cite{usgeologicalsurveyMineralCommoditySummaries2025}, refining percentages from the IEA Global Critical Minerals Outlook 2024\cite{internationalenergyagencyGlobalCriticalMinerals2024}, and cathode and EV percentages from the IEA Global EV Outlook 2024\cite{internationalenergyagencyGlobalEVOutlook2024}. Note that lithium mining totals in the United States are not disclosed by USGS, but are reported in other sources\cite{howellNevadaHasLoads2024}.}
    \label{fig:04-li_distribution}
\end{figure}

Global lithium production currently primarily draws on brine and hard rocks resources, with sedimentary clay deposits representing a significant but underdeveloped third source. Spodumene hard rock deposits, mined primarily in Australia and refined in China, currently represent the largest share of overall lithium production, and sit toward the lower-to-middle of the cost curve at roughly \$5,000---8,000/t lithium carbonate equivalent (LCE)~\cite{flemingLongrunAvailabilityMineral2024, wesselkaemperOptimizationProcessingCosts2025}. Continental brine deposits found in desert salars in Chile and Argentina represent the lowest-cost production globally, with high brine grade and high aridity enabling solar evaporation-based concentration at costs as low as \$3,000---5,000/t LCE~\cite{flemingLongrunAvailabilityMineral2024}. However, with the high water usage of these resources in water-stressed regions \cite{sunRobustAssessmentsLithium2025, halkesLifeCycleAssessment2024}, alternate `direct lithium extraction' (DLE) technologies are emerging as a lower-footprint alternative for processing these brine-based resources, though these are currently only operating in China and at relatively limited scale. A range of other `unconventional' resources, such as geothermal brines, jadarite, zinnwaldite, and others, currently contribute negligibly to global supply. 

Of these alternative resources, sedimentary clay deposits, most prominently found in the western United States, represent the largest U.S. undeveloped unconventional lithium source. Largely because of these clays, the United States held about 12\% of world resources in 2023 (Fig. \ref{fig:04-li_distribution}), which has since increased to 20\% as of 2026 \cite{usgeologicalsurveyMineralCommoditySummaries2026}. However, regulatory timelines, depressed prices, uncertainty in financing, and community and indigenous opposition have led to delays in developing these resources; currently, only two projects have been permitted: Thacker Pass and Rhyolite Ridge, both clay-based resources. With the latter suffering from a collapse in funding~\cite{haasRhyoliteRidgeLithium2025}, Thacker Pass remains the only U.S. project under construction, with a target operational date of 2028 - more than two decades after significant exploration activity started in 2007 \cite{alonzoEvenLithiumPrices2025, NolanTestimony2025, calaway2025AnnualReport2025}. Its target annual production capacity of 160,000 tonnes lithium carbonate equivalent (LCE) by 2044~\cite{sgscanadaincSK1300Technical2024} would be roughly 10\% of world production in 2025~\cite{usgeologicalsurveyMineralCommoditySummaries2026}. Given that the Silver Peak evaporative brine facility in Nevada has been the only lithium producer in the United States since the beginning of the 21st century, with a current nameplate capacity of roughly 6,000 tonnes, this marks a major increase in U.S. production. For perspective, if all Nevada clay projects operated by publicly listed companies were to come online (including Thacker Pass), that could result in a total of 320,000 tonnes LCE of annual production~\cite{sgscanadaincSK1300Technical2024, fluorenterprisesRhyoliteRidgeLithiumBoron2022, dyerAmendedResourceEstimate, laneNI43101Technical2024, riordanTonopahLithiumClaims2023, samariMineralResourceEstimate2024}. Thus, while various Direct Lithium Extraction projects in the United States are under construction, their relatively limited total annual production capacity of 40,000 tonnes~\cite{worleyltdParadoxLithiumProject, gallardoSouthWestArkansas2025, brushNi43101Technical2023} and relatively low reported operating costs (\$5,600-7,400 per tonne LCE \cite{fitzgeraldCanEmergingIndustrial2025}) point towards an initial focus on better understanding the competitiveness of more expensive U.S. clay resources. 

On the other hand, China's domestic lithium production has grown roughly tenfold between 2017 and 2025, from about 37,000 to an estimated 330,000 tonnes LCE, driven primarily by rapid expansion of low-grade \textit{lepidolite} hard rock mica deposits~\cite{benchmarkteamRestartCATLsJianxiawo2026, usgeologicalsurveyMineralCommoditySummaries2026}. At industry-reported costs of roughly \$10,000---16,000/t LCE \cite{liuChinaLithiumBoom2024, atkinsonLithiumPriceReversion2023, jamasmieCATLCutLithium2024}, or roughly the 90th percentile of the global cost curve~\cite{sulcaLithiumPricesWhere2024}, lepidolite is the lithium source most sensitive to price fluctuations and the relevant marginal producer against which alternative suppliers must compete. Lepidolite production has contracted during periods of low prices, with disruptions to lepidolite supply from Jiangxi province repeatedly triggering significant volatility in Chinese lithium futures markets \cite{liChinasCATLSuspends2025, benchmarkteamRestartCATLsJianxiawo2026}. While the dynamics for these projects are examined in more detail in Section~\ref{sec:04-jianxiawo}, this price-driven entry and exit behavior, combined with the scale of existing installed capacity, suggests that lepidolite is likely to remain the relevant marginal producer as demand grows, making it an appropriate competitive benchmark. A key question for countries seeking to reduce their dependence on incumbent supply chains is therefore whether domestic production can undercut this marginal cost and compete on the global market.

From a competitive market-based supply security perspective, the relevant threshold is whether a project's average lifetime operating cost is competitive with the current marginal producer: the price at which an incumbent's production becomes uneconomic sets the ceiling that a new entrant must undercut to capture market share \cite{tiltonMineralEconomicsPolicy2016}. While `competitiveness' can be defined from a financial perspective, in terms of the expected return relative to risk and alternative uses of capital, typically evaluated via net present value (NPV) or internal rate of return (IRR) metrics \cite{pindyckIrreversibilityUncertaintyInvestment1990, brennanEvaluatingNaturalResource1985}, it is this long-run operating cost competitiveness, rather than project-level returns or the exact cost curve of non-marginal production, that determines whether domestic supply development could durably compete on the global scale. 

Production characteristics such as costs and emissions are not separable from the resource itself: the chemistry and geology of the deposit, regardless of its resource type, determine the extraction and processing pathway, which in turn drives both the cost structure and also externalities such as emissions and environmental damages. This means that geographically shifting production is not simply a matter of replicating existing processes, but also that these production-resource pairings could inherently provide advantages or disadvantages. This resource-pathway relationship creates opportunities for policy to shape supply chain configurations in ways that could simultaneously address security, cost, and environmental objectives, rather than trading one off against another \cite{riofrancosSecuritySustainabilityNexus2023, haddadHowMakeLithium2023, buschEffectsDemandRecycling2025}. On the other hand, communicating the national security and domestic sourcing benefits of decarbonization projects can help navigate local public opposition \cite{ujiPursuingDecarbonizationNational2023}. The extent to which such `win-win-win' outcomes are achievable depends on the relative cost and environmental impact profiles (e.g. emissions) of competing pathways across production locations, which this work aims to characterize.

Assessing whether these outcomes are achievable requires characterization of the economic and emissions profiles of these competing pathways. Yet, the literature remains incomplete, particularly for the Chinese lepidolite operations that define the current benchmark for marginal production. While a substantial body of literature has examined models of lithium production through life cycle assessment with varying levels of detail (see the review of over 40 such studies in Appendix \ref{app:04-methods-litreview}), process-based technoeconomic analyses are comparatively sparse (e.g.  \cite{wesselkaemperOptimizationProcessingCosts2025, warrenTechnoEconomicAnalysisLithium2021}), and studies that jointly evaluate both dimensions across multiple resource types are rarer still, with only \citet{mousavinezhadNextLithiumBoom2025}, \citet{gaoLithiumExtractionHard2023}, and \citet{tianIntegratedEnvironmentalEconomic2026} to our knowledge. \citet{mousavinezhadNextLithiumBoom2025} is the only study to evaluate multiple production pathways along the full mine-to-product scope in both dimensions, but it focuses on four projects representing US-based pathways (U.S. evaporative brine, U.S. direct lithium extraction (DLE) from brine, U.S. clay, and spodumene mined in Canada and processed in the US). 

Lepidolite in particular has received relatively limited attention in the lithium literature, and results are notably varied: \citet{liuCarbonFootprintBatteryGrade2025} estimated its global warming potential (GWP) to be the highest among production pathways in China at roughly 24.5 tonnes \ce{CO2eq}/t \ce{Li2CO3}. In contrast, \citet{wangLifeCycleEnvironmental2026} estimated a much higher GWP impact of lepidolite at 66.0 tonnes \ce{CO2eq}/t \ce{Li2CO3}. Lepidolite's production costs are also relatively poorly characterized in the open literature.  \citet{gaoLithiumExtractionHard2023} report that the production cost of lithium carbonate from lepidolite concentrate via the sulfate roasting method was about \$4,500 per tonne \ce{Li2CO3}, excluding the cost of lepidolite concentrate production. \citet{tianIntegratedEnvironmentalEconomic2026} do not report production costs directly; instead, they analyze the profitability of joint lithium and byproduct production from lepidolite concentrate, using plant-level inventories compiled from environmental impact assessment reports. Multiplying their reported inventories by their reported prices yields an implied operating cost of roughly \$9,800/t \ce{Li2CO3}. This reconstructed cost includes a reported purchased concentrate cost, but excludes labor, maintenance, and capital, and is gross of byproduct credits. \citet{tianIntegratedEnvironmentalEconomic2026} also conduct a life cycle analysis, but allocate burdens across multiple coproducts and report results only graphically, and thus thierr results cannot be re-expressed on a common basis with this study. Outside of these two peer-reviewed studies, news reports and grey literature have identified lepidolite production costs to vary widely between \$10,000 and \$16,000 per tonne \cite{atkinsonLithiumPriceReversion2023, liuChinaLithiumBoom2024, jamasmieCATLCutLithium2024}. We plot the \citet{tianIntegratedEnvironmentalEconomic2026} cost and grey literature values, together with literature estimates for the Silver Peak and Thacker Pass operations, in Fig.~\ref{fig:04-costemissions}, noting that each estimate comes from a different study with different scopes and assumptions.

As a step towards addressing these literature gaps, we model the production costs and cradle-to-gate greenhouse gas emissions of three lithium projects representing distinct resource-pathway combinations: U.S. brine-evaporation, U.S. clay-sulfuric acid, and Chinese lepidolite-sulfate roasting. We define flowsheets for each project as shown in Fig. \ref{fig:04-step_breakdown}, based on publicly available project disclosures, patents, and other literature, and we build bottom-up mass-balanced `process-based models based on these flowsheets. We aim to assess whether the cost and emissions advantages suggested by resource type hold, and what this implies for the competitiveness of lithium production in the United States relative to lepidolite in China as the current marginal producer. 

\begin{figure}
    \centering
    \altfig{Three process flow diagrams showing the step-by-step production sequences for each lithium project modeled. Panel a shows Silver Peak, an evaporative brine operation. Panel b shows Thacker Pass, a clay-based operation. Panel c shows Jianxiawo, a lepidolite hard-rock operation. Each diagram shows production steps as rounded boxes grouped within shaded facility boundaries, with inward arrows indicating key reagent inputs and outward arrows indicating coproducts or waste streams at each step.}
        \includegraphics[width=0.76\linewidth]{figures/04-stepwise.png}
    \altfigend
    \caption{Process overview diagrams for each project modeled as representatives of lithium production pathways: (a) U.S. evaporative brine (Silver Peak); (b) U.S. clays (Thacker Pass); (c) Chinese lepidolite (Jianxiawo)}
    \footnotesize{Each subfigure shows the sequence of production steps, key reagents consumed (arrows inward) and coproducts or waste streams produced (arrows outward). Each shaded box in the background represents a facility, where production steps (rounded boxes on top of shaded boxes) are co-located, with the same sitewide operating parameters.}
    \label{fig:04-step_breakdown}
\end{figure} 


\section{Results and Discussion}\label{sec:04-results}
\begin{figure}[!h]
    \centering
    \altfig{Scatter plot with production cost in dollars per tonne of lithium carbonate equivalent on the x-axis and life cycle greenhouse gas emissions in tonnes \ce{CO2}-equivalent per tonne lithium carbonate on the y-axis. Three projects are plotted: Silver Peak at $\sim$\$5,500 and 6.0 tonnes \ce{CO2eq}, Thacker Pass at $\sim$\$8,700 dollars and 16.7 tonnes \ce{CO2eq}, and Jianxiawo at $\sim$\$12,100 an 30.3 tonnes \ce{CO2eq}. Each project has a darker shaded rectangle showing the sensitivity range from this study and a lighter rectangle spanning the range of values reported in the literature, with individual literature data points marked as lines. A dark grey spur extends from the Jianxiawo midpoint showing modeled cost reduction (approximately 470 dollars less) under a hazardous waste non-compliance scenario. A similar spur also extends rightward from the Thacker Pass midpoint showing the modeled cost increase (approximately \$7,200 more) under a sulfur supply-shock scenario.}
        \includegraphics[width=\textwidth]{figures/04-costemissions.png}
    \altfigend
    \caption[Production cost (\$/t \ce{Li2CO3}) versus life cycle greenhouse gas emissions (t \ce{CO2eq}/t \ce{Li2CO3}) for three key lithium projects representing distinct resource-pathway combinations.]{Production cost (\$/t \ce{Li2CO3}) versus life cycle greenhouse gas emissions (t \ce{CO2eq}/t \ce{Li2CO3}) for three key lithium projects representing distinct resource-pathway combinations: U.S. evaporative brine (Silver Peak), U.S. clay (Thacker Pass), and Chinese lepidolite (Jianxiawo)}
    \small{Darker shaded rectangles span the range from this study, with the central point indicating the study midpoint; these ranges reflect variation in the three largest sources of cost and emissions uncertainty for each project, identified through sensitivity analysis (Fig. \ref{fig:04-silverpeak-sensitivity}, Fig. \ref{fig:04-thackerpass-sensitivity}, Fig. \ref{fig:04-jianxiawo-sensitivity}). The complete uncertainty range incorporating all parameters simultaneously is shown in Fig. \ref{fig:c-extendedresults}; the ranges shown here are intended to represent plausible operating variation rather than absolute worst-case bounds. The dark grey spur extending from the Jianxiawo midpoint shows the modeled cost reduction under a hazardous waste non-compliance scenario, in which lepidolite roasting residues are disposed of as non-hazardous waste rather than under current Chinese regulations (see Section~\ref{sec:04-results}). A similar spur also extends rightward from the Thacker Pass midpoint showing the modeled cost increase (approximately \$7,200 more) under a sulfur supply-shock scenario, equivalent to the price shock seen due to the closure of the Strait of Hormuz in early 2026. Lighter rectangles span the range of values reported across grey and peer-reviewed literatures, with each line representing one data point: A: \citet{srkconsultingSECTechnicalReport2023a}; B: \citet{sgscanadaincSK1300Technical2024}; C: \citet{ambroseUnderstandingFutureLithium2020a}; D: \citet{mousavinezhadNextLithiumBoom2025}; E: \citet{flemingLongrunAvailabilityMineral2024}; F: \citet{wesselkaemperOptimizationProcessingCosts2025}; G: \citet{tianIntegratedEnvironmentalEconomic2026}; H: \citet{chenLithiumSpotPrice2023}; I: \citet{atkinsonLithiumPriceReversion2023}; J: \citet{liuChinaLithiumBoom2024}; K: \citet{liLithiumWeaknessPushes2024}; L: \citet{sunRobustAssessmentsLithium2025}; M: \citet{manjongLifeCycleModelling2021}; N: \citet{ambroseUnderstandingFutureLithium2020b}; O: \citet{schenkerCurrentFutureImpacts2025}; P: \citet{khakmardanLifeCycleAssessment2024}; Q: \citet{iyerLifecycleAnalysisLithium2024}; R: \citet{royLithiumClayAssessing2024}; S: \citet{wangLifeCycleEnvironmental2026};  T: \citet{liuCarbonFootprintBatteryGrade2025}. See Table \ref{tab:c1-litreview-matrix} for a summary of all literature studies.}
    \label{fig:04-costemissions}
\end{figure}

Our results indicate that US-based lithium production from evaporative brine and sulfuric acid-based clay pathways have lower production costs and life cycle emissions than production from Chinese-based lepidolite sulfate roasting pathways. The Silver Peak project occupies the lowest-cost, lowest-emission frontier across both our estimates and the literature, at an estimated cost of about \$5,500 [\$4,500-\$6,800] per tonne LCE, and 6.0 [4.7-7.3] tonnes \ce{CO2eq} per tonne LCE. Modeled production costs and emissions from the clays in the generic version of the Thacker Pass flowsheet are closer to modeled production from lepidolite from Jianxiawo in comparison, though lower in both midpoint costs (\$8,700 [\$6,700-\$12,000] vs. \$12,100 [\$11,200-\$13,900]) and emissions (16.7 [13.5-19.4] tonnes \ce{CO2eq} vs. 30.3 [24.5-40.4] tonnes \ce{CO2eq}). 

Our modeled ranges for all three projects are broadly consistent with reported literature values (Fig. \ref{fig:04-costemissions}), meaning our internally coherent combined cost and emissions model structure described in Section~\ref{sec:04-scopemethod} is also externally coherent. The generally wide ranges for lepidolite across both the literature and our study also reflects the data scarcity in reported production costs and emissions for this pathway, underscoring the importance of better understanding this major source of lithium production.

From this, we see that Silver Peak's cost and emissions advantage is largely robust across all modeled ranges, but the ranges aiming to represent plausible operating variation for Thacker Pass and Jianxiawo partially overlap in terms of cost. This is further mediated by two scenarios: where hazardous wastes generated by Jianxiawo's production process are treated as if they were non-hazardous under regulatory non-compliance, which slightly increases Jianxiawo's cost-competitiveness, and sulfur price volatility affecting Thacker Pass operating costs, substantially decreasing that project's viability. 

In the former case, we consider a scenario in which impurity sludge is handled as non-hazardous waste. Under this scenario, tailings disposal costs fall to approximately \$640/t, a reduction of roughly forty percent. Thus, even under non-compliance, treatment of waste and residue management (waste rock piling, tailings, and wastewater treatment) still remain a substantial cost driver, because the flotation and classification tailings streams are generated at very large volumes relative to product output, in spite of low direct unit disposal costs. This tailings cost is notable in the context of a sometimes prevailing assumption that Chinese environmental regulations permit structurally lower production costs than Western jurisdictions: hazardous waste compliance at Jianxiawo adds a cost burden largely absent from the U.S. projects modeled here, though this is not the dominant cost, and the fundamentals of processing low-grade ore set a floor that no regulatory environment can fully offset. 

In the other case, we consider the impact of short-term price spikes in sulfuric acid, in light of geopolitical conflict and security concerns that motivate domestic lithium production broadly. Sulfur is primarily produced as a byproduct of petroleum refining and natural gas `sweetening', and thus its price is set by these industries. Sulfur supply is thus nearly inelastic to sulfur demand, and lithium production accounts for a negligible share of global consumption, leaving lithium producers that depend on sulfuric acid as price-takers in a market they cannot influence. We model a supply-shock scenario at \$1.06/kg under the closure of the Strait of Hormuz~\cite{silveiralimaGlobalSulfurPrices2026}. Applied to the midpoint cost scenario, the direct cost of sulfur rises to \$9,200/t \ce{Li2CO3}, with a corresponding increase in production costs to \$15,900/t \ce{Li2CO3}. This is above Jianxiawo's conservative bound, but the gap would narrow somewhat since an increase in sulfur costs would likely raise Jianxiawo's sulfate reagent costs, though through weaker, less direct channels. Unlike hazardous waste classification at Jianxiawo, which sets a cost floor, a price spike of this magnitude is unlikely to be permanent across the entire life-of-mine. Still, this sensitivity dampens Thacker Pass's global competitiveness if sulfur prices remain elevated for extended periods of time. 

On the life-of-mine basis used throughout, the emissions ranking is unambiguous: Thacker Pass's modeled emissions fall entirely below Jianxiawo's, and the two ranges are disjoint under both the representative and full uncertainty analyses. The cost ranking is favorable but narrow, with Thacker Pass's conservative estimate undercutting Jianxiawo's midpoint by roughly one percent, a margin within the modeling uncertainties discussed below. We explore the sensitivity of this analysis to project-specific assumptions in the sections below.

\subsection{Cost and emissions of producing lithium via U.S. evaporative brine}\label{sec:04-silverpeak}

The Silver Peak project, located in the historic silver mining town of the same name in southwestern Nevada, is the only active operating lithium producing facility within the United States. Currently owned by Albemarle, this project has produced lithium since 1966 through a traditional evaporative brine, batched-chemical process, with a current production target of roughly 6000 tonnes of lithium carbonate each year \cite{srkconsultingSECTechnicalReport2023a}. In this pathway, lithium-bearing brine is pumped from subsurface aquifers and concentrated through solar evaporation in large, shallow ponds before downstream chemical purification yields battery-grade \ce{Li2CO3} (Fig. \ref{fig:04-step_breakdown}a). Most desert salar brines are composed of sodium chloride solutions with lithium, magnesium, calcium and potassium ions \cite{garrettHandbookLithiumNatural2004}; as water evaporates from brines, many of these ions naturally precipitate from the solution. However, magnesium and lithium are typically soluble, and magnesium removal is typically the most chemically demanding purification step for lithium brines \cite{torresSalarAtacamaLithium2024, flemingLongrunAvailabilityMineral2024}. As opposed to South American salars such as Salar de Atacama ($\sim$2,000 mg/L), Silver Peak's brine resource is relatively dilute at approximately 95 \si{mg/L}, requiring extensive evaporation before chemical processing can begin; in our model, approximately 99\% of water in the brine evaporates as the brine concentrates to 6,750 PPM strong brine. However, Silver Peak's relatively low magnesium content allows for an all-chemical, batched precipitation process with lime and soda ash, a simpler production method than solvent extraction or ion exchange processes necessary for other salar brines~\cite{srkconsultingSECTechnicalReport2023a, flexerLithiumRecoveryBrines2018}. While these reagents end up as the single largest contributor on both cost and emission axes, Silver Peak's cost and emissions profile is still the most favorable of the three projects modeled here and thus globally competitive. The structural reliance on solar evaporation for lithium concentration shifts a substantial share of its cost and emissions onto pond maintenance and upkeep, at the cost of requiring significant water losses to the environment. 

Evaluating the step-by-step cost breakdown (Fig. \ref{fig:04-silverpeak-summary}a), we see that three steps dominate: batch treatment ($\sim$\$1,910/t), evaporation ponds ($\sim$\$1,340/t), and brine pumping ($\sim$\$980/t). The evaporation pond cost is predominantly labor and the "Other" category, here represented primarily by maintenance operations on the large pond footprint. Pond-level pumping electricity is assumed to be included with the pumping step, another major cost contributor. Batch treatment and lithium carbonate precipitation ($\sim$\$845/t) are the two steps with significant material costs, and soda ash is the primary reagent cost driver at both. Brine softening consumes more soda ash than carbonate precipitation does, as the calcium introduced by the lime dose (used to remove contained magnesium within the concentrated brine) must itself be removed. Reagents account for roughly 28\% of operating cost, of which soda ash alone is 24\%, so the pathway is neither purely reagent-driven nor insensitive to global commodity markets, though the pond capital and labor base remains the single largest non-reagent exposure.

From an emissions perspective, batch chemical treatment ($\sim$2.8 t \ce{CO2eq}/t), brine pumping ($\sim$2.0 t \ce{CO2eq}/t), and lithium carbonate precipitation ($\sim$1.1 t \ce{CO2eq}/t) are by far the dominant emissions steps (Fig. \ref{fig:04-silverpeak-summary}b). The brine pumping is Scope 2, driven by electricity consumption for pumping using an emissions factor based on the WECC's Northwest Power Pool (NWPP) grid \cite{usepaSummaryData2025}, while the other two are Scope 3, from the embodied emissions of the lime and soda ash used in the process. While soda ash production in the United States is generally cleaner ($\sim$0.3-0.7 t \ce{CO2eq}/t \ce{Na2CO3}) due to its direct production from Trona rather than the heat-intensive Solvay process used in China ($\sim$0.7-1.2 t \ce{CO2eq}/t \ce{Na2CO3}) \cite{wangGreenhouseGasesRegulated2024, scottCanSyntheticSoda2023, royLithiumClayAssessing2024}, the heavy use of lime for magnesium removal ($\sim$0.83-1.7 t \ce{CO2eq}/t \ce{CaO} at 1.2 tonnes \ce{CaO} per tonne \ce{Li2CO3}) and soda ash to remove the calcium introduced by the lime and precipitate lithium carbonate (5.1 tonnes \ce{Na2CO3} per tonne \ce{Li2CO3}) results in significant impacts from the embodied emissions associated with the production of these upstream inputs. However, this emissions profile does also indicate that pumping is by far the most significant electricity consumer, and that decarbonization of the Nevada grid would have a direct and substantial impact on Silver Peak's emissions intensity.

\begin{figure}
    \centering
    \altfig{Two-panel stacked bar chart for the Silver Peak brine project. Panel a shows production cost in dollars per tonne lithium carbonate broken down by process step. Panel b shows greenhouse gas emissions in tonnes \ce{CO2}-equivalent per tonne lithium carbonate broken down by the same steps. Error bars on each step show the sensitivity envelope from varying machine design parameters and material input costs between conservative and optimistic values.}
        \includegraphics[width=1\linewidth]{figures/04-Silver Peak-Steps-Summary.png}
    \altfigend
    \caption{Estimated U.S. evaporative brine costs (a) and emissions (b) by process step.}
    \small{Error bars represent a sensitivity envelope formed by simultaneously applying conservative and optimistic values for all machine design parameters and material input costs that define scenario variations; they do not reflect variation in process chemistry, reagent stoichiometry, or steps for which no scenario variation is defined (e.g., tailings management). As a result, the envelope is wider for steps with more varied parameters and narrower or absent for steps with few or none. Error bars should therefore be interpreted as indicative of the dominant sources of operating uncertainty rather than as statistical confidence intervals.}
    \label{fig:04-silverpeak-summary}
\end{figure}

\begin{figure}
    \centering
    \altfig{Two-panel tornado chart for the Silver Peak project. Panel a shows the top five inputs with the greatest influence on operating cost, centered on $\sim$\$5,500 per tonne lithium carbonate. Panel b shows the top five inputs with the greatest influence on \ce{CO2} emissions, centered on approximately 6.0 tonnes \ce{CO2}-equivalent per tonne lithium carbonate. Bars extend left for optimistic and right for conservative input values, showing the magnitude of each input's effect on the total.}
        \includegraphics[width=1\linewidth]{figures/04-Silver Peak-Sensitivity.png}
    \altfigend
    \caption{Tornado sensitivity plots for operating costs (a) and carbon dioxide emissions (b) for lithium production via U.S. evaporative brine.}
    \small{The top five inputs with the greatest influence on each outcome are shown. The central vertical line indicates the midpoint estimate ($\sim$\$5,500/t \ce{Li2CO3} and $\sim$6.0 t \ce{CO2eq}/t \ce{Li2CO3}, respectively). Bars extend left and right in optimistic and conservative using estimates of input costs or impact factors for materials and operating designs for machine block models, respectively. Machine inputs represent design scenarios rather than continuous parameter variation, with labels representing the most significant design variation relating to the measured output; machine labels represent actual use or throughput-adjusted consumption at each scenario. Inputs whose scaled variation produced limited cost or emissions impact are not labeled.}
    \label{fig:04-silverpeak-sensitivity}
\end{figure}

The cost and emissions sensitivities shown in Fig. \ref{fig:04-silverpeak-sensitivity} tell somewhat different stories. For costs, the evaporation pond capital and operating assumptions and the soda ash price nearly tied at the top --- \$969/t and \$961/t of range respectively --- followed by brine pump costs at \$376/t, meaning pond sizing and reagent pricing have comparable uncertainty ranges, even though maintenance and labor costs are relatively small (0.3-0.5\% of total capital costs as annual maintenance operations costs, i.e. 8-12 laborers managing salt removal and dike upkeep). For emissions, the embodied emissions of soda ash dominate, with its uncertainty spanning roughly 2.0 t \ce{CO2eq}/t. With lime having the third largest uncertainty range, emissions uncertainties are primarily driven by upstream supply chain assumptions rather than on-site operations. Thus, while soda ash's price and embodied emissions move both the cost of production and overall emissions, other variables tend to more strongly affect only one or the other, and policymakers may be able to achieve cost and emissions objectives through distinct levers. 

Several modeling simplifications limit this study. Most substantively, several machine blocks (particularly thickening and dewatering, and washing and purification) are parameterized for throughputs characteristic of much larger operations (i.e. Thacker Pass and Jianxiawo); at Silver Peak's modest production scale of about 6,000 t \ce{Li2CO3} per year, these steps currently contribute negligibly to modeled costs and emissions. Additionally, other modeling simplifications are made given limited site-specific engineering data, such as incorporating pond circulation pumping costs into brine extraction pumping costs, and combining the liming pre-treatment step into the batch chemical purification step. That said, the relative consistency between our modeled range and the literature range on both cost and emissions axes, and the relative narrowness of that literature range compared to the other two projects, suggests that these simplifications do not substantially distort the overall cost and emissions competitiveness picture for a project that has been operating and reported on for nearly six decades.

With Silver Peak's relatively limited resources, its limited ability to expand due to water withdrawal limitations, and the overall lack of similar brine resources in the United States that can be exploited through the traditional evaporative pathway \cite{wesselkaemperOptimizationProcessingCosts2025}, clay resources represent a credible pathway for expanding domestic lithium production capacity, and are the next focus of this analysis. 

\subsection{Cost and emissions of producing lithium via U.S. claystone}\label{sec:04-thackerpass}

The Thacker Pass project, located in Humboldt County, Nevada, is a pre-production, under-construction lithium clay deposit operated by Lithium Americas Corporation. The project targets production of lithium from hectorite clay, a magnesium-rich sedimentary mineral in which lithium and magnesium occupy the same structural positions in the clay lattice \cite{castorLithiumRichClaystoneMcDermitt2020, zhaoRecentAdvancesLithium2023}. Unlike hard rock ores such as spodumene or lepidolite, which require high-temperature roasting or calcination to break down the mineral structure before lithium can be extracted \cite{gaoLithiumExtractionHard2023}, the clay structure at Thacker Pass is amenable to direct acid leaching at moderate temperatures, avoiding the most energy-intensive step of those pathways \cite{royLithiumClayAssessing2024}. However, because sulfuric acid dissolves both lithium and magnesium simultaneously, large quantities of magnesium enter the leach solution and must subsequently be removed, making magnesium management the central processing challenge of the clay pathway — analogous to its role in brine processing, but at substantially higher concentrations. 

As described in Section~\ref{sec:04-scopemethod}, this project is modeled using a generic acid-leach clay flowsheet rather than the project-specific Thacker Pass design. The project-specific design for the Thacker Pass project incorporates a magnesium crystallization route that enables recycling of magnesium compounds as reagents, an important optimization given the magnesium-rich composition of the ore body, at the cost of additional energy usage. Lithium Americas' latest project disclosure has defined a staged production ramp-up across four phases targeting 40,000, 80,000, 120,000, and 160,000 t \ce{Li2CO3} per year from a resource averaging 2,538 PPM lithium \cite{sgscanadaincSK1300Technical2024}. This analysis models Thacker Pass at 66,800 t \ce{Li2CO3} per year, consistent with the production volume assumed in the prior feasibility disclosure \cite{rothFeasibilityStudyNational2022}, to facilitate direct comparison with previous studies (e.g. \cite{wesselkaemperOptimizationProcessingCosts2025, royLithiumClayAssessing2024}). Thus, while this is not the lifetime average across the staged production plan, we aim to represent a useful intermediate snapshot of the project's operating economics and externalities.

\begin{figure}
    \centering
    \altfig{Two-panel stacked bar chart for the Thacker Pass clay project. Panel a shows production cost in dollars per tonne lithium carbonate broken down by process step. Panel b shows greenhouse gas emissions in tonnes \ce{CO2}-equivalent per tonne lithium carbonate broken down by the same steps. Error bars on each step show the sensitivity envelope from varying machine design parameters and material input costs between conservative and optimistic values. On the sulfur burning step in both panels, the exported electricity credit is drawn as a hatched segment overlaid on the right-hand end of the bar, so that the bar tip reads the gross figure before the credit and a black diamond marks the net figure after it.}
        \includegraphics[width=1\linewidth]{figures/04-Thacker Pass-Steps-Summary.png}
    \altfigend
    \caption{Estimated U.S. clay costs (a) and emissions (b) by process step.}
    \small{Error bars represent a sensitivity envelope formed by simultaneously applying conservative and optimistic values for all machine design parameters and material input costs that define scenario variations; they do not reflect variation in process chemistry, reagent stoichiometry, or steps for which no scenario variation is defined (e.g., tailings management). As a result, the envelope is wider for steps with more varied parameters and narrower or absent for steps with few or none. Error bars should therefore be interpreted as indicative of the dominant sources of operating uncertainty rather than as statistical confidence intervals. Sulfur burning also includes a credit rather than only a cost burden because the on-site sulfuric acid plant exports power to the rest of the facility; the exported electricity credit is drawn as a hatched segment overlaid on the right-hand end of the bar, so the bar tip reads the gross figure before the credit, and the black diamond marks the net figure after it. Under optimistic assumptions the credit exceeds the embodied emissions of the purchased sulfur and the step is net-negative at $\sim$-0.33 t \ce{CO2eq}/t, which is why its error bar in (b) crosses zero while its bar does not.}
    \label{fig:04-thackerpass-summary}
\end{figure}

The step-level cost breakdown (Fig. \ref{fig:04-thackerpass-summary}a) shows that costs are primarily composed of three key blocks: the hydrometallurgical circuit from acid leaching through lithium carbonate precipitation ($\sim$\$3,340/t), of which impurity removal alone is \$2,000/t; on-site sulfuric acid production at \$2,120/t,
driven almost entirely by the cost of sulfur (\$1,980/t); and waste handling with a total cost of \$1,570/t, split between the waste rock pile (\$800/t) and tailings management (\$770/t). The direct mining operation is smaller than the waste it produces, with blasting and loading together at \$780/t. Additionally, as sulfuric acid production from sulfur combustion is highly exothermic, the heat can be captured to produce electricity. The value of the power exported, calculated as a direct credit at the local electricity price (i.e. negative power consumption), is $\sim$\$370/t \ce{Li2CO3} at the midpoint (\$260 conservative / \$480 optimistic), a significant credit to the cost of lithium carbonate production (4.3\% (2.0\% to 7.2\%) of cost), as shown in Fig. \ref{fig:04-thackerpass-summary}a as the hashed segment of the sulfur burning step bar. 

The emissions breakdown (Fig. \ref{fig:04-thackerpass-summary}b) mirrors this concentration, with most emissions coming from the embodied emissions of reagents, but with a difference in ordering: impurity removal is the single largest emissions contributor at $\sim$12.1 t \ce{CO2eq}/t \ce{Li2CO3}, or 72\% of the total, despite being 23\% of the cost, reflecting the high embodied emissions of impurity-removal reagents (primarily lime). The sulfuric acid route has the opposite relationship:  sulfur burning and acid leaching together account for $\sim$2.7 t \ce{CO2eq}/t ($\sim$16.5\% of emissions) against 24\% of cost. However, because burning sulfur releases no carbon and the exported power displaces grid electricity, a credit worth 1.3 t \ce{CO2eq}/t at the midpoint can be attributed to this step, as seen with the hashed bar (Fig. \ref{fig:04-thackerpass-summary}b). Scope 2 emissions appear in electricity-intensive mechanical steps such as crushing and attrition scrubbing, while Scope 1 emissions are largely confined to mining operations (diesel for vehicles), but both are relatively minor. This emissions structure implies that decarbonization of the Thacker Pass pathway depends primarily on reducing the embodied emissions of key reagents, especially lime and soda ash, rather than decarbonizing on-site electricity or fuel use. 

That dependence is concentrated in a single reagent -- lime -- and it is also the reason our cost estimate sits above the disclosed one. Our generic flowsheet results in an estimated midpoint cost of $\sim$\$8,700/t \ce{Li2CO3}, exceeding the reported \$8,040/t life-of-mine cost of the latest disclosure \cite{sgscanadaincSK1300Technical2024}, the
\$7,200/t of the 2022 disclosure \cite{rothFeasibilityStudyNational2022}, and the \$6,300--\$6,990 reported elsewhere. As noted in
Section~\ref{sec:04-scopemethod}, the generic flowsheet fully precipitates magnesium with quicklime, consuming approximately 9,900 kg \ce{CaO}/t \ce{Li2CO3}, compared to 2,630 kg/t in the project-specific magnesium crystallization design. That difference results in an additional \$1,280/t LCE of cost and 8.46 t CO2eq/t LCE of emissions in our model, which would simultaneously reduce the single largest cost and emissions step by about two-thirds. Of course, the magnesium crystallization route would be a major cost and emissions driver as well, with high electricity and steam requirements needed to evaporate water from the process slurry to induce magnesium precipitation instead of the chemical route, so the net benefit of this switch would be muted. So, while this all-chemicals route is more flexible and can be generically applied to other clay production routes, refining the model to incorporate this project-specific route is an avenue for potential future work on the competitiveness of Thacker Pass specifically. 

\begin{figure}
    \centering
    \altfig{Two-panel tornado chart for the Thacker Pass project. Panel a shows the top five inputs with the greatest influence on operating cost, centered on $\sim$\$8,700 per tonne lithium carbonate. Panel b shows the top five inputs with the greatest influence on \ce{CO2} emissions, centered on approximately 16.7 tonnes \ce{CO2}-equivalent per tonne lithium carbonate. Bars extend left for optimistic and right for conservative input values, showing the magnitude of each input's effect on the total.}
        \includegraphics[width=1\linewidth]{figures/04-Thacker Pass-Sensitivity.png}
    \altfigend
    \caption{Tornado sensitivity plots for operating costs (a) and carbon dioxide emissions (b) for lithium production via U.S. clays.}
    \small{The top five inputs with the greatest influence on each outcome are shown. The central vertical line indicates the midpoint estimate ($\sim$\$8,700/t \ce{Li2CO3} and $\sim$16.7 t \ce{CO2eq}/t \ce{Li2CO3}, respectively). Bars extend left and right in optimistic and conservative using estimates of input costs or impact factors for materials and operating designs for machine block models, respectively. Machine inputs represent design scenarios rather than continuous parameter variation, with labels representing the most significant design variation relating to the measured output; machine labels represent actual use or throughput-adjusted consumption at each scenario. Inputs whose scaled variation produced limited cost or emissions impact are not labeled.}
    \label{fig:04-thackerpass-sensitivity}
\end{figure}

The dominant source of cost uncertainty in the Thacker Pass Project is the price of sulfur (Fig. \ref{fig:04-thackerpass-sensitivity}a), modeled with a midpoint of \$0.228/kg~\cite{intratecSulfurPriceHistorical} and bounded by extremes of \$0.0748~\cite{mullinsTampaMoltenSulphur2024,usgeologicalsurveyMineralCommoditySummaries2025} and \$0.4734~\cite{usgeologicalsurveyMineralCommoditySummaries2023} per kilogram of sulfur, based on reported contract settlements from the third quarter of 2024 and second quarter of 2022, respectively. This midpoint sulfur cost implies \$80/t \ce{H2SO4} on an operating basis and \$91/t
including annualized capital, in the same band as the \$88/t on-site figure implied by the 2022
project disclosure \citep{rothFeasibilityStudyNational2022}, even though the plant's capital cost,
staffing and equipment are constructed independently from~\cite{davenportSulfuricAcidManufacture2006,
ulrichGuideChemicalEngineering1984}. At a consumption of 8,680 kg sulfur per tonne \ce{Li2CO3}, this swings the total sulfur cost from \$650/t \ce{Li2CO3} to \$4,110/t \ce{Li2CO3}. Beyond sulfuric acid, soda ash and lime introduce secondary reagent cost uncertainty of similar magnitude to one another.

Process water costs represent the second largest source of cost uncertainty. Water in arid, remote Nevada is relatively expensive compared to other mining regions, and our assumed range of \$0.50---\$5.00/\si{m^3} produces meaningful cost variation at the volumetric consumption implied by the generic flowsheet. That consumption is likely overestimated: our model reports a net freshwater makeup of 13.5 million \si{m^3/yr} at 66,800 \si{t/yr} against 16.5 million \si{m^3/yr} at 160,000 \si{t/yr} in the Phase 4---5 disclosure \cite{sgscanadaincSK1300Technical2024}. The gross circulating draw is 43.4 million \si{m^3/yr}, of which 69\% is credited as recycle. We model recycling as a structural make-up within three circuits using generalized dewatering and mother liquor steps, rather than the project-specific zero-liquid discharge conditions prescribed in the disclosure and permits. Furthermore, ore is reported to be slurried on-site prior to transport to the processing facility in the latest disclosure \cite{sgscanadaincSK1300Technical2024}; we instead model slurrying after the crushing step, in the attrition scrubbing step, which likely understates both transport costs and the water balance at the mine face. Given local opposition to the project's water usage \cite{alonzoenvironmentQuestionsWaterRights2025}, explicit modeling of water recycling loops is a necessary refinement before drawing conclusions about Thacker Pass's water intensity relative to the disclosed design, but our model represents water demand for a baseline, partially optimized generic plant.

The emissions sensitivity (Fig. \ref{fig:04-thackerpass-sensitivity}b) is driven by the embodied emissions factors of the reagents, with sulfur and lime carrying comparable ranges of $\sim$2.3 and $\sim$2.1 t \ce{CO2eq}/t, ahead of soda ash. Thus, while lime contributes less to cost uncertainty, it has a much higher impact on emissions. Further, because approximately two-thirds of lime's embodied emissions are process emissions chemically unavoidable from calcination \cite{lavegliaDecarbonizingLimebasedConstruction2023}, reducing the clay pathway's emissions exposure requires either substituting lime in the impurity removal step with process innovation (e.g. the magnesium recycling process described above) or sourcing lime from producers deploying carbon capture. This makes the clay pathway's emissions trajectory dependent on upstream industrial decarbonization decisions in the lime sector in a way that is structurally distinct from the electricity-dominated emissions profile of the brine pathway. 

Taken together, the generic flowsheet used here is conservative relative to the disclosed Thacker Pass design in both cost and emissions, primarily due to the absence of the magnesium crystallization route. The choice to incorporate these production steps into the Thacker Pass-specific flowsheet appears to be beneficial not only in terms of economics, but also carbon emission externalities. Despite these limitations, the modeled results suggest that Thacker Pass remains broadly competitive with Chinese lepidolite production even under conservative assumptions --- a finding that is likely to strengthen if the model were refined toward the project-specific design.

\subsection{Cost and emissions of producing lithium via Chinese lepidolite}\label{sec:04-jianxiawo}

The Jianxiawo lepidolite project, located in Yichun, Jiangxi Province, is operated by Contemporary Amperex Technology Co. Limited (CATL), currently the world's largest battery manufacturer. As China's largest lithium source, Jianxiawo provides around 3\% of global supply and 14\% of China's supply \cite{liChinasCATLSuspends2025, andrezejewskiResourcesEnergyQuarterly2025, usgeologicalsurveyMineralCommoditySummaries2026}, making it the single most consequential source of marginal lithium supply on the global cost curve.

Lepidolite is a fluorine-rich lithium mica mineral with a low bulk lithium grade, with Jiangxi lepidolite ore at approximately 0.27---0.4\% \ce{Li2O} \cite{sulcaLithiumPricesWhere2024}, or roughly 1250 to 1850 PPM Li.  Lepidolite's complex mineralogy requires high-temperature sulfate roasting to liberate the lithium. The high fluorine content in lepidolite can also generate hydrofluoric acid as a byproduct during roasting, necessitating additional reagents (e.g. lime) and generating hazardous waste streams that contribute to both costs and environmental liabilities \cite{yanExtractionLithiumLepidolite2012, dengSystematicReviewSulfate2026}. But this low grade is precisely what has allowed lepidolite projects to rapidly expand in recent years: mines operating under local kaolinite permits were not required to register as lithium mines as long as stated ore grades remained below 0.4\% \ce{Li2O}\cite{zhangChinaHasDomestic2025}. Our model deliberately assumes the upper bound of this range (0.4\%~\ce{Li2O}), to represent the most favorable ore grade at which these mines could operate under the kaolinite permit threshold, and therefore a conservative estimate of lepidolite's cost disadvantage relative to U.S. sources.

CATL's vertical integration, which spans lithium mining through cell manufacturing, is a critical context for interpreting Jianxiawo's economics: the mine has continued operating during periods when lepidolite production costs exceeded spot lithium prices because downstream refining, cathode active material, and battery segments remained profitable \cite{jamasmieCATLCutLithium2024}. Yet the decision to idle the mine between September 2024 and February 2025, despite this integrated structure \cite{benchmarkteamCATLSuspendsOperations}, suggests that lepidolite's standalone production costs exceeded the internal transfer price at which CATL's downstream operations could profitably absorb the material — providing indirect evidence that the operation is genuinely near marginal. The standalone costs modeled here can therefore be interpreted as a meaningful benchmark for the cost at which lepidolite production enters and exits the global supply curve.

\begin{figure}
    \centering
    \altfig{Two-panel stacked bar chart for the Jianxiawo lepidolite project. Panel a shows production cost in dollars per tonne lithium carbonate broken down by process step. Panel b shows greenhouse gas emissions in tonnes \ce{CO2}-equivalent per tonne lithium carbonate broken down by the same steps. Error bars on each step show the sensitivity envelope from varying machine design parameters and material input costs between conservative and optimistic values. A red vertical marker on the tailings bar indicates the modeled cost under a hazardous waste non-compliance scenario.}
        \includegraphics[width=1\linewidth]{figures/04-Jianxiawo-Steps-Summary.png}
    \altfigend
    \caption{Estimated Chinese lepidolite costs (a) and emissions (b) by  step.}
    \small{Error bars represent a sensitivity envelope formed by simultaneously applying conservative and optimistic values for all machine design parameters and material input costs that define scenario variations; they do not reflect variation in process chemistry, reagent stoichiometry, or steps for which no scenario variation is defined (e.g., tailings management). As a result, the envelope is wider for steps with more varied parameters and narrower or absent for steps with few or none. Error bars should therefore be interpreted as indicative of the dominant sources of operating uncertainty rather than as statistical confidence intervals. The red vertical marker on the tailings bar indicates the modeled tailings cost under a hazardous waste non-compliance scenario (see Section~\ref{sec:04-results}).} 
    \label{fig:04-jianxiawo-summary}
\end{figure}

The step-level cost breakdown (Fig. \ref{fig:04-jianxiawo-summary}a) is defined by a single dominant feature: the sulfate roasting step accounts for approximately \$5,400/t \ce{Li2CO3} (44\% of total modeled cost), comprising material costs (primarily sodium sulfate), utility costs from high-temperature kiln operation (primarily natural gas combustion), and maintenance costs for the capital intensive, corrosion-prone roasting train. The second block is waste handling, at \$2,820/t LCE (23\% of cost): the waste rock pile at \$1,480/t, tailings disposal at \$1,100/t, and wastewater treatment at \$235/t. Tailings are expensive because impurity sludge contains fluoride and trace heavy metals (thallium in particular) that are precipitated out of the leach solution; this stream would be classified as hazardous under Chinese regulations, requiring secure landfill disposal~\cite{nationaldevelopmentandreformcommission2024}. The waste rock pile is also expensive, because the very low-grade ore ($\leq 0.4\%$ \ce{Li2O}) requires immense amounts of material to be moved. The impact of classifying this hazardous material as nonhazardous is discussed in Section~\ref{sec:04-results}. Outside of these two major categories, lithium carbonate precipitation contributes a further $\sim$\$960/t in material costs (soda ash), while all other steps are comparatively minor. 

The emissions breakdown (Fig. \ref{fig:04-jianxiawo-summary}b) is similarly dominated by the roasting step at $\sim$19.5 t \ce{CO2eq}/t \ce{Li2CO3}, but with an important structural distinction from the Thacker Pass pathway: roasting emissions are divided roughly equally between Scope 1 direct combustion emissions from high-temperature kiln operation and Scope 3 embodied emissions from sulfate roasting salts (sodium sulfate, potassium sulfate) and lime \cite{yanExtractionLithiumLepidolite2012, tranChapter3Lithium2015}. This means that even full decarbonization of the natural gas combustion component would leave roughly half of roasting emissions unaddressed, and that emissions reductions require both fuel switching and upstream reagent supply chain decarbonization simultaneously. Outside of the crucial sulfate roasting step, lithium carbonate precipitation is the only other notable emissions contributor at $\sim$4.2 t \ce{CO2eq}/t, driven by the embodied emissions of soda ash. It is important to note that this step likely has a higher environmental impact in China than in the United States, given that the soda ash used in this step is likely produced from the heat-intensive Solvay process ($\sim$0.7-1.2 t \ce{CO2eq}/t \ce{Na2CO3}) \cite{wangGreenhouseGasesRegulated2024, scottCanSyntheticSoda2023, royLithiumClayAssessing2024}. All other steps have smaller impacts in comparison.  

\begin{figure}
    \centering
    \altfig{Two-panel tornado chart for the Jianxiawo project. Panel a shows the top five inputs with the greatest influence on operating cost, centered on $\sim$\$12,100 per tonne lithium carbonate. Panel b shows the top five inputs with the greatest influence on \ce{CO2} emissions, centered on approximately 30.3 tonnes \ce{CO2}-equivalent per tonne lithium carbonate. Bars extend left for optimistic and right for conservative input values, showing the magnitude of each input's effect on the total.}
        \includegraphics[width=1\linewidth]{figures/04-Jianxiawo-Sensitivity.png}
    \altfigend
    \caption{Tornado sensitivity plots for operating costs (a) and carbon dioxide emissions (b) for lithium production via Chinese lepidolite.}
    \small{The top five inputs with the greatest influence on each outcome are shown. The central vertical line indicates the midpoint estimate ($\sim$\$12,100/t \ce{Li2CO3} and $\sim$30.3 t \ce{CO2eq}/t \ce{Li2CO3}, respectively). Bars extend left and right in optimistic and conservative using estimates of input costs or impact factors for materials and operating designs for machine block models, respectively. Machine inputs represent design scenarios rather than continuous parameter variation, with labels representing the most significant design variation relating to the measured output; machine labels represent actual use or throughput-adjusted consumption at each scenario. Inputs whose scaled variation produced limited cost or emissions impact are not labeled.}
    \label{fig:04-jianxiawo-sensitivity}
\end{figure}

The cost sensitivity for the Jianxiawo lepidolite pathway (Fig. \ref{fig:04-jianxiawo-sensitivity}a) is shaped by three large contributing variables. Soda ash price produces the widest range at approximately \$965/t, reflecting both the large stoichiometric requirement of soda ash for lithium carbonate precipitation and the wide spread between Chinese Solvay-process and U.S. trona-derived soda ash prices. Sodium sulfate price and the sulfate roasting train design (the latter driven primarily by varying combustion efficiency assumptions in the roasting process) follow at approximately \$940/t and \$800/t, respectively. This is structurally different from Thacker Pass, where a single reagent dominated, and implies that cost reduction at Jianxiawo requires simultaneous optimization across the roasting circuit and multiple reagent supply chains rather than a single-lever intervention. Lime and process water contribute secondary uncertainty of smaller magnitudes. 

In contrast to the cost sensitivity, where uncertainty is distributed roughly equally across three drivers, emissions uncertainty is dominated by a single factor: sodium sulfate embodied emissions span $\sim$28---36 t \ce{CO2eq}/t \ce{Li2CO3} (Fig. \ref{fig:04-jianxiawo-sensitivity}b). The next largest drivers are paired with roughly equal ranges: the roasting machine energy intensity and soda ash follow with $\sim$4 t \ce{CO2eq}/t, while the high-energy grinding mill and lime span about $\sim$1 t \ce{CO2eq}/t. As a result, there is one primary lever for reducing emissions for the Jianxiawo flowsheet, but the mixture of embodied emissions in material inputs and the high energy requirements of machine processes requires a multi-faceted approach for total emissions reduction. 

Several limitations specific to the Jianxiawo model warrant elaboration. Most fundamentally, lepidolite processing is considerably less documented in the open literature than either brine or clay pathways: the sulfate roasting flowsheet is reconstructed from patents filed in China and lab-scale studies \cite{yanExtractionLithiumLepidolite2012, luongFactorsAffectingExtraction2013, parkEffectiveLithiumExtraction2023, dengSystematicReviewSulfate2026, nanMethodsExtractingLithium2022, nanMethodsExtractingLithium2019, zouMethodEfficientLithium2024}, given the lack of S-K 1300 or NI 43-101-equivalent technical report disclosure requirements for Chinese-listed mining companies. This data gap is itself a finding, reflecting the limited public transparency of Chinese lithium production operations, but also means that process parameters such as reagent ratios, roasting temperatures, and conversion efficiencies carry greater epistemic uncertainty than projects owned by U.S. or Canadian-listed companies. 

Similarly, the tailings model is relatively simple: impurity sludge is assigned a hazardous disposal cost based on published Chinese regulatory gate fees \cite{nationaldevelopmentandreformcommission2024}, but the volume and composition of sludge generated are derived from reagent
stoichiometry rather than measured plant data. The reagent allowance, sludge
tonnage, cake solids fraction, and lithium loss at this step all descend from a single declared figure carried at one significant figure, and this remains the single largest epistemic uncertainty in the Jianxiawo estimate.

Additionally, the production volume modeled here, based on the most commonly reported production figure of 46,000 tonnes LCE, is less than half of the highest reported capacity of 111,000 t LCE \cite{benchmarkteamRestartCATLsJianxiawo2026, andrezejewskiResourcesEnergyQuarterly2025}, and thus carries more uncertainty than projects owned by U.S. or Canadian-listed companies, which are anchored to stock market-required engineering disclosures. Given the presence of economies of scale in the roasting circuit, cost and emissions estimates are sensitive to this assumption in ways that are difficult to bound without plant-level data. Finally, as noted above, the standalone costs modeled here do not capture the cost-sharing or transfer-pricing arrangements that may apply within CATL's vertically integrated operations, and thus may not fully reflect the true `marginal' lithium producer in China.

Our bottom-up model suggests that neither the cost nor the emissions profile of Jianxiawo is competitive with Silver Peak. The overlap with Thacker Pass in the cost space is attributable to conservative generic flowsheet assumptions applied to Thacker Pass, and the most favorable ore grade permissible under the kaolinite permit threshold (0.4\% \ce{Li2O}) for Jianxiawo, suggesting the true gap between U.S. clay and Chinese lepidolite production is likely wider than our midpoint estimates indicate. At the same time, our process-based estimates fall between figures reported in the literature and industry reports \cite{gaoLithiumExtractionHard2023, atkinsonLithiumPriceReversion2023, liuChinaLithiumBoom2024}. As a result, lepidolite's cost and emissions disadvantage, while real and structural, may be somewhat less severe than what has been commonly reported. Both conclusions point in the same direction: more detailed, process-grounded modeling of Chinese lepidolite operations is needed to bound these estimates, and represents a priority for future work.

\section{Limitations}\label{sec:04-limitations}

Several important scope limitations of the present analysis should be addressed in future work, beyond those that were discussed for each individual project, such as the generic clay flowsheet and lime consumption at Thacker Pass (Section~\ref{sec:04-thackerpass}), the roasting flowsheet and sludge chain at Jianxiawo (Section~\ref{sec:04-jianxiawo}), and machine
blocks throughput parameterization for Silver Peak (Section~\ref{sec:04-silverpeak}). The following discussion applies across all models. 

First, mine closure and reclamation costs are not modeled here. These costs can be substantial, particularly for open-pit mining operations like Thacker Pass and Jianxiawo, and their exclusion likely understates the full lifecycle cost of all three projects — though the relative magnitude of this omission may differ across jurisdictions given differences in reclamation requirements between the United States and China.

Second, this analysis reports production costs at the mine gate and does not model cost, insurance, and freight (CIF) delivery costs to end markets. Chinese producers may retain a competitive advantage in delivered cost to Asian battery manufacturers even where their production costs are higher, given proximity to demand centers and established logistics infrastructure. U.S. producers serving domestic battery manufacturers could benefit from localized supply chains, but establishing these logistics channels from rural Nevada to battery manufacturing plants primarily located in the U.S. automotive corridor \cite{chengUSIndustrialPolicy2024} may be more costly in comparison to China. A full delivered-cost analysis, incorporating transportation, tariffs, and insurance, would provide a more complete picture of competitive positioning under different trade policy scenarios.

Further, water consumption is an important dimension that has been partially addressed in this analysis, under several engineering assumptions. Water flows are explicitly tracked and dosed at each production step for all three projects, and recycling is credited wherever the returning stream is feasible from an engineering perspective. For example, structural water make-up is credited in the Thacker Pass circuit and the Jianxiawo beneficiation circuit, so 69\% and 39\% of each project's gross circulating draw is met by recycle rather than fresh water. Within the Jianxiawo beneficiation circuit, 60\% of the water demand is credited as recycled. A similar credit cannot be applied  in the water leach circuit, since every stream leaving it carries either lithium not yet recovered or contaminants mobilized from the ore; that circuit alone accounts for 58\% of the project's fresh water draw. More generally, many assumptions were made regarding the solids fractions of output waste streams, and these quantities, in addition to secondary quantities such as seepage and entrainment at tailings facilities, constrain the amount of water that is reusable or requires further treatment. 

Additionally, tailings and residue management costs are modeled using a staged mechanical dewatering mechanism combined with project-specific disposal routes~\cite{sgscanadaincSK1300Technical2024,srkconsultingSECTechnicalReport2023a} and gate fees for hazardous and non-hazardous waste~\cite{nationaldevelopmentandreformcommission2024}. Each stage is parameterized by a fixed cost per tonne of dry solids and a marginal cost per tonne of water removed, calibrated against published dewatering costs and energy intensities and estimations from the grey literature and disclosures~\cite{carneiroConceptualCostComparison2018, klohncrippenbergerStudyTailingsManagement2017, usepaBiosolidsTechnologyFact2000, usepaBiosolidsTechnologyFact2003, usepaBiosolidsTechSheetPlatePress2000, crystalFilterPressedDryStacking2018, smithEnergyConventionalWater2017, sgscanadaincSK1300Technical2024, silvercorpmetalsinc.TailingsStorageFacility2025}. While the dominant cost drivers are directly observed, particularly the hazardous disposal burden at Jianxiawo, some parameters within the model remain as order-of-magnitude placeholders, including on-site tailings facility and wastewater treatment operating costs. These are unlikely to affect the paper's comparative conclusions materially, given that disposal costs at the U.S. projects are small relative to reagent costs, but represent a priority refinement for site-specific validation. More substantively, the environmental impacts of waste disposal, outside of greenhouse gas emissions from electricity usage, are not currently modeled. For the ore-based pathways, where large volumes of leach residue and impurity sludge are generated and need to be dewatered, this omission likely modestly understates total Scope 1 and 2 emissions as a fraction of total emissions, though the dominant emissions drivers identified in the results (reagent embodied emissions and fuel combustion) are unaffected by this gap.

Beyond these refinements to existing models, several extensions would broaden the contribution of this work. The three projects modeled here are representative but not comprehensive: several other clay projects exist in the United States at various stages of development \citep{wesselkaemperOptimizationProcessingCosts2025}, and other lepidolite and spodumene operations in China could be modeled. In particular, the brine-direct lithium extraction (DLE) pathway, with projects under development in both the United States and China, appears unlikely to be marginal at current prices but would be important for characterizing the full domestic cost curves of both countries and understanding how the competitive landscape shifts under higher-price scenarios. Altogether, these additions would enable a richer analysis of how cost curves shift under different demand growth scenarios and how individual projects enter and exit production as prices move, providing an empirical foundation for a more complete treatment of supply curve dynamics.

All of these extensions described above, like this current analysis, are static in nature, characterizing production costs and emissions under a given set of assumptions and uncertainties. However, with project-level cost curves established for a broader set of pathways, they can serve as inputs to dynamic models that can capture firm- and market-level responses to policy interventions, technological learning, and shifting market conditions. The progression from static project baselines to dynamic supply curve modeling is therefore a natural next step, and the framework developed here is designed with that extension in mind.

\section{Conclusions, Policy Implications, Future Work}

%No need to repeat this yet again
%This study developed bottom-up technoeconomic and life cycle assessment models for three lithium production projects representing distinct resource-pathway combinations: Silver Peak (evaporative brine, United States), Thacker Pass (sulfuric acid-leach clay, United States), and Jianxiawo (sulfate roasting lepidolite, China). 
Our results demonstrate that the cost and emissions profiles of lithium production depend strongly on the resource from which it is extracted: the thermochemical and hydrometallurgical requirements of each pathway create structural differences in production costs, emissions intensities, and sensitivity to price volatility that are challenging to overcome with project-specific optimizations. Under central assumptions, both U.S. projects produce lithium carbonate at a lower cost and life cycle emissions than the Chinese lepidolite operation that appears to set the marginal price on the global supply curve, with evaporative brine being distinctly favorable in both dimensions and claystone remaining broadly competitive even under conservative generic flowsheet assumptions. Crucially, while Silver Peak brine has limited ability to expand due to limited land and water rights, the opening of the Thacker Pass clay project, which when operating at maximum production capacity could supply 10\% of current global lithium demand \cite{usgeologicalsurveyMineralCommoditySummaries2026}, could establish the United States as a major producer of lithium as demand continues to grow. 

The United States holds some structural natural resource advantages relevant to lithium production costs and emissions, though these should not be overstated given China's extensive and low-cost chemical processing industry. One example is soda ash: U.S. producers source it from mined trona in Wyoming's Green River Basin, the world's largest natural soda ash deposit \cite{usgeologicalsurveyMineralCommoditySummaries2026}, at substantially lower cost and emissions intensity than the Solvay process predominant in China \cite{wangGreenhouseGasesRegulated2024, scottCanSyntheticSoda2023, royLithiumClayAssessing2024}, conferring a structural advantage across all three pathways modeled here. More broadly, the abundant solar and wind resources of the arid Southwest where U.S. lithium deposits are concentrated, and relatively low natural gas prices represent potential further advantages as energy costs become more significant at larger production scales. These advantages could be offset, due to structural dependencies that run the other way: Thacker Pass's largest cost driver is sulfur, consuming 8,680 kg of sulfur per tonne \ce{Li2CO3}. With sulfur as a byproduct of petroleum refining and gas processing whose price U.S. lithium producers cannot influence, the prices observed following the closure of the Strait of Hormuz in early 2026 cause Thacker Pass's modeled cost to rise to \$15,900/t \ce{Li2CO3},
above Jianxiawo's conservative bound. Thus, a domestic pathway pursued for supply security can itself be exposed to a globally traded feedstock, transferring the risk from one type of material vulnerability to another.

On the other hand, a common assumption in supply chain security discourse is that Chinese production benefits from lower environmental regulatory costs. Our results complicate this narrative: hazardous waste disposal requirements for lepidolite roasting residues contribute approximately \$1,100/t \ce{Li2CO3} (about 9\%) to Jianxiawo's modeled costs — a burden largely absent from the U.S. projects, which remains notable even under a hypothetical non-compliance scenario in which hazardous tailings were treated as non-hazardous (\$640/t \ce{Li2CO3}). Furthermore, there is evidence that Chinese environmental compliance costs for lepidolite operations are likely to increase rather than decrease: as noted in Section~\ref{sec:04-jianxiawo}, many lepidolite mining permits in Jiangxi Province were originally issued as kaolinite permits, and ongoing regulatory tightening around this permitting gap \cite{benchmarkteamCATLSuspendsOperations, zhangChinaHasDomestic2025} could raise operating costs and restrict capacity in ways not captured in current cost curve estimates.

These project-level baselines also could provide the foundation for a systematic analysis of policy interventions. Building on the models developed here, future work could analyze the relative competitiveness of each pathway under existing policies such as including production tax credits under the IRA and OBBBA, Made in China 2025 industrial subsidies, and Chinese export restrictions, or prospective interventions such as a carbon border adjustment mechanism, currently being implemented in Europe. Such analyses would allow policymakers to identify which interventions most efficiently close the competitiveness gap between U.S. and Chinese production and under what conditions domestic lithium supply development becomes self-sustaining without ongoing policy support.

More broadly, this framework for integrated process-level technoeconomic and life cycle emissions assessment across geographic production pathways does not need to be specific to lithium. While the specific projects and economic landscape analyzed here reflect a particular moment in the geopolitics of critical minerals, the framework itself is material-agnostic: wherever production costs and environmental outcomes are jointly shaped by resource type and production location, the same approach could apply. Applied to other critical mineral systems, this approach can contribute to a more rigorous and consistent basis for evaluating supply chain security, environmental performance, and the effectiveness of industrial policy across the materials and manufacturing systems underpinning the energy transition.

\section{Methodology}\label{sec:04-scopemethod}

This work evaluates and compares the production costs and cradle-to-gate greenhouse gas emissions of three lithium projects representing distinct resource-pathway combinations, using a transparent, empirically-grounded technoeconomic-life cycle assessment framework and consistent assumptions across geographical locations. Given that the emissions literature is substantially richer and more methodologically mature than the technoeconomic literature for lithium, consistency between modeled emissions profiles and reported life cycle assessment values can thus provide an independent cross-check on the underlying process model, and therefore the cost estimates derived from it. 

In order of increasing complexity of mineral resource and processing pathway, we define flowsheets for each project as shown in Fig. \ref{fig:04-step_breakdown}, based on publicly available project disclosures, patents, and other literature. In particular, Silver Peak and Thacker Pass are owned and operated by Albemarle and Lithium Americas, respectively, two companies that are publicly traded on U.S. stock exchanges. As a result, under the SEC's S-K 1300 \cite{SECgovModernizationProperty} regulation (built on Canada's NI 43-101 regulation), they are required to disclose mineral resources and reserves through reports written by an independent Qualified Person, from which high-level flowsheets can be determined \cite{srkconsultingSECTechnicalReport2023a, sgscanadaincSK1300Technical2024}. On the other hand, publicly available detailed representations of lepidolite projects in China are limited, given the lack of similar disclosure requirements in China. We focus on Jianxiawo (Yichun, China, \ref{fig:04-step_breakdown}c), the largest lepidolite project in China by production volume: its integration with CATL, the world's largest battery manufacturer, means that its operational changes, including production suspensions and restarts, attract substantial coverage in industry press and commodity market reporting, providing more contextual data than is available for most Chinese lepidolite operations (e.g.~\cite{liChinasCATLSuspends2025,jamasmieCATLCutLithium2024,benchmarkteamRestartCATLsJianxiawo2026,benchmarkteamCATLSuspendsOperations}). Furthermore, its usage of the sulfate roasting pathway, the dominant processing route for lepidolite in China~\cite{dengSystematicReviewSulfate2026}, makes it representative of the pathway, beyond being the largest individual project.

The primary methodology we employ in this analysis is a series of process-based technoeconomic models, from which environmental impacts can be estimated and policy interventions can be applied. The modeling framework is structured hierarchically across four levels of resolution, enabling consistent comparison across projects while preserving sensitivity to pathway-specific technology requirements. At the \textit{supply chain} level, facilities are linked by \textit{transportation routes} into a directed network from raw material extraction to refined \ce{Li2CO3} product. Each \textit{facility} aggregates co-located production steps operating under shared geographic and operational parameters (e.g. electricity grid composition, labor costs, utility prices, and plant schedule), ensuring that location-specific cost and emissions factors are applied consistently. Within each facility, \textit{production steps} represent individual unit operations (e.g., acid leaching, impurity removal, lithium carbonate precipitation), each defined by a primary input and output stream, secondary reagent inputs, and coproduct or waste streams routed to defined sinks. At each step, three material volume quantities are tracked: the primary input, the process basis (the internal scaling quantity, such as \si{m^3} slurry processed or tonnes solids crushed), and the primary output, which are related using two conversion factors (from input to step basis and step basis to output) that accommodate the diversity of natural reporting units across feedstock types and process steps. Reagent consumption is calculated as a function of constituent concentration in the process stream, allowing stoichiometric or empirical dosing ratios to be applied directly to the relevant constituent rather than to bulk throughput. Each production step is parameterized using \textit{machine blocks}, which specify equipment sizing, capital cost scaling relationships, utility intensities, and labor requirements. This separation between process chemistry (defined at the production step level) and equipment characteristics (defined at the machine block level) allows alternative equipment configurations or technology variants to be substituted without altering the underlying mass balance, enabling systematic sensitivity and scenario analysis. Costs and emissions are computed at the production step level and aggregated upward; because material flows are propagated from a terminal production target back through the full process chain while taking into account process unit conversions and yield rates, throughput, reagent consumption, and equipment sizing at every step are directly linked to the final output volume. For further detail on this hierarchical modeling system and a formal mathematical formulation, see Appendix \ref{app:04-methods-system}. 

Seven overarching facility types are defined, as shown in Fig. \ref{fig:04-step_breakdown}: brine evaporation, blast mining, pregnant leach solution (PLS) production from batch chemical treatment, PLS production from sulfuric acid leaching, PLS production from sulfate roasting and water leaching, lithium carbonate production (which is parameterized differently depending on whether the PLS is a chloride or sulfate-based solution), and a sulfuric acid plant specifically for the Thacker Pass claystone route. Each of these facilities consist of multiple production steps, with some steps shared between pathways, for a total of thirty two production steps modeled (Fig. \ref{fig:04-step_breakdown}). Production volumes are modeled at approximately 6,000 t \ce{Li2CO3}/yr for Silver Peak~\cite{srkconsultingSECTechnicalReport2023a}, 66,800 t/yr for Thacker Pass (consistent with the 2022 feasibility disclosure \cite{rothFeasibilityStudyNational2022} and other literature studies~\cite{wesselkaemperOptimizationProcessingCosts2025, royLithiumClayAssessing2024}), and 46,000 t LCE/yr for Jianxiawo (the most commonly cited production figure~\cite{flemingLongrunAvailabilityMineral2024, andrezejewskiResourcesEnergyQuarterly2025}). All three production estimates are stated as tonnes of packaged product at an assumed $\leq$0.4 wt\% moisture specification, equivalent to 0.996 t contained \ce{Li2CO3} per tonne of stated capacity.
 
For each pathway, we compute average lifetime operating costs of lithium carbonate production in USD per metric tonne \ce{Li2CO3}. As discussed in the introduction, this metric is more appropriate for assessing whether a project can sustainably compete with current or future marginal producers than a full NPV or IRR analysis. For each production step, operating expenditures (OPEX) are computed as the sum of variable costs (materials, labor, utilities) and fixed operating costs (maintenance and overhead), expressed per unit of annual \ce{Li2CO3} production. Capital expenditures (CAPEX), such as equipment, buildings, and auxiliary systems, are estimated using installation cost factors and equipment scaling relationships. For continuous processes, overall equipment costs are scaled with throughput using power law exponents (e.g. the six-tenths rule) while batch processes are expanded by adding parallel lines \cite{petersPlantDesignEconomics1991, ausimmMonograph27Cost, ulrichGuideChemicalEngineering1984}. These capital costs are not included in the primary cost comparison given the focus on operating competitiveness, but are reported for reference and cross-comparison with the literature. 

These results are extended to capture the environmental impacts of production, where impact factors for reagent use, energy consumption, and utilities are multiplied by the mass and energy flows through the system (see Appendix \ref{app:04-methods-system-emissions}). Scope 1 emissions result from specified process emissions (which are limited in our three cases) or direct combustion of fuel, primarily for heating process throughputs. Scope 2 emissions result from electricity usage, which is then multiplied by the local grid emissions factor (we use average emission rates here, but marginal rates would be appropriate for a full consequential analysis~\cite{nationalacademiesofscienceengineeringandmedicineAssessmentTechnologiesImproving2021}). Scope 3 emissions are calculated by multiplying each reagent's annual consumption by its cradle-to-gate emissions factor, reflecting the upstream production burden of that reagent as a purchased commodity input. This methodology is applied uniformly across all pathways, with geography-specific emissions factors applied where production location affects this upstream burden (e.g., soda ash from trona mining in the United States versus the Solvay process in China). 

For both costs and emissions, conservative, midpoint, and optimistic values are determined for each machine block design and material input cost, yielding sensitivity ranges that represent plausible operating variation across the dominant uncertainty drivers for each project. Process chemistry (e.g. reagent stoichiometry, conversion factors, and yield rates) is held fixed within each pathway, since these are generally determined by the underlying mineralogy and cannot be varied independently of the process design. Similarly, utility prices, labor rates, and land costs are not varied in the sensitivity analysis, as these are set by location rather than by operational choices, and their cross-pathway differences are instead captured through the location-specific parameters applied at the facility level.

While Silver Peak is modeled using flowsheets derived from project-specific disclosures due to its unique geological resource, the clay pathway is represented using a generic acid-leach flowsheet grounded in the fundamental process chemistry of Nevada sedimentary clay deposits and heavily informed by design of Thacker Pass, rather than following the Thacker Pass flowsheet exactly. Several other clay projects in Nevada share similar mineralogy and processing fundamentals (e.g. Rhyolite Ridge~\cite{fluorenterprisesRhyoliteRidgeLithiumBoron2022}, Bonnie Claire~\cite{samariMineralResourceEstimate2024}, Tonopah Flats~\cite{dyerAmendedResourceEstimate}, TLC~\cite{riordanFalchaniLithiumProject2024}) and a generic flowsheet provides a more transferable baseline for this class of resource than a project-specific model would. The key distinction is in impurity management: because sulfuric acid dissolves both lithium and magnesium simultaneously from magnesium-rich hectorite, the generic flowsheet precipitates all magnesium with quicklime, whereas the Thacker Pass project-specific design employs a magnesium crystallization route that recycles magnesium compounds as reagents, reducing lime consumption to approximately 2,630~kg/t at the cost of higher thermal energy demand~\cite{sgscanadaincSK1300Technical2024}. This omission renders the generic flowsheet as conservative relative to the Thacker Pass design specifically, but ideally representative of what a first-generation clay operation without project-specific optimization would require. The cost and emissions implications of this choice are discussed in Section~\ref{sec:04-thackerpass}.

Chinese lepidolite is likewise is modeled based on a generic sulfate roasting-water leach flowsheet based on assumptions from CATL's Jianxiawo project, as well as patents~\cite{nanMethodsExtractingLithium2019,nanMethodsExtractingLithium2022,zouMethodEfficientLithium2024} and previous studies~\cite{yanExtractionLithiumLepidolite2012, luongIronSulphateRoasting2014, parkEffectiveLithiumExtraction2023, dengSystematicReviewSulfate2026}, given that there are no publicly available disclosures from CATL comparable to the SK-1300 and NI 43-101 reports available for the US projects. Thus, we consider the standard commercial route for lepidolite-based lithium production in the Jiangxi region, using froth flotation, grinding, roasting at 800--950$^\circ$C with \ce{Na2SO4}, \ce{K2SO4}, and \ce{CaO}, water leaching, staged impurity removal, and carbonate precipitation~\ref{fig:04-step_breakdown}. Reagent ratios are taken from bench-scale tests on flotation concentrate~\cite{yanExtractionLithiumLepidolite2012, luongIronSulphateRoasting2014} and are not validated against plant-level consumption. Ore grade is set at the 0.4\%~\ce{Li2O} upper bound of reported Jiangxi grades~\cite{sulcaLithiumPricesWhere2024}, while standalone costs exclude any transfer-pricing advantage within CATL's vertically integrated operations. The cost and emissions implications of these choices are discussed in Section~\ref{sec:04-jianxiawo}

An extensive expanded description of our methodology is provided in the Supplemental Information, and the full open-source spreadsheet model is available for download at XYZ.

\section{Acknowledgements}
This research was supported by the Carnegie Mellon University College of Engineering Moonshot Award for ``Engineering Competitiveness: Critical Technologies, Supply Chains, and Infrastructure,'' the National Science Foundation Award # 2241237, and the National Science Foundation Graduate Research Fellowship under Grant # DGE2140739. We would like to acknowledge J.P. Helveston, A. Reamer, J.P. MacDuffie, K. Haverkamp, E. Wilson, E. Spiller, A. Ramji, A. Haddad, A. Yao, M. Mauter, W. Peng, M. Davidson, S. Davis, and E. Olivetti for their feedback and perspectives at various conferences and presentations; and the Fuchs, Karplus (LEO) and Michalek (VEG) research groups as well as the 2024-2025 Resources for the Future Critical Minerals Research Laboratory for their feedback throughout the development of this work. We would also like to acknowledge the use of large language models to edit and refine some of the text and code in this work, though the authors reviewed all outputs and take full responsibility for all final text.

\clearpage

\bibliographystyle{unsrtnat}
\bibliography{thesis}

\clearpage

\section{Supplementary Appendix}\label{app:04}

\subsection{Literature Review}\label{app:04-litreview}

Many studies of the lithium system and broader lithium battery system have been conducted, evaluating all aspects from costs, environmental impacts of many types, security, process technology choices, and more. Our study aims to answer questions about the onshoring of lithium production in the United States, and compares against currently marginal production -- lepidolite resources in China. We suggest that this security and policy-focused framing combined with engineering details provides unique insights in comparison to a very extensive set of previous studies, compiled primarily from the last five years, as shown in Table \ref{tab:c1-litreview-matrix}. 

In particular, the two dimensions that most distinguish this work from others are the detailed process modeling and connection between cost and emissions modeling, as very few papers do both. More importantly, even among such papers, open, examinable models are not published, making conclusions difficult to evaluate and extend. There are limitations to our scope. For example, we focus on three lithium production steps in this paper, and further modeling of other pathways would be necessary for a complete competitive picture.  Multi-material papers do exist with some process-level modeling, usually of downstream steps (e.g. battery manufacturing), but modeling of specific critical minerals production pathways such as lithium is often treated at a lower resolution. Additionally, on the emissions side, this paper only covers greenhouse gas emissions. Other impact categories, such as air quality considerations and water usage, are important elements of the lithium sustainability question, but our framework is extensible to these in future work. Furthermore, studies (e.g. \citet{thomsenMineManufacturerAssessing2025}) have indicated that the impacts of transportation on the critical mineral and battery supply chains are undercounted. Although the impacts were found to be relatively negligible in this study given the constrained geographical scope, extensions of this framework to multi-country production networks would necessitate more detailed modeling of these transportation impacts that many studies miss. Finally, the approach towards uncertainty in this study should be noted. A significant limitation in LCA studies (and to a lesser extent, TEA studies) is that they make a singular point estimate, without considering operating and input data uncertainty. Some studies do make this consideration, as seen in the table below, with treatments ranging from full Monte Carlo and probabilistic uncertainty methods to parametric variation to simple scenario-based analyses. Probabilistic methods can provide a rich picture of output distributions, but require well-characterized input distributions, which can be challenging to measure for emerging lithium production pathways. Parametric variations are more tractable, but applying relatively arbitrary, uniform swings (e.g. $\pm$20\%, $\pm$50\%) without grounding such variation in what is actually possible can lead to overemphasis on variables that have little real leverage and underweighting those that do. Thus, we use empirically estimated conservative, midpoint, and optimistic estimates for triangular sensitivity analysis to represent genuine operating uncertainty. 

% OF NOTE :
% We note that the lepidolite GWP estimate of 66.0~t~\ce{CO2eq}/t~\ce{Li2CO3} reported by \citet{wangLifeCycleEnvironmental2026} is substantially higher than other lepidolite estimates in the literature, including our own. Examination of the underlying inventory indicates that the reported comminution electricity consumption is approximately an order of magnitude higher than values reported elsewhere for comparable grinding duties - namely, roughly 400 kWh/tonne ore versus 35 kWh/tonne reported in our model and 20-50 kWh/tonne reported in other hard rock lithium ore studies. We plot this value for completeness and transparency but note that it likely reflects an inventory error rather than a genuine process characteristic of lepidolite production.

Further detailed discussion of our choice of system boundaries and modeling scope is given in the following sections. 


%% ---- Color definitions ----
\definecolor{cellcolor}{HTML}{3D3D3D}
\definecolor{greenrow}{HTML}{D9EAD3}
\definecolor{ourstudybg}{HTML}{FFF2CC}
\definecolor{keyheader}{HTML}{DAE8FC}
\definecolor{sectionrow}{HTML}{F3F3F3}
\definecolor{codelabel}{HTML}{999999}
 
%% ---- Shorthand commands ----
\providecommand{\X}{\cellcolor{cellcolor}}
\providecommand{\anonstrut}{\rule{0pt}{5.5pt}}
 
%% ---- Custom narrow centered column type ----
\newcolumntype{D}{>{\centering\arraybackslash}p{0.45cm}}
 
%% ---- Shared column header block (reused in all three sidewaystable floats) ----
\newcommand{\LRheader}{%
\toprule
 & \multicolumn{14}{c}{\textbf{Scope}}
 & \multicolumn{9}{c}{\textbf{Outcomes}}
 & \multicolumn{3}{c}{\textbf{Sensitivity}} \\
\cmidrule(lr){2-15}\cmidrule(lr){16-24}\cmidrule(lr){25-27}
 & \multicolumn{2}{c}{\parbox{1.0cm}{\centering\tiny Material\\Scope}}
 & \multicolumn{6}{c}{\parbox{2.8cm}{\centering\tiny Material Production\\Pathway Considered}}
 & \multicolumn{2}{c}{\parbox{0.9cm}{\centering\tiny Process\\Scope}}
 & \multicolumn{2}{c}{\parbox{1.0cm}{\centering\tiny Geographic\\Variation}}
 & \multicolumn{2}{c}{\parbox{0.9cm}{\centering\tiny Trans-\\portation}}
 & \multicolumn{3}{c}{\parbox{1.4cm}{\centering\tiny Emissions}}
 & \multicolumn{3}{c}{\parbox{1.4cm}{\centering\tiny Cost\\Modeling}}
 & \multicolumn{3}{c}{\parbox{1.4cm}{\centering\tiny Key\\Contributions}}
 & \multicolumn{3}{c}{\parbox{1.4cm}{\centering\tiny Sensitivity\\Analysis}} \\
\cmidrule(lr){2-3}\cmidrule(lr){4-9}\cmidrule(lr){10-11}\cmidrule(lr){12-13}
\cmidrule(lr){14-15}\cmidrule(lr){16-18}\cmidrule(lr){19-21}
\cmidrule(lr){22-24}\cmidrule(lr){25-27}
\textbf{Reference} &
  \rotatebox{90}{\parbox{3.8cm}{\raggedright\tiny Lithium only}} &
  \rotatebox{90}{\parbox{3.8cm}{\raggedright\tiny Includes other minerals/materials}} &
  \rotatebox{90}{\parbox{3.8cm}{\raggedright\tiny Brine-Traditional}} &
  \rotatebox{90}{\parbox{3.8cm}{\raggedright\tiny Brine-DLE}} &
  \rotatebox{90}{\parbox{3.8cm}{\raggedright\tiny Spodumene}} &
  \rotatebox{90}{\parbox{3.8cm}{\raggedright\tiny Lepidolite}} &
  \rotatebox{90}{\parbox{3.8cm}{\raggedright\tiny Clay}} &
  \rotatebox{90}{\parbox{3.8cm}{\raggedright\tiny Other (inc.\ recycling / circular / emerging)}} &
  \cellcolor{keyheader}\rotatebox{90}{\parbox{3.8cm}{\raggedright\tiny Production Capacity}} &
  \cellcolor{keyheader}\rotatebox{90}{\parbox{3.8cm}{\raggedright\tiny Li Process Modeling}} &
  \rotatebox{90}{\parbox{3.8cm}{\raggedright\tiny Comprehensive (or an attempt to be)}} &
  \rotatebox{90}{\parbox{3.8cm}{\raggedright\tiny Case Studies / Representative Subset}} &
  \rotatebox{90}{\parbox{3.8cm}{\raggedright\tiny Costs}} &
  \rotatebox{90}{\parbox{3.8cm}{\raggedright\tiny Emissions}} &
  \rotatebox{90}{\parbox{3.8cm}{\raggedright\tiny GHGs}} &
  \rotatebox{90}{\parbox{3.8cm}{\raggedright\tiny Local Air Pollution}} &
  \rotatebox{90}{\parbox{3.8cm}{\raggedright\tiny Other (e.g.\ Water)}} &
  \cellcolor{keyheader}\rotatebox{90}{\parbox{3.8cm}{\raggedright\tiny Process-Based}} &
  \rotatebox{90}{\parbox{3.8cm}{\raggedright\tiny Regression-Based}} &
  \rotatebox{90}{\parbox{3.8cm}{\raggedright\tiny Other}} &
  \cellcolor{keyheader}\rotatebox{90}{\parbox{3.8cm}{\raggedright\tiny Supply Chain Security}} &
  \cellcolor{keyheader}\rotatebox{90}{\parbox{3.8cm}{\raggedright\tiny Discussion of Policy Implications}} &
  \cellcolor{keyheader}\rotatebox{90}{\parbox{3.8cm}{\raggedright\tiny Models Policy Intervention}} &
  \rotatebox{90}{\parbox{3.8cm}{\raggedright\tiny Monte Carlo / Probabilistic}} &
  \rotatebox{90}{\parbox{3.8cm}{\raggedright\tiny Parametric Variation / Numerical}} &
  \cellcolor{keyheader}\rotatebox{90}{\parbox{3.8cm}{\raggedright\tiny Triangular (Low/Med/High) / Specific Cases}} \\
\midrule}
 
 
%% ============================================================
%%  PART 1a — Our study + Main LCA/TEA papers M1–M21
%% ============================================================

\begin{sidewaystable}
\singlespacing
\centering
\renewcommand{\arraystretch}{0.75}
\caption[Literature Review Scope and Attribute Matrix.]{%
  Literature Review Scope and Attribute Matrix.
  Filled cells indicate the attribute is present.
  \colorbox{ourstudybg}{The yellow row} is the present study;
  \colorbox{greenrow}{green rows} are notable related papers.
  \colorbox{keyheader}{Blue-tinted headers} indicate key contribution dimensions.
  Grey codes index other studies; see \ref{tab:c2-litreview-legend} for the legend.}
\label{tab:c1-litreview-matrix}
\alttbl{Literature review matrix cross-referencing approximately 100 studies against 26 attribute columns spanning four groups: resource type and geography (lithium, other minerals, brine, hard rock, clay, lepidolite, DLE, multiple geographies), analysis type (TEA, LCA, cost model, process model, sensitivity analysis, energy and water, reagent detail), scope dimensions (single project, comparative, policy, supply chain, trade or geopolitics, security or vulnerability), and output characteristics (country-level data, time series, open data or code, peer reviewed, novel contribution). The table is split across four pages covering LCA and TEA papers M1 through M47, security and policy papers S1 through S19, and technical papers T1 through T26. The present study's row is highlighted in yellow and covers the widest combination of attributes. Notable related papers are highlighted in green.}

\footnotesize
\setlength{\tabcolsep}{2pt}
\begin{tabular}{p{4.5cm}*{26}{D}}
\LRheader
\specialrule{0.8pt}{2pt}{0pt}
\rowcolor{ourstudybg}\textbf{Our study} & \X &  & \X &  &  & \X & \X &  & \X & \X &  & \X & \X & \X & \X &  &  & \X &  &  & \X & \X &  &  &  & \X \\
\specialrule{0.8pt}{0pt}{4pt}
\anonstrut\textcolor{codelabel}{\scalebox{0.55}{\textsf{M1}}} & \X &  & \X &  & \X &  &  &  &  &  &  &  &  &  & \X &  & \X &  &  &  &  & \X &  &  &  &  \\
\anonstrut\textcolor{codelabel}{\scalebox{0.55}{\textsf{M2}}} & \X &  & \X &  & \X &  &  & \X &  & \X &  & \X & \X & \X &  &  &  &  &  &  &  &  &  & \X &  &  \\
\anonstrut\textcolor{codelabel}{\scalebox{0.55}{\textsf{M3}}} & \X &  & \X &  & \X & \X &  &  & \X &  &  & \X &  &  &  &  &  &  & \X &  &  &  &  &  &  & \X \\
\rowcolor{greenrow}\citet{ambroseUnderstandingFutureLithium2020b} & \X &  & \X &  & \X & \X &  &  & \X & \X &  & \X &  &  & \X & \X & \X &  &  & \X &  &  &  &  &  & \X \\
\anonstrut\textcolor{codelabel}{\scalebox{0.55}{\textsf{M4}}} & \X &  & \X & \X & \X &  & \X &  & \X &  & \X &  &  &  &  &  &  &  & \X &  &  & \X &  &  & \X &  \\
\rowcolor{greenrow}\citet{chordiaLifeCycleEnvironmental2022} & \X &  & \X &  & \X &  &  &  &  & \X &  & \X &  &  & \X &  & \X &  &  &  & \X &  &  &  &  & \X \\
\anonstrut\textcolor{codelabel}{\scalebox{0.55}{\textsf{M5}}} & \X &  &  &  & \X &  &  &  & \X & \X &  & \X & \X & \X & \X & \X & \X &  &  &  &  &  &  & \X &  &  \\
\anonstrut\textcolor{codelabel}{\scalebox{0.55}{\textsf{M6}}} & \X &  & \X & \X &  &  &  &  &  & \X &  & \X &  &  & \X & \X & \X &  &  &  &  &  &  &  & \X &  \\
\anonstrut\textcolor{codelabel}{\scalebox{0.55}{\textsf{M7}}} & \X &  &  & \X &  &  &  &  & \X &  &  & \X &  &  &  &  &  &  &  & \X &  & \X & \X &  & \X &  \\
\rowcolor{greenrow}\citet{flemingLongrunAvailabilityMineral2024} & \X &  & \X & \X & \X & \X & \X & \X &  &  & \X &  &  &  &  &  &  &  & \X &  &  &  &  &  &  &  \X \\
\rowcolor{greenrow}\citet{gaoLithiumExtractionHard2023} & \X &  &  &  & \X & \X &  &  &  & \X &  &  &  &  &  & \X & \X & \X &  &  &  &  &  &  & \X &  \\
\anonstrut\textcolor{codelabel}{\scalebox{0.55}{\textsf{M8}}} & \X &  & \X & \X &  &  &  &  &  & \X &  & \X &  &  & \X & \X & \X &  &  &  &  &  &  &  &  &  \\
\anonstrut\textcolor{codelabel}{\scalebox{0.55}{\textsf{M9}}} &  & \X & \X & \X & \X &  &  &  &  &  &  & \X &  &  & \X &  & \X &  &  &  & \X & \X & \X &  &  &  \\
\anonstrut\textcolor{codelabel}{\scalebox{0.55}{\textsf{M10}}} & \X &  & \X &  &  &  &  &  &  & \X & \X &  & \X & \X & \X & \X & \X &  &  &  &  &  &  &  &  &  \\
\anonstrut\textcolor{codelabel}{\scalebox{0.55}{\textsf{M11}}} & \X &  &  & \X &  &  &  & \X &  & \X &  & \X &  &  & \X & \X & \X & \X &  &  &  &  &  &  & \X & \X \\
\anonstrut\textcolor{codelabel}{\scalebox{0.55}{\textsf{M12}}} & \X &  & \X &  & \X &  &  &  &  & \X &  & \X & \X &  &  &  &  &  &  &  &  &  &  &  &  &  \\
\anonstrut\textcolor{codelabel}{\scalebox{0.55}{\textsf{M13}}} &  & \X & \X & \X & \X &  &  &  &  & \X &  & \X &  & \X & \X &  &  &  &  &  &  &  &  & \X &  & \X \\
\anonstrut\textcolor{codelabel}{\scalebox{0.55}{\textsf{M14}}} & \X &  & \X & \X &  &  &  &  &  & \X &  &  &  &  & \X & \X & \X &  &  &  &  &  &  &  &  &  \\
\anonstrut\textcolor{codelabel}{\scalebox{0.55}{\textsf{M15}}} & \X &  &  & \X &  &  & \X &  &  &  &  &  &  &  &  &  &  &  &  &  &  &  &  &  &  &  \\
\anonstrut\textcolor{codelabel}{\scalebox{0.55}{\textsf{M16}}} & \X &  & \X & \X &  &  & \X &  &  & \X &  & \X &  & \X & \X & \X & \X &  &  &  &  &  &  &  &  & \X \\
\anonstrut\textcolor{codelabel}{\scalebox{0.55}{\textsf{M17}}} & \X &  &  &  & \X &  &  &  &  & \X &  & \X &  & \X & \X &  &  &  &  &  &  & \X &  &  &  & \X \\
\anonstrut\textcolor{codelabel}{\scalebox{0.55}{\textsf{M18}}} & \X &  & \X &  & \X &  &  &  &  & \X &  & \X &  & \X & \X & \X & \X &  &  &  &  &  &  &  &  & \X \\
\anonstrut\textcolor{codelabel}{\scalebox{0.55}{\textsf{M19}}} &  & \X & \X &  & \X &  &  &  &  &  &  & \X &  &  & \X &  &  &  &  &  & \X &  &  &  &  & \X \\
\anonstrut\textcolor{codelabel}{\scalebox{0.55}{\textsf{M20}}} & \X &  & \X &  & \X &  &  &  &  & \X &  & \X &  & \X & \X & \X & \X &  &  &  &  &  &  &  &  & \X \\
\rowcolor{greenrow}\citet{khakmardanComparativeLifeCycle2023} &  & \X & \X &  & \X &  & \X & \X &  & \X &  & \X &  & \X & \X &  & \X &  &  &  &  &  &  &  &  &  \\
\anonstrut\textcolor{codelabel}{\scalebox{0.55}{\textsf{M21}}} & \X &  &  &  & \X &  &  &  &  & \X &  & \X &  &  & \X & \X & \X &  &  &  &  &  &  &  &  &  \\
\bottomrule
\end{tabular}
\end{sidewaystable}
\tagmcend
 
 
%% ============================================================
%%  PART 1b — Main LCA/TEA papers M22–M47
%% ============================================================

\tagmcbegin{artifact}
\begin{sidewaystable}
\ContinuedFloat
\singlespacing
\centering
\renewcommand{\arraystretch}{0.75}
\caption[]{Continued.}
\alttbl{Continuation of the literature review matrix, covering LCA and TEA papers M22 through M47}
\footnotesize
\setlength{\tabcolsep}{2pt}
\begin{tabular}{p{4.5cm}*{26}{D}}
\LRheader
\anonstrut\textcolor{codelabel}{\scalebox{0.55}{\textsf{M22}}} & \X &  &  &  & \X &  &  &  &  & \X &  & \X &  &  & \X & \X & \X &  &  &  &  &  &  &  &  &  \\
\anonstrut\textcolor{codelabel}{\scalebox{0.55}{\textsf{M23}}} & \X &  &  &  & \X &  &  &  &  & \X &  & \X &  &  & \X & \X & \X &  &  &  &  &  &  &  &  & \X \\
\rowcolor{greenrow}\citet{liuCarbonFootprintBatteryGrade2025} & \X &  &  & \X & \X & \X &  & \X & \X & \X & \X &  &  &  & \X &  &  &  &  &  &  &  &  &  &  & \X \\
\anonstrut\textcolor{codelabel}{\scalebox{0.55}{\textsf{M24}}} &  & \X &  &  &  &  &  & \X &  &  & \X &  &  &  & \X &  &  &  &  &  &  & \X & \X &  &  & \X \\
\anonstrut\textcolor{codelabel}{\scalebox{0.55}{\textsf{M25}}} &  & \X & \X &  & \X &  &  & \X &  &  &  & \X &  & \X & \X &  &  &  &  &  &  &  &  &  &  &  \\
\anonstrut\textcolor{codelabel}{\scalebox{0.55}{\textsf{M26}}} &  & \X & \X &  & \X &  &  &  &  & \X &  & \X &  &  & \X &  &  &  &  &  &  &  &  &  & \X &  \\
\anonstrut\textcolor{codelabel}{\scalebox{0.55}{\textsf{M27}}} & \X &  &  & \X &  &  &  &  &  & \X &  & \X &  &  & \X & \X & \X &  &  &  &  &  &  &  & \X &  \\
\rowcolor{greenrow}\citet{mousavinezhadNextLithiumBoom2025} & \X &  & \X & \X & \X &  & \X &  & \X & \X &  & \X &  &  & \X &  &  &  &  &  & \X & \X &  &  &  &  \\
\anonstrut\textcolor{codelabel}{\scalebox{0.55}{\textsf{M28}}} & \X &  &  &  &  &  & \X &  &  & \X &  & \X &  &  & \X &  & \X &  &  &  &  &  &  &  &  &  \\
\anonstrut\textcolor{codelabel}{\scalebox{0.55}{\textsf{M29}}} & \X &  &  & \X &  &  &  &  &  & \X &  & \X &  &  & \X & \X & \X &  &  &  &  &  &  &  & \X &  \\
\rowcolor{greenrow}\citet{peiselerCarbonFootprintDistributions2024} &  & \X & \X &  & \X &  &  & \X & \X &  & \X &  &  &  & \X &  &  &  & \X &  &  & \X &  & \X &  & \X \\
\anonstrut\textcolor{codelabel}{\scalebox{0.55}{\textsf{M30}}} & \X &  &  &  &  &  & \X &  &  & \X &  & \X &  & \X & \X & \X & \X &  &  &  &  &  &  &  &  &  \\
\anonstrut\textcolor{codelabel}{\scalebox{0.55}{\textsf{M31}}} & \X &  & \X & \X &  &  &  &  &  & \X &  & \X &  &  & \X & \X & \X &  &  &  &  &  &  &  &  &  \\
\anonstrut\textcolor{codelabel}{\scalebox{0.55}{\textsf{M32}}} & \X &  &  & \X &  &  &  &  &  & \X &  &  &  &  & \X &  &  &  &  &  &  &  &  &  &  &  \\
\anonstrut\textcolor{codelabel}{\scalebox{0.55}{\textsf{M33}}} & \X &  & \X & \X &  &  &  &  &  & \X & \X &  &  &  & \X & \X & \X &  &  &  & \X &  &  &  & \X &  \\
\anonstrut\textcolor{codelabel}{\scalebox{0.55}{\textsf{M34}}} &  & \X & \X & \X &  &  &  &  &  & \X &  & \X &  &  &  &  & \X &  &  &  &  &  &  & \X &  &  \\
\anonstrut\textcolor{codelabel}{\scalebox{0.55}{\textsf{M35}}} & \X &  &  &  & \X &  &  &  &  & \X & \X &  &  & \X & \X & \X & \X &  &  &  &  &  &  &  & \X &  \\
\anonstrut\textcolor{codelabel}{\scalebox{0.55}{\textsf{M36}}} & \X &  &  &  & \X &  &  &  &  & \X &  &  &  &  & \X &  &  & \X &  &  &  &  &  &  &  &  \\
\anonstrut\textcolor{codelabel}{\scalebox{0.55}{\textsf{M37}}} & \X &  & \X &  & \X &  &  & \X &  & \X &  &  &  &  & \X & \X & \X &  &  &  &  &  &  &  &  & \X \\
\anonstrut\textcolor{codelabel}{\scalebox{0.55}{\textsf{M38}}} & \X &  & \X &  & \X & \X & \X &  &  & \X &  & \X & \X &  &  &  &  &  &  & \X &  &  &  &  &  &  \\
\rowcolor{greenrow}\citet{sunRobustAssessmentsLithium2025} & \X &  & \X & \X & \X & \X & \X &  & \X & \X & \X &  &  &  & \X &  & \X &  &  &  & \X & \X &  &  &  &  \\
\anonstrut\textcolor{codelabel}{\scalebox{0.55}{\textsf{M39}}} &  & \X & \X &  & \X & \X &  &  &  &  &  &  & \X & \X & \X &  &  &  &  & \X &  &  &  &  &  &  \\
\anonstrut\textcolor{codelabel}{\scalebox{0.55}{\textsf{M40}}} & \X &  & \X & \X & \X & \X & \X &  &  & \X &  & \X &  &  &  &  &  &  &  & \X &  &  &  &  &  &  \\
\anonstrut\textcolor{codelabel}{\scalebox{0.55}{\textsf{M41}}} & \X &  & \X & \X & \X & \X &  &  &  &  & \X & \X &  &  & \X &  & \X &  &  & \X &  &  &  &  &  &  \\
\anonstrut\textcolor{codelabel}{\scalebox{0.55}{\textsf{M42}}} & \X &  &  & \X &  &  &  &  &  & \X &  & \X &  &  &  &  & \X &  &  &  &  &  &  &  &  &  \\
\anonstrut\textcolor{codelabel}{\scalebox{0.55}{\textsf{M43}}} & \X &  &  &  & \X & \X &  &  &  & \X &  & \X & \X & \X & \X & \X & \X & \X &  &  &  &  &  &  & \X &  \\
\anonstrut\textcolor{codelabel}{\scalebox{0.55}{\textsf{M44}}} & \X &  &  & \X &  &  &  & \X & \X & \X &  &  & \X &  &  &  &  & \X &  &  &  &  &  &  &  &  \\
\rowcolor{greenrow}\citet{wesselkaemperOptimizationProcessingCosts2025} & \X &  & \X & \X & \X &  & \X &  & \X & \X &  & \X & \X &  &  &  &  & \X &  & \X & \X & \X &  &  & \X &  \\
\anonstrut\textcolor{codelabel}{\scalebox{0.55}{\textsf{M45}}} &  & \X & \X &  & \X &  & \X & \X & \X & \X &  & \X &  &  &  &  &  & \X &  & \X & \X & \X &  &  & \X &  \\
\anonstrut\textcolor{codelabel}{\scalebox{0.55}{\textsf{M46}}} & \X &  & \X & \X & \X & \X &  & \X &  & \X &  & \X &  &  &  &  & \X &  &  &  &  &  &  &  & \X &  \\
\anonstrut\textcolor{codelabel}{\scalebox{0.55}{\textsf{M47}}} & \X &  & \X & \X & \X & \X & \X & \X &  &  & \X & &  &  &  &  &  & \X &  &  &  &  &  &  & \X &  \\
\bottomrule
\end{tabular}
\end{sidewaystable}
 \tagmcend
 
%% ============================================================
%%  PART 1c — Security & Policy papers S1–S19
%% ============================================================

\tagmcbegin{artifact}
\begin{sidewaystable}
\ContinuedFloat
\singlespacing
\centering
\renewcommand{\arraystretch}{0.75}
\caption[]{Continued.}
\alttbl{Continuation of the literature review matrix, covering security and policy papers S1 through S19}
\footnotesize
\setlength{\tabcolsep}{2pt}
\begin{tabular}{p{4.5cm}*{26}{D}}
\LRheader
\multicolumn{27}{l}{\cellcolor{sectionrow}\textit{\small Security \& Policy Papers}} \\
\anonstrut\textcolor{codelabel}{\scalebox{0.55}{\textsf{S1}}} & \X &  & \X &  & \X & \X & \X & \X &  &  & \X &  &  &  &  &  &  &  &  & \X &  & \X &  &  &  & \X \\
\anonstrut\textcolor{codelabel}{\scalebox{0.55}{\textsf{S2}}} &  & \X &  &  &  &  &  &  &  &  &  &  &  &  & \X &  &  &  &  &  & \X &  &  &  & \X &  \\
\anonstrut\textcolor{codelabel}{\scalebox{0.55}{\textsf{S3}}} &  & \X &  & \X &  &  &  & \X &  &  &  &  &  &  &  &  &  &  &  &  & \X & \X &  &  &  &  \\
\anonstrut\textcolor{codelabel}{\scalebox{0.55}{\textsf{S4}}} &  & \X &  &  &  &  &  &  &  &  &  &  &  &  &  &  &  &  &  &  & \X & \X &  &  &  &  \\
\anonstrut\textcolor{codelabel}{\scalebox{0.55}{\textsf{S5}}} & \X &  & \X &  & \X &  &  & \X &  &  & \X &  &  &  &  &  &  &  &  &  & \X &  &  & \X &  &  \\
\anonstrut\textcolor{codelabel}{\scalebox{0.55}{\textsf{S6}}} &  & \X &  &  &  &  &  &  &  &  & \X &  &  &  &  &  &  &  &  &  & \X &  &  &  &  & \X \\
\anonstrut\textcolor{codelabel}{\scalebox{0.55}{\textsf{S7}}} & \X &  & \X &  & \X &  &  &  &  &  &  & \X &  &  &  &  &  &  &  &  & \X & \X &  &  &  &  \\
\anonstrut\textcolor{codelabel}{\scalebox{0.55}{\textsf{S8}}} &  & \X &  &  &  &  &  &  &  &  & \X &  &  &  &  &  &  &  &  &  & \X &  &  &  &  &  \\
\rowcolor{greenrow}\citet{haddadHowMakeLithium2023} & \X &  & \X & \X & \X &  &  &  &  &  & \X &  &  &  & \X & \X & \X &  &  &  & \X & \X &  &  &  & \X \\
\anonstrut\textcolor{codelabel}{\scalebox{0.55}{\textsf{S9}}} &  & \X &  &  &  &  &  &  &  &  & \X &  &  &  &  &  &  &  &  &  & \X &  &  &  &  &  \\
\anonstrut\textcolor{codelabel}{\scalebox{0.55}{\textsf{S10}}} &  & \X &  &  &  &  &  &  &  &  &  &  &  &  &  &  &  &  &  &  & \X & \X &  &  &  &  \\
\anonstrut\textcolor{codelabel}{\scalebox{0.55}{\textsf{S11}}} & \X &  &  &  &  &  &  &  &  &  & \X &  &  &  &  &  &  &  &  &  & \X &  &  &  &  &  \\
\anonstrut\textcolor{codelabel}{\scalebox{0.55}{\textsf{S12}}} &  & \X &  &  &  &  &  &  &  &  &  & \X &  &  &  &  &  &  &  &  & \X &  &  &  &  &  \\
\anonstrut\textcolor{codelabel}{\scalebox{0.55}{\textsf{S13}}} & \X &  & \X &  &  &  &  &  &  &  &  & \X &  &  &  &  &  &  &  &  & \X &  &  &  &  &  \\
\anonstrut\textcolor{codelabel}{\scalebox{0.55}{\textsf{S14}}} & \X &  & \X &  & \X & \X &  &  &  &  & \X &  &  &  &  &  &  &  &  &  & \X &  &  &  &  &  \\
\anonstrut\textcolor{codelabel}{\scalebox{0.55}{\textsf{S15}}} & \X &  & \X &  & \X & \X & \X & \X &  &  & \X &  &  &  &  &  &  &  &  &  & \X &  & \X &  &  &  \\
\anonstrut\textcolor{codelabel}{\scalebox{0.55}{\textsf{S16}}} &  & \X &  &  &  &  &  &  &  &  &  & \X &  &  &  &  &  &  &  &  & \X &  &  &  &  &  \\
\anonstrut\textcolor{codelabel}{\scalebox{0.55}{\textsf{S17}}} &  & \X &  &  &  &  &  &  &  &  & \X &  &  &  &  &  &  &  &  &  & \X &  &  &  &  &  \\
\anonstrut\textcolor{codelabel}{\scalebox{0.55}{\textsf{S18}}} &  & \X &  &  &  &  &  &  &  &  & \X &  &  &  &  &  &  &  &  &  & \X &  &  &  &  &  \\
\anonstrut\textcolor{codelabel}{\scalebox{0.55}{\textsf{S19}}} & \X &  &  &  &  &  & \X &  &  &  &  & \X &  &  &  &  &  &  &  &  & \X & \X &  &  &  &  \\
\bottomrule
\end{tabular}
\end{sidewaystable}
\tagmcend
 
 
%% ============================================================
%%  PART 1d — Technical papers T1–T26
%% ============================================================

\tagmcbegin{artifact}
\begin{sidewaystable}
\ContinuedFloat
\singlespacing
\centering
\renewcommand{\arraystretch}{0.75}
\caption[]{Continued.}
\alttbl{Continuation of the literature review matrix, covering technical papers T1 through T26}
\footnotesize
\setlength{\tabcolsep}{2pt}
\begin{tabular}{p{4.5cm}*{26}{D}}
\LRheader
\multicolumn{27}{l}{\cellcolor{sectionrow}\textit{\small Technical Papers (no costs)}} \\
\anonstrut\textcolor{codelabel}{\scalebox{0.55}{\textsf{T1}}} & \X &  & \X &  & \X & \X & \X &  &  & \X &  &  &  &  &  &  &  &  &  &  &  &  &  &  &  &  \\
\anonstrut\textcolor{codelabel}{\scalebox{0.55}{\textsf{T2}}} & \X &  &  &  &  &  &  &  &  & \X &  &  &  &  &  &  &  &  &  &  &  &  &  &  &  &  \\
\anonstrut\textcolor{codelabel}{\scalebox{0.55}{\textsf{T3}}} & \X &  &  &  &  &  &  & \X &  & \X &  &  &  &  & \X & \X & \X &  &  &  &  &  &  &  &  &  \\
\anonstrut\textcolor{codelabel}{\scalebox{0.55}{\textsf{T4}}} & \X &  &  &  &  & \X &  &  &  & \X &  &  &  &  &  &  &  &  &  &  &  &  &  &  &  &  \\
\anonstrut\textcolor{codelabel}{\scalebox{0.55}{\textsf{T5}}} & \X &  & \X & \X &  &  &  &  &  & \X &  &  &  &  &  &  &  &  &  &  &  &  &  &  &  &  \\
\anonstrut\textcolor{codelabel}{\scalebox{0.55}{\textsf{T6}}} & \X &  & \X & \X &  &  &  &  &  &  & \X &  &  &  &  &  & \X &  &  &  &  & \X &  &  &  &  \\
\anonstrut\textcolor{codelabel}{\scalebox{0.55}{\textsf{T7}}} &  & \X & \X &  & \X & \X & \X & \X &  & \X & \X &  &  &  &  &  &  &  &  &  &  &  &  &  &  &  \\
\anonstrut\textcolor{codelabel}{\scalebox{0.55}{\textsf{T8}}} & \X &  &  &  &  &  & \X &  &  & \X &  &  &  &  &  &  &  &  &  &  &  &  &  &  &  &  \\
\anonstrut\textcolor{codelabel}{\scalebox{0.55}{\textsf{T9}}} & \X &  &  &  & \X &  &  &  &  &  &  &  &  &  &  &  &  &  &  &  &  &  &  &  &  &  \\
\anonstrut\textcolor{codelabel}{\scalebox{0.55}{\textsf{T10}}} & \X &  &  &  & \X & \X &  & \X &  & \X &  &  &  & \X &  &  & \X &  &  &  &  &  &  &  &  &  \\
\anonstrut\textcolor{codelabel}{\scalebox{0.55}{\textsf{T11}}} & \X &  &  &  & \X &  &  &  &  & \X &  &  &  &  &  &  &  &  &  &  &  &  &  &  &  &  \\
\anonstrut\textcolor{codelabel}{\scalebox{0.55}{\textsf{T12}}} & \X &  &  &  & \X &  &  &  &  & \X &  &  &  &  &  &  &  &  &  &  &  &  &  &  &  &  \\
\anonstrut\textcolor{codelabel}{\scalebox{0.55}{\textsf{T13}}} & \X &  & \X & \X & \X & \X & \X & \X &  & \X &  &  &  &  &  &  &  &  &  &  &  &  &  &  &  &  \\
\anonstrut\textcolor{codelabel}{\scalebox{0.55}{\textsf{T14}}} & \X &  & \X & \X & \X &  &  & \X &  & \X &  & \X &  &  &  &  &  &  &  &  &  &  &  &  &  &  \\
\anonstrut\textcolor{codelabel}{\scalebox{0.55}{\textsf{T15}}} & \X &  & \X &  & \X & \X & \X & \X &  & \X & \X & \X &  &  &  &  &  &  &  &  &  &  &  &  &  & \X \\
\anonstrut\textcolor{codelabel}{\scalebox{0.55}{\textsf{T16}}} & \X &  & \X & \X &  &  &  & \X &  &  & \X & \X &  &  &  &  &  &  &  &  &  &  &  &  &  &  \\
\anonstrut\textcolor{codelabel}{\scalebox{0.55}{\textsf{T17}}} &  & \X & \X &  &  &  &  &  &  &  &  &  &  &  & \X & \X & \X &  &  &  &  &  &  &  & \X &  \\
\anonstrut\textcolor{codelabel}{\scalebox{0.55}{\textsf{T18}}} & \X &  &  &  &  & \X &  &  &  & \X &  &  &  &  &  &  &  &  &  &  &  &  &  &  &  &  \\
\anonstrut\textcolor{codelabel}{\scalebox{0.55}{\textsf{T19}}} & \X &  & \X &  &  &  &  &  &  & \X &  &  &  &  &  &  &  &  &  &  &  &  &  &  &  &  \\
\anonstrut\textcolor{codelabel}{\scalebox{0.55}{\textsf{T20}}} & \X &  & \X & \X &  &  &  & \X &  & \X &  & \X &  &  &  &  &  &  &  &  &  &  &  &  &  &  \\
\anonstrut\textcolor{codelabel}{\scalebox{0.55}{\textsf{T21}}} & \X &  & \X & \X & \X & \X & \X & \X &  & \X &  & \X &  &  &  &  &  &  &  &  &  &  &  &  &  &  \\
\anonstrut\textcolor{codelabel}{\scalebox{0.55}{\textsf{T22}}} & \X &  & \X &  &  &  &  &  &  & \X &  & \X &  &  &  &  &  &  &  &  &  &  &  &  &  &  \\
\anonstrut\textcolor{codelabel}{\scalebox{0.55}{\textsf{T23}}} & \X &  & \X &  &  &  &  &  &  & \X &  &  &  &  &  &  &  &  &  &  &  &  &  &  &  &  \\
\anonstrut\textcolor{codelabel}{\scalebox{0.55}{\textsf{T24}}} & \X &  &  &  &  & \X &  &  &  & \X &  &  &  &  &  &  &  &  &  &  &  &  &  &  &  &  \\
\anonstrut\textcolor{codelabel}{\scalebox{0.55}{\textsf{T25}}} & \X &  &  &  &  & \X & \X &  &  & \X &  &  &  &  &  &  &  &  &  &  &  &  &  &  &  &  \\
\anonstrut\textcolor{codelabel}{\scalebox{0.55}{\textsf{T26}}} & \X &  &  & \X &  &  &  &  &  &  &  &  &  &  &  &  &  &  &  &  &  &  &  &  &  &  \\
\bottomrule
\end{tabular}
\end{sidewaystable}
\tagmcend

\clearpage

%% ============================================================
%%  PART 2 — CODE LEGEND  (portrait, single page)
%% ============================================================
\captionsetup[longtable]{singlelinecheck=false}
{
\singlespacing
\setlength{\LTpre}{0pt}
\setlength{\LTpost}{4pt}

\small
\setlength{\tabcolsep}{4pt}
\begin{longtable}{p{0.9cm}>{\raggedright\arraybackslash}p{4cm}@{\hskip 0.5cm}p{0.9cm}>{\raggedright\arraybackslash}p{4cm}@{\hskip 0.5cm}p{0.9cm}>{\raggedright\arraybackslash}p{4cm}}
\caption[Extended Lithium Literature Review Code Legend.]{Extended Lithium Literature Review Code Legend. \\ Codes prefixed: \textsf{M} = main LCA/TEA papers (\emph{n}=47);
\textsf{S} = security \& policy papers (\emph{n}=19);
\textsf{T} = technical papers, no costs (\emph{n}=26).}
\label{tab:c2-litreview-legend}\\
\toprule
\textbf{Code} & \textbf{Reference} &
\textbf{Code} & \textbf{Reference} &
\textbf{Code} & \textbf{Reference} \\
\midrule
\endfirsthead
\toprule
\textbf{Code} & \textbf{Reference} &
\textbf{Code} & \textbf{Reference} &
\textbf{Code} & \textbf{Reference} \\
\midrule
\endhead
\bottomrule
\endfoot
\textcolor{codelabel}{\textsf{\textbf{M1}}} & \citet{agusdinataSocioenvironmentalImpactsLithium2018} & \textcolor{codelabel}{\textsf{\textbf{M32}}} & \citet{schenkerLithiumGeothermalBrines2024} & \textcolor{codelabel}{\textsf{\textbf{S16}}} & \citet{renMultiplexCharacteristicsVulnerability2024} \\
\textcolor{codelabel}{\textsf{\textbf{M2}}} & \citet{alessiaChallengesSustainableLithium2021} & \textcolor{codelabel}{\textsf{\textbf{M33}}} & \citet{schenkerCurrentFutureImpacts2025} & \textcolor{codelabel}{\textsf{\textbf{S17}}} & \citet{sunGlobalCompetitionLithiumIon2021} \\
\textcolor{codelabel}{\textsf{\textbf{M3}}} & \citet{ambroseUnderstandingFutureLithium2020a} & \textcolor{codelabel}{\textsf{\textbf{M34}}} & \citet{schombergExtendedLifeCycle2021} & \textcolor{codelabel}{\textsf{\textbf{S18}}} & \citet{tanMappingCobaltLithium2024} \\
\textcolor{codelabel}{\textsf{\textbf{M4}}} & \citet{buschEffectsDemandRecycling2025} & \textcolor{codelabel}{\textsf{\textbf{M35}}} & \citet{setuEnvironmentalImpactAssessment2026} & \textcolor{codelabel}{\textsf{\textbf{S19}}} & \citet{vivodaLithiumNexusEnergy2024} \\
\textcolor{codelabel}{\textsf{\textbf{M5}}} & \citet{fengEnvironmentalImpactsLithium2024} & \textcolor{codelabel}{\textsf{\textbf{M36}}} & \citet{mousavinezhadSustainableLithiumProduction2025} & \textcolor{codelabel}{\textsf{\textbf{T1}}} & \citet{averillREVIEWEXTRACTIVEPROCESSES1978} \\
\textcolor{codelabel}{\textsf{\textbf{M6}}} & \citet{fengEnvironmentalImpactsLithium2025} & \textcolor{codelabel}{\textsf{\textbf{M37}}} & \citet{stampEnvironmentalImpactsTransition2012} & \textcolor{codelabel}{\textsf{\textbf{T2}}} & \citet{battagliaRecoveryLithiumCarbonate2022} \\
\textcolor{codelabel}{\textsf{\textbf{M7}}} & \citet{fitzgeraldCanEmergingIndustrial2025} & \textcolor{codelabel}{\textsf{\textbf{M38}}} & \citet{sterbaLithiumMiningAccelerating2019} & \textcolor{codelabel}{\textsf{\textbf{T3}}} & \citet{bonkExAnteComparativeLife2022} \\
\textcolor{codelabel}{\textsf{\textbf{M8}}} & \citet{halkesLifeCycleAssessment2024} & \textcolor{codelabel}{\textsf{\textbf{M39}}} & \citet{thomsenMineManufacturerAssessing2025} & \textcolor{codelabel}{\textsf{\textbf{T4}}} & \citet{dengSystematicReviewSulfate2026} \\
\textcolor{codelabel}{\textsf{\textbf{M9}}} & \citet{hannaEnvironmentalImpactsReshoring2025} & \textcolor{codelabel}{\textsf{\textbf{M40}}} & \citet{tranChapter3Lithium2015} & \textcolor{codelabel}{\textsf{\textbf{T5}}} & \citet{ehrenProcessProducingLithium2018} \\
\textcolor{codelabel}{\textsf{\textbf{M10}}} & \citet{heEnvironmentalLifeCycle2025} & \textcolor{codelabel}{\textsf{\textbf{M41}}} & \citet{vannesteTechnoEconomicReviewCurrent2026} & \textcolor{codelabel}{\textsf{\textbf{T6}}} & \citet{flexerLithiumRecoveryBrines2018} \\
\textcolor{codelabel}{\textsf{\textbf{M11}}} & \citet{huangLifeCycleAssessment2021} & \textcolor{codelabel}{\textsf{\textbf{M42}}} & \citet{veraEnvironmentalImpactDirect2023} & \textcolor{codelabel}{\textsf{\textbf{T7}}} & \citet{garrettHandbookLithiumNatural2004} \\
\textcolor{codelabel}{\textsf{\textbf{M12}}} & \citet{ibarra-gutierrezProjectEconomicsLithium2021} & \textcolor{codelabel}{\textsf{\textbf{M43}}} & \citet{wangLifeCycleEnvironmental2026} & \textcolor{codelabel}{\textsf{\textbf{T8}}} & \citet{haddadElectrochemicalLithiumExtraction2024} \\
\textcolor{codelabel}{\textsf{\textbf{M13}}} & \citet{istrateDecarbonizingLithiumionBattery2024} & \textcolor{codelabel}{\textsf{\textbf{M44}}} & \citet{warrenTechnoEconomicAnalysisLithium2021} & \textcolor{codelabel}{\textsf{\textbf{T9}}} & \citet{kunduRecoveryLithiumSpodumenebearing2023} \\
\textcolor{codelabel}{\textsf{\textbf{M14}}} & \citet{iyerLithiumProductionNorth2022} & \textcolor{codelabel}{\textsf{\textbf{M45}}} & \citet{wesselkaemperPrimaryBatteryMaterial2025} & \textcolor{codelabel}{\textsf{\textbf{T10}}} & \citet{liRecoveryLithiumMineral2019} \\
\textcolor{codelabel}{\textsf{\textbf{M15}}} & \citet{iyerLithiumProductionNorth2023} & \textcolor{codelabel}{\textsf{\textbf{M46}}} & \citet{yelatontsevProcessingLithiumOres2021} & \textcolor{codelabel}{\textsf{\textbf{T11}}} & \citet{liuNovelApproachesLithium2019} \\
\textcolor{codelabel}{\textsf{\textbf{M16}}} & \citet{iyerLifecycleAnalysisLithium2024} & \textcolor{codelabel}{\textsf{\textbf{M47}}} & \citet{zhouImplicationsUncertaintiesGlobal2026} & \textcolor{codelabel}{\textsf{\textbf{T12}}} & \citet{liuProcessAnalysisStudy2021} \\
\textcolor{codelabel}{\textsf{\textbf{M17}}} & \citet{iyerLifecycleGreenhouseGas2026} & \textcolor{codelabel}{\textsf{\textbf{S1}}} & \citet{andradeExplorationProductionUnderstanding2024} & \textcolor{codelabel}{\textsf{\textbf{T13}}} & \citet{liuReviewLithiumExtraction2023} \\
\textcolor{codelabel}{\textsf{\textbf{M18}}} & \citet{jiangEnvironmentalImpactsLithium2020} & \textcolor{codelabel}{\textsf{\textbf{S2}}} & \citet{ballingerVulnerabilityElectricVehicle2019} & \textcolor{codelabel}{\textsf{\textbf{T14}}} & \citet{marcinovLithiumProductionRecovery2023} \\
\textcolor{codelabel}{\textsf{\textbf{M19}}} & \citet{kallitsisThinkGlobalAct2024} & \textcolor{codelabel}{\textsf{\textbf{S3}}} & \citet{bridgeLithiumBrexitGlobal2023} & \textcolor{codelabel}{\textsf{\textbf{T15}}} & \citet{meshramExtractionLithiumPrimary2014} \\
\textcolor{codelabel}{\textsf{\textbf{M20}}} & \citet{kellyEnergyGreenhouseGas2021} & \textcolor{codelabel}{\textsf{\textbf{S4}}} & \citet{buschEffectsDemandRecycling2025} & \textcolor{codelabel}{\textsf{\textbf{T16}}} & \citet{munkLithiumBrinesAGlobal2016} \\
\textcolor{codelabel}{\textsf{\textbf{M21}}} & \citet{khakmardanLifeCycleAssessment2024} & \textcolor{codelabel}{\textsf{\textbf{S5}}} & \citet{calisaya-azpilcuetaAssessmentSupplyChain2020} & \textcolor{codelabel}{\textsf{\textbf{T17}}} & \citet{notterContributionLiIonBatteries2010} \\
\textcolor{codelabel}{\textsf{\textbf{M22}}} & \citet{khakmardanConstructingLifeCycle2025} & \textcolor{codelabel}{\textsf{\textbf{S6}}} & \citet{chengElectricVehicleBattery2024} & \textcolor{codelabel}{\textsf{\textbf{T18}}} & \citet{parkEffectiveLithiumExtraction2023} \\
\textcolor{codelabel}{\textsf{\textbf{M23}}} & \citet{lappalainenEnvironmentalImpactsLithium2024} & \textcolor{codelabel}{\textsf{\textbf{S7}}} & \citet{egbueCriticalIssuesSupply2012} & \textcolor{codelabel}{\textsf{\textbf{T19}}} & \citet{perezMethodProductionBattery2014} \\
\textcolor{codelabel}{\textsf{\textbf{M24}}} & \citet{llamas-orozcoEstimatingEnvironmentalImpacts2023} & \textcolor{codelabel}{\textsf{\textbf{S8}}} & \citet{greitemeierChinasHoldLithiumion2025} & \textcolor{codelabel}{\textsf{\textbf{T20}}} & \citet{stringfellowTechnologyRecoveryLithium2021} \\
\textcolor{codelabel}{\textsf{\textbf{M25}}} & \citet{machalaLifeCycleComparison2025} & \textcolor{codelabel}{\textsf{\textbf{S9}}} & \citet{helbigSupplyRisksAssociated2018} & \textcolor{codelabel}{\textsf{\textbf{T21}}} & \citet{swainRecoveryRecyclingLithium2017} \\
\textcolor{codelabel}{\textsf{\textbf{M26}}} & \citet{manjongLifeCycleModelling2021} & \textcolor{codelabel}{\textsf{\textbf{S10}}} & \citet{henselChallengesOpportunitiesGlobal2023} & \textcolor{codelabel}{\textsf{\textbf{T22}}} & \citet{torresSalarAtacamaLithium2024} \\
\textcolor{codelabel}{\textsf{\textbf{M27}}} & \citet{mousavinezhadEnvironmentalImpactAssessment2024} & \textcolor{codelabel}{\textsf{\textbf{S11}}} & \citet{jinChinasLithiumSupply2023} & \textcolor{codelabel}{\textsf{\textbf{T23}}} & \citet{wilkomirskyProductionLithiumCarbonate1999} \\
\textcolor{codelabel}{\textsf{\textbf{M28}}} & \citet{mousavinezhadSustainableLithiumProduction2025} & \textcolor{codelabel}{\textsf{\textbf{S12}}} & \citet{manleyModelAssessIndustry2022} & \textcolor{codelabel}{\textsf{\textbf{T24}}} & \citet{yanExtractionLithiumLepidolite2012} \\
\textcolor{codelabel}{\textsf{\textbf{M29}}} & \citet{nikfarUnlockingSustainableLithium2025} & \textcolor{codelabel}{\textsf{\textbf{S13}}} & \citet{miattoRiseFallAmerican2020} & \textcolor{codelabel}{\textsf{\textbf{T25}}} & \citet{zhaoRecentAdvancesLithium2023} \\
\textcolor{codelabel}{\textsf{\textbf{M30}}} & \citet{royLithiumClayAssessing2024} & \textcolor{codelabel}{\textsf{\textbf{S14}}} & \citet{ouyangSystematicRisksGlobal2024} & \textcolor{codelabel}{\textsf{\textbf{T26}}} & \citet{wangLifeCycleEnvironmental2026} \\
\textcolor{codelabel}{\textsf{\textbf{M31}}} & \citet{schenkerRegionalizedLifeCycle2022} & \textcolor{codelabel}{\textsf{\textbf{S15}}} & \citet{prinaceraiGeographyControlDeep2024} &  &  \\
\end{longtable}

}

%% Belt-and-suspenders: explicitly restore global spacing in case anything
%% above emitted a \global spacing command that escaped the group.
\onehalfspacing

\subsubsection{Firm-level process modeling, costs, and economic incentives}\label{app:04-litreview-incentives}

At a theoretical level, the structure of global value chains and which locations become developed are driven by costs borne by firms, compared to their competitiveness with current and future prices. Comparative cost advantages exist between locations that depend on `factor endowments' (e.g., labor, capital, resources) \cite{ricardoForeignTrade1817, leamerHeckscherOhlinModelTheory1995}, mediated by regulation. For a commodity, the structure of the market is generally determined by the price of the commodity, with mines `turning on' and `turning off' when the price of the commodity exceeds their marginal operating cost, mediated by uncertainty in future prices. While lithium is not a pure commodity in the sense that significant amounts of production are under contract, we can model `competitiveness' from a national security perspective at a first order by understanding the relative operating production costs of differing production pathways.

We do so via technoeconomic modeling: a generalized set of methods to determine production costs as a function of location and production pathway (technology choice). While many applications of technoeconomic methods have been applied to battery and electric vehicle manufacturing (e.g. \cite{duffnerBatteryCostModeling2020, saktiTechnoeconomicAnalysisOptimization2015, orangiTechnoEconomicModelBenchmarking2022}), far fewer works have evaluated production costs through the upper and midstream battery supply chain, particularly in the context of lithium.

This gap may stem from the ongoing pace of technological advancements in lithium extraction and refining, as well as the wide range of technologies at varying stages of maturity, from well-established methods like spodumene processing to emerging processes such as direct lithium extraction (DLE). The diversity of technological readiness, combined with limited data availability and rapidly shifting market conditions, makes it challenging to develop standardized cost models across the upper and midstream supply chain. As such, many papers in the literature provide reviews of bench-scale or even industrial-scale processes from a scientific, technological perspective (see the non-comprehensive 20 studies listed under `Technical Review Papers' in \ref{tab:c1-litreview-matrix}), but do not fully consider the economic dimension.

For those that do, two primary methods emerge: regression-based models and process-specific models.

Regression-based cost modeling can give a general sense of costs given a certain project's characteristics. Three studies use regression models to predict costs without considering emissions: \citet{ambroseUnderstandingFutureLithium2020a}, \citet{flemingLongrunAvailabilityMineral2024}, and \citet{zhouImplicationsUncertaintiesGlobal2026}. All of these studies rely on white and grey literature databases, often S\&P Global Capital IQ cost data, estimating average production costs from input variables for each project (e.g., deposit grade and location, details in company reports, and evaporation rate for brine deposits). These studies are limited when production in different countries and pathways needs to be understood at a detailed level; in particular, the effect of targeted policy interventions cannot be easily well-understood (e.g., reducing cost of capital, subsidies for specific material inputs, etc.). Furthermore, regression-based models suffer when input data is sparse, which is the general case for the lithium supply chain, and cannot adequately predict the costs of developing technologies (e.g., direct lithium extraction). Even some mature projects have production costs that appear to be notably overstated or understated, for example \$10,300 for Silver Peak \cite{zhouImplicationsUncertaintiesGlobal2026, ambroseUnderstandingFutureLithium2020a} versus the \$5,520 midpoint of our model and \$4,600---6,750 range reported in other literature, or \$1,200 for lithium production from lepidolite in Yichun versus the \$12,140 midpoint of our model and \$10,000-16,680 range reported in other literature. Thus, while regression-based cost modeling may be sufficient for certain questions, we suggest that an engineering-calibrated model structure may result in more realistic outcomes. 

Numerous studies, especially in the last five years (e.g. \cite{chordiaLifeCycleEnvironmental2022, schenkerLithiumGeothermalBrines2024, khakmardanComparativeLifeCycle2023, iyerLifecycleAnalysisLithium2024, kellyEnergyGreenhouseGas2021, mousavinezhadEnvironmentalImpactAssessment2024, jiangEnvironmentalImpactsLithium2020, nikfarUnlockingSustainableLithium2025}), have analyzed lithium production pathways from an environmental perspective, especially given lithium's reputation for high water usage in sensitive areas, and the general understanding that mineral processing requires significant amounts of energy and chemical resources. The large majority of this literature originates from the Life Cycle Assessment (LCA) methodology, which not only considers direct and indirect greenhouse gas (GHG) and particulate emissions but also additional environmental impact factors, such as water impacts, ecotoxicity, radiative damage, and many others. These models typically build detailed models of processes throughout the included system (`foreground'), while relying on external models and measurements of the system that is not explicitly modeled (`background').

This literature can be limited because:
\begin{itemize}
    \item Production in specific locations with firm-specific pathways are modeled, yet measurements from these studies are often used to represent the entire industry without regard for these variations in geography and technology \cite{sunRobustAssessmentsLithium2025, kallitsisThinkGlobalAct2024};
    
    \item Measurements, calculated as emissions factors (emissions per unit output), are necessarily `attributional' and do not capture the marginal change in emissions if production were to change \cite{nationalacademiesofscienceengineeringandmedicineAssessmentTechnologiesImproving2021}; and

    \item Studies are conducted independently, without necessarily trying to have the same scope, so there is no guarantee of equal comparison across studies \cite{bergersonLifeCycleAssessment2020}. (Most studies use the same functional unit of one kilogram or metric ton of lithium carbonate, but boundaries of what are included in the foreground versus background systems are not always specified clearly or explicitly disagree with other studies.)
\end{itemize}

Of particular importance is the difference in emissions across geographical areas. \citet{kallitsisThinkGlobalAct2024} highlight this across battery supply chains generally, finding that variation in production location across Chinese and North American states and provinces, as well as 19 other countries, can result in a factor of approximately 4 difference between the carbon footprints of different battery material manufacturing processes. \citet{sunRobustAssessmentsLithium2025} highlight this finding for lithium in particular: ``environmental impact intensities of battery-grade \ce{Li2CO3}. production from various mine sites differ between 4 times for GHG emissions to 2885 times for land use.'' Indeed, when we compile reported emissions from 103 pathways across 19 studies, we find widely different emissions factors across geographies and production pathways (Fig. 1). However, due to the depth to which these studies analyze different production pathways, they can provide a foundation upon which process-based integrated technoeconomic--environmental models can be built.

\subsubsection{Combining economic and environmental dimensions of lithium production}\label{app:04-litreview-tealca}
Several examples within the battery supply chain have combined dimensions of technoeconomic and life-cycle assessment, particularly in the context of battery manufacturing; a non-exhaustive search of recent literature (within five years) in the battery manufacturing space turned up 6 models \cite{daiEverBattClosedloopBattery2019, chenTechnoeconomicAnalysisLithiumion2023, guvenLifecycleAssessmentLifecycle2023, gutschCostsCarbonFootprint2024, maPathwayDecisionsReuse2024, orangiBottomupFrameworkInvestigate2024}. However, far fewer have focused on the lithium system specifically, and those that have have been limited to a few production pathways. These studies are limited by only modeling a subset of all technical production pathways, with varying levels of detail about cost and environmental impact, focusing on individual pathways (mineral resources \cite{gaoLithiumExtractionHard2023, liRecoveryLithiumMineral2019} and geothermal brines \cite{huangLifeCycleAssessment2021}). However, as previously discussed, not only can we not necessarily compare analyses or interventions across these studies due to possible inconsistencies in system boundaries, functional unit definitions, and other assumptions \cite{bergersonLifeCycleAssessment2020}, but standard solar evaporative brines and clay resources (which are significant in the United States) are not modeled.

\subsubsection{Empirical Technological Context and Systems Scoping}\label{app:04-litreview-context}
Two primary compounds are used in lithium-ion batteries: lithium hydroxide and lithium carbonate. Lithium hydroxide is primarily used when producing nickel-rich cathodes (e.g., NMC, NCA) due to its lower decomposition temperature. The complex integration of multiple ions, particularly nickel, requires high-purity inputs and precise temperatures to achieve the proper structure. Lithium carbonate, which is cheaper, is used for LFP and LCO batteries, due to the less stringent thermal and structural requirements for producing those chemically simpler cathodes. Other lithium compounds (lithium fluoride, lithium chloride) are used for the lithium in the electrolytes.

In our analysis, we initially limit our scope to lithium carbonate, as is standard in both industry (e.g. \cite{tradingeconomics.comLithiumPriceChart, fastmarketsLithiumCarbonate99}) and literature (e.g. \cite{flemingLongrunAvailabilityMineral2024, usgeologicalsurveyMineralCommoditySummaries2026, vikstromLithiumAvailabilityFuture2013}). Discussion of the implications of this scope limitation are included in the Limitations section.

There are three major sources of lithium: brines, hard rock ores, and clays. Each of these can be broken down into additional subcategories: salar/continental brines and geothermal brines; spodumene, lepidolite, and other minerals; hectorite, zinnwaldite, jadarite, and others. Each of these categories has a relatively standard extraction and refining process, though new technologies (such as Direct Lithium Extraction (DLE)) are being developed. While much more detail could be incorporated due to inter-firm differences in production methods, the specific chemical composition of each mineral resource, and nuances in daily operations, we show in \ref{fig:c-pathways} a simplified block diagram of primary production pathways for the production of lithium carbonate from different lithium-bearing resources. This includes the pathways described in the main body of the paper, namely Silver Peak, Thacker Pass, and Jianxiawo, but these models can be adapted with other localized parameters that can allow for comparison with other possible production pathways involving other pathways in China and the United States, pathways in non-aligned countries, multi-nation production pathways, or other outside options.

\begin{figure}
    \centering
    \altfig{Flowchart showing five lithium production pathways organized in three rows: material source at top, processing in the middle, and primary product at bottom. From left to right the sources are: brine evaporative (Chile, Argentina, USA) flowing through pumping and pond evaporation to batch precipitation or solvent extraction polishing, producing lithium carbonate; brine direct lithium extraction (China, USA, Canada, Argentina) flowing through pumping and adsorption to ion exchange or solvent extraction; spodumene ore (Australia, China, Brazil, Canada, Ghana, USA) flowing through drill and blast mining and optional beneficiation to sulfuric acid roasting and water leaching; lepidolite ores (China) flowing through drill and blast mining to sulfate roasting and water leaching; and lithium clays (USA, Mexico, Peru) flowing through surface excavation and comminution to direct acid leaching or sulfate roasting and water leaching. Established processes are shown with solid-line boxes, emerging processes with dashed-line boxes, and processes not in scope are greyed out.}
        \includegraphics[width=1\linewidth]{figures/c-pathways.png}
    \altfigend
    \caption[{Existing and future lithium production pathways}]{Existing (solid lines) and future (dashed lines) lithium production pathways}
    \label{fig:c-pathways}
\end{figure}

\clearpage
\subsection{Limitations with Scoping and Measurement} \label{app:04-limitations}

Several limitations arise due to the scope of this work, which we explore below. 

\subsubsection{Focus on Lithium Carbonate for Batteries} \label{app:04-limitations-carbonate}

This study primarily focuses on lithium demand from the battery sector, which represents approximately 88\% of total projected demand \cite{usgeologicalsurveyMineralCommoditySummaries2026}. While this approach excludes smaller markets such as glass manufacturing, lubricants, or pharmaceuticals, the overwhelming dominance of battery demand means that the findings should still capture the primary driver of overall future lithium demand.

Furthermore, we are presently only considering lithium carbonate as the primary product of comparison. Limiting the scope to lithium carbonate provides a more consistent basis for comparing different production pathways, as all resource types can be processed into lithium carbonate. In contrast, lithium hydroxide is currently primarily derived from lithium carbonate, but could be produced directly from lithium resources in a variety of generally more novel, less mature processes, making direct technoeconomic comparisons more complex. As mentioned above, this scoping limitation is standard in both industry and literature, as some benchmark is needed for price comparisons and academic studies. Future work could explicitly model lithium hydroxide production and assess how technology, geography, and policy interventions impact the relative competitiveness of lithium carbonate versus lithium hydroxide pathways, and subsequent implications for meeting consumer preferences.

\subsubsection{Static versus Dynamic Modeling}\label{app:04-limitations-static}

This analysis employs a static technoeconomic modeling framework. While this limits the ability to represent technological improvements, shifts in overarching market conditions, or responses by firms to interventions occurring in other locations, the static approach provides a transparent, point-in-time comparison of pathways under consistent assumptions. However, we argue that this type of analysis offers the perspective of a first-order response from lithium producing firms, who by and large would react to policies by making changes to their baseline business-as-usual production. These baselines (and first-order responses to policy) can be useful for identifying which production pathways and regions are currently competitive or vulnerable to policy interventions.

\subsubsection{Technological Maturity}\label{app:04-limitations-maturity}

Lithium production pathways vary significantly in technological maturity, from well-established methods like spodumene conversion and evaporative brine concentration to emerging technologies such as direct lithium extraction (DLE) or lithium recovery from clay deposits. While uncertainties in cost and performance for newer technologies are unavoidable, including these pathways provides valuable insight into the potential role they could play in diversifying supply chains. Highlighting these uncertainties under best practices for emerging technologies \cite{woods-robinsonControversyConsensusCommon2025, bredehoeft47R11CostEstimate2020} and recognizing that the study will have limitations can make the study useful for identifying where further technological development or data collection could reduce uncertainties-and vulnerabilities.

\subsubsection{Focus on Representative, Significant Pathways Rather Than Comprehensiveness}\label{app:04-limitations-representative}

Our analysis aims to construct representative models for each major lithium production pathway across key production pathways, accounting for the most significant current and future pathways identified by the supply chain security and materials demand literature, rather than attempting an exhaustive accounting of every facility or process variant. This approach is necessitated by the substantial variability in cost and emissions data, which stems from differences in process design, energy sources, and operational practices, which are difficult to quantify precisely. While alternative pathways can be expanded upon later with project-specific data, we focus on the pathways define the current state of production in the United States and the marginal production worldwide. 

\clearpage
\subsection{Additional Results and Model Outputs}\label{app:04-moreresults}

\begin{figure}
    \centering
    \altfig{Scatter plot with production cost in dollars per tonne lithium carbonate on the x-axis (0 to 16,000) and life cycle greenhouse gas emissions in tonnes \ce{CO2}-equivalent per tonne on the y-axis (0 to roughly 70). Three projects are shown as colored rectangles: U.S. Evaporative Brine at Silver Peak, Nevada in light blue in the lower left around 4,400 to 7,000 dollars and 4.7 to 8.5 tonnes \ce{CO2eq}; U.S. Clays at Thacker Pass, Nevada in dark blue in the center around 5,900 to 13,300 dollars and 12 to 20.5 tonnes \ce{CO2eq}; and Chinese Lepidolite at Jianxiawo, Jiangxi in red in the upper right around 8,000 to 14,000 dollars and 15 to 65 tonnes \ce{CO2eq}. Darker shaded rectangles show this study's range with a central dot at the midpoint. Lighter shaded rectangles span the full range of literature values, with individual data points marked by letters corresponding to cited sources. Literature points are distinguished as peer-reviewed (dashed borders), grey literature (dotted borders), or disclosure-reported values (solid borders). The lepidolite pathway shows the widest uncertainty in both cost and emissions.}
        \includegraphics[width=1\linewidth]{figures/c-fullres.png}
    \altfigend
    \caption[Production cost versus life cycle greenhouse gas emissions for three key lithium projects representing distinct resource-pathway combinations, with full uncertainty ranges.]{Production cost (\$/t \ce{Li2CO3}) versus life cycle greenhouse gas emissions (t\ce{CO2eq}/ /t \ce{Li2CO3}) for three key lithium projects representing distinct resource-pathway combinations, with full uncertainty ranges.}
    \label{fig:c-extendedresults}
\end{figure}

\ref{fig:c-extendedresults} shows the simultaneous impact of all of the best-case and worst-case scenarios across model input variation - machine design and material costs. Comparing this figure to \ref{fig:04-costemissions}, the top 3 restricted ranges capture the majority of the variation for Silver Peak and Thacker Pass, where a small number of drivers dominate; more of the literature range is also covered by this broader interpretation of model inputs. For Jianxiawo, the full uncertainty range is considerably wider in both dimensions, reflecting the larger number of parameters carrying meaningful uncertainty when lepidolite process data is reconstructed from patents and lab-scale studies rather than engineering disclosures. Notably, while the ranges for the Silver Peak project remain relatively separated from the other two projects, the full ranges for Thacker Pass and Jianxiawo overlap more substantially. However, two features of the analysis limit the scope of this overlap. First, it requires simultaneously pessimistic assumptions for the clay pathway and optimistic assumptions for lepidolite --- a combination that is unlikely to hold in practice given that the dominant cost drivers for each pathway are largely independent of one another. Second, the Thacker Pass range shown here reflects the generic flowsheet and high reagent costs, which is already conservative relative to the project-specific design; the overlap would likely narrow or disappear if the magnesium crystallization route and no consideration of offsite sulfuric acid production costs were incorporated. 

Across all three projects, the conservative bounds are asymmetrically larger than the optimistic bounds relative to the midpoint, reflecting the fact that cost and emissions overruns are structurally easier to achieve than equivalent improvements, particularly where reagent consumption or equipment sizing are at the high end of their plausible ranges.

\begin{table}[htbp]
\caption[Modeled cost structure and key input intensities for the three focus projects.]{Modeled cost structure and key input intensities for the three focus projects, at the midpoint scenario. Cost-structure rows sum to the production cost; differences of one unit are rounding.}
\label{tab:04-summary-majornumbers}
\centering
\small
\begin{tabular}{lccc}
\toprule
 & Silver Peak & Thacker Pass & Jianxiawo \\
\midrule
Production volume (t \ce{Li2CO3}/yr)              & 6,000  & 66,800 & 46,000 \\
\textbf{Production cost (\$/t \ce{Li2CO3})}       & \textbf{5,520} & \textbf{8,703} & \textbf{12,140} \\
\textbf{GHG emissions (t \ce{CO2eq}/t)}           & \textbf{6.01}  & \textbf{16.67} & \textbf{30.26} \\
\midrule
\multicolumn{4}{l}{\textit{Cost structure (\$/t \ce{Li2CO3})}} \\
\quad Reagents                                    & 1,518 & 5,126 & 5,192 \\
\quad Energy and utilities                        & 585   & 83    & 2,423 \\
\quad Labor                                       & 2,402 & 1,326 & 428 \\
\quad Waste, tailings and haulage                 & 26    & 1,565 & 3,108 \\
\quad Maintenance and overhead                    & 989   & 603   & 988 \\
\midrule
\multicolumn{4}{l}{\textit{Key inputs (per t \ce{Li2CO3})}} \\
\quad Electricity consumed (\si{\kilo\watt\hour}) & 6,288 & 2,170 & 4,742 \\
\quad Electricity exported, on-site acid plant (\si{\kilo\watt\hour}) & --- & $-$4,137 & --- \\
\quad Purchased process water (\si{\cubic\metre}) & 0.09  & 203   & 408 \\
\quad Lime (kg)                                   & 1,162\textsuperscript{a} & 9,924 & 4,684 \\
\quad Soda ash (kg)                               & 5,057 & 3,625 & 3,511 \\
\quad Pathway reagent (kg)                        & ---   & sulfur 8,680 & \ce{Na2SO4} 15,612 \\
\bottomrule
\end{tabular}

\vspace{2pt}
\footnotesize
\textsuperscript{a}Dosed as \ce{CaO}; purchased as quicklime and slaked on site.
\end{table}

Table \ref{tab:04-summary-majornumbers} also summarizes key costs 

\clearpage
\subsection{Detailed Methods Writeup}\label{app:04-methods}

Lithium production systems span heterogeneous resource types (brines, hard rock, claystone, geothermal fluids), processing pathways (physical beneficiation, leaching, solvent extraction, precipitation), and geographic supply chains. These systems differ materially in ore grade, lithium concentration, reagent intensity, energy demand, logistics requirements, and capital intensity. Direct comparison across projects is therefore non-trivial. We use Technoeconomic Assessment (TEA) combined with elements of Life Cycle Analysis (LCA) to identify the `competitiveness' of these projects.

The purpose of this framework is to:
\begin{itemize}
    \item Standardize comparison across lithium production pathways using a consistent mass and energy accounting structure.
    \item Link chemistry-dependent material intensity to equipment sizing, enabling realistic scaling of capital and operating costs.
    \item Enable parametric sensitivity analysis (e.g., grade, recovery, transport distance, energy price).
    \item Provide modular extensibility, such that new technologies, detailed process modeling, or site-specific data can be incorporated without restructuring the full model.
    \item Support implementation and evaluation of engineering-constrained policy interventions, including measures such as carbon pricing, water-use and environmental restrictions, production-based subsidies, and domestic production requirements, where process-level detail is required to understand the impact of such interventions.
\end{itemize}

\subsubsection{Methodological Literature Review}\label{app:04-methods-litreview}

Across the recent lithium supply chain modeling literature, various levels of abstraction are engaged, to answer differing research and policy questions. Many papers cover multiple resource pathways (traditional brines, direct lithium extraction, spodumene, clays), but far fewer simultaneously analyze key parameters around costs and emissions while also being open and transparent about their model forms and assumptions.

Regression and reduced-form cost models \citep[e.g.,][]{ambroseUnderstandingFutureLithium2020b, flemingLongrunAvailabilityMineral2024, buschEffectsDemandRecycling2025} infer costs from project descriptors and focus on outcomes, rather than process models and manufacturing choices. These models can allow for near-comprehensive coverage of projects, but make counterfactual policy analysis challenging: if a policy changes reagent prices, electricity grid compositions, recovery, financing terms, or the like, these regression models can only respond indirectly at best.

Process-specific, bottom-up technoeconomic assessment (TEA) and life cycle assessment (LCA) models are on the opposite end: they typically represent a pathway in detail, but often only a specific project or set of projects. The LCA and environmental impacts literature is particularly robust, with over 30 articles since 2020 discussing various environmental impacts (usually, but not exclusively, greenhouse gases) but these are not necessarily consistent with each other given differing system boundaries. Technoeconomic studies are much fewer in number, but their lack of clarity on specific engineering modeling assumptions and lack of publicly available, interactive models makes it challenging to assess conclusions and apply further policy interventions.

This work contributes an openly examinable, modular, process-explicit, machine-level production modeling framework, combining both technoeconomic and life-cycle assessment principles within a single accounting framework. This level of bottom-up detail and open modeling has to the best of the authors' awareness not been applied in the context of lithium production and its distinct production pathways.

This methodology is inspired by parallel methodologies from process-based cost modeling and physical input-output tables, with sufficient detail to remain responsive to policy-relevant engineering variables, while avoiding the data intensity, engineering complexity, and narrow scopes associated with high-fidelity process simulation methods such as Aspen Plus or IDAES \cite{leeIDAESProcessModeling2021}. With each step centered around target production volumes, linked together with key operational parameters such as yield rates, reagent intensity, and conversion factors, capital and operating costs remain causally linked across production steps, facilities, and supply-chain representations. By operating at an intermediate resolution, the framework enables standardized comparisons across lithium pathways, while maintaining extensibility to emerging technologies and site-specific data.

\clearpage
\subsection{Modeling System}\label{app:04-methods-system}

\paragraph{System Architecture and Module Hierarchy}\label{app:04-methods-system-hierarchy}\leavevmode\par

The modeling framework is structured hierarchically into four module types:

\begin{description}
    \item[Production Step] Production steps are the fundamental unit operation or process block in which material transformation occurs, in the style of physical input-output analysis. Production steps can be high-level representative blocks, or replaced with sets of specific step-by-step breakdowns, or even integrated with detailed physics-based models.

    \item[Facility] A facility module is a collection of production steps co-located at a single site, operated as an integrated processing system. Production steps in the same facility are subject to the same geographical (e.g.\ electricity grid, labor costs) and operational context (e.g.\ working days per year, shifts per day).

    \item[Transportation] Transportation route modules are representations of discrete movements of material between facilities, composed of transportation legs to account for the possibility of multiple transportation steps between facilities (e.g.\ truck to barge to rail).

    \item[Supply Chain] Supply chain modules are network-level representations linking facilities, with transportation legs connecting them as appropriate. This level also aggregates the statistics from the facility, transportation, and production step level.
\end{description}

\paragraph{Production Step Formulation}\label{app:04-methods-system-prodstep}\leavevmode\par

Each production step is represented as a structural transformation consisting of a primary input, secondary inputs (reagents, utilities, consumables), a process unit (black-box transformation of inputs to outputs), a primary output, and secondary outputs (byproducts, wastes, emissions). For modeling tractability, each production step is constrained to one primary input stream and one primary output stream, enabling consistent mass accounting and simplified scaling relationships.

\medskip
\noindent\textbf{Modeling of Materials: Inputs, Outputs, Conversion Factors, and Units.}\quad

The model explicitly defines three material volume quantities at each production step: primary inputs (e.g., tonnes ROM ore, \si{m^3} brine extracted), the process basis (e.g., \si{m^3} slurry processed, tonnes solids milled, \si{m^3} solution reacted), and primary outputs (e.g., \si{m^3} slurry, tonnes \ce{Li2CO3} product, tonnes concentrate). These quantities are not inherently equivalent, nor in the same units. Equipment sizing and capital cost estimations are based on the process unit throughput, which serves as the internal scaling metric at each step.

To accommodate the diversity of natural reporting units across process steps and feedstock types, we define two conversion factors to mediate between the three volume quantities: 

\begin{description}[noitemsep, topsep=0pt]
    \item[\textbullet\enskip Input-to-Process Conversion Factor (IPCF)]: Converts the primary input from its natural reporting units (e.g., tonnes ore, \si{m^3} brine, tonnes slurry) into standardized process units.
    
    \item[\textbullet\enskip Process-to-Output Conversion Factor (POCF)]: Converts standardized process units into primary output units (e.g., tonnes \ce{Li2CO3}, tonnes concentrate, \si{m^3} solution).
\end{description}

These factors are designed to be flexible, allowing for ad hoc adjustments such as volume changes due to reagent addition or recycling loops in the process model. As such, care must be taken to ensure that units are consistent across the three volume quantities (step input, step basis, and step output). Although the model architecture is designed to allow for multiple primary inputs and outputs, the present implementation models a single primary input, representing the output of the previous step. The IPCF then converts this volume into the process basis units without further modifications to its composition. 

Secondary inputs, such as reagents and additives, are dosed relative to either the process volume or the concentration of a target constituent, expressed in relative mass units (e.g. kg reagent per kg constituent, or kg per tonne process basis). They are assumed not to materially affect the bulk process volume; where reagents do significantly affect process volume, this is instead absorbed into the IPCF. With the addition of these secondary inputs, the primary and secondary inputs become the process basis, but no interaction is assumed yet; the transformation of the process basis into the primary output occurring at the POCF interface is where any compositional changes are assumed, treated as a black-box operation. 

Material constituents are tracked for each input, process basis, and output volume, recorded as mass fractions relative to the process unit (e.g., \si{kg/m^3}, kg/kg, etc.). Constituent concentrations $c_\text{constituent}$ are always expressed as weights (usually parts per million (PPM)) relative to the constituent-bearing fraction of the volume, such as mg Li per kg dissolved solids in a slurry mixture, or mg Li per kg \ce{Li2CO3} in a lithium cake. This convention ensures that constituent concentrations are invariant to dilution or concentration of the carrying phase; in this way, when units change, the underlying concentrations of the tracked constituent carrying phase do not need to vary. However, understanding the absolute constituent mass per material flow unit is important, for example to size reagent doses; often, reagents are intended to target a particular constituent of the process stream (e.g., lithium, magnesium, calcium, etc.) rather than the entirety of the production stream. Thus, each volume has an associated Constituent Conversion Factor (CCF), which converts from the constituent units (always mass units) to the volume unit. Reagent use $M_r$ (also always in mass units) is therefore modeled as:

\begin{equation}
    M_r = r_\text{reagent} \cdot c_\text{constituent} \cdot \text{CCF}
    \label{eq:c-reagentmass}
\end{equation}

where $r_\text{reagent}$ is the reagent ratio in kg reagent per kg constituent, $c_\text{constituent}$ is the constituent concentration (in this work, always reported as ppm, or mg/kg).

The process basis volume is determined by the volume defining basis, which may either be the process step basis itself or a constituent. The latter is required when equivalent outputs demand different process volumes depending on feed composition. For example, all else equal, more volume of lower lithium content brine will need to be processed than higher lithium content brine, to yield the same mass of lithium carbonate product. In such cases, the process basis is sized by the ratio of constituent mass per process unit to constituent mass per output unit. CCFs are essential to this calculation: since constituent mass must be conserved across this process black box boundary (subject to yield losses), both the process and output concentrations must be converted to the same absolute units before the sizing ratio is formed. This is particularly important when the constituent-bearing material changes physical identity between steps, often resulting in significantly different concentration values --- for example, dissolved lithium in concentrated brine becoming lithium in lithium carbonate crystals during precipitation --- where concentration values reference different physical bases and cannot be directly compared without conversion through their respective CCFs. 

Thus, with a defined output production target, we can determine the process basis volume $Q_{\text{process}}$ and input to each step:
\begin{equation}
    \sum \text{Inputs} \times \text{IPCF} = Q_{\text{process}} \times r_{\text{yield}} 
    \times r_{\text{conv}} = \frac{\sum \text{Outputs}}{\text{POCF}}
\end{equation}
where $r_{\text{yield}}$ is the bulk yield of the process volume (assuming a fraction of everything is lost proportionally). The conversion ratio $r_{\text{conv}}$ is equal to one if the volume-defining basis is the process step basis itself. However, when sizing is constituent driven, it is given by:
\alteq{Conversion ratio r-conv equals c-in times the basis capacity conversion factor, divided by c-out times the output capacity conversion factor times the process operating capacity factor.}
\begin{equation}
    r_{\text{conv}} = \frac{c_{\text{in}} \cdot \text{CCF}_{\text{basis}}}{c_{\text{out}} \cdot \text{CCF}_{\text{output}} \cdot \text{POCF}}
\end{equation}
\alteqend
representing the ratio of absolute constituent mass per process basis unit to absolute constituent mass per output unit, with the POCF normalizing for the unit conversion between process basis and output. 

Steps are linked sequentially by passing each step's output as the input to the succeeding step, such that a single terminal production target propagates upstream through the full process chain, determining throughput, input demand, and output volume at every step. With process basis volumes established, equipment can be sized, allowing capital and operating costs to be estimated.

\paragraph{Modeling of Operations}\leavevmode\par

The operations modeled in each production step are derived from standard Process Based Cost Modeling \citep{kirchainCostModelingMaterials2001, fuchsStrategicMaterialsSelection2008}. Given the target throughput, we must consider the time utilization of each machine. Some machines are sized for continuous throughput, whereas others are based on batch operations, in which case throughput can be calculated by time per batch. In either case, with target production set to an annual basis, we can calculate the required line time needed to meet a given production target. This quantity can be compared with line time available (LTA), which accounts for plant schedule and downtime, to determine how many machines to run simultaneously, or how much the continuous process needs to be oversized relative to the baseline.

In the present implementation, facility-level operating assumptions define the plant schedule:

\begin{description}[noitemsep, topsep=0pt]
    \item[\textbullet\enskip Working days per year] ($d_{\text{py}}$)
    \item[\textbullet\enskip Shifts per day] ($s_{\text{pd}}$)
    \item[\textbullet\enskip Hours per shift] ($h_{\text{ps}}$)
    \item[\textbullet\enskip Unpaid breaks] ($u_b$)
    \item[\textbullet\enskip Paid breaks] ($p_b$)
    \item[\textbullet\enskip Scheduled maintenance] ($s_{\text{cm}}$)
    \item[\textbullet\enskip Unplanned downtime] ($u_{\text{pd}}$)
\end{description}
From these, the model computes:
\begin{align}
    \text{PLT} &= d_{\text{py}} \times s_{\text{pd}} \times (h_{\text{ps}} - u_b) \\
    \text{WLT} &= d_{\text{py}} \times s_{\text{pd}} \times (h_{\text{ps}} - u_b - p_b)
\end{align}
where PLT denotes paid line time and WLT denotes worked line time. Each production step then computes its own LTA by adjusting worked line time for step-specific downtime:
\alteq{Line time available equals working line time multiplied by one minus unplanned downtime fraction minus scheduled changeover and maintenance fraction.}
\begin{equation}
    \text{LTA} = \text{WLT} \times (1 - u_{\text{pd}} - s_{\text{cm}})
\end{equation}
\alteqend
This separation preserves the process-based cost modelling distinction between paid labor time and productive equipment time, allowing labor cost and machine capacity to scale differently.

\noindent\textbf{Batch Processes.}\quad
For batch operations, throughput is governed by discrete processing cycles characterized by a base volume per batch ($V_{\text{base}}$), a batch cycle time ($t_{\text{cycle}}$), and a batch setup time ($t_{\text{setup}}$). The effective production rate per machine is:
\alteq{Base volume divided by the sum of cycle time and setup time.}
\begin{equation}
    \frac{V_{\text{base}}}{t_{\text{cycle}} + t_{\text{setup}}}
\end{equation}
\alteqend
Required line time (LTR) is thus:
\alteq{Line time required equals the step throughput Q-step divided by the base volume per cycle-plus-setup time.}
\begin{equation}
    \text{LTR} = \frac{Q_{\text{step}}}{V_{\text{base}} / (t_{\text{cycle}} + t_{\text{setup}})}
\end{equation}
\alteqend
and machines required are:
\alteq{Number of machines equals the ceiling of line time required divided by line time available if a dedicated line, or simply line time required divided by line time available otherwise.}
\begin{equation}
    N_{\text{machines}} = 
    \begin{cases}
        \lceil \text{LTR} / \text{LTA} \rceil & \text{if dedicated line} \\
        \text{LTR} / \text{LTA} & \text{otherwise}
    \end{cases}
\end{equation}
\alteqend
Labor requirements are scaled analogously: if labor is dedicated, headcount is rounded up to the nearest integer; otherwise, fractional allocation is permitted. Utilities (electricity, gas, water, steam, compressed air) scale linearly with required line time. Thus, in batch systems, capacity expansion occurs discretely via additional parallel lines, labor and utilities scale proportionally to processing time, and setup time explicitly penalizes low throughput operation.

\medskip
\noindent\textbf{Continuous Processes.}\quad

For continuous operations, throughput is parameterized relative to a reference design capacity ($V_{\text{base}}$). Rather than accumulating discrete cycle time, scaling is expressed as a volume ratio:
\alteq{Volume ratio r-vol equals the step throughput Q-step divided by the product of base design capacity V-base and line time available.}
\begin{equation}
    r_{\text{vol}} = \frac{Q_{\text{step}}}{V_{\text{base}} \times \text{LTA}}
\end{equation}
\alteqend
In contrast to batch operations, only a single line is assumed ($N_{machines} = 1$). Total scaled equipment cost scales with a power-law exponent ($\alpha_{\text{equip}}$), consistent with standard chemical engineering scaling laws. 

\alteq{Scaled equipment cost equals base equipment cost multiplied by the volume ratio raised to the equipment scaling exponent alpha-equip.}
\begin{equation}
    C_{\text{equip,scaled}} = C_{\text{equip,base}} \cdot r_{\text{vol}} ^ {\alpha_{\text{equip}}}
\end{equation}
\alteqend

Utilities (electricity, natural gas, diesel, etc.), measured in per-hour quantities, scale with the same throughput exponent ($\alpha_{\text{equip}}$), multiplied by the operating time ($LTA$).

\begin{equation}
    Q_u = Q_{u,\text{base}} \cdot r_{\text{vol}} ^ {\alpha_{\text{equip}}} \cdot \text{LTA}
\end{equation}

Labor ($L$) scales with its own exponent ($\alpha_{\text{labor}}$), with dedicated labor treated in the same manner as with batch processes. This formulation reflects the economic reality that continuous systems (e.g., leach trains, evaporators, large thickeners) are typically resized rather than duplicated, and exhibit sub-linear capital scaling consistent with the six-tenths rule \citep{petersPlantDesignEconomics1991}. 

\begin{equation}
    L = L_{base} \cdot r_{\text{vol}} ^ {\alpha_{\text{labor}}} \cdot \text{LTA}
\end{equation}

Unlike batch systems, setup time does not appear explicitly; instead, downtime and availability are incorporated directly into available operating time, i.e. the line time required is equal to the total line time available

\medskip
\noindent\textbf{Location-Based Operating Parameters}\quad

In this analysis, the United States (Nevada) and China (Jiangxi) share the same scheduling parameters: 350 days/year, three 8-hour shifts, 5\% unplanned downtime, and 8\% scheduled maintenance. This reflects the assumption that equipment-driven availability is primarily a function of process technology rather than location; the principal location-dependent differentiation is instead captured through factor costs (labor, utilities, capital) as described below. Actual plant availability data for Chinese lepidolite operations are not publicly available to ground a location-specific adjustment.

\paragraph{Modeling of Costs}\label{app:04-methods-system-costs}\leavevmode\par

Following operational scaling, costs are computed at the level of each Production Step and then aggregated to the facility level. The cost model preserves the canonical PBCM separation between variable costs, which scale with throughput, and fixed costs, which reflect annualized capital and overhead charges, but also allows for aggregation of costs across the more policy- and investment-relevant split of operational expenditures (OPEX) and capital expenditures (CAPEX).

For each step, total annual cost is decomposed as:
\begin{align}
    C_{\text{step}} = \, & C_{\text{materials}} + C_{\text{labor}} + C_{\text{utilities}} + 
    C_{\text{machine}} + C_{\text{tooling}} \nonumber \\
    & + C_{\text{building}} + C_{\text{aux}} + C_{\text{maintenance}} + C_{\text{overhead}}
\end{align}

These are grouped as follows:

\begin{description}[noitemsep, topsep=0pt]
    \item[\textbullet\enskip Variable]: Materials, Labor, Utilities
    \item[\textbullet\enskip Fixed]: Machine, Tooling, Buildings, Auxiliary Equipment, Maintenance, Overhead
    \item[\textbullet\enskip OPEX]: Materials, Labor, Utilities, Maintenance, Overhead
    \item[\textbullet\enskip CAPEX]: Machine, Tooling, Buildings, Auxiliary Equipment
\end{description}

\medskip
\noindent\textbf{Material Costs.}\quad

Material costs are derived from reagent usage determined during chemistry propagation. Annual reagent cost is computed from the usage quantity (Eq. \ref{eq:c-reagentmass}) multiplied by the unit cost:

\begin{align}
    C_r &= M_r \times \text{unit material cost}
\end{align}

Only secondary inputs (i.e., reagents introduced at the step) are assigned costs. Primary inputs are treated as upstream transfers and are costed in their originating steps or facilities. This structure ensures full traceability between stoichiometric assumptions, impurity removal efficiencies, yield losses, and economic impact.

\medskip
\noindent\textbf{Labor Costs.}\quad

Labor requirements are determined during operational scaling, as described above. Annual labor cost is then computed as:
\begin{equation}
    C_{\text{labor}} = w \times L \times \text{PLT}
\end{equation}
where $w$ is the wage (including benefits), $L$ is required labor headcount, and $PLT$ is paid line time per year. Labor cost is tied to paid time rather than productive time, preserving the PBCM-based distinction between compensated hours and machine availability.

\medskip
\noindent\textbf{Utility Costs.}\quad

Utility consumption (electricity, natural gas, process water, cooling water, steam, compressed air) is determined during operational scaling as described above. Utility cost is then:
\begin{equation}
    C_{\text{utility}} = \sum_{u} Q_u \cdot P_u
\end{equation}
where $Q_u$ is quantity consumed and $P_u$ is the facility- (location-) specific price, allowing explicit representation of location-dependent utility prices.

\medskip
\noindent\textbf{Machine and Tooling Capital Costs.}\quad

The number of machines is determined as described in the operational scaling section above. As production is modeled on an annual basis, we compute the annualized machine cost using a capital recovery factor ($CRF$), and a local machine cost factor $CF_{\text{local}}$ where appropriate:
\alteq{Capital recovery factor CRF as a function of discount rate r and lifetime n equals r times the quantity one plus r raised to the n, divided by the quantity one plus r raised to the n minus one.}
\begin{equation}
    \text{CRF}(r, n) = \frac{r(1+r)^n}{(1+r)^n - 1}
    \label{eq:c-crf}
\end{equation}
\alteqend
\begin{equation}
    C_{\text{machine,annual}} = N_\text{{machines}} \cdot \text{CRF} \cdot C_{\text{equip,scaled}} \cdot CF_{\text{local}}
\end{equation}
If machine lifetime differs from facility lifetime, a step-specific CRF is applied. The local machine cost factor accounts for differences in machine costs across locations. Tooling costs are defined and annualized analogously; if no separate lifetime is provided, tooling costs are recovered over the facility-level capital recovery period. For this analysis, inclusion of measurements of $CF_{\text{local}}$ are not yet made, which is an area for future work.

\medskip
\noindent\textbf{Building Cost.}\quad

Building cost scales with equipment footprint: in batch systems, footprint is proportional to number of machines; in continuous systems, footprint scales by a throughput exponent. Annualized building cost is:
\begin{equation}
    C_{\text{building}} = \text{BCRF} \cdot P_{\text{build}} \cdot A
\end{equation}
where $P_{\text{build}}$ is building cost per \si{m^2}, $A$ is required footprint, and BCRF is the building capital recovery factor, calculated in a similar manner to \ref{eq:c-crf}, but instead using a building lifetime (default 25 years) instead of $n$.

\medskip
\noindent\textbf{Auxiliary Equipment, Maintenance, Overhead, and Other Costs.}\quad

The cost of additional auxiliary equipment not included in the model blocks and maintenance are modeled as fractions of machine capital. For auxiliaries, since they are capital costs, we multiply by the annualized machine cost:
\begin{equation}
    C_{\text{aux}} = f_{\text{aux}} \cdot C_{\text{machine,annual}}
\end{equation}

In contrast to auxiliary equipment, since maintenance costs are often measured as a fraction of total installed machine costs~\cite{petersPlantDesignEconomics1991, ausimmMonograph27Cost}, we multiply here by the scaled machine costs:
\begin{equation}
    C_{\text{maint}} = f_{\text{maint}} \cdot C_{\text{equip,scaled}}
\end{equation}

Fixed overhead is applied to the sum of machine, tooling, building, auxiliary, and maintenance costs:
\begin{equation}
    C_{\text{overhead}} = f_{\text{over}} \cdot 
    (C_{\text{machine,annual}} + C_{\text{tool}} + C_{\text{building}} + C_{\text{aux}} + 
    C_{\text{maint}})
\end{equation}

Additionally, an optional additional catch-all OPEX adder is defined similar to the fixed overhead costs, to account for any further additional machine or process-specific costs. In this present analysis, this adder is set to zero for all machine blocks. 

\begin{equation}
    C_{\text{opex\_excess}} = f_{\text{opex\_excess}} \cdot 
    (C_{\text{machine,annual}} + C_{\text{tool}} + C_{\text{building}} + C_{\text{aux}} + 
    C_{\text{maint}})
\end{equation}

\medskip
\noindent\textbf{Location-Based Operating Costs}\quad

Several cost parameters used in this analysis above are location-dependent. The locational input data defines unit prices for each utility: electricity (\texttt{elec\_price}, \$/kWh), natural gas (\texttt{gas\_price}, \$/MMBTU), diesel (\texttt{diesel\_price}, \$/L), propane (\texttt{propane\_price}, \$/L), cooling water (\texttt{cool\_water\_price}, \$/\si{m^3}), steam (\texttt{steam\_price}, \$/kg), and compressed air (\texttt{comp\_air\_price}, \$/\si{Nm^3}). These are the $P_u$ terms in the utility cost equation. The direct labor wage including benefits ($w$, \$/hr) and building cost per unit area ($P_{\text{build}}$, \$\si{m^2}) are also location-specific.

Capital cost parameters set at the location level include the discount rate (\texttt{dr}), capital recovery period (\texttt{crp}), building recovery period (\texttt{brp}), and the machine cost factor (\texttt{machine\_cost\_factor}, i.e.\ $CF_{\text{local}}$). Fixed cost fractions for auxiliary equipment (\texttt{aux\_equip}), maintenance (\texttt{maint}), and overhead (\texttt{fixed\_over}) are also defined per location, as are tailings disposal costs for nonhazardous and hazardous waste streams.

The most consequential differences are in labor (\$50/hr in the U.S.\ vs.\ \$10/hr in China), electricity (\$0.09/kWh in the U.S.\ vs.\ \$0.075/kWh in China), and natural gas (\$5.04/MMBTU in the U.S.\ vs.\ \$8.50/MMBTU in China). The machine cost factor is set at 1.0 for the U.S.\ and 0.75 for China, reflecting lower equipment procurement costs in China, though this parameter is acknowledged as an area requiring further empirical grounding.

In addition to location-level parameters, reagent unit costs and embodied emissions factors are defined per project in the project-specific input data, with conservative, midpoint, and optimistic scenarios. These capture supply-chain-specific procurement conditions (e.g., sulfuric acid cost in Nevada vs.\ Jiangxi, soda ash sourced domestically vs.\ imported) that cannot be reduced to a single location profile. During scenario evaluation, the \texttt{update\_materials} routine applies the appropriate scenario rank to all facility material cost and impact entries, enabling the triangular sensitivity analysis described earlier.

\paragraph{Modeling of Emissions}\label{app:04-methods-system-emissions}\leavevmode\par

Environmental impacts are computed in parallel with financial costs in each Production Step. Where possible, physically modeled flows (utilities consumed and coproducts routed to sinks) are aggregated, and location-specific impact factors are applied to calculate their environmental impact, in line with standard methodologies from process-based, bottom-up studies in the Life Cycle Assessment (LCA) literature.

Environmental impacts are divided into two primary channels: utility-related impacts (energy and resource consumption), and sink-related impacts (coproducts and waste streams). For each step, impacts are reported as utility impacts (by utility type), sink impacts (by sink and coproduct), Scope 1 and Scope 2 totals, and total step impacts. Facility and supply chain totals are obtained by summing across all steps. Since these emissions are tightly coupled to operational modeling, we can not only observe the environmental impact of process assumptions and design choices, but also allow for relatively straightforward extensions to externality monetization and the broader connection between environmental and economic outcomes.

\medskip
\noindent\textbf{Coproduct and Waste Impacts (Scope 1).}\quad

The model tracks secondary outputs from each step, representing coproducts or waste streams routed to facility-defined sinks (e.g., landfill, wastewater treatment, air). Each secondary output includes a conversion factor from step production volume, an annual volume, and an assigned sink. During volume throughput calculations, coproduct volumes are computed consistently with yield and throughput.

Impact factors for sinks can be defined in two ways: sink-level factors (uniform across all coproducts), or coproduct-specific factors within a sink, allowing differentiated treatment (e.g., solid tailings vs.\ brine discharge to wastewater). For a coproduct $c$ routed to sink $s$, impacts are computed as:
\begin{equation}
    I_{c,j} = V_c \cdot f_{s,c,j}
\end{equation}
where $V_c$ is annual coproduct volume and $f_{s,c,j}$ is the impact factor for that sink-coproduct pair and impact category $j$. Tailings sinks carry a factor built from the electricity the sink draws: 
\begin{equation}
    f_{s,j} = e_s \cdot \rho_{\mathrm{DS}}(s_{\mathrm{in}}) \cdot \phi_j
\end{equation}
with $e_s$ the sink's specific energy in \si{kWh/t} dry solids, $\rho_{\mathrm{DS}}$ the dry solids per cubic meter at the registration solids fraction, and $\phi_j$ the facility's grid intensity in category $j$. Because this draw is purchased electricity, it is booked to Scope 2 rather than Scope 1. The remaining sink terms are undefined and assumed to evaluate to zero, such as diesel, off-site haulage and treatment at gated sinks, impoundment construction, and land use. Given these waste emissions are expected to be small relative to the Scope 2 and Scope 3 terms that dominate all three pathways, this omission is unlikely to affect the comparative results materially, but represents an area for future refinement. 

\medskip
\noindent\textbf{Utility-Related Impacts (Scope 1 and 2).}\quad

Utility consumption is computed during operational scaling, as described above, yielding annual quantities of electricity (kWh), natural gas (MMBTU), propane (L), diesel (L), cooling water (\si{m^3}), steam (kg), and compressed air (\si{Nm^3}). Each utility $u$ is associated with an impact factor vector $\{f_{u,j}\}$, where $j$ indexes environmental impact categories (e.g., \ce{CO2}, water use, \ce{NOx}, \ce{SO2}). While most of these impacts occur on-site (e.g. combustion emissions from diesel, natural gas, propane) and thus are considered Scope 1, usage of offsite electricity is considered Scope 2. In this analysis, we primarily consider greenhouse gas impacts, though the framework is easily extendable to other impact categories. Utility-related impacts for a step are computed as:
\begin{equation}
    I_{u,j} = Q_u \cdot e_{u,j}
\end{equation}
where $Q_u$ is annual utility consumption and $e_{u,j}$ is the impact factor for category $j$. Impact factors for each of these utilities are specified in \ref{tab:c-emission-factors}. 

\medskip
\noindent\textbf{Upstream Material Use Impacts (Scope 3).}\quad

Secondary inputs (reagents) consumed in each production step carry embodied emissions from their own upstream production and transport. For each reagent $r$ consumed at step $i$, the Scope 3 impact is computed as:

\begin{equation}
    I_{r,j} = M_r \cdot e_{r,j}
\end{equation}

where $M_r$ is the annual mass consumption of reagent $r$ (Eq. \ref{eq:c-reagentmass}), and $e_{r,j}$ is the cradle-to-gate emissions factor for reagent $r$ and impact category $j$ (e.g., kg \ce{CO2eq}/kg reagent).

Reagents are treated as purchased commodity inputs unless otherwise specified; their embodied emissions reflect the upstream production processes of the dominant manufacturing route in the relevant geography. For example, soda ash (\ce{Na2CO3}) consumed at U.S. facilities is assigned an emissions factor reflecting trona mining and calcination ($\sim$ 0.3--0.7 t \ce{CO2eq}/t \ce{Na2CO3}), while soda ash consumed at Chinese facilities is assigned a higher factor reflecting the energy-intensive Solvay process ($\sim$0.7--1.2 t \ce{CO2eq}/t \ce{Na2CO3})~\cite{wangGreenhouseGasesRegulated2024, scottCanSyntheticSoda2023, royLithiumClayAssessing2024}. Soda ash is an intentionally produced commodity chemical rather than a waste product of another process; its emissions factor therefore reflects the full cradle-to-gate burden of dedicated soda ash manufacturing. Similarly, lime (\ce{CaO}) is assigned an emissions factor reflecting limestone calcination, where approximately two-thirds of emissions are unavoidable process emissions from carbonate decomposition and the remainder reflects kiln fuel combustion~\cite{lavegliaDecarbonizingLimebasedConstruction2023}. Emissions factors for all reagents are sourced from published life cycle inventory databases and literature, and are summarized below in \ref{tab:c-emission-factors}.

\medskip
\noindent\textbf{Defined Emissions Impacts.}\quad

A table summarizing all emissions factors, by location and material, with parametric variation where possible, is found below in \ref{tab:c-emission-factors}. This table includes emission intensities of chemical reagent and energy inputs (kg\,\ce{CO2eq}\,kg$^{-1}$ for solids/liquids; kg\,\ce{CO2eq}\,kWh$^{-1}$ for electricity; kg\,\ce{CO2eq}\,MMBtu$^{-1}$ for natural gas; kg\,\ce{CO2eq}\,m$^{-3}$ for water; kg\,\ce{CO2eq}\,kg$^{-1}$ for steam; kg\,\ce{CO2eq}\,L$^{-1}$ for liquid fuels) under optimistic, midpoint, and conservative scenarios. For soda ash, separate rows reflect U.S. trona-based production (applied to Nevada scenarios) and Chinese Solvay-process production (applied to Jiangxi scenarios); electricity grid intensity is location-specific; all other inputs adopt world-average supply-chain factors. Inputs lacking scenario variation use a single point estimate spanning all three columns. Superscripts correspond to the notes below the table.

\begin{table}[htbp]
\caption{Emission intensities of chemical reagent and energy inputs}
\label{tab:c-emission-factors}
\begin{ThreePartTable}
  \footnotesize
  \setlength{\tabcolsep}{6pt}
  \begin{tabular*}{\textwidth}{@{\extracolsep{\fill}} l l r r r @{}}
    \toprule
    \multirow{2}{*}[-3pt]{Input}
      & \multirow{2}{*}[-3pt]{\makecell[l]{Geographic\\scope}}
      & \multicolumn{3}{c}{Emission intensity} \\
    \cmidrule(l){3-5}
      & & \multicolumn{1}{c}{Optimistic}
        & \multicolumn{1}{c}{Midpoint}
        & \multicolumn{1}{c}{Conservative} \\
    \midrule
    %%------------------------------------------------------------
    %% BLOCK 1: Chemical reagents with opt/mid/con variation
    %%------------------------------------------------------------
    \multicolumn{5}{@{}l}{\textit{Chemical reagents --- with scenario range}} \\[2pt]
    Flocculant (generic)\tnote{$\dagger$}
      & World
      & 2.85\tnote{a}    & 3.914\tnote{a}   & 5.38\tnote{a}   \\[2pt]
    Lime
      & World
      & 0.83\tnote{b}    & 1.20\tnote{a}    & 1.70\tnote{a,e} \\[2pt]
    Limestone
      & World
      & 0.00225\tnote{a} & 0.004\tnote{a}   & 0.006\tnote{b}  \\[2pt]
    \multirow{2}{*}{Soda ash}
      & USA
      & 0.30\tnote{c}    & 0.50\tnote{c}    & 0.70\tnote{c}   \\
      & China
      & 0.70\tnote{c}    & 1.20\tnote{a}    & 1.80\tnote{e}   \\[2pt]
    Sodium sulfate
      & World
      & 0.10\tnote{a}    & 0.25\tnote{a}    & 0.60\tnote{f}   \\[2pt]
    Sulfuric acid
      & World
      & 0.0509\tnote{a}  & 0.0969\tnote{a}  & 0.14\tnote{b}   \\[6pt]
    Sulfur
      & World
      & 0.16\tnote{g}     & 0.30\tnote{g}     & 0.43\tnote{g}   \\[2pt]
    %%------------------------------------------------------------
    %% BLOCK 2: Chemical reagents — single point estimate
    %%------------------------------------------------------------
    \multicolumn{5}{@{}l}{\textit{Chemical reagents --- single point estimate}} \\[2pt]
    Explosives
      & World
      & \multicolumn{3}{c}{2.5\tnote{b}}                       \\[2pt]
    Filter aid
      & World
      & \multicolumn{3}{c}{0.4\tnote{a}}                       \\[2pt]
    Wash / process water
      & World
      & \multicolumn{3}{c}{0}                                   \\[6pt]
    %%------------------------------------------------------------
    %% BLOCK 3: Energy and utility inputs
    %%------------------------------------------------------------
    \multicolumn{5}{@{}l}{\textit{Energy and utility inputs}} \\[2pt]
    \multirow{2}{*}{Electricity grid}
      & Nevada, USA    & \multicolumn{3}{c}{0.1932\tnote{h}}   \\
      & Jiangxi, China & \multicolumn{3}{c}{0.5836\tnote{i}}   \\[2pt]
    Natural gas
      & World
      & \multicolumn{3}{c}{53.06\tnote{d}\textsuperscript{,\,}%
        \raisebox{0.5ex}{\scriptsize$\star$}}                   \\[2pt]
    Steam
      & World
      & \multicolumn{3}{c}{0.15\tnote{j}}                      \\[2pt]
    Diesel
      & World
      & \multicolumn{3}{c}{2.68\tnote{d}}                      \\[2pt]
    Propane
      & World
      & \multicolumn{3}{c}{1.51\tnote{d}}                      \\
    \bottomrule
  \end{tabular*}
  \begin{tablenotes}[flushleft]
    \footnotesize
    \item[\textit{a}]\hspace{0.05cm} ecoInvent 3.8, as compiled in \citet{royLithiumClayAssessing2024}.
    \item[\textit{b}]\hspace{0.05cm} GREET model (Argonne National Laboratory) \cite{wangGreenhouseGasesRegulated2024}.
    \item[\textit{c}]\hspace{0.05cm} U.S. soda ash production is primarily trona-based;
          reported emission intensities span
          0.3--0.7 kg \ce{CO2eq} kg$^{-1}$ \citep{scottCanSyntheticSoda2023}.
          The China optimistic value (0.70) equals the upper end of the
          U.S. trona range, used as a conservative floor for
          Solvay-process production.
    \item[\textit{d}]\hspace{0.05cm} EPA Center for Corporate Climate Leadership,
          \textit{GHG Emission Factors Hub}.
          Natural gas: \ce{CO2} only (53.06 kg \ce{CO2} MMBtu$^{-1}$);
          diesel: 10.21 kg \ce{CO2} gal$^{-1}$ converted to
          2.68 kg \ce{CO2eq} L$^{-1}$;
          propane: 5.72 kg \ce{CO2} gal$^{-1}$ converted to
          1.51 kg \ce{CO2eq} L$^{-1}$ \citep{usepaEmissionFactorsGreenhouse2025}.
    \item[\textit{e}]\hspace{0.05cm} Estimated conservative upper bound; not directly reported in a single source.
    \item[\textit{f}]\hspace{0.05cm} Mannheim-process production; the high-temperature, fuel-oil-fired kiln step is identified as the dominant direct emission source \citep{farahaniAssessingEnergyLoad2024}.
    \item[\textit{g}]\hspace{0.05cm} Elemental sulfur is recovered as a byproduct of petroleum refining and gas sweetening, so its cradle-to-gate factor is governed by the allocation rule applied upstream rather than by measurement. Values are derived from this table's sulfuric acid factors at the stoichiometric sulfur-to-acid mass ratio (0.327), and are corroborated by a published cradle-to-gate factor of 0.18 kg \ce{CO2eq}\,kg$^{-1}$ \citep{bafaInformationsblattCO2FaktorenVersion2026}. 
    \item[\textit{h}]\hspace{0.05cm} EPA eGRID summary data for the Nevada grid
          subregion \citep{usepaSummaryData2025}.
    \item[\textit{i}]\hspace{0.05cm} China Ministry of Ecology and Environment,
          provincial grid emission factor for Jiangxi \citep{ministryofecologyandenvironmentGuanYuFaBu2023NianDianLiErYangHuaTanPaiFangYinZiDeGongGaoAnnouncementRelease2025}.
    \item[\textit{j}]\hspace{0.05cm} Raw assumption; no published source.
    \item[$\dagger$] 
          $\star$ Natural gas factor is \ce{CO2} only; \ce{CH4} and
          \ce{N2O} contributions are negligible for combustion at this scale but could be added if needed.
  \end{tablenotes}
\end{ThreePartTable}
\end{table}

\paragraph{Modeling of Water}\label{app:04-methods-system-water}\leavevmode\par

Water is a key input for mining operations, serving as a key constraint on operations and flowsheet design, regulatory compliance, and local social license to operate~\cite{schombergExtendedLifeCycle2021, riofrancosSecuritySustainabilityNexus2023}. In our model, it is not only an operating cost, both as an input but also as an output that needs to be treated as a discharge, but also a heavy determinant of flowsheet designs, given that most chemical processing steps operate at differing solids fractions, with differing water purchase needs and dewatering of waste streams for disposal. Furthermore, water usage represents the sharpest point of divergence among the three modeled projects in regulatory terms, with Thacker Pass under a zero-discharge water pollution control permit and Silver Peak returning its streams to the pond system, while Jianxiawo carries no equivalent requirements. Because rigorous site-wide closed water balances rely on highly variable, site-specific parameters that cannot be uniformly calibrated across all three facilities, we do not attempt a strict closed balance. Water is instead tracked at the three points where it carries a price or a consequence: where it is purchased as make-up, where it is recovered and returned to a circuit, and where it leaves the flowsheet as a coproduct routed to a sink. 

Water consumed at a step is dosed as a reagent against the step basis and split into purchased \texttt{process\_water} and \texttt{recycled\_process\_water}, the latter carried at zero cost. The split is set by the circuit's physical losses rather than assumed: for a circuit of gross water demand $G$ per unit of step basis,
\begin{equation}
    M = W_{\text{imp}} + W_{\text{cake}} + W_{\text{bleed}} + W_{\text{mill}}, \qquad R = G - M
\end{equation}

\noindent where $M$ is the purchased make-up, $R$ the recycled draw, and the loss terms are the water retained in tailings at their as-placed solids fraction, the water leaving in a filter cake, the deliberate circuit bleed, and mill discharge. The returns available to $R$ are the thickener overflow and the water removed by the dewatering trains at the tailings sinks, $\Delta W_k = 1/s_{k-1} - 1/s_k$ per tonne of dry solids as defined in \ref{app:04-methods-facilitychoices-tailings}.

\medskip
\noindent\textbf{Wastewater Treatment.}\quad

Coproducts that are neither returned to a circuit nor placed with the solids report to a wastewater sink, priced per cubic meter of stream received:
\begin{equation}
    C_{\text{ww}} = \sum_{c} V_{c} \, p_{j(c)}
\end{equation}
where $V_c$ is the annual routed volume and $p_j$ the unit price in jurisdiction $j$. Following the tailings sink convention, each sink is specified at one location jurisdiction. Both wastewater sinks are set at \$1/\si{\cubic\meter}. This approximation is assumed for Jianxiawo based on reported costs for wastewater treatment at Chinese industrial parks, which span roughly that order of magnitude~\cite{huStudyDevelopmentPerformance2019, lyuExploringCostWastewater2020}. Given no wastewater is modeled in Nevada for its two projects, the assumed price is untested; it is retained so any possible non-zeroed routing is priced rather than silently marked as free.

\begin{table}[htbp]
\centering
\small
\caption[Coproduct streams routed to the \texttt{wastewater\_discharge} sinks.]{Coproduct streams routed to the \texttt{wastewater\_discharge} sink across all modeled pathways. All routings are within the Jianxiawo model; the Nevada projects route no volume to these sinks.}
\label{tab:wastewater-routings}
\begin{tabular}{llll}
\toprule
Pathway sheet & Step & Coproduct & Project(s) \\
\midrule
Solution Processing--Sulfate & \texttt{thicken\_dewater} & \texttt{pressed\_water}      & Jianxiawo \\
Solution Processing--Sulfate & \texttt{wash\_purif\_step}& \texttt{wash\_water}         & Jianxiawo \\
Sulfate Roast--Water Leach   & \texttt{dewatering}       & \texttt{press\_filtrate}     & Jianxiawo \\
Solution Processing--Sulfate & \texttt{precipitation\_step} & \texttt{mother\_liquor}      & Jianxiawo \\
\bottomrule
\end{tabular}

\vspace{2pt}
\end{table}

\paragraph{Machine Model Blocks}\label{app:04-methods-system-blocks}\leavevmode\par

Each Production Step is further parameterized using one of several possible machine blocks, to represent alternative equipment configurations or operational archetypes. While reagents and coproducts are handled at the Production Step level, machine-specific characteristics are aggregated into blocks which can then be scaled as appropriate, allowing for heterogeneity in operational performance while preserving consistent chemistry and throughput logic at the step level. As a result, all operational variables other than materials --- labor, utilities, machine, tooling, building, auxiliaries, maintenance, and overhead --- are handled by machine model blocks. By default, each step uses a specified \texttt{machine\_default} block, but alternative blocks may be substituted to represent technology variants or scenario analyses (e.g., tornado analysis).

\paragraph{Facility Aggregation}\label{app:04-methods-system-facilities}\leavevmode\par

A facility aggregates multiple production steps in series or parallel configuration. Intermediate material transfers within a facility are assumed to occur without modeled transportation losses unless explicitly specified (e.g., evaporation ponds, pumping). Facility-level metrics (CAPEX, OPEX, power consumption, water use, emissions) are calculated as the sum of constituent production steps.

\paragraph{Transportation Modeling}\label{app:04-methods-system-transport}\leavevmode\par

Transportation legs connect facilities across the supply chain; multiple legs can be combined into a route to account for multi-modal transport (e.g., spodumene shipped via rail, then barge, then truck). For each leg, transported mass or volume equals the primary output stream quantity. With the exception of international marine freight, transportation is assumed to be round-trip, accounting for empty return:
\begin{equation}
    \text{Total Transport} = 2 \times \text{Distance} \times \text{Payload}
\end{equation}
Energy use and emissions are derived from transport work and mode-specific intensity factors. Transportation costs and emissions are roughly estimated from road freight literature \cite{mulhollandCO2EmissionsTrucks2023} rather than mine haul references. However, given mine-site haul distances of less than 10 km at both Thacker Pass and Jianxiawo, haulage contributes negligibly to total modeled cost and emissions. 

\paragraph{Supply Chain Integration}\label{app:04-methods-system-sc}\leavevmode\par

The supply chain model links facilities and transportation legs into a directed network from extraction to refined product delivery.

\clearpage
\subsubsection{Facility-level Modeling Choices}\label{app:04-methods-facilitychoices}

Review of publicly available NI 43-101 and S-K 1300 technical reports and feasibility studies reveals that lithium projects are typically reported as a sequence of high-level process areas (`flowsheets') rather than individual machines. These areas are often labeled inconsistently across projects, even within the same resource pathways, presenting a distinct modeling challenge. Directly reproducing project-specific structures would limit cross-project comparability, while fully disaggregating to machine-level details for each project would require a level of engineering and economic understanding beyond what is presented in such disclosures. Given the supply chain, facility, production step, and machine block representations described above, we leverage our understanding of each pathway to create common representations of pathway segments, aggregating commonly sequenced production steps into facility representations.

Each of the pathway segment descriptions below represents a modeled facility in our supply chains. Design decisions are described at this hierarchical level, and the blocks necessary to compose each facility or pathway segment are specified. Engineering details, capital cost aggregations, and reagent usage are described within each of the production step and machine block models themselves.

\paragraph{Pumping and Evaporation Ponds (Generic Brine Pathway)}\label{app:04-methods-facilitychoices-pumpponds}\leavevmode\par

This is the generic pathway for solar evaporation-based processes. The pumping section is also relevant as the initial step for direct lithium extraction (DLE).

Lithium concentrations in brines are typically on the order of hundreds to a few thousand ppm, meaning that millions of cubic meters of brine must be extracted, transported, and concentrated to produce a relatively small mass of final product. As a result, hydraulic transport and surface area requirements are the primary determinants of capital and operating costs. Different projects have differing hydrogeologic (e.g.\ brine source depths, rock porosity) and climatic parameters (e.g.\ evaporation rate) that can be further developed as needed, so this model aims to preserve sensitivity to these variables while remaining sufficiently transferable across both traditional salar-based production and DLE-based systems.

This facility is composed of the following Production Steps:

\begin{description}[noitemsep, topsep=0pt]
    \item[\textbullet\enskip\texttt{Pumping:}] Extraction and hydraulic transport of lithium brines. 
    \item[\textbullet\enskip\texttt{Evaporation:}] Solar concentration of extracted brine.
\end{description}

\medskip
\noindent\textbf{Pumping.}\quad

The pumping step extracts lithium-bearing brine from subsurface aquifers and delivers it to the evaporation pond system. The primary input is raw brine in the ground (\texttt{raw\_brine\_from\_ground}), expressed in liters per year, which passes through the step without compositional transformation. As a result, the step basis is brine pump throughput, with an input-to-process conversion factor (IPCF) of 1.0. With a process-to-output conversion factor (POCF) of 1.0, it exits as pumped brine (\texttt{pumped\_brine}) of identical volume and chemistry. The step therefore functions as a pure transport and energy accounting block: its primary purpose in the model is to assign pumping-related capital and electricity costs and emissions to the extraction stage, rather than to modify the brine composition.

The step is parameterized by the full multi-ion brine chemistry at the point of extraction: constituents Li, K, Mg, Ca, Na, B, Cl, and \ce{SO4} are tracked, all expressed in ppm (mg/kg brine). These concentrations are set from the project-specific data, and for Silver Peak the dilute feed brine is specified at approximately 90 ppm Li, 190 ppm Mg, 450 ppm Ca, and 49,610 ppm Na, as identified in \ref{app:04-methods-modelblocks-oreconc}. The brine density is specified as 1.058 \si{g/cm^3} and is used as the constituent conversion factor (CCF), converting from ppm (kg constituent per kg brine) to the volumetric step basis (liters pumped per year). Because the output chemistry depends on the input chemistry - the pumped brine inherits the composition of the feed - the \texttt{primary\_output\_1\_chemistry\_dependence} flag is set to \texttt{True}, ensuring that any upstream modifications to feed composition (e.g., sensitivity analysis on brine grade) propagate correctly to downstream steps.

No reagents are consumed at this step, and no coproducts are generated. Electricity consumption for pumping is the sole utility input and is the dominant cost and emissions driver of this step; its magnitude is set by the machine block parameters rather than by the process chemistry (see the \texttt{brine\_pump} model block, \ref{app:04-methods-blockdefs-brinepump}). 

\medskip
\noindent\textbf{Evaporation.}\quad

The evaporation step concentrates the pumped brine through solar evaporation until it reaches the target strong brine chemistry required for downstream chemical purification. The process basis is pond area in \si{m^2}, and the step is constituent-driven: the volume-defining basis is the Li constituent of the input brine stream, meaning that the required pond area scales with the lithium mass that must be processed to meet the annual production target, rather than with bulk brine volume directly. This is the correct formulation because lower-grade brines require proportionally larger pond areas to yield the same mass of lithium product.

The CCF for the step basis converts from ppm Li (mg Li per kg brine) to \si{m^2} of pond area, and is defined as brine density multiplied by evaporation rate: for Silver Peak, with an evaporation rate of 2.44~mm/day and brine density of 1.058 \si{g/cm^3}, this yields a CCF of 9.4225 $\times 10^{-4}$ kg brine/(\si{m^2} $\cdot$ year) per ppm. Pond area follows as $A = \dot m_{\text{Li}}/(c \cdot \text{CCF})$, so required area scales inversely with evaporation rate. The output is concentrated lithium brine (\texttt{concentrated\_lithium\_brine}), whose target chemistry is specified in the project-specific data as the strong brine composition at the point of handover to downstream batch chemical purification. For Silver Peak, this is approximately 6,750 ppm Li, 14,250 ppm Mg, and 140,000 ppm Na at a density of 1.25 \si{g/cm^3}~\cite{srkconsultingSECTechnicalReport2023a}. The difference between input and output Li concentrations (90 ppm vs. 6,750 ppm) implies over 99\% evaporative water loss. The pond recovery efficiency parameter (0.51 for Silver Peak~\cite{srkconsultingSECTechnicalReport2023a}) accounts for lithium losses during evaporation attributable to co-precipitation and brine seepage.

Evaporated water is tracked as a coproduct routed to the \texttt{Air} sink, representing water vapor loss to the atmosphere. No reagents are consumed at this step. Unlike the pumping step, the chemistry of the output is not flagged as depending on the input chemistry, since the output concentration is set directly from the project-specific intermediate concentration target rather than computed from a reaction stoichiometry; this means that the model treats the evaporation step as a concentration target to be reached rather than a direct chemical transformation to be simulated. Capital and operating costs are driven by pond area and are parameterized through the \texttt{evaporation\_pond} model block (Appendix \ref{app:04-methods-blockdefs-evappond}).

\paragraph{Blasting and Excavating}\label{app:04-methods-facilitychoices-blastexcav}\leavevmode\par

This is the generic pathway for mining in the traditional sense, applied to ore-based extraction of spodumene, lepidolite, and clays. Physical fragmentation and movement of ore (and surrounding material) is strongly dependent on deposit characteristics: hard rock ores (spodumene, lepidolite) typically require drilling and blasting to generate moveable useful material, whereas clay deposits are often weakly bound and may be excavated without blasting. At the system level, blasting and excavation are aggregated to capture the primary energy and cost drivers associated with rock breaking and material movement, without specifying detailed mine parameters such as bench height, block size, or exact fleet composition. Hauling is treated as a project-specific transportation step connecting the mine site to the initial processing plant.

This facility is composed of the following blocks:

\begin{description}[noitemsep, topsep=0pt]
    \item[\textbullet\enskip\texttt{Blasting}] (optional): Drilling and blasting for hard rock fragmentation.
    \item[\textbullet\enskip\texttt{Excavation/Loading:}] Excavation and loading of fragmented or unconsolidated material.
\end{description}

\medskip
\noindent\textbf{Blasting.}\quad
The blasting step represents drilling and explosive fragmentation of in-situ rock prior to excavation. The primary input is \texttt{raw\_material} (in-situ rock, in metric tonnes), with an IPCF of 1.0, and the step basis is \texttt{rock\_blast\_throughput} in metric tonnes. The step passes material through without compositional transformation; output chemistry is flagged as dependent on input chemistry, propagating the initial ore grade set from the project-specific data. The CCF is 0.001 kg ore per (tonne-ppm), the standard conversion for ppm-in-mass-basis constituents to absolute kg per tonne.

The sole reagent is explosives, dosed against the step basis (total rock moved, including waste) using the powder factor drawn from the project-specific data --- for Jianxiawo, 0.25 kg explosives per tonne rock~\cite{darlingSMEMiningEngineering2011}; for Thacker Pass, 0.10 kg per tonne~\cite{sgscanadaincSK1300Technical2024}. Explosives are not a constituent-elimination reagent; their function is fragmentation, not chemical removal of any tracked species. No coproducts are generated at this step. The blasting step is bypassed entirely when the project-specific powder factor is absent or marked as `--', as is the case for projects where unconsolidated material requires no blasting (the no-blasting path routes directly to \texttt{excav\_load}).

The initial ore chemistry is set in \texttt{run\_scenarios.py} by loading the project-specific initial concentration and applying a dilution correction before passing it to the blasting step. If the reported grade is a ROM (run-of-mine) grade, the in-situ grade is back-calculated as $c_{\text{in-situ}} = c_{\text{ROM}} / (1 - d)$, where $d$ is the dilution factor; if the reported grade is an in-situ grade, the ROM grade is computed as $c_{\text{ROM}} = c_{\text{in-situ}} \times (1 - d)$. This ensures that the chemistry entering the blasting step reflects true in-situ conditions, while the chemistry exiting the excavation step reflects post-dilution ROM conditions as they would be observed at the processing plant gate.

\medskip
\noindent\textbf{Excavation/Loading.}\quad

The excavation and loading step represents mechanical fragmentation, loading, and preparation of blasted (or directly excavated) material for haulage to the processing facility. The primary input is \texttt{blasted\_rock} (from the blasting step) or \texttt{raw\_material} directly (in the no-blasting case), in metric tonnes. The step basis is \texttt{loose\_rock\_throughput} in tonnes of loose rock, with IPCF of 1.0. The output is \texttt{rom\_ore\_feed}, expressed in cubic meters of loose ore for haulage; it represents both the fraction of moved material that constitutes ore rather than waste, and the conversion from the mass basis to the volumetric basis on which haulage is sized.

The POCF for the output and the waste coproduct conversion factor are not fixed, but are set programmatically from project-specific mining factors:
\begin{align}
	\rho_{\text{loose}} &= \frac{\rho_{\text{bank}}}{S} \\
	\text{POCF}_{\text{rom}} &= \frac{1}{(1 + r_{\text{strip}})\,\rho_{\text{loose}}} \\
	f_{\text{waste}} &= \frac{r_{\text{strip}} + l}{1 + r_{\text{strip}}}
\end{align}
where $r_{\text{strip}}$ is the strip ratio (tonnes waste rock per tonne ore), $l$ is the loss fraction, $\rho_{\text{bank}}$ is the in-situ bulk density of the deposit, and $S$ is the bulking (swell) factor relating loose volume to in-situ volume. The POCF performs two roles at once: it selects the fraction of total material moved that becomes usable ore after accounting for dilution by accidental waste inclusion, and it converts that ore from the mass basis to the loose volume on which haulage and downstream crushing are sized. Thus, the POCF carries units of $\si{m^3}$ of loose ore per tonne of material moved. The waste remainder, $f_{\text{waste}}$, is routed as a coproduct to the \texttt{waste\_pile} sink on the (unchanged) tonnage basis. This sink carries an estimated cost of \$1.50 per tonne of waste, covering haulage
to the storage facility, dumping and reclamation, calculated from the Thacker Pass disclosure's waste-and-topsoil mining cost of \$2.71/t 
\cite{sgscanadaincSK1300Technical2024} less the estimated cost of drill and blast mining calculated by this model (\$1.24/t). The same cost is applied Jianxiawo as an engineering assumption, since no equivalent cost was found for Chinese mines. 

The yield rate on \texttt{rom\_ore\_feed} is set to $(1-l)$, capturing the ore fraction lost during the excavation and loading process independently of the strip ratio. Because $\rho_{\text{loose}}$ appears in the POCF here and again as the crushing IPCF with the opposite dimension, it cancels from the material balance, with our swell factor only governing haulage volume, not mine tonnage.

Swell factors are taken as $S = 1.20$ for Thacker Pass claystone and $S = 1.75$ for Jianxiawo pegmatite-hosted lepidolite, corresponding to published bulking factors of 20--40\% for clay and 75--80\% for granite~\cite{theengineeringtoolboxSoilRockBulk2009}. With $\rho_{\text{bank}} = 1.9$ $\si{t/m^3}$, the Thacker Pass loose density is $1.9/1.20 = 1.58$ $\si{t/m^3}$, consistent with the ore bulk density for transport of 1.6 $\si{t/m^3}$ given in the 2022 feasibility disclosure~\cite{rothFeasibilityStudyNational2022}.

The ROM ore chemistry (\texttt{rom\_comp}) is set as a target composition on the \texttt{rom\_ore\_feed} output directly, reflecting post-dilution grade. For projects where ROM grade is reported directly (\texttt{ROM feed = True}), the ROM composition is used as-is; for projects where in-situ grade is reported, the dilution correction described above is applied. The step CCF is $0.001$ kg ore per (tonne-ppm), consistent with blasting. No reagents are consumed at this step. This production step uses the \texttt{excavation\_loading} machine block (Appendix \ref{app:04-methods-blockdefs-excavload}).

\paragraph{Batch Chemical Purification (Silver Peak Method)}\label{app:04-methods-facilitychoices-batchchem}\leavevmode\par

This lithium production pathway is noted for its relative simplicity, being the first developed for lithium extraction from brines; production at Silver Peak, Nevada began in 1966. The method operates primarily on brines with low Mg:Li mass ratios, given the relatively high selectivity of magnesium for carbonate ions. Magnesium, calcium, and other impurities are removed as part of earlier concentration steps, either through selective precipitation during solar evaporation or by chemical treatment with lime and soda ash prior to lithium carbonate precipitation.

Six studies model the batch chemical purification method \citep{ambroseUnderstandingFutureLithium2020b, wesselkaemperOptimizationProcessingCosts2025, fengEnvironmentalImpactsLithium2025, sunRobustAssessmentsLithium2025, mousavinezhadNextLithiumBoom2025, schenkerCurrentFutureImpacts2025}, in which batches of slaked lime (\ce{Ca(OH)2}) and sodium carbonate (\ce{Na2CO3}) are added to concentrated lithium brines to precipitate lithium carbonate. \citet{mousavinezhadNextLithiumBoom2025} identify the Silver Peak process as using a resin-based ion exchange process, likely following the model structure of \citet{schenkerCurrentFutureImpacts2025}. While the latter do not specify what they mean by `ion exchange' in their model, no mention of resin-based ion exchange methods appears in the other studies or in the company's NI 43-101 disclosure itself. As a result, our model is based primarily on the reported steps in the Silver Peak disclosure and the model structure reported in \citet{wesselkaemperOptimizationProcessingCosts2025}, and is composed of the following blocks: 

\begin{description}[noitemsep, topsep=0pt]
    \item[\textbullet\enskip\texttt{batch\_treatment:}] Batch treatment of concentrated brine with lime and 
    soda ash.
    \item[\textbullet\enskip\texttt{sls\_step:}] Solid--liquid separation and filtration.
\end{description}

\medskip
\noindent\textbf{Batch Treatment.}\quad

The batch treatment step precipitates dissolved magnesium and calcium impurities from concentrated lithium brine, leaving lithium in solution. The primary input is \texttt{concentrated\_lithium\_brine} from the evaporation step, with an IPCF of 1.0. The step basis is \texttt{concentrated\_brine\_throughput} in liters per batch, with a CCF of $1.25 \times 10^{-6}$ kg/L reflecting the density of concentrated brine and converting ppm constituent concentrations to absolute kg per liter. The output is \texttt{treated\_slurry}, passed to the separation step with a POCF of 1.0 (volume is conserved) and output chemistry dependent on input chemistry.

Three reagents are dosed against constituent targets:
\begin{description}[noitemsep, topsep=0pt]
	\item[\textbullet\enskip Slaked lime \ce{Ca(OH)2} targeting Mg] at a ratio of 3.658 kg \ce{Ca(OH)2} per kg Mg, with 95\% elimination of Mg. (Note this is dosed as 2.769 kg of quicklime that is slaked). This ratio is derived from the stoichiometric molecular weight ratio $\mathrm{MW}(\ce{Ca(OH)2})/\mathrm{MW}(\ce{Mg}) = 74.09/24.305$, multiplied by an operational excess factor of 1.2 to account for reaction inefficiencies. The key reaction is \ce{Mg^{2+} + Ca(OH)2 -> Mg(OH)2 v + Ca^{2+}}; precipitated \ce{Mg(OH)2} reports to the solids phase, while the calcium introduced by the lime remains in solution and must be subsequently removed.
	\item[\textbullet\enskip Soda ash \ce{Na2CO3} targeting Ca (calcium natively present in the brine)] at a ratio of 3.173 kg \ce{Na2CO3} per kg Ca natively present in the brine, with 90\% elimination of Ca. This is derived from MW(\ce{Na2CO3})/MW(Ca) $\times 1.2 = 105.988/40.078 \times 1.2$, the stoichiometric requirement plus one operational excess factor. The reaction is \ce{Ca^{2+} + Na2CO3 -> CaCO3 v + 2Na^+}.
	\item[\textbullet\enskip Soda ash \ce{Na2CO3} as a secondary Ca removal dose, dosed against Mg] at a ratio of 6.279 kg \ce{Na2CO3} per kg Mg, with 0\% elimination of Mg. Because the lime reaction exchanges one mole of dissolved \ce{Ca^{2+}} for every mole of \ce{Mg^{2+}} precipitated, the calcium burden created by liming scales with the magnesium in the feed rather than with the calcium in it, and is dosed against Mg as its stoichiometric proxy. This is derived from MW(\ce{Na2CO3})/MW(Mg) $\times 1.2 \times 1.2 = 105.988/24.305 \times 1.44$, the product of the soda ash requirement per kg Ca and the Ca introduced per kg Mg dosed with lime, each carrying one excess factor. This assumes the full lime dose, including its operational excess, reaches solution as \ce{Ca^{2+}}. 
\end{description}

No coproducts are generated at this step; precipitated solids remain suspended in the treated slurry and are separated in the subsequent step. This step uses the \texttt{batch\_treatment\_generic} machine block (Appendix \ref{app:04-methods-blockdefs-batchtreat}).

\medskip
\noindent\textbf{Separation.}\quad

The separation step clarifies the treated slurry by removing precipitated \ce{Mg(OH)2} and \ce{CaCO3} solids, yielding a purified lithium-bearing solution (\texttt{pls\_post\_polishing}) suitable for downstream lithium carbonate precipitation. The primary input is \texttt{treated\_slurry} in liters, converted to the step basis (\texttt{mixture\_being\_separated} in \si{m^3/h}) via an IPCF of $10^{-3}$ (L to \si{m^3}). The CCF is $1.25 \times 10^{-6}$ kg/L, consistent with the upstream step. The output \texttt{pls\_post\_polishing} has a POCF of 1000 (\si{m^3} to L) and a yield rate of 0.995, reflecting a 0.5\% lithium loss to solids during filtration.

One reagent is consumed: a filter aid (precoat, e.g.\ diatomaceous earth) dosed at 0.005 kg per \si{m^3} of treated slurry against the step basis. This value is a rough engineering estimate; a suitable published reference for diatomaceous earth dosing in mineral brine clarification has not been identified and represents a low-priority refinement.

The residue coproduct is routed to the \texttt{mgcake\_stack\_nv\_35} sink, representing an on-site filter cake at approximately 35\% solids fraction. The coproduct conversion factor of 0.2434 \si{m^3} wet solids per \si{m^3} treated slurry is derived from first-principles stoichiometry. The batch treatment reagent doses (3.658 kg \ce{Ca(OH)2} per kg Mg at 95\% removal; 3.173 kg \ce{Na2CO3} per kg Ca at 90\% removal) yield 40.6 g/L \ce{Mg(OH)2} and 74.4 g/L \ce{CaCO3} dry precipitate at the concentrated brine chemistry (Mg: 14,250 ppm; Ca: 1,688 ppm; $\rho$ = 1.25 kg/L), totaling 115.0 g/L dry solids. At a filter press cake of 35 wt\% solids and wet density of 1.35 \si{g/cm^3}, this yields 0.2434 \si{m^3} wet cake per \si{m^3} slurry. This estimate is cross-validated against Silver Peak's disclosed lime consumption of 1.3 t \ce{CaO} per tonne \ce{Li2CO3}. The residue composition is assumed to be as 25\% \ce{Mg(OH)2}, 25\% \ce{CaCO3}, and 50\% water by volume. We assume no volume loss here, as water lost is replenished from water from tailings slurry dewatering and recycling from downstream lithium carbonate production steps, though a full water balance is a key future area for improvement. 

This step uses the \texttt{thickener\_generic} machine block (Appendix \ref{app:04-methods-blockdefs-thickener}).

\paragraph{Clay Comminution and Acid Leach}\label{app:04-methods-facilitychoices-clayleach}\leavevmode\par

Clay-based lithium sources allow for relatively straightforward chemical extraction from fine-grained matrices, and comminution serves primarily to control input particle sizes. Even weakly bound sediments, as in projects such as Rhyolite Ridge, Tonopah Flats, and Bonnie Claire, are typically delivered from the mine as run-of-mine fragments and therefore require simple primary crushing for size control prior to slurry preparation \citep{khakmardanLifeCycleAssessment2024, mousavinezhadNextLithiumBoom2025}. For harder claystone systems such as Thacker Pass, Clayton Valley, and TLC, additional mechanical processing, modeled here as attrition scrubbing, is required to break down aggregated particles prior to leaching.

Classification and screening are key operations in clay flowsheets to control slurry particle size and ensure consistent leaching performance. Rather than specifying individual screening or hydrocyclone devices, these operations are abstracted into a generic classifier block parameterized by cut size and efficiency. Given the importance of water management in Nevada-based claystone systems, a thickening step is included following classification to regulate solids concentration and recycle process water.

Clay-origin slurries are processed with sulfuric acid leaching, after which the leached slurry undergoes multistage solid--liquid separation, typically represented as a counter-current decantation (CCD) system, to recover pregnant leach solution and minimize soluble lithium losses to residue. Water recovery and solution concentration are key drivers at this stage and materially affect both reagent consumption and downstream purification requirements. Staged neutralization and impurity removal steps are required to control iron, aluminum, magnesium, calcium, and residual acidity prior to lithium recovery. While specific precipitation sequences vary across project disclosures, these operations perform the same functional role: conditioning the sulfate solution to within acceptable impurity tolerances. The resulting purified lithium sulfate solution is then processed into lithium carbonate.

Of note is that the last three steps of this process are not included in the Thacker Pass project flowsheet, where a magnesium recycling route was adopted in which magnesium salts are removed via evaporative crystallization -- consuming fewer reagents at the cost of substantially higher thermal energy demand. Our generic model instead fully precipitates all magnesium with quicklime, causing total quicklime consumption to reach approximately 9,900 kg \ce{CaO}/t \ce{Li2CO3}, compared to the disclosed 2,630 kg/t. This addition of quicklime also necessitates addition of sodium carbonate to polish the additional calcium, significantly increasing costs. As the Thacker Pass process composition does not specify constituent ions, the modeled quicklime dose is calibrated against the disclosed magnesium removal duty rather than any assay-derived stoichiometry. 

This facility is composed of the following blocks:

\begin{description}[noitemsep, topsep=0pt]
    \item[\textbullet\enskip\texttt{Crushing:}] Primary crushing for particle size control.
    \item[\textbullet\enskip\texttt{Attrition Scrubbing:}] Attrition scrubbing for harder claystone systems.
    \item[\textbullet\enskip\texttt{Classification:}] Classification and screening for slurry particle size control.
    \item[\textbullet\enskip\texttt{Thickening:}] Thickening to regulate solids concentration and recycle process water.
    \item[\textbullet\enskip\texttt{Acid Leaching:}] Sulfuric acid leaching of classified clay slurry.
    \item[\textbullet\enskip\texttt{Neutralization:}] Staged neutralization, removing Fe, Al, and others, via pH adjustment and precipitation. 
    \item[\textbullet\enskip\texttt{Multistage Primary Counter-Current Decantation:}] Primary counter-current decantation for pregnant leach solution recovery.
    \item[\textbullet\enskip\texttt{Impurity Removal:}] 
    Chemical removal of impurities (Mg, Ca, etc.) via pH adjustment and precipitation
\end{description}

\medskip
\noindent\textbf{Crushing.}\quad

The crushing step reduces run-of-mine (\texttt{rom\_ore\_feed}) claystone to a particle size suitable for slurry preparation and attrition scrubbing. The primary input is ROM ore in $\si{m^3}$ of loose rock, converted to the step basis of \texttt{crushed\_ore\_throughput} in tonnes via an IPCF equal to the loose ore density $\rho_{\text{loose}} = \rho_{\text{bank}}/S$, which for Thacker Pass is $1.9/1.20 = 1.58$ $\si{t/m^3}$. This value is loaded dynamically from project-specific data rather than hardcoded, and is the same $\rho_{\text{loose}}$ applied in the excavation and loading POCF, so the two cancel and the mine tonnage is independent of the swell factor. The input CCF is set alongside the IPCF as $\text{IPCF} \times \text{step CCF}$, so that constituent grades pass through the mining-to-comminution boundary unchanged. The step CCF is 0.001 kg ore per (tonne-ppm), the standard conversion for mass-basis ppm constituents. Only Li is tracked as a constituent through the mining and comminution circuit, sufficient to propagate the ore grade forward to sizing calculations in downstream steps.

No reagents are consumed at this step. A \texttt{rejected\_coarse} coproduct stream routed to the \texttt{rock\_pile\_nv\_100} sink is defined in the data schema with a conversion factor of zero, reflecting the modeling assumption that no coarse material is rejected to waste at the crushing stage. All material is assumed to be crushed until in-spec for feed into attrition scrubbing in the following step. 

This step uses the \texttt{crushing\_cycle} machine block (Appendix \ref{app:04-methods-blockdefs-crushing}).

\medskip
\noindent\textbf{Attrition Scrubbing.}\quad

The attrition scrubbing step deagglomerates crushed claystone particles and produces a pumpable slurry at a target solids concentration suitable for classification. The primary input is \texttt{crushed\_ore} in tonnes, converted to the step basis (\texttt{scrubbed\_slurry\_throughput}, in \si{m^3} of slurry at approximately 30 wt\% solids) via an IPCF of 3.40. This factor accounts for both the mass-to-volume conversion of crushed ore into a slurry and an assumed 20\% contribution of coarse recycled material from the downstream classification step: for a target slurry at 30 wt\% solids using a grain density of 2.6 \si{t/m^3} and water density of 1.0 \si{t/m^3}, the mixture density is $\rho_{\text{mix}} = 1/(w_s/\rho_s + (1-w_s)/\rho_w)$, giving a step CCF of $w_s \cdot \rho_{\text{mix}} / 10^3 \approx 3.68 \times 10^{-4}$ kg solids per \si{m^3} per ppm. Including the 20\% recycle flow from classification yields the effective IPCF of approximately 3.40.

The output \texttt{scrubbed\_slurry} is a diluted slurry at approximately 20 wt\% solids, reflecting water addition during scrubbing. The POCF of 1.613 converts from the 30 wt\% step basis to the 20 wt\% output basis, derived as the ratio of the inverse CCF values at the two solids fractions: $\text{POCF} = (1/(0.2 \cdot \rho_{20\%})) / (1/(0.3 \cdot \rho_{30\%}))$. The output CCF is correspondingly $w_s \cdot \rho_{\text{mix}} / 10^3 \approx 2.28 \times 10^{-4}$ kg solids per \si{m^3} per ppm at 20 wt\% solids.

Two reagents are dosed against the step basis volume: a slurry dispersant at 0.001 \si{kg/m^3} of slurry (a rough engineering estimate for clay/silt deflocculation, with no constituent elimination), and recycled process water at 1.177 \si{m^3/m^3} of slurry, dosed at zero cost. This dose is the step's full water demand net of the classifier recycle: the water content of the scrubbed slurry output, $0.912 \times 1.613=1.472$ \si{m^3} per \si{m^3} of step basis, less the $0.2 \times 1.613 \times 0.912=0.294$ \si{m^3} returning with the 20\% coarse recycle at the same density. This is equivalent to 91.2 \si{m^3} of water in total per 100 \si{m^3} of scrubbed slurry output, of which 18.2 \si{m^3} arrives with the classifier recycle and 73.0 \si{m^3} is dosed. The dose is drawn from the beneficiation circuit's internal returns (thickener overflow and the water reclaimed from the coarse-gangue tailings dewatering), rather than purchased. Water not reclaimed from the tailings is treated as lost, as described under the tailings sinks; the circuit's raw make-up, which covers that loss, enters at classification. No coproducts are generated at this step.

This step uses the \texttt{attrition\_scrubber} machine block (Appendix \ref{app:04-methods-blockdefs-attritionscrubber}).

\medskip
\noindent\textbf{Classification.}\quad
The classification step separates fine-grained particles suitable for acid leaching from coarser material that is recycled back to attrition scrubbing, abstracting the function of hydrocyclones and/or screens into a single sizing operation. This is the beneficiation step of the clay pathway: the coarse fraction is both recirculated for further attrition and, in part, rejected to tailings, so the fines advancing to thickening carry an upgraded lithium grade relative to the crushed ore. The primary input is \texttt{scrubbed\_slurry} at 20 wt\% solids, with an IPCF of 1.0 (the step basis is the same volumetric unit as the incoming slurry) and a CCF of $2.28 \times 10^{-4}$ kg solids per \si{m^3} per ppm, consistent with the output of the preceding step.

The output \texttt{classified\_slurry} leaves at approximately 15 wt\% solids, carrying an effective POCF of $0.695$ built from three factors: the density ratio $1.380$ converting the 20 wt\% step basis to the 15 wt\% output, the exit deflation $1-0.20=0.80$, the fresh-feed fraction of classifier throughput, since 20\% of throughput recirculates to attrition scrubbing, and the mass rejection ratio that carries the beneficiation upgrade, formed from the ratio of the initial ROM grade to upgraded leach-feed grade ($f/c=2,538/4,033=0.629$). The first two are held in the spreadsheet cell; the third is applied at runtime from the project's ROM and leach feed grades so that the cell does not go stale if either grade moves. The output CCF is $1.65 \times 10^{-4}$ kg solids per \si{m^3} per ppm, based on a 15 wt\% solids mixture. 

The yield rate is set to 0.92, matching the 8\% loss to coarse gangue mineralization reported in the Thacker Pass disclosure~\cite{sgscanadaincSK1300Technical2024}. The solids split between classified ore and rejected coarse materials is a derived quantity rather than a declared one, and must be read against a stated basis. Per tonne of fresh crushed solids entering the circuit, the two-product formula gives a concentrate mass fraction of $Y=R \cdot f/c=0.92 \times 2,538/4,033=0.579$, with the complementary $0.421$ rejected as coarse gangue at $2,538 \times 0.08/0.421=482$ ppm Li. Per unit of the classifier's own step basis, which is inflated by the circulating load to $1/0.80=1.25$ times the fresh solids, the same split reads 0.463 forwarded, 0.337 rejected, and 0.200 recirculated to \texttt{scrubbing}, summing to unity. The reject is the disclosed life-of-mine figure, taken from Table 13-4 of the 2024 disclosure \cite{sgscanadaincSK1300Technical2024}, which reports 611.8 Mt of dry leach ore against 1,056.7 Mt of dry ore mined, a rejection of $1-611.8/1056.7=0.421$ of fresh solids. This matches with the statement in Section 14.4.2.2 of the same disclosure noting that up to 46\% of the ore fed to the process will be rejected during classification. 

One coproduct is generated: \texttt{rejected\_coarse\_gangue}, a wet reject leaving at the step-basis density of approximately 20 wt\% solids, with a feed-basis conversion factor of $(1-Y) \times 0.80=0.273$ \si{m^3} per \si{m^3} of step basis. For the output, dilution from 20 to 15 wt\% requires 0.200 \si{m^3} of water per \si{m^3} of step basis. Classification is also where the beneficiation circuit's raw make-up enters, sized at runtime as the water leaving the circuit in solid streams, amounting to 0.109 \si{m^3} per \si{m^3} of step basis. The balance of the dilution requirement, subsequently 0.091 \si{m^3} per \si{m^3} at runtime, is drawn from the thickener overflow and is dosed as recycled water at zero cost, as is the whole of the attrition scrubbing demand. 

This step uses the \texttt{classifier\_generic} machine block (Appendix \ref{app:04-methods-blockdefs-classifier}).

\medskip
\noindent\textbf{Thickening.}\quad

The thickening step concentrates the classified slurry from approximately 15 wt\% solids to approximately 55 wt\% solids, preparing a dense leach feed for the acid leaching circuit and recovering process water in the thickener overflow for recycle. The primary input is \texttt{classified\_slurry} with an IPCF of 1.0 and an input CCF of $1.65 \times 10^{-4}$ kg solids per \si{m^3} per ppm, consistent with the output of the classification step ($w_s \cdot \rho_{\text{mix}} / 10^3$ at 15 wt\% solids, grain density 2.6 \si{t/m^3}).

The output \texttt{thick\_leach\_feed} is at approximately 55 wt\% solids, with a POCF of 0.199 converting from the 15 wt\% step basis to the 55 wt\% output, derived as the ratio of inverse CCFs at the two solids fractions. The output CCF is $8.31 \times 10^{-4}$ kg solids per \si{m^3} per ppm at 55 wt\% solids. The lithium grade is unchanged across this step; the upgrade is assumed to be carried entirely by classification, and thickening is a pure densification at constant grade.

One reagent is consumed: a polymer flocculant dosed at $2.6 \times 10^{-4}$ \si{kg/m^3} of slurry feed against the step basis, with no constituent elimination. This dosage is calibrated to match the reported combined flocculant consumption across the thickener and decanter circuits in the Thacker Pass disclosure of approximately 260 g/t \cite{sgscanadaincSK1300Technical2024}. The thickener overflow (\texttt{thickener\_overflow}) is a coproduct carrying 0.801 \si{m^3} per \si{m^3} of thickening basis, the water entering with the 15 wt\% feed less the water leaving in the 55 wt\% underflow, and is routed to the \texttt{process\_water\_recycle} sink at zero cost, supplying the whole of water demand for the attrition scrubbing step and the balance of classification dilution.

This step uses the \texttt{thickener\_generic} machine block (Appendix \ref{app:04-methods-blockdefs-thickener}).

\medskip
\noindent\textbf{Acid Leaching.}\quad

The acid leaching step dissolves lithium and other soluble species from the thickened claystone slurry using sulfuric acid, producing a leach slurry (\texttt{leached\_slurry}) containing a pregnant leach solution rich in \ce{Li2SO4} alongside gangue residue, with the sulfuric acid being produced on-site at the sulfuric acid plant (Appendix \ref{app:04-methods-facilitychoices-sap}). The primary input is \texttt{thick\_leach\_feed} at approximately 55 wt\% solids, converted to the step basis (\texttt{leach\_slurry\_throughput} at approximately 34 wt\% solids) via an IPCF of 1.748, reflecting the dilution of the thickened feed upon addition of acid solution and makeup water. This IPCF is derived as the ratio of the inverse CCFs at 55 wt\% and 34 wt\% solids, with the 34 wt\% stream taken acid-inclusive: its ore loading is $1/(1/2.6 + (1/0.34 - 1 - 0.49)/1 + 0.49/1.84) = 476$ \si{kg/m^3}, against $0.55 \cdot \rho_{\text{55\%}} = 831$ \si{kg/m^3} for the thickened feed, using grain density 2.6 \si{t/m^3} and an acid density of 1.84 \si{t/m^3}. The input CCF is $8.31 \times 10^{-4}$ kg solids per \si{m^3} per ppm (at 55 wt\% solids), while the step and output CCF is $4.76 \times 10^{-4}$ kg ore per \si{m^3} per ppm at 34 wt\% solids, acid-inclusive. 

The principal reaction is the dissolution of lithium from the clay lattice:
\vspace{-1em}
\alteq{Lithium oxide combined with silicon dioxide and water in solid form reacts with aqueous sulfuric acid to produce aqueous lithium sulfate, solid silicon dioxide, and water.}
\begin{equation}
\ce{Li2O \cdot nSiO2 \cdot mH2O (s) + H2SO4 (aq) -> Li2SO4 (aq) + nSiO2 (s) + (m+1)H2O}
\end{equation}
\alteqend

Other gangue elements (e.g. Al, Mg, Ca, Fe) form soluble sulfates or secondary solids; their removal is handled in downstream impurity removal steps.  

Two reagents are dosed against the step basis:
\begin{description}[noitemsep, topsep=0pt]
\item[\textbullet\enskip Sulfuric acid] at 233.1 kg \ce{H2SO4}/\si{m^3} leach slurry against the step basis volume. The dose is encoded on the acid charge rather than on the slurry, as 0.49 kg \ce{H2SO4} per kg of ore solids, the disclosure's specified design consumption~\cite{sgscanadaincSK1300Technical2024, rothFeasibilityStudyNational2022}; the slurry density therefore cancels and the dose does not drift when the leach basis is re-encoded. There is no constituent elimination (acid consumption drives dissolution but the Li constituent is tracked separately via yield rate).
\item[\textbullet\enskip Recycled process water] at at 0.301 \si{m^3/m^3} of leach slurry against the step basis, representing the additional dilution water needed to transition from the 55 wt\% thickened feed to the 34 wt\% leach slurry operating density. Per the Thacker Pass disclosure the leach feed is repulped with CCD wash solution rather than raw water~\cite{sgscanadaincSK1300Technical2024}, so the dose is taken as an internal return at zero cost; the corresponding raw water demand sits in the CCD wash-water dose, which is sized as the net water demand of the leach--CCD circuit and therefore covers it.
\end{description}

The output yield rate is 0.92, the disclosure's expected average lithium extraction in the acid leach circuit~\cite{sgscanadaincSK1300Technical2024}. Because no mass leaves the step --- the solid residue remains entrained in the leach slurry until the CCD circuit --- this recovery is a constituent efficiency rather than a volume split: the POCF is set to $1/0.92=1.087$ so that $y \cdot \text{POCF}$ returns the physical volume, and the unrecovered lithium is carried in the downstream grade anchor, leaving the flowsheet once in the CCD filter cake. No coproducts are generated at this step.

This step uses the \texttt{acid\_leach\_circuit} machine block (Appendix \ref{app:04-methods-blockdefs-acidleach}).

\medskip
\noindent\textbf{Neutralization.}\quad

The neutralization step consumes residual free acid in the leached slurry and initiates precipitation of iron and aluminum impurities prior to solid--liquid separation in the CCD circuit. The primary input is \texttt{leached\_slurry} at approximately 34 wt\% solids, with the input, step, and output CCFs all equal to $4.76 \times 10^{-4}$ kg ore per (\si{m^3} $\cdot$ ppm), computed acid-inclusively as $1/(1/\rho_s + (1/w_s - 1 - D)/\rho_w + D/\rho_{\ce{H2SO4}})/10^3$ at 34 wt\% solids, grain density 2.6 \si{t/m^3} and acid dose $D$ = 0.49 kg/kg ore. 

Two reagents are dosed against the step basis:
\begin{description}[noitemsep, topsep=0pt]
\item[\textbullet\enskip Limestone (\ce{CaCO3})] at 25.93 \si{kg/m^3} of leach slurry for bulk acid neutralization to pH $\sim$3.5, calibrated to match the reported limestone consumption rate of approximately 2.95 t/t \ce{Li2CO3} in the Thacker Pass disclosure~\cite{sgscanadaincSK1300Technical2024}. The addition of limestone is assumed to not significantly change the throughput volume. No constituent elimination is modeled at this step; the gypsum precipitated here remains in the slurry and reports to the CCD filter cake through \texttt{s\_cake\_CCD\_TP}; \ce{Fe(OH)3} is not tracked as mass.
\item[\textbullet\enskip Flocculant] at 0.03 \si{kg/m^3} of leach slurry, as a settling aid for the nascent precipitates ahead of the CCD. This value is a rough estimate.
\end{description}

The output \texttt{neutralized\_pls} passes to the CCD circuit with a POCF of 1.0 and a yield of 1.0.

This step uses the \texttt{staged\_neutralization} machine block (Appendix \ref{app:04-methods-blockdefs-neutralization}).

\medskip
\noindent\textbf{Counter-Current Decantation.}\quad

The Counter-Current Decantation (CCD) step separates the pregnant leach solution from the solid leach residue through multi-stage counter-current decantation, recovering dissolved lithium from entrained solution while minimizing lithium losses to the filter cake. The primary input is \texttt{neutralized\_pls} at approximately 34 wt\% solids, with input and step CCFs of $4.76 \times 10^{-4}$ kg solids per \si{m^3} per ppm, consistent with the neutralization step output.

The output \texttt{washed\_pregnant\_leach\_solution} has a POCF of 1.23, reflecting the volumetric expansion from a 34 wt\% solids slurry feed to a clarified PLS at 16.1 wt\% total dissolved solids (178.3 g/L) and density $\rho_{PLS}=1/(0.1607/2.6+0.8393/1.0)=1.110$ \si{t/m^3}. The physical PLS volume is fixed by the wash-water balance: each \si{m^3} of CCD feed carries 690.35 kg of water and 254.89 kg of dissolved salts, the latter being the charged acid less the 10.9\% fixed as gypsum at neutralization, paired with its cation; 377.60 kg of solids report to the cake (333.00 kg of leach residue plus 44.60 kg of that gypsum), taking 231.43 kg of liquor with them at the 62 wt\% cake specification. Water used for washing is dosed at 3.3 times the cake liquor water~\citet{steinCalculateCountercurrentWashing1990}, or 0.641 \si{m^3} per \si{m^3} of feed, which sets the liquor at 16.1 wt\% TDS and leaves 1,355 kg of PLS liquor, or 1.22 \si{m^3}. The dosed raw draw of 0.916 \si{m^3/m^3} exceeds this last-stage wash because it also covers
the leach repulp requirement.

The yield rate is 0.99, the reported 0.5--1.5\% lithium loss to the leach residue~\cite{sgscanadaincSK1300Technical2024}; because this recovery is a constituent efficiency and no volume leaves with it, the tabulated POCF is $1.2057/0.99=1.2179$ so that $y \cdot \text{POCF}$ returns the physical volume. The output CCF shifts to $\rho_{PLS}10^3=1.115 \times 10^{-3}$ kg per \si{m^3} per ppm, reflecting the transition from a solids-containing slurry basis to a PLS density basis.

One reagent is consumed: flocculant at 0.04 \si{kg/m^3} of CCD feed. This value is a rough estimate; typical polymer flocculant doses in CCD circuits for sulfate leach slurries range from approximately 20--80 g/t of solids processed~\cite{sgscanadaincSK1300Technical2024}, and the modeled dose is broadly consistent with the lower end of this range given the solids loading implied by the 34 wt\% feed. Wash water is dosed explicitly at 0.942~\si{m^3} per \si{m^3} of CCD feed, sized as the water leaving in the PLS and cake less the water entering with the thickened leach feed. This draw covers the leach repulp requirement as well as the CCD's own, since the leach is repulped with CCD wash solution and takes 0.3011~\si{m^3/m^3} as an internal return at zero cost.

The filter cake coproduct (\texttt{filter\_cake}) is routed to the \texttt{claycake\_stack\_nv\_62} sink with a conversion factor of 0.3644 \si{m^3} per \si{m^3} of CCD feed. 70\% of the ore solids are assumed to survive dissolution, leading to 333 kg of leach residue from the 475.72 kg of ore solids. 32.33 kg of gypsum is also produced from the limestone dosed at neutralization. At a 62 wt\% permit specification, these 365.33 kg of cake solids is 589.24 kg of cake, which at the corresponding mixture density of 1.617 \si{t/m^3} occupies 0.3644 \si{m^3}.

This step uses the \texttt{multi\_stage\_ccd\_thickening} machine block (Appendix \ref{app:04-methods-blockdefs-ccd}).

\medskip
\noindent\textbf{Impurity Removal.}\quad

The impurity removal step precipitates dissolved magnesium and calcium from the washed PLS, producing a clarified lithium sulfate solution (\texttt{pls\_post\_polishing}) suitable for soda ash precipitation of lithium carbonate. This step represents the most significant divergence from the Thacker Pass project-specific flowsheet, which uses magnesium crystallization, ion exchange, and zero-liquid-discharge evaporation rather than chemical precipitation, as discussed in Section~\ref{sec:04-scopemethod}. As such, the reagent doses, sludge volumes, and cost implications of this generic all-chemical route are approximate but conservative, and thus should be revisited with a more detailed mass balance in future work.

The primary input is \texttt{washed\_pregnant\_leach\_solution} at approximately 178.3 g/L total dissolved solids (TDS) and density $\rho_\text{PLS} \approx 1.11$ \si{t/m^3}. The input, step, and output CCFs are all $\rho_\text{PLS} \times 10^{-3} = 1.11 \times 10^{-3}$~kg/(\si{m^3}$\cdot$ppm), converting lithium concentration in ppm (mg Li/kg PLS) to kg~Li per \si{m^3} of PLS. The POCF is 1000, converting from the \si{m^3} PLS step basis to the litre output units of \texttt{pls\_post\_polishing}. The yield rate is 0.99. Because the lithium grade is unchanged across the step, this figure is simultaneously the volume split and the lithium recovery: 1\% of the feed solution leaves the step as residual PLS entrained in the washed sludge cake.

Three reagents and a wash-water stream drive the impurity removal step; all are dosed against the step basis (\si{m^3} PLS):

\begin{description}[noitemsep, topsep=0pt]

\item[\textbullet\enskip Lime (\ce{CaO})] at 71.4 \si{kg/m^3} PLS. The dose is stoichiometric on the dissolved sulfate the liquor carries out of the CCD circuit; the acid charge fixes the dissolved sulfate at \texttt{Dose\_H2SO4\_TP} $\times (1 - \texttt{f\_Neut\_TP}) = 0.437$ kg \ce{H2SO4} per kg of ore solids, of which 85.4\% reports to the PLS, giving 178.3 \si{g/L} TDS at \texttt{TDS\_TP} = 0.161. One mole of \ce{CaO} neutralizes one mole of that sulfate regardless of whatever cation carries it; the dose is taken on the precipitating share of the sulfate only. Assuming the cation is \ce{Mg^{2+}}:
%
\begin{equation*}
  \ce{MgSO4 + CaO + H2O -> Mg(OH)2 v + CaSO4 v}
\end{equation*}
%
No excess factor is applied. The entire dissolved cation load is already inside \texttt{salts\_CCD\_TP}, so an allowance for residual iron and aluminum on top of it would double-count; charge balance on the two host phases cited above places roughly 90\% of the acid-consumed cation charge on species requiring hydroxide precipitation in any case. This corresponds to approximately 9,900 kg \ce{CaO}/t \ce{Li2CO3}, substantially exceeding the Thacker Pass disclosed consumption of 2,630 kg/t \cite{sgscanadaincSK1300Technical2024}. This is due to the modeling counterfactual; the latter uses a project-specific magnesium crystallization route, which removes the same magnesium by evaporative crystallization and recycles it, using a significant amount of electricity, while this generic all-chemical route instead pays with highly increased lime consumption. 

\item[\textbullet\enskip Soda ash (\ce{Na2CO3})] at 2.6 \si{kg/m^3} PLS for calcium polishing. Because the lime dose is stoichiometric with the dissolved sulfate, essentially all dosed calcium leaves as gypsum; what remains dissolved is a small residual that must be removed to avoid contaminating the downstream \ce{Li2CO3} product. Soda ash precipitates it as \ce{CaCO3}:
\alteq{Aqueous calcium ion reacts with sodium carbonate to produce calcium carbonate precipitate and two sodium ions.}
\begin{equation*}
  \ce{Ca^{2+} + Na2CO3 -> CaCO3 v + 2Na^+}
\end{equation*}
\alteqend
The residual is a declared allowance of 1 \si{kg/m^3} dissolved \ce{Ca}, stated at one significant figure because no assay or disclosure fixes it for this liquor. It is bracketed below by gypsum saturation in pure water, 2.1 \si{g/L} \ce{CaSO4.2H2O} or 0.49 \si{kg/m^3} as \ce{Ca}, and by \ce{Ca(OH)2} saturation at the liming pH, roughly 0.9 \si{kg/m^3}; the dose follows as $1.0 \times (105.99/40.08) = 2.6$ \si{kg/m^3}. Unlike the lime dose, this duty is concentration-driven, so it dependent on the liquor volume basis and does not scale with the lime dose.

\item[\textbullet\enskip Flocculant] at 0.1 \si{kg/m^3} PLS as a settling aid. This is a rough estimate.

\item[\textbullet\enskip Wash water] at 0.3470 \si{m^3/m^3} PLS, raw process water. The sludge is filtered and the cake washed to displace entrained pregnant solution: the retained wash equals the cake liquor (0.3570 \si{m^3}) less the residual PLS (0.0100 \si{m^3}), holding the lithium loss at the 1\% yield figure. Single perfect displacement is assumed, so the dosed volume is the minimum consistent with the cake moisture.

\end{description}

The impurity sludge coproduct (\texttt{impurity\_sludge}) is routed to the \texttt{sludge\_stack\_nv\_50} sink as a filter cake at 50 wt\% solids. The conversion factor is derived from the reagent stoichiometry: per \si{m^3} of PLS the precipitated solids comprise $178.3 \times 58.32/120.37 = 74.3$ kg \ce{Mg(OH)2}, 
$(71.4 \times 40.08/56.08 - 1.0) \times 172.17/40.08=215.1$ kg gypsum (all dosed calcium less the declared 1.0 \si{kg/m^3} residual), $2.6 \times 100.09/105.99 = 2.5$ kg \ce{CaCO3}, and 0.1 kg flocculant, totaling 292.0 \si{kg/m^3}. Divided by the dry-solids loading of a 50 wt\% cake, $0.50 \times 1/(0.50/2.4 + 0.50/1.0) \times 10^3 = 705.9$ \si{kg/m^3} at an assumed composite precipitate grain density of 2.4 \si{t/m^3}, this gives 0.414 \si{m^3} wet cake per \si{m^3} PLS. The cake solids fraction is an engineering assumption rather than a disclosed value: a plate-and-frame press on a sludge that is roughly three-fourths gypsum by mass reaches this range, though the gelatinous \ce{Mg(OH)2} fraction degrades filtration, and 0.45--0.55 brackets the plausible outcome. Dry solids tonnage is invariant to this choice; only the retained water and the dewatering duty move with it. The step closes volumetrically as $0.99 + 0.414 = 1 + 0.282 + 0.122$, the final term being the volume of the precipitated solids, whose mass enters through the reagent doses and the dissolved magnesium and sulfate rather than as tracked volume.

This step uses the \texttt{staged\_neutralization} machine block (Appendix \ref{app:04-methods-blockdefs-neutralization}); the same generic machine block is used for both the neutralization and impurity removal steps, as both involve reagent addition to a slurry in agitated tanks. 

\paragraph{Lepidolite Comminution and Processing}\label{app:04-methods-facilitychoices-lepidproc}\leavevmode\par

Our model is based primarily on the methods described in \citet{dengSystematicReviewSulfate2026} and \citet{yanExtractionLithiumLepidolite2012}, as the sulfate roasting--water leach pathway is identified as the most common approach for processing lepidolite \citep{gaoLithiumExtractionHard2023, parkEffectiveLithiumExtraction2023, nanMethodsExtractingLithium2019, nanMethodsExtractingLithium2022, zouMethodEfficientLithium2024}.

Hard rock ores such as lepidolite require mechanical size reduction and beneficiation prior to chemical processing. Run-of-mine lepidolite material is crushed to approximately 150--250 $\mu$m, classified by size, recrushed if necessary, and then beneficiated with froth flotation. The flotation step is critical for this pathway: at typical Jiangxi ore grades of less than 0.4\% \ce{Li2O} ($\sim$ 1,860 ppm Li), direct roasting of ROM material would require processing enormous reagent and energy quantities per unit of lithium product. Froth flotation using amine collectors and phosphate depressants selectively recovers lepidolite minerals into a concentrate at approximately 1.63\% \ce{Li2O} ($\sim$ 7,576 ppm Li)~\cite{garrettHandbookLithiumNatural2004, jiaoKeyTechnologiesDevelopment2024}, reducing the mass throughput by approximately 60\%~\citep{jiaoKeyTechnologiesDevelopment2024}. After flotation, the throughput is ground to a product size of approximately 50--75 $\mu$m. This latter grinding step uses high-energy mechanical activation milling, which introduces structural defects in the mica lattice to facilitate the solid-state reactions during subsequent roasting. This product is then dewatered prior to entering the roasting circuit.

Following comminution, lepidolite is processed via a sulfate roasting--water leach pathway. Due to the significant presence of fluorine, acid-, chlorine-, or autoclave-based methods can produce corrosive mixtures and/or toxic gases; sulfate-based methods instead allow for the formation of precipitated fluorides (e.g., \ce{CaF2}, \ce{AlF3}), at the expense of more costly sulfate reagents. The roasting step, conducted at approximately 800--950$^\circ$C, converts lithium in tightly bound mica lattices to water-soluble lithium sulfate. The resulting calcine is then subjected to water leaching, producing a pregnant leach solution containing lithium, other dissolved metals, and residual sulfate species. Impurities are removed through staged neutralization and pH adjustment to precipitate aluminum, iron, magnesium, and other species. Solid--liquid separation then yields a clarified lithium solution suitable for downstream carbonate precipitation via soda ash addition.

This facility is composed of the following blocks:

\begin{description}[noitemsep, topsep=0pt]
    \item[\textbullet\enskip\texttt{Crushing:}] Primary and secondary size reduction.
    \item[\textbullet\enskip\texttt{Classification:}] Classification and screening to separate out particle sizes for recrushing. 
    \item[\textbullet\enskip\texttt{Flotation:}] Froth flotation for lepidolite concentration and separation from gangue.
    \item[\textbullet\enskip\texttt{Thickening:}] Thickening of the froth concentrate to grinding feed density, recovering process water.
    \item[\textbullet\enskip\texttt{Grinding:}] Mechanical activation milling using a high-energy ball mill grinder to target product size.
    \item[\textbullet\enskip\texttt{Dewatering:}] Dewatering of activated concentrate prior to roasting.
    \item[\textbullet\enskip\texttt{Sulfate Roasting:}] Sulfate roasting at 800--950$^\circ$C.
    \item[\textbullet\enskip\texttt{Water Leaching:}] Water leaching of calcine, filtration, and cake wash to produce pregnant leach solution.
    \item[\textbullet\enskip\texttt{Impurity Removal and Neutralization:}] Staged neutralization and pH adjustment for impurity 
    precipitation.
    \item[\textbullet\enskip\texttt{Clarification:}] Clarification of lithium solution for 
    downstream processing.
\end{description}

\medskip
\noindent\textbf{Crushing.}\quad

The crushing step reduces ROM lepidolite ore (\texttt{rom\_ore\_feed}) from run-of-mine fragments to a particle size of approximately 150--250~$\mu$ suitable for wet classification and flotation. The crusher operates in closed circuit with the classifier: oversize is recirculated for re-crushing, so it processes fresh feed plus the circulating load. At an assumed circulating load of 200\% (recycle $= 2\times$ fresh feed; see Classification), crusher throughput is $(1+\mathrm{CL}) = 3\times$ the fresh ROM rate. The primary input is fresh ROM ore in \si{m^3}, converted to \texttt{crushed\_ore\_throughput} in tonnes via an IPCF of 8.1 (the lepidolite grain density (2.7 \si{t/m^3}) scaled by the circulating-load factor of 3). This sizes the crusher on the recirculating throughput while drawing only fresh ROM from mining. Only Li is tracked through the comminution circuit; no reagents are consumed. The output \texttt{crushed\_ore} carries its full discharge (fresh plus recycle) to classification with a POCF of 1.0 and a yield rate of 0.99, the once-through lithium recovery (1\% loss to fines and dust routed to solid tailings).

This step uses the \texttt{crushing\_cycle} machine block (Appendix \ref{app:04-methods-blockdefs-crushing}).

\medskip
\noindent\textbf{Classification}\quad

The classification step separates fine lepidolite particles from coarser gangue in preparation for flotation and production of an enriched lepidolite concentrate, closing the crushing circuit and splitting its feed into a fine fraction that reports forward and a coarse oversize that recirculates to crushing. At an assumed circulating load of 200\% (cf.\ the 204\% cyclone case of \citet{rejaImprovingClosedCircuit2022}), the classifier processes $3\times$ the forward fines, of which the fraction $1/(1+\mathrm{CL}) = 1/3$ advances.

The primary input is \texttt{crushed\_ore} in tonnes, converted to the step basis (\texttt{classified\_slurry\_throughput} in \si{m^3} of slurry at approximately 30 wt\% solids) via an IPCF of 2.7037, derived as $1/(w_s \cdot \rho_{\text{mix}})$ at 30 wt\% solids with grain density 2.7 \si{t/m^3}: $\rho_{\text{mix}} = 1/(0.3/2.7 + 0.7/1.0) = 1.233$ \si{t/m^3}, giving $1/(0.3 \times 1.233) = 2.704$. This is unchanged from the crusher output, as the classifier receives the crusher's full discharge at the normal 30 wt\% slurry density. The step CCF is $3.70 \times 10^{-4}$ kg solids per (\si{m^3} $\cdot$ ppm), computed as $w_s \cdot \rho_{\text{mix}}/10^3$.

The POCF is $1/3 = 0.3333$, which sizes the classifier on the full recirculating throughput while forwarding only the once-through fines. The forwarded \texttt{classified\_slurry} remains at approximately 30 wt\% solids with an output CCF of $3.70 \times 10^{-4}$, its lithium grade unchanged from the feed at approximately 1,858 ppm Li (0.40\% \ce{Li2O}) --- a size-only separation in which slimes are removed at feed grade. The yield rate is 0.93, the once-through lithium recovery (7\% loss to classification slimes), consistent with the 5--8\% range for Jiangxi lepidolite desliming~\cite{zhangInnovativeRecoveryLepidolite2025}. Process water is dosed against the step basis at 0.288 \si{m^3/m^3} (accounting for deflation by the circulating-load factor of 3). That total is split at runtime between purchased and recycled water. The raw make-up is sized structurally as the water leaving the circuit, at 0.098 \si{m^3/m^3}, roughly split as 0.084 impounded with the classification slimes and flotation tailings, 0.005 leaving with the roaster feed cake and evaporated at roasting, and 0.009 bled from the circuit as press filtrate, with a small amount left with the coproduct mill discharge. The balance, 0.190 \si{m^3/m^3}, is dosed as \texttt{recycled\_process\_water} at zero cost and is met by the thickener overflow together with decant reclaimed from the tailings storage facility. The reclaim depends on the solids fraction at which the tailings are finally placed, which are assumed as 0.55. %The tailings filtration plant implied by placement above the as-discharged 30 wt\% is not itself represented in the model. <-- 
Water not reclaimed is treated as lost, as described under the tailings sinks.

One coproduct is generated: the classification slimes (\texttt{classification\_slimes}), 7\% of feed solids, routed to the \texttt{gangue\_impound\_cn\_30} sink (nonhazardous wet reject at $\sim$30 wt\% solids) with a corresponding conversion factor of 0.02333 ($0.07 \times 1/3$) of the step-basis volume, as the once-through slimes tonnage is preserved despite the circulating-load-inflated basis; the recirculating oversize is an internal circuit stream, not a tailings stream. This step uses the \texttt{classifier\_generic} machine block (Appendix \ref{app:04-methods-blockdefs-classifier}).

\medskip
\noindent\textbf{Flotation}\quad

The flotation step concentrates the deslimed lepidolite fines into a
lithium-enriched concentrate by froth flotation, and is the key grade-concentrating operation in this pathway. The step is constituent-driven: the volume-defining basis is the Li constituent of the input stream, so the required flotation circuit volume scales with the lithium mass to be processed rather than with bulk throughput. The primary input is \texttt{classified\_slurry} at approximately 30 wt\% solids, with an IPCF of 1.0 and a step CCF of $3.70 \times 10^{-4}$ kg solids per (\si{m^3} $\cdot$ ppm).
 
The output \texttt{flotation\_froth} is a lepidolite concentrate that leaves flotation at an estimated 30 wt\% solids with a POCF of 1.0 and an output CCF of $3.70 \times 10^{-4}$ kg constituent per ($\si{m^3} \cdot$ ppm). This 30 wt\% reflects the typical density of a concentrate slurry, typically 20–35 wt\% solids depending on froth depth, wash water, and drainage time. No densification occurs here; dewatering of the froth to grinding-feed density is performed by the downstream thickening step. The step performs the key grade concentration, from 1,858 ppm Li (0.40\% \ce{Li2O}) in the deslimed feed to 7,576 ppm Li (1.63\% \ce{Li2O}) in the concentrate --- a $4.08\times$ grade upgrade carried by the concentrate constituent grade.  The yield rate is 0.75, the flotation lithium recovery, taken at the conservative end of the 75--82\% range reported for industrial-scale desliming--flotation circuits processing low-grade Jiangxi ore~\cite{jiaoKeyTechnologiesDevelopment2024, zhangInnovativeRecoveryLepidolite2025}.

Two reagents are dosed against the step basis volume: depressants (sodium hexametaphosphate) at 0.1 \si{kg/m^3}, which suppress non-lepidolite silicate minerals to improve selectivity; and collectors (dodecylamine) at 0.12 \si{kg/m^3}, which selectively adsorb onto lepidolite surfaces to render them hydrophobic and recoverable by froth flotation~\cite{jiaoKeyTechnologiesDevelopment2024}. None of these reagents eliminate tracked constituents from the process stream given their function for physical separation.

One coproduct is generated: the flotation tailings (\texttt{flotation\_tailings}). The tailings fraction follows from the flotation mass balance via the two-product formula, which gives the concentrate mass yield as
\[
	Y_{\text{conc}} = R_{\text{flot}} \cdot \frac{f}{c} = 0.75 \frac{1,858}{7,576} = 0.1839,
\]
where $R_{\text{flot}} = 0.75$ is the flotation lithium recovery and $f = 1,858$ ppm and $c = 7,576$ ppm \ce{Li} (0.40\% and 1.63\% \ce{Li2O}) are the feed and concentrate grades. Thus 18.39\% of the flotation-feed solids report to the concentrate and the remaining $1 - 0.1839 = 0.8161$ (81.61\%) report to the tailings. Since the tailings and feed are both taken at approximately 30 wt\% solids, the volume fraction equals the mass fraction, giving a conversion factor of 0.8161 of the step-basis volume. The tailings are routed to the \texttt{gangue\_impound\_cn\_30} sink (nonhazardous wet reject at $\sim$30 wt\% solids).

This step uses the \texttt{flotation\_generic} machine block (Appendix \ref{app:04-methods-blockdefs-flotation}).

\medskip
\noindent\textbf{Thickening.}\quad

The thickening step dewaters the flotation froth from pulp density to grinding-feed density, recovering process water as a recyclable overflow. The primary input is \texttt{flotation\_froth} at approximately 30 wt\% solids, with an IPCF of 1.0 and a step CCF of $3.70 \times 10^{-4}$ kg solids per (\si{m^3} $\cdot$ ppm).

The output \texttt{thick\_grind\_feed} is the thickened concentrate at approximately 60 wt\% solids, with a POCF of 0.3836 and an output CCF of $9.64 \times 10^{-4}$ kg constituent per (\si{m^3} $\cdot$ ppm). The POCF follows from the same-solids volume ratio between the two densities: at 30 wt\% the slurry carries $w_s \cdot \rho_{\text{mix}} = 0.30 \times 1.233 = 370$ kg solids/\si{m^3}, and at 60 wt\% it carries $0.60 \times 1.607 = 964$ kg solids/\si{m^3}, so a fixed solids mass occupies $370/964 = 0.3836$ of its original volume. The lithium grade is unchanged at $7,576$ ppm Li and the yield rate is 1.0; the separation is water-only, so the mass-basis composition is conserved and the CCF propagates with it.

One reagent is dosed against the step basis: flocculant at 0.005 \si{kg/m^3} slurry ($\sim$13.5 g/t solids) to promote settling, estimated from similar lithium-based concentrate flotation circuits \cite{opokuOverviewCoarseParticle2025, tramflocPolymersWaterWastewater2013}. 

One coproduct is generated: the thickener overflow (\texttt{thickener\_overflow}), consisting of recovered process water, with a conversion factor of 0.6164 of the step-basis volume --- the complement of the POCF ($0.3836 + 0.6164 = 1.0000$) --- routed to the \texttt{process\_water\_recycle} sink at zero cost, where it forms the largest single return to the classification makeup recycle stream.

This step uses the \texttt{thickener\_generic} machine block (Appendix \ref{app:04-methods-blockdefs-thickener}).

\medskip
\noindent\textbf{Grinding.}\quad

The grinding step reduces the lepidolite concentrate to a fine product size of approximately 50--75 $\mu$ via high-energy mechanical activation milling, breaking down the mica crystal structure to yield an activated concentrate that is reactive in subsequent sulfate roasting. This mechanical activation is a distinguishing feature of the lepidolite pathway --- the fine grinding increases the surface area and introduces structural defects in the mica lattice that facilitate the solid-state reactions during roasting~\citep{dengSystematicReviewSulfate2026, yanExtractionLithiumLepidolite2012}.

The primary input is \texttt{thick\_grind\_feed} at 60 wt\% solids, with an IPCF of 1.0 (the step basis is the same volumetric unit). The step and input CCFs are both 0.0009643 kg solids per (\si{m^3} $\cdot$ ppm), consistent with the 60 wt\% slurry density basis. No reagents are consumed. The output \texttt{ground\_slurry} carries a yield rate of 0.995, reflecting a 0.5\% loss during milling. This loss is represented as the coproduct \texttt{mill\_discharge}. 

This step uses the \texttt{high\_energy\_grinding\_mill} machine block (Appendix \ref{app:04-methods-blockdefs-hegm}).

\medskip
\noindent\textbf{Dewatering.}\quad

The dewatering step removes water from the ground lepidolite concentrate slurry to produce a filter cake suitable for high-temperature sulfate roasting. Roasting requires a low-moisture feed; excess water would consume thermal energy in evaporation rather than driving the sulfate activation reactions, and could also cause handling problems in the roasting train. The primary input is \texttt{ground\_slurry} at 60 wt\% solids, with input and step CCFs of $9.64 \times 10^{-4}$ kg constituent per (\si{m^3} $\cdot$
ppm), derived as $w_s \cdot \rho_{\text{mix}} / 10^3$ at 60 wt\% solids with grain density 2.7 \si{t/m^3}: $\rho_{\text{mix}} = 1/(0.6/2.7 + 0.4/1.0) = 1.607$ \si{t/m^3}.

The output \texttt{roaster\_feed} is an 80 wt\% solids filter cake, expressed in tonnes rather than \si{m^3}, with a POCF of 1.205 converting from \si{m^3} of 60 wt\% slurry to tonnes of 80 wt\% cake: $\text{POCF} = w_{s,\text{in}} \cdot \rho_{\text{mix,in}} / w_{s,\text{out}} = 0.6 \times 1.607 / 0.8$. The output CCF is $8.0 \times 10^{-4}$ kg constituent per (tonne $\cdot$ ppm), reflecting the mass basis of the cake: $w_{s,\text{cake}} / 10^3 = 0.8 / 10^3$. No reagents are consumed. The press filtrate (\texttt{press\_filtrate}) is generated at 0.402 \si{m^3/m^3} of input slurry and routed to the \texttt{wastewater\_discharge\_cn} sink. This is modeled as the beneficiation circuit's bleed rather than a recycle. Studies on spodumene flotation~\cite{sauberEffectRecycledProcess2025} indicate that recirculating process water from flotation would accumulate dissolved metal ions (e.g. Al, Fe, Mg) alongside residual collector reagents, suppressing collector adsorption and degrading recovery, and this is assumed to extend to lepidolite flotation and its higher complexity ore. The filtrate is the most concentrated point in the circuit and thus the most appropriate bleed location; we assume its volume is set by the dewatering mass balance rather than by a contaminant limit, so the bleed rate (3.0\% of circuit water) is assumed rather than derived. The press filtrate quantity is consistent with the reduction in water content from 0.643 \si{t/m^3} at 60 wt\% solids to 0.241 \si{t/m^3} at 80 wt\% solids. Yield rate is 1.0.

This step uses the \texttt{mechanical\_dewatering} machine block (Appendix \ref{app:04-methods-blockdefs-dewater}).

\medskip
\noindent\textbf{Sulfate Roasting.}\quad

The sulfate roasting step thermally activates the dewatered lepidolite concentrate by reacting it with sodium sulfate, potassium sulfate, and lime at 800--950$^\circ$
C, converting lithium from the tightly bound mica lattice into water-soluble lithium sulfate. The central reaction is:
\alteq{Solid lepidolite consisting of lithium oxide, aluminum oxide, three silicon dioxide units, and potassium fluoride reacts with sodium sulfate to produce solid lithium sulfate, solid aluminum oxide, solid silicon dioxide, and potassium fluoride as a gas or solid.}
\begin{equation}
	\ce{(Li2O)(Al2O3)(3SiO2)(KF)(s) + Na2SO4 -> Li2SO4 (s) + Al2O3 (s) + 3SiO2 (s) + KF (g/s)}
\end{equation}
\alteqend
with fluorine species captured by lime addition as \ce{CaF2} and \ce{AlF3} precipitates rather than volatilizing as toxic \ce{HF} gas.

The primary input is \texttt{roaster\_feed} in tonnes of 80 wt\% solids cake, with an IPCF of 1.6, representing the total roasting mix mass per tonne of input cake after addition of all reagents and residual moisture. This factor accounts for: dry ore solids (0.8 t per tonne cake), sodium sulfate addition (0.5 $\times$ 0.8 = 0.4 t), lime addition (0.15 $\times$ 0.8 = 0.12 t), potassium sulfate addition (0.1 $\times$ 0.8 = 0.08 t), and residual moisture (0.2 t). The input and step CCF are both $5.0 \times 10^{-4}$ kg constituent per (tonne $\cdot$ ppm): $w_{s,\text{cake}} / (G_{18} \times 10^3) = 0.8 / (1.6 \times 10^3)$.

Three reagents are dosed against the \texttt{roasting\_mix} step basis:
\begin{description}[noitemsep, topsep=0pt]
\item[\textbullet\enskip Sodium sulfate (\ce{Na2SO4})] at 250 kg/tonne roasting mix (equivalent to 0.5 t per tonne dry lepidolite concentrate), derived from \citet{yanExtractionLithiumLepidolite2012}.
\item[\textbullet\enskip Lime (\ce{CaO})] at 75 kg/tonne roasting mix (equivalent to 0.15 t per tonne dry concentrate). The base value of 0.1 t/t from \citet{yanExtractionLithiumLepidolite2012} is increased by 50\% following \citet{luongFactorsAffectingExtraction2013} to ensure complete fluoride capture, suppressing \ce{HF} volatilization.
\item[\textbullet\enskip Potassium sulfate (\ce{K2SO4})] at 50 kg/tonne roasting mix (equivalent to 0.1 t per tonne dry concentrate), derived from \citet{yanExtractionLithiumLepidolite2012}, to supplement the alkali sulfate flux and improve lithium extraction efficiency.
\end{description}
The output \texttt{roasted\_mix} is the calcine after roasting, with water evaporated during the process. The POCF is 0.875, derived as the ratio of the calcine mass (excluding evaporated water: $0.8 \times (1 + 0.5 + 0.15 + 0.1) = 1.4$ t) to the total roasting mix mass (1.6 t). The output CCF is $5.71 \times 10^{-4}$ kg constituent per (tonne $\cdot$ ppm): $w_{s,\text{cake}} / (0.8 \times (1 + 0.5 + 0.15 + 0.1) \times 10^3)$. A fluorine offgas coproduct (\texttt{fluorine\_offgas}) is defined with a conversion factor of zero, reflecting the assumption that lime additions during roasting fix essentially all fluorine as solid precipitates before volatilization.
A second coproduct, \texttt{evaporated\_water}, carries the cake moisture entering with the 80 wt\% roaster feed, at 0.125 tonnes per tonne of roasting mix ($0.2/1.6$), and is routed to the \texttt{atmosphere} sink. Closure on the mass basis is exact: $0.875+0.125=1.000$.

This step uses the \texttt{sulfate\_roasting\_train\_lepidolite} machine block (Appendix \ref{app:04-methods-blockdefs-sulfateroasting}).

\medskip
\noindent\textbf{Water Leaching.}\quad

The water leaching step dissolves the lithium sulfate and other soluble species from the roasted calcine into aqueous solution, producing a pregnant leach solution for downstream impurity removal. The primary input is \texttt{roasted\_mix} in tonnes of calcine, converted to the step basis (\texttt{leaching\_slurry} at approximately 20 wt\% solids) via an IPCF of 4.370, derived as $1/(w_s \cdot \rho_{\text{mix}})$ at 20 wt\% solids with assumed calcine grain density of 2.7 \si{t/m^3}: $\rho_{\text{mix}} = 1/(0.2/2.7 + 0.8/1.0) = 1.144$ \si{t/m^3}, giving $1/(0.2 \times 1.144) = 4.370$. The input CCF is $5.71 \times 10^{-4}$ kg constituent per (tonne $\cdot$ ppm), consistent with the roasting output. The step CCF is $1.31 \times 10^{-4}$ kg constituent per (\si{m^3} $\cdot$ ppm), derived as the product of the 20 wt\% slurry density ($w_s \cdot \rho_{\text{mix}} / 10^3 = 0.229$ kg/L) and the mass fraction of original ore solids in the calcine ($1 / (1 + 0.5 + 0.15 + 0.10) = 0.571$), which corrects for the dilution of the ore-derived Li quantity by reagent-derived solids in the calcine.

Process water is dosed at 0.983 \si{m^3/m^3} of leach slurry basis, composed of two components. The leach water, $(1 - w_s) \cdot \rho_{\text{mix}} / \rho_{\text{water}} = 0.8 \times 1.144 / 1.0 = 0.915$ \si{m^3/m^3}, is basis-inclusive, since the step basis is defined after the water is added. The cake wash, $V_\text{cakeliq} - r_\text{res} = 0.0782 - 0.0100 = 0.0682$ \si{m^3/m^3}, is additive, added at the filter to displace the liquor held in the cake.

The POCF is 0.9335, the liquor volume per unit slurry volume less the declared residual retained after washing, $P - r_\text{res} = 0.9435 - 0.0100$. The output CCF is $1.05 \times 10^{-3}$ kg constituent per (\si{m^3} $\cdot$ ppm), the pregnant leach solution density $\rho_\text{PLS}/10^3$ on the washed liquor. The yield rate is 0.90, the combined roast, leach, and wash lithium extraction reported for a \ce{Na2SO4}/\ce{K2SO4}/CaO roast~\cite{gaoLithiumExtractionHard2023, zhaoRecentAdvancesLithium2023}. The unextracted lithium is unconverted mica that never becomes water-soluble \ce{Li2SO4}, and reports to the coproduct residue. The leach residue coproduct (\texttt{leach\_residue}) is routed to the \texttt{slag\_stack\_cn\_65} sink with a conversion factor of 0.1347 \si{m^3/m^3} of leach slurry, the insoluble grain volume plus the retained liquor ($0.0565 + 0.0782$). Because the wash adds water, product and coproduct sum to $1 + 0.0682 = 1.0682$. Residue constituents include Li, Al, Fe, Si, and \ce{CaF2}, classified as non-hazardous gangue.

This step uses the \texttt{water\_leach} machine block (Appendix \ref{app:04-methods-blockdefs-waterleach}).

\medskip
\noindent\textbf{Impurity Removal and Neutralization.}\quad

The impurity removal step precipitates calcium, residual aluminium and fluoride, magnesium, and trace heavy metals from the pregnant leach solution via staged pH adjustment, producing a clarified lithium sulfate solution suitable for downstream carbonate precipitation. The primary input is \texttt{pregnant\_leach\_solution}, a clarified liquor. Input, step, and output CCFs are all equal to $1.05 \times 10^{-3}$ kg constituent per (m$^{3} \cdot$ppm), the pregnant leach solution density  $\rho_{PLS}/10^3$, derived from the dissolved load at 76.8 g/kg. 
 
Three reagents are dosed against the step basis volume:
 
\begin{description}[noitemsep, topsep=0pt]
\item[\textbullet\enskip Neutralizing agent] at 4 \si{kg/m^3} of PLS, modeled as a generic mixed reagent representing a combination of lime, limestone, and/or NaOH for staged pH adjustment. The dosage is a rough estimate, given that no measured PLS impurity assay is available for this pathway. While there is the known duty of polishing of reagent-derived calcium -- the roast dose of 0.15~t \ce{CaO} per tonne of dry concentrate supplies 107.2~kg~Ca/t, of which 19.6~kg/t is consumed capturing fluoride as \ce{CaF2}, with the balance entering the leach liquor -- the rest is intended to capture heavy metals and impurities such as thallium. This uncertainty in both composition and neutralizing agent combination warrants further refinement against measured PLS impurity data.
\item[\textbullet\enskip Sulfuric acid] at 1 \si{kg/m^3} of PLS for pH re-adjustment after primary neutralization, to avoid overshooting the target pH and precipitating lithium-bearing compounds. This is also a rough estimate.
\item[\textbullet\enskip Sulfide precipitant] at 0.1 \si{kg/m^3} of PLS, representing sodium hydrosulfide or equivalent for thallium and residual heavy-metal polishing. The dose is a practical, assumed minimum metering rate, as no heavy metal assay is currently available for this pathway. 
\end{description}

The output \texttt{neutralized\_pls} passes to clarification with a POCF of 1.0 and a yield rate of 1.0: precipitation occurs in the agitated reactor train, but no solid--liquid separation takes place there, so no material leaves the step. The precipitated solids are carried forward in the neutralized stream and are separated at clarification.

This step uses the \texttt{staged\_neutralization} machine block (Appendix \ref{app:04-methods-blockdefs-neutralization}).

\medskip
\noindent\textbf{Clarification.}\quad

The clarification step performs a final solid--liquid separation on the neutralized PLS to remove residual precipitates and suspended solids, yielding a clarified lithium solution (\texttt{pls\_post\_polishing}) suitable for soda ash precipitation. The primary input is \texttt{neutralized\_pls}, with input and step CCFs of $1.05 \times 10^{-3}$ kg constituent per (m$^3 \cdot$ppm), consistent with the upstream step. One reagent is consumed: flocculant at 0.005 \si{kg/m^3} of PLS, estimated from similar lithium-based concentrate flotation circuits \cite{opokuOverviewCoarseParticle2025, tramflocPolymersWaterWastewater2013}.
 
The output \texttt{pls\_post\_polishing} is expressed in liters, with a POCF of 1000 (\si{m^3} to L) and an output CCF of $1.05 \times 10^{-6}$ kg constituent per (L$\cdot$ppm), derived by dividing the \si{m^3}-basis CCF by 1000. This step performs the solid--liquid separation for the neutralization circuit: precipitation occurs upstream in the agitated reactor train, and the resulting solids are separated here.

The impurity sludge coproduct is derived from an assumed reagent stoichiometry rather than from a PLS assay, which is not available for this pathway. It assumed to be primarily formed of calcium from upstream targeted fluorine removal, but also contains heavy metal contaminants such as thallium, rubidium, and cesium, which make this sludge hazardous. Polishing of dissolved calcium proceeds as
\begin{equation}
	\ce{Na2CO3(aq) + Ca^{2+}(aq) -> CaCO3(s) v + 2Na^{+}(aq)},
\end{equation}
so precipitated solids per unit reagent are $M_{\ce{CaCO3}}/M_{\ce{Na2CO3}} = 0.944$. The assumed reagent dosage $D_\text{neut} = 4$ \si{kg/m^3} gives a dry-solids loading of
\begin{equation}
	\mathrm{DS} = D_\text{neut}\cdot\frac{M_{\ce{CaCO3}}}{M_{\ce{Na2CO3}}}
	            = 4.0 \times 0.944 = 3.777 \si{kg/m^3}\text{ of PLS}.
\end{equation}
The sludge is placed as a filter cake at $s = 0.50$, whose ideal dry-solids loading is 729.7 \si{kg/m^3} at a grain density of 2.7 \si{t/m^3}, giving a conversion factor of $5.18 \times 10^{-3}$ \si{m^3} of cake per \si{m^3} of PLS, routed to the \texttt{sludge\_gate\_cn\_50} sink.

The yield rate of 0.996 follows directly from this. At $s = 0.50$ the cake retains one kilogram of liquor per kilogram of dry solids, and that liquor is at PLS grade, so the lithium lost is exactly the liquor volume fraction withdrawn:
\begin{equation}
	y = 1 - \frac{\mathrm{DS}\,(1-s)/s}{\rho_\text{PLS}}
	  = 1 - \frac{3.777}{1050.8} = 0.996.
\end{equation}

This step uses the \texttt{solid\_liquid\_separation} machine block (Appendix \ref{app:04-methods-blockdefs-sls}).

\paragraph{Pregnant Leach Solution to Lithium Carbonate (Common to All Pathways)}\label{app:04-methods-facilitychoices-pls}\leavevmode\par

This section applies to the conversion of purified lithium-bearing solution to lithium carbonate product \citep{srkconsultingSECTechnicalReport2023a, chordiaLifeCycleEnvironmental2022}. Although upstream chemistry differs between chloride-based brine pathways (\ce{LiCl}) and sulfate-based ore pathways (\ce{Li2SO4}), the downstream precipitation and solids handling configuration is broadly similar.

In both cases, lithium carbonate is produced by reaction of dissolved lithium with soda ash. Chloride systems generate sodium chloride as a byproduct while sulfate systems generate sodium sulfate; these differences primarily affect salt management and recycle strategy rather than equipment configuration. The core precipitation, thickening, washing, and drying sequence remains functionally consistent across pathways.

Precipitation is conducted at elevated temperatures to improve reaction kinetics and reduce lithium carbonate solubility. Thermal requirements are modeled generically, with pathway-specific parameterization to account for differences in reaction temperature, capturing energy demand without prescribing a specific crystallizer architecture. Following soda ash precipitation, lithium carbonate is separated from mother liquor via thickening and filtration. Washing and drying complete product finishing. Rather than specifying individual equipment types, these operations are represented as aggregated solids-handling steps parameterized by recovery efficiency, moisture content, and recycle fraction.

This facility is composed of the following blocks, with pathway-specific parameterization depending on whether the upstream pathway is chloride- or sulfate-based:

\begin{description}[noitemsep, topsep=0pt]
    \item[\textbullet\enskip\texttt{Lithium Carbonate Precipitation:}] Soda ash precipitation, including preheating and cooling.
    \item[\textbullet\enskip\texttt{Thickening and Dewatering:}] Thickening and dewatering of lithium carbonate solids.
    \item[\textbullet\enskip\texttt{Washing and Purification:}] Washing and purification, to remove residual Na, K, Mg. 
    \item[\textbullet\enskip\texttt{Drying and Packaging:}] Drying and packaging of final product.
\end{description}

\medskip
\noindent\textbf{Lithium Carbonate Precipitation.}\quad

Lithium carbonate is produced by reacting dissolved lithium with soda ash (\ce{Na2CO3}):
\begin{equation}
    2\,\text{Li}^+ + \ce{Na2CO3} \rightarrow \ce{Li2CO3(s)} + 2\,\text{Na}^+
\end{equation}
In chloride-based brine systems the byproduct is NaCl, whereas in sulfate-based ore systems it is \ce{Na2SO4}; these differences affect downstream salt management but not the precipitation equipment configuration. Both pathways use the same underlying production step, with pathway-specific parameterization for the input chemistry, soda ash dosage, and yield rate.
This step uses the \texttt{precipitator} machine block (Appendix \ref{app:04-methods-blockdefs-precipitator}).

\medskip
\noindent\textit{Chloride pathway.}\quad
 
The primary input is \texttt{pls\_post\_polishing} (concentrated lithium brine) in liters, with the step constitutent-driven on the Li constituent such that the IPCF is 1.0 by definition --- the step volume scales with the ratio of lithium mass in output volume to input PLS, rather than the bulk PLS volume. The input and step CCFs are both $1.25 \times 10^{-6}$ kg Li per (L $\cdot$ ppm), reflecting the density of the concentrated PLS (1.25 kg/L) divided by $10^6$ to convert from ppm (mg/kg) to kg/L. The POCF is $10^{-3}$ (L to \si{m^3}). The output \texttt{li2co3\_slurry} is a 20 wt\% \ce{Li2CO3} slurry with density $\rho_\text{mix} = 1/(0.2/2.11 + 0.8/1.0) = 1.118$ \si{t/m^3}, giving an output CCF of $0.2 \times 1.118 / 10^3 = 2.24 \times 10^{-4}$ kg constituent per (m$^3 \cdot$ppm). The yield rate is 0.97, reflecting a 3\% lithium loss to mother liquor and wash streams, consistent with reported losses from the precipitation circuit in the Silver Peak disclosure \cite{srkconsultingSECTechnicalReport2023a}.
 
Soda ash is dosed at 11.5 kg \ce{Na2CO3} per kg Li, targeting the Li constituent, at $1.51\times$ excess against the stoichiometric $M_{\ce{Na2CO3}}/(2M_{\ce{Li}}) = 7.635$. This corresponds to 2.26 t \ce{Na2CO3} per tonne \ce{Li2CO3} produced, consistent with the ratio described in the flowsheet of the Silver Peak disclosure \cite{srkconsultingSECTechnicalReport2023a}.
 
\medskip
\noindent\textit{Sulfate pathway.}\quad
 
The primary input is \texttt{pls\_post\_polishing} in liters, with the step constituent-driven on the Li constituent and IPCF of 1.0 by definition. The input and step CCFs are pathway-specific, reflecting the different clarified liquors the two sulfate-route projects deliver to precipitation. Both take the form $\rho_\text{PLS}/10^6$. For Jianxiawo this is $1.05 \times 10^{-6}$ kg Li per (L$\cdot$ppm) on a 7.5 wt\% (81.0~g/L) TDS solution basis, and for Thacker Pass $1.115 \times 10^{-6}$ on a 16.8 wt\% (187.3 g/L) basis, in both cases carried unchanged from the upstream clarification or CCD circuit.

The precipitation step combines the inputs with soda ash to form lithium carbonate compounds, with a POCF of $10^{-3}$ (L to \si{m^3}). The non-crystallized liquor leaves the step as a \texttt{mother\_liquor} coproduct with a conversion factor (set at runtime on both branches) as the basis volume less the forwarded crystal slurry, $10^{-3} - y \, c_\text{in}\,\text{CCF}_\text{in} / (c_\text{out}\,\text{CCF}_\text{out})$, tracking both the upstream grade and CCF pair. For Thacker Pass, this is $9.67 \times 10^{-4}$ \si{m^3} per liter of PLS basis or approximately 8.8 million \si{m^3/yr}, and for Jianxiawo it is $9.80 \times 10^{-4}$, approximately 10.0 million \si{m^3/yr}. The sink differs by project. Thacker Pass operates under a zero-discharge water pollution control permit, so the liquor is routed to \texttt{process\_water\_recycle} at zero cost and the routing closes the water balance only. Jianxiawo carries no such obligation, so its liquor is routed to \texttt{wastewater\_discharge\_cn}. In neither case is lithium assumed to be extractable from the stream, with it carrying the lithium not recovered at precipitation (0.098 \si{kg~Li/m^3} at Jianxiawo and 0.045 \si{kg~Li/m^3} at Thacker Pass). Because the step is constituent-driven on Li, the yield is the lithium recovery and the liquor carries exactly the complement. The yield rate differs by project and is set from the \texttt{Carbonate Recovery Factor} declared per project. Thacker Pass takes 0.97, scaled to reach roughly the same amount of losses reported in the 2024 disclosure~\cite{sgscanadaincSK1300Technical2024}. Jianxiawo takes 0.90, derived from lab-reported recovery rates for sulfate liquor solutions~\cite{liuProcessAnalysisStudy2021}, since no project-specific recovery is disclosed. 
 
Soda ash is dosed at 16.55 kg \ce{Na2CO3} per kg Li, targeting the Li constituent, a $2.17\times$ excess against the stoichiometric $M_{\ce{Na2CO3}}/(2M_{\ce{Li}}) = 7.635$ kg \ce{Na2CO3} per kg Li. This corresponds to 3.51 t \ce{Na2CO3} per tonne \ce{Li2CO3} produced. The excess is calibrated to the Thacker Pass disclosure~\cite{sgscanadaincSK1300Technical2024}. The higher excess relative to the chloride route is assumed reasonable justified by the elevated relative ionic strength of sulfate-based PLS. The number derived is consistent with the Thacker Pass disclosure~\cite{sgscanadaincSK1300Technical2024}; we assume it is also reasonable for a lepidolite-sourced sulfate PLS.

\medskip
\noindent\textbf{Thickening and Dewatering.}\quad

The thickening and dewatering step concentrates the \ce{Li2CO3} slurry from approximately 20 wt\% solids to approximately 70 wt\% solids via gravity settling and mechanical pressing, recovering mother liquor for recycle or discharge. The primary input is \texttt{li2co3\_slurry} at 20 wt\% solids, with input and step CCFs of $2.24 \times 10^{-4}$ kg constituent per (m$^3 \cdot$ppm), derived as $w_s \cdot \rho_\text{mix} / 10^3$ at 20 wt\% solids with \ce{Li2CO3} grain density 2.11 \si{t/m^3}: $\rho_\text{mix} = 1/(0.2/2.11 + 0.8/1.0) = 1.118$ \si{t/m^3}.
 
The output \texttt{li2co3\_cake} is at 70 wt\% solids with density $\rho_\text{mix} = 1/(0.7/2.11 + 0.3/1.0) = 1.583$ \si{t/m^3}, giving an output CCF of $0.7 \times 1.583 / 10^3 = 1.108 \times 10^{-3}$ kg constituent per (m$^3 \cdot$ppm). The POCF is $0.202$, derived as $(0.2 \times \rho_{20\%}) / (0.7 \times \rho_{70\%}) = (0.2 \times 1.118) / (0.7 \times 1.583)$, representing the volumetric contraction from 20 wt\% to 70 wt\% solids. One reagent is consumed: flocculant at 0.005 \si{kg/m^3} of slurry feed, estimated from similar lithium-based concentrate flotation circuits \cite{opokuOverviewCoarseParticle2025, tramflocPolymersWaterWastewater2013}. 
 
The mother liquor overflow is treated differently by pathway. For the chloride (brine) route, the overflow is assumed to be recycled back to the precipitation reactor or the evaporation pond system, consistent with the Silver Peak  disclosure \cite{srkconsultingSECTechnicalReport2023a}; the conversion factor is therefore set to zero and no wastewater treatment cost is incurred. For the sulfate route, the overflow leaves at a conversion factor of $\approx 0.80$ \si{m^3/m^3} of slurry feed, consistent with the calculated volume fraction removed during thickening ($\sim$0.80 \si{m^3/m^3}), and its sink is set per project in the same way as the mother liquor. Thacker Pass operates under a zero-discharge water pollution control permit, with water recycling being a central design feature of that flowsheet~\cite{sgscanadaincSK1300Technical2024}, so the overflow is routed to \texttt{process\_water\_recycle} at zero cost. That routing closes the water balance only: the stream is filtrate from a lithium carbonate slurry rather than clean water, so no credit is taken against any make-up demand. Jianxiawo carries no such obligation, so its overflow is routed to \texttt{wastewater\_discharge\_cn}.

This step uses the \texttt{mechanical\_dewatering} machine block (Appendix \ref{app:04-methods-blockdefs-dewater}).

\medskip
\noindent\textbf{Washing and purification.}\quad

The washing and purification step removes residual mother liquor and soluble impurities from the \ce{Li2CO3} filter cake by counter-washing with process water, improving product purity. The primary input is \texttt{li2co3\_cake} at 70 wt\% solids, with input, step, and output CCFs all equal to $1.108 \times 10^{-3}$ kg constituent per (m$^3 \cdot$ppm), consistent with the thickening output. The POCF and yield rate are 1.0; washing does not change the cake volume or lithium content materially at this stage.
 
Process water is consumed as a wash input; the dosage is a rough estimate. The wash liquor is handled by pathway: for the chloride route it is recycled back to the precipitation reactor or pond system per the Silver Peak environmental disclosure~\cite{srkconsultingSECTechnicalReport2023a}, with conversion factor zero; for the sulfate route, it leaves at $0.1$ \si{m^3/m^3} of cake, reflecting the pass-through of wash water after contact with the cake, with the sink set per project: \texttt{process\_water\_recycle} at Thacker Pass under the zero-discharge permit, \texttt{wastewater\_discharge\_cn} at Jianxiawo. As with the thickener overflow, the Thacker Pass routing is a zero-cost closure term and is not credited to any make-up demand.

This step uses the \texttt{wash\_purify\_combo} machine block (Appendix \ref{app:04-methods-blockdefs-washpurify}).

 
\medskip
\noindent\textbf{Drying and packaging.}\quad

The drying and packaging step removes residual moisture from the washed \ce{Li2CO3} cake and prepares the final battery-grade product. The primary input is \texttt{lithium\_carbonate\_cake} in \si{m^3} at 70 wt\% solids, converted to the step basis in tonnes via an IPCF of $1.583$ \si{t/m^3}, equal to the cake density $\rho_\text{mix} = 1/(0.7/2.11 + 0.3/1.0)$; this represents the theoretical maximum density of a 70 wt\% \ce{Li2CO3} cake and likely slightly overestimates the actual bulk density in practice. The input CCF is $1.108 \times 10^{-3}$ kg constituent per (m$^3 \cdot$ppm), consistent with the upstream steps. The step CCF is $7.0 \times 10^{-4}$ kg constituent per (tonne$\cdot$ppm), derived as $w_s / 10^3 = 0.7 / 10^3$, reflecting the mass fraction of \ce{Li2CO3} in the cake on a tonne basis.
 
The output \texttt{lithium\_carbonate\_product} is packaged product in tonnes at the $\leq$0.4 wt\% moisture specification, with a POCF of $0.7028$, the product mass per tonne of 70 wt\% cake at that specification ($0.7/0.996 = 0.7028$), and an output CCF of $9.96 \times 10^{-4}$ kg constituent per (tonne$\cdot$ppm), the 996 kg of \ce{Li2CO3} contained in each tonne of shipped product. Closure on the cake mass basis is exact ($0.7028 + 0.2972 = 1.000$), with the balance leaving as evaporated moisture.

This step uses the \texttt{drying\_packaging\_generic} machine block (Appendix \ref{app:04-methods-blockdefs-drypack}).

\paragraph{Tailings and Residue Management}\label{app:04-methods-facilitychoices-tailings}\leavevmode\par

Tailings and residue management are key contributors to the cost and environmental impacts of minerals extraction. We model this step differentially across each pathway type, given spodumene, lepidolite, clay, and brine-hosted deposits each have their own typical mineralizations and byproducts \citep{khakmardanComparativeLifeCycle2023, chordiaLifeCycleEnvironmental2022, setuEnvironmentalImpactAssessment2026, royLithiumClayAssessing2024}. Tailings management practices vary across projects, ranging from conventional slurry impoundments to filtered dry-stack systems, and depend on factors such as water availability, seismic risk, climatic conditions, and regulatory constraints. Ore-based processing generates solid residues from comminution, beneficiation, and leaching operations. These residues must be dewatered, transported, and permanently stored in accordance with site-specific geotechnical and environmental requirements. Furthermore, in the production of lithium from pregnant leach solutions, additional residues are generated, which need to be managed. 

Because tailings handling manages coproducts rather than primary outputs, it is not modeled as a machine block within the \texttt{ProductionStep} paradigm. Each tailings handling block is instead registered as a \textit{sink} carrying a small set of physical declarations, with every coproduct row on pathway sheets naming the sink its stream reports to. For each tailings sink, it is defined by one material class, one placement method, and one jurisdiction. Sinks are named \texttt{\{material\}\_\{method\}\_\{jurisdiction\}\_\{s\_in\}}, where the trailing number is the solids fraction at which the stream leaves its production step.

Each sink declares the solids fraction at which the stream arrives ($s_{\text{in}}$), an ordered mechanical dewatering train, the solids fraction reached in place after consolidation or air-drying ($s_{\text{final}}$), the grain density $\rho_s$, the placement method, the jurisdiction, and $x$, the mass fraction of the dry solids that is gelatinous metal hydroxide. The train is a list of $(\text{mechanism}, s_{\text{out}})$ pairs drawn from \textit{thicken}, \textit{vacuum}, and \textit{press}, and may be empty where the upstream flowsheet already delivers the stream at a solids fraction where the sink does not process further. Where the emitting machine block contains no separation unit, the sink additionally declares $s_{\text{feed}}$, the solids fraction at which the stream physically leaves the step; the train then begins at the density $s_{\text{feed}}$ while $s_{\text{in}}$ remains the basis on which the routed volume and the unit cost are both expressed. This keeps the dewatering duty priced at the machine that performs it without disturbing the volumetric closure of the emitting step.

Placement is one of \textit{impound} (hydraulic deposition behind containment), \textit{stack} (spread, compact, no containment), or \textit{offsite\_gate} (a per-wet-tonne tariff with no on-site storage term). As hydroxides are gelatinous, they cannot achieve the same dewatered solids fractions as crystalline tailings, and thus $x$ is used to determine the extent of dewatering that can be applied to hydroxide mixtures.  The attainable ceiling blends on retained water, since each fraction of the solids holds its own:

\begin{equation}
    m_w^{\text{ceil}}(x) = x\,m_w^{\text{hyd}} + (1-x)\,m_w^{\text{cry}},
    \qquad
    s_{\text{ceil}} = \frac{1}{1 + m_w^{\text{ceil}}},
    \qquad
    m_w(s) = \frac{1-s}{s}
    \label{eq:c-tailings-ceiling}
\end{equation}

\noindent A declared $s_{\text{out}}$ exceeding $s_{\text{ceil}}$ for its mechanism and material raises an error rather than being priced, so a stream cannot be dewatered in the model by a machine that could not do it. As noted in \ref{tab:c-tailings-stages}, the pure-material ceilings are 0.55, 0.75, and 0.85 for thickening, vacuum filtration, and pressure filtration of free-filtering mineral solids \cite{crystalFilterPressedDryStacking2018, klohncrippenbergerStudyTailingsManagement2017}, and 0.15, 0.25, and 0.55 for gelatinous hydroxide \cite{usepaBiosolidsTechnologyFact2000, usepaBiosolidsTechSheetPlatePress2000, usepaBiosolidsTechnologyFact2003}. The unit handling cost, in dollars per tonne of dry solids, is then:

\begin{equation}
    c_i = \underbrace{\sum_{k \in \text{train}} \big( C_{m,k} + c_{w,k}\,\Delta W_k \big)}_{\text{dewatering}}
        + \underbrace{a_r}_{\text{placement}}
        + \underbrace{t_r}_{\text{facility operation}}
        + \underbrace{\pi_j / \rho_{\text{dry}}}_{\text{storage capacity}}
        + p_e \bigg( \sum_{k \in \text{train}} \big( E_{m,k} + e_{w,k}\,\Delta W_k \big) + b_r \bigg)
    \label{eq:c-tailings}
\end{equation}

\noindent where $\Delta W_k = 1/s_{k-1} - 1/s_k$ is the water removed by stage $k$ in tonnes per tonne of dry solids with $s_0= s_{\text{feed}}$ where declared and $s_0=s_{\text{in}}$ otherwise, $(C_{m,k}, E_{m,k})$ are the cost and electricity charged per tonne of dry solids passing through the machine, $(c_{w,k}, e_{w,k})$ are the cost and electricity charged per tonne of water it removes, $(a_r, b_r)$ are the placement cost and electricity intensity for method $r$, $t_r$ is the facility operating cost for that method, $\pi_j$ is the cost of purchasing a cubic meter of storage capacity in jurisdiction $j$, $\rho_{\text{dry}}$ is the in-place dry solids density in \si{t/m^3}, and $p_e$ is the facility-specific electricity price. On the \textit{offsite\_gate} route the storage and facility-operation terms are replaced by $g_j / s_{\text{out}}$, converting a gate fee quoted per wet tonne to a dry-tonne basis at the train's exit solids fraction. The total annual cost is the coproduct volume routed to the sink multiplied by $c_i \cdot \delta(s_{\text{in}}, \rho_s)$, with $\delta(s, \rho_s) = s / \big(s/\rho_s + (1-s)\big)$ the dry solids carried per cubic meter of the stream as it leaves the step. Because the routed volume and the unit cost share the $s_{\text{in}}$ basis, unit costs expressed per cubic meter of stream are not comparable across sinks arriving at different solids fractions; per tonne of dry solids and per cubic meter of stored volume are the comparable units, and both are reported below.

The two-term stage cost is what makes the model path-dependent in the right way. Equipment is sized on the dry solids it must handle per hour (e.g. settling area, plate area, cycle count) and that capacity is paid for regardless of how wet the feed is; this is $C_m$. What the machine then does is remove water, and the marginal cost of that removal rises as the residual water becomes more tightly held; this is $c_w$. Thus, under the present form the cost rises monotonically with the target and falls as the feed arrives drier.

Specific energy is charged per stage rather than per tonne of water alone because the thermodynamic minimum is not the binding constraint. Expelling filtrate against a 1~MPa differential requires theoretical minimum of 0.278~kWh per tonne of water, roughly two orders of magnitude below observed plate-and-frame duty; the balance is cycle overhead (hydraulics, cloth wash, blow-down, dead time), which scales with machine cycles and therefore with dry-solids throughput. The same comparison explains why no thermal drying stage is modeled: evaporating water costs 628~kWh/t of water against roughly 100 kWh/t dry solids (DS) mechanically for the equivalent duty, so residual water above the pressure-filtration ceiling is left to evaporate at no energy cost.

Stage coefficients are carried separately for free-filtering mineral solids and for gelatinous hydroxide, and are mass-blended at $x$ for mixed streams. Table~\ref{tab:c-tailings-stages} gives the attainable ceiling and the four coefficients for each mechanism and material.

\begin{table}[htbp]
\centering
\small
\begin{tabular}{llrrrrr}
\toprule
& & & \multicolumn{2}{c}{Cost} & \multicolumn{2}{c}{Electricity} \\
\cmidrule(lr){4-5}\cmidrule(lr){6-7}
Mechanism & Material & $s_{\text{ceil}}$ & $C_m$ & $c_w$ & $E_m$ & $e_w$ \\
& & & (\$/t DS) & (\$/t \ce{H2O}) & (kWh/t DS) & (kWh/t \ce{H2O}) \\
\midrule
Thickening          & mineral   & 0.55 & 0.10 & 0.03 & 0.40 & 0.15 \\
Vacuum filtration   & mineral   & 0.75 & 0.30 & 0.30 & 0.80 & 1.20 \\
Pressure filtration & mineral   & 0.85 & 0.50 & 1.20 & 1.00 & 2.00 \\
\midrule
Gravity thickening  & hydroxide & 0.15 & 0.20 & 0.50 & 2.00 & 3.00 \\
Belt pressing       & hydroxide & 0.25 & 2.00 & 9.50 & 6.00 & 6.00 \\
Plate-and-frame     & hydroxide & 0.55 & 3.00 & 8.50 & 15.00 & 30.00 \\
\bottomrule
\end{tabular}
\caption{Attainable solids ceilings and stage cost and electricity coefficients by mechanism and material. Ceilings for mixed streams follow Equation~\ref{eq:c-tailings-ceiling}; coefficients are mass-blended at $x$.}
\label{tab:c-tailings-stages}
\end{table}


Two patterns in the table drive the model's behavior. Within each material, the water-basis coefficient $c_w$ rises by roughly an order of magnitude from gravity settling to pressure filtration, reflecting the increasing difficulty of removing progressively more tightly held water. Between materials, every coefficient is higher for hydroxide, most markedly for the pressing stages, where the compressibility of the cake governs both cycle time and reagent demand.

The mineral coefficients are calibrated so that the assembled trains reproduce published whole-system energy intensities: a thickened circuit returns 0.85 kWh/t DS (dry solids) and a thickened-plus-filtered circuit 3.16 kWh/t DS, against 1.63 and 2.65 kWh/t DS reported by \citet{carneiroConceptualCostComparison2018}, whose totals additionally include deposition and are charged separately here as placement. The hydroxide water-basis coefficients are not free parameters: inverting the cumulative U.S. EPA biosolids cost ladder gives a belt-press increment of $27.0/2.667 = 10.1$ and a plate-and-frame increment of $21.6/2.182 = 9.9$ dollars per tonne of water~\citep{usepaBiosolidsTechnologyFact2000, usepaBiosolidsTechSheetPlatePress2000}, and that two mechanically distinct machine classes price water almost identically is the reason a two-parameter form fits the published ladder at all. The plate-and-frame electricity is consistent with the approximately 100 kWh/t DS reported at $s = 0.5$ by \citet{smithEnergyConventionalWater2017}.

Mass-blending the coefficients at $x$ understates the duty for mixed streams, since the compressible fraction also degrades the filterability of the crystalline fraction; the blended value is therefore a floor. The dewatering and placement electricity is passed to the emissions accounting framework as a sink impact factor and booked to Scope 2. Stack placement carries no electrical duty, so sinks that perform no mechanical dewatering contribute nothing to this term.

Storage is costed on the volumetric capacity a stream consumes, rather than as a flat tonnage adder, and at the density the material reaches once in place. The quantity $\pi_j$ is the price of a discrete tranche of storage capacity, which a project must purchase repeatedly as the facility fills, rather than a one-time balance-sheet entry. For Jiangxi, $\pi = 2.82$~\$/\si{m^3} is taken from the first phase of Silvercorp's Ying tailings storage facility No.~3, which cost USD~28M and carries approximately 9.92~\si{Mm^3} of capacity~\cite{silvercorpmetalsinc.TailingsStorageFacility2025}; that facility is a lead--zinc--silver flotation impoundment in Henan, applied here to Jiangxi lepidolite as a Chinese analogue with publicly available data. For Nevada, $\pi = 0.53$~\$/\si{m^3} is derived from the Thacker Pass disclosure: sustaining capital of \$561.7M expands the clay tailings filter stack by 13.0~\si{Mm^2} of footprint over the life of mine, and the facility stores 1,124.8~\si{Mm^3} across a total footprint of 13.836~\si{Mm^2} including the initial 83.6~ha, giving $(561.7/13.0)/(1,124.8/13.836) = 0.53$~\$/\si{m^3}~\cite{sgscanadaincSK1300Technical2024}. The 5.3$\times$ ratio between the two is physically coherent: an impoundment buys a dam, a liner, underdrainage and a decant structure, whereas a dry stack buys footprint preparation.

Placement is charged at \$0.43/t dry solids (DS) with 0.8 kWh/t dry solids (DS) for slurry deposition. For stacking, a single jurisdiction-neutral figure of \$1.73/t DS is applied across all projects, derived by apportioning the \$2.00--3.50/t DS all-in dry-stack operating envelope reported by \citet{crystalFilterPressedDryStacking2018} using the 63\% transport-and-placement share of filtered-tailings operating cost reported by \citet{carneiroConceptualCostComparison2018}, giving a range of \$1.26--2.21/t DS. Facility operating costs, such as surveillance, decant and seepage pumping, water management, and embankment raising, are charged at \$0.43/t DS for impoundments and \$0.055/t DS for stacks, obtained by differencing the same two sources; the impoundment value is an order-of-magnitude placeholder with a derivation attached. Waste rock piles pay placement but not facility operation costs, being dumped rather than enclosed in engineered containment. Disposal gate fees for hazardous material in China are drawn from China's National Development and Reform Commission for hazardous secure landfill in Jiangxi province~\cite{nationaldevelopmentandreformcommission2024}.

The water reclaimed from the model is generally routed towards recycled process water, assuming that it does not contain contaminants that would make recycling not possible. Water reclaimed at a sink in excess of the recycled demand of the circuit it serves would have to be discharged, but no such surplus arises in any modeled project. At Jianxiawo the beneficiation make-up is sized as the available return, so mechanical reclaim and recycled draw are equal by construction on each side. At Thacker Pass, the recycled draws exceed the returns admissible to them once the constituent-bearing streams routed to \texttt{process\_water\_recycle} are excluded, so there is no surplus to route. At Silver Peak, the water returned from the magnesium cake press is routed to the pond circuit and is not carried as a separate stream. Of note, the water released as placed tailings settle from their as-placed solids fraction toward $s_{\text{final}}$ (consolidation drainage) is treated as a loss rather than a reclaim, being too diffuse across a deposit to collect. The water balance therefore credits the circuit only at the train's exit solids fraction, and $s_{\text{final}}$ enters the cost only through the in-place density where one is disclosed.


\medskip
\noindent\textbf{Type 1: Evaporative Residue (Brines)}\quad

Waste streams from brine production primarily consist of salts generated through precipitation during the evaporative phase, and brine-derived chemical residues generated during concentration, chemical treatment, and selective lithium extraction. Solid volumes are from evaporation low relative to ore-based pathways, but still significant given the volume of brine that needs to be processed per tonne output, with a strong dependence on brine composition, especially the ratio of Magnesium to Lithium in the feed. 

Firstly, as brine is concentrated across a series of solar ponds, major salts precipitate sequentially in approximate order of decreasing solubility, and are harvested periodically from the pond floors \citep{flexerLithiumRecoveryBrines2018, garrettHandbookLithiumNatural2004, becerraTechnicalReportSummary2024}. Depending on the brine composition, this order is typically halite (\ce{NaCl}), gypsum (\ce{CaSO4 * 2H2O}), sylvite (\ce{KCl}), and carnallite (\ce{KCl * MgCl2 * 6H2O}). These evaporation salt precipitates are physically stable and chemically non-reactive; at operations they are typically restacked adjacent to ponds, and thus incur primarily mechanical haulage costs. We thus model these as primary maintenance costs of the evaporation ponds, and do not explicitly model them using this tailings model.

The chemical residues, primarily precipitated \ce{Mg(OH)2} and \ce{CaCO3}, leave the batch treatment circuit at $s = 0.35$ and are routed to \texttt{mgcake\_stack\_nv\_35}, an on-site nonhazardous stack~\cite{srkconsultingSECTechnicalReport2023a}. Roughly half the dry solids are assumed to be gelatinous \ce{Mg(OH)2}, so $x = 0.50$, which places the stream's thickening ceiling at $s = 0.236$ and its vacuum ceiling at $s = 0.375$, both below the solids fraction at which it arrives. The mechanism is therefore selected by the material rather than declared: the sink buys pressure filtration only, taking the stream from $s = 0.35$ to $s = 0.50$ against a blended press ceiling of 0.668. That stage removes 0.857 t of water per tonne of dry solids. The stream is assumed to consolidate to $s = 0.65$ in place, which is not charged.

\medskip
\noindent\textbf{Type 2: Leach Residue (Clay)}\quad

Clay processing generates three solid waste streams. Firstly, the coarse gangue rejected at classification leaves at $s = 0.20$ and is routed to \texttt{gangue\_stack\_nv\_20}, where it is dry-stacked; the water released is reclaimed and credited against the beneficiation circuit's raw make-up at classification. Being a free-filtering silicate sand, $x = 0$, and reaching a stackable 85 wt\% solids from a 20 wt\% slurry requires two machines in series: a thickener to $s = 0.55$, at the ceiling of gravity dewatering, followed by a pressure filter to $s = 0.85$. The two stages remove 3.18 and 0.64~t of water per tonne of dry solids respectively, and cost \$0.20 and \$1.27/t DS, so that 83\% of the water leaves in the cheap stage while 87\% of the cost is incurred in the expensive one. Regardless, this is a necessary process given the strict water usage limits set upon the Thacker Pass project \cite{lithiumnevadacorp.FactSheetThacker2022}. 


Secondly, the impurity sludge from neutralization and impurity removal is modeled as a filter cake at $s = 0.50$ and routed to \texttt{sludge\_stack\_nv\_50}, an on-site lined containment facility rather than off-site disposal \cite{sgscanadaincSK1300Technical2024}. Its composition follows from the reagent doses at the step: the 71.4 \si{kg/m^3} lime dose precipitates 215.1 \si{kg/m^3} of gypsum, 74\% of the 292.0 \si{kg/m^3} of solids formed, while the magnesium removed from the liquor gives 74.3 \si{kg/m^3} of \ce{Mg(OH)2} and the soda ash and flocculant contribute the remaining 2.6 \si{kg/m^3}. Gypsum is coarse and free-filtering while the rest is not, giving $x = 0.26$ and a blended pressure-filtration ceiling of 0.744. The \texttt{staged\_neutralization} block that emits this stream is an agitated reactor train and contains no separation unit, so the sludge leaves the step as a single slurry at $s_{\text{feed}}\approx0.2$, obtained from the step's own mass balance as 292.0 kg of precipitate in 1,466 kg of slurry per \si{m^3} of feed solution. The sink therefore carries the separation itself, thickening to $s=0.30$ and then pressing to $s=0.50$, below the blended ceilings of 0.325 and 0.744 respectively. The two stages remove 1.69 and 1.33 t of water per tonne of dry solids and cost \$0.38 and \$5.28/t dry solids (DS), requiring 2.3 and 17.0 \si{kWh/t} DS respectively. The declared 0.50 remains the fraction at which the cake is placed and stored, so the volume routed from the step and the water credited back to the plant are both unchanged; what the declaration adds is the duty of producing that cake, previously carried nowhere in the flowsheet.

Lastly, the coarse CCD filter cake is composed of spent claystone, carrying little residual lithium or reagents, and it leaves the CCD step as thickened underflow at $s_{\text{feed}} = 0.33$ and is pressed to $s = 0.62$, the solids fraction specified by the Nevada water permit~\cite{lithiumnevadacorp.FactSheetThacker2022}. The press is specified in the \texttt{multi\_stage\_ccd\_thickening} machine block, but the duty is charged at \texttt{claycake\_stack\_nv\_62} so that it scales on the water removed. Placement and storage occur at the in-place dry density of 0.914 \si{t/m^3} reported in the permit. That density implies a void ratio of 1.84 and a degree of saturation of 65\%, which is where filtered tailings usually sit, as they are placed unsaturated at somewhere between 50 and 75\% saturation~\citep{crystalFilterPressedDryStacking2018}. All told this stream costs \$4.91/t dry solids (DS) to handle, above the coarse gangue at \$3.83 and well below the impurity sludge at \$9.94, but it is the largest of the three by tonnage, at 42.96 t DS per tonne of lithium carbonate product.

\medskip
\noindent\textbf{Type 3: Complex Reactive Ore Residue (Lepidolite)}\quad

Lepidolite processing generates three residue streams with distinct handling requirements. The impurity sludge from the clarification step --- dominated by calcium carbonate from polishing of reagent-derived calcium --- is classified as hazardous waste under Chinese regulations due to elevated concentrations of thallium, rubidium, and cesium~\cite{sulcaLithiumPricesWhere2024, fisherScrutinisingLithiumTechnology2022}. This stream is modeled at $s=0.50$, and routed to an off-site hazardous secure landfill with a gate fee of \$280/t wet (\textasciitilde\$560/t dry solids (DS) at this solids fraction) \cite{nationaldevelopmentandreformcommission2024}. The hazardous classification follows from rubidium and cesium arriving with the liquor retained in the cake, as well as thallium precipitated into the cake by the sulfide addition at neutralization. This stream is routed to \texttt{sludge\_gate\_cn\_50}. It is entirely gelatinous hydroxide ($x = 1$) arriving at $s = 0.50$, marginally below the 0.55 ceiling of plate-and-frame pressing on such a material, so the sink declares an empty train; the press producing a 50 wt\% hydroxide cake is not priced in this sink and its cost, which the dewatering ladder would put among the highest in the model, is carried nowhere in the flowsheet.

The classification slimes (\texttt{classification\_slimes}) and the flotation tailings (\texttt{flotation\_tailings}) are both non-hazardous gangue rejects, each a dilute slurry at approximately $s = 0.30$, and both are routed to the same on-site conventional wet tailings storage facility (sink \texttt{gangue\_impound\_cn\_30}), which does not require lining under Chinese standards for non-hazardous lepidolite gangue. The sink buys a thickener to $s = 0.50$, removing 1.33 t of water per tonne of dry solids.

Finally, the water leach residue is filtered and washed within the leach step and arrives at $s = 0.65$, routed to on-site dry stacking as \texttt{slag\_stack\_cn\_65}. The sink declares $s_\text{in} = s_\text{final} = 0.65$ and an empty train, and pays only placement, facility operation, and storage. The ground-concentrate mill discharge is routed to \texttt{conc\_stack\_cn\_60} on the same basis. Fluoride in the leach residue is stabilized as \ce{CaF2} by the 50\% lime excess dosed at roasting, and residues from analogous calcium-sulfate roast--water leach processes meet general solid-waste standards rather than hazardous ones, so no gate fee is applied.

\paragraph{Sulfuric Acid Plant}\label{app:04-methods-facilitychoices-sap}\leavevmode\par

This pathway describes an on-site production facility for sulfuric acid from elemental sulfur, which is included in the design for the Thacker Pass site \citep{sgscanadaincSK1300Technical2024}. Two primary reasons motivate the on-site manufacturing of sulfuric acid. Firstly, the very high quantity of sulfuric acid required (0.49 tonnes per tonnes solid in the leaching process) would require many shipments of possibly hazardous, corrosive acid to a remote Nevada processing location, leading to significantly increased costs and reliance on external sulfuric acid sources; shipment of sulfur by itself is not only a smaller physical quantity but also avoids such hazards. Furthermore, sulfur burning and conversion is a very exothermic process, with recoverable high-grade heat that generates a significant amount of on-site electricity.

The plant used in this model is follows a sulfur-burning double-contact double-absorption train flowsheet using the equipment and energy balances of \citet{davenportSulfuricAcidManufacture2006}, which is based on a 4,400 t/d \ce{H2SO4} train system. It is decomposed into two production steps, split at the point where the carrier changes from a gas stream to liquid acid:

\begin{description}[noitemsep, topsep=0pt]
    \item[\textbullet\enskip\texttt{Sulfur burning:}] Sulfur combustion, catalytic \ce{SO2} conversion, waste-heat recovery and power generation.
    \item[\textbullet\enskip\texttt{Acid absorption:}] \ce{SO3} absorption into circulating acid and dilution to product strength.
\end{description}

Constituents on the \texttt{Sulfuric Acid Plant} sheet are declared as mass fractions rather than using the PPM convention used elsewhere in the model, as the tracked material is the bulk of the stream at every stage with no constituent parts. 

\medskip
\noindent\textbf{Sulfur Burning, Conversion and Power Recovery.}\quad

The step basis is \texttt{burner\_throughput} in kilograms of sulfur burned, chosen because every downstream quantity in the plant (e.g. acid mass, reaction heat, gas volume, catalyst charge) can be determined proportionally to the input sulfur quantity. The step evaluates
\begin{equation}
    \ce{S + O2 -> SO2}; \qquad \ce{SO2 + 1/2 O2 ->[V2O5] SO3}
\end{equation}
over a four-bed 3--1 double-contact converter at 99.9\% overall conversion \cite{davenportSulfuricAcidManufacture2006}. This step uses the \texttt{sulfur\_burner} machine block (Appendix \ref{app:04-methods-blockdefs-sulfurburner}).

Elemental sulfur is dosed as a reagent at 1.0 kg S per kilogram of \texttt{burner\_throughput} dosed against 1 unit of combustion air as the primary input; this is a modeling convenience to allow us to model the sulfur with a purchase price and embodied impact factor. Vanadium pentoxide catalyst is dosed at $9.354 \times 10^{-5}$ kg per kilogram of \texttt{burner\_throughput}, derived from a four-bed charge of 690~\si{m^3} at 550 \si{kg/m^3}, replaced at 10\% per year against an annual burn of 60,000 kg S/h over 6,762 estimated operating hours, where the 60,000 kg S/h is roughly approximated from the 4,400 t/d \ce{H2SO4} design target. Air and stack gas are outside the tracked mass balance.

The output is \texttt{so3\_gas}, forwarded to acid absorption. The conversion factor from step basis to output units is
\[
    \frac{M_{\ce{SO3}}}{M_{\ce{S}}} = \frac{80.064}{32.065} = 2.4968
\]
kilograms of \ce{SO3} per kilogram of sulfur burned, and the output carries sulfur as a constituent at a mass fraction of $M_{\ce{S}}/M_{\ce{SO3}} = 0.4005$. The yield rate is 0.999\cite{davenportSulfuricAcidManufacture2006}, and the complementary 0.001 leaves as a \texttt{stack\_gas} coproduct routed to the \texttt{atmosphere} sink.

\medskip
\noindent\textbf{Acid Absorption.}\quad
 
The step basis is \texttt{acid\_plant\_throughput} in kilograms of \ce{H2SO4} at 100 wt\%, and the step evaluates
\begin{equation}
    \ce{SO3 + H2O -> H2SO4}
\end{equation}
This step uses the \texttt{acid\_absorption} machine block (Appendix \ref{app:04-methods-blockdefs-acidabsorption}).
 
At this step, the \ce{SO3} entering absorption is converted to the sulfuric acid basis at $M_{\ce{H2SO4}}/M_{\ce{SO3}} = 1.2250$, with sulfur retained as the tracked constituent at the same 0.4005 mass fraction it carried out of the burner. The output \texttt{sulfuric\_acid\_product} is assumed with absorption treated as going to completion with 100\% yield. 

Absorption make-up water is dosed stoichiometrically at
\[
    \frac{M_{\ce{H2O}}}{M_{\ce{H2SO4}}} \times 10^{-3} = \frac{18.015}{98.079} \times 10^{-3} = 1.8368 \times 10^{-4}
\]
cubic metres per kilogram of product acid. This is the water chemically bound in the product and does not include tower circulation, which is assumed to be an internal recycle and carries no make-up demand.

Taken together the two carrier changes fix the plant's overall stoichiometry at $M_{\ce{S}}/M_{\ce{H2SO4}} = 0.3269$ kilograms of sulfur per kilogram of acid, which the 0.999 converter yield relaxes slightly to 0.3273. At the Thacker Pass midpoint, this produces 1.909 Mt of acid per year from 624.9 kt of sulfur, which corresponds to a demand of 28.58 t \ce{H2SO4} and 9.35 t S per tonne of lithium carbonate, an amount very close to the 9.15 t S/t LCE reported in the project disclosure \cite{sgscanadaincSK1300Technical2024}. 

\medskip
\noindent\textbf{Heat recovery and power.}\quad
 
The combustion and conversion of sulfur to sulfuric acid release approximately 4,035 MJ per tonne of acid, roughly three quarters of the total reaction enthalpy of the plant, at temperatures high enough to raise superheated steam. Absorption releases a further 1,350 MJ per tonne at approximately 100\textdegree C; this heat is low-grade and is not recovered. The recovered high-grade heat is passed through a condensing turbogenerator. The plant consumes part of that output itself, primarily in the main air blower that pushes combustion air through the furnace and converter, with the remainder offsetting grid electricity the lithium plant would otherwise purchase. The net position is booked as a negative electricity consumption on the \texttt{sulfur\_burner} block (Appendix \ref{app:04-methods-blockdefs-sulfurburner}).
 
Because the credit is booked at the local electricity price, the acid plant's net power position propagates directly into the project's cost and scope 2 emissions. At the midpoint the plant is a net exporter at 4,449~kWh per tonne of lithium carbonate, worth \$400/t \ce{Li2CO3} and $-1.42$~t \ce{CO2eq}/t. This recorded as an avoided-grid-purchase credit and not a reduction in the embodied burden of the acid itself; the two are reported separately rather than netted as they represent different life cycle scopes. 

\clearpage
\subsubsection{Machine Model Block Methodology}\label{app:04-methods-machineblocks}

\paragraph{Model Block Data Aggregation Methodology}\label{app:04-methods-machineblocks-aggregation}\leavevmode\par

This study employs a structured bottom-up aggregation methodology to translate heterogeneous sources of industry data (technical reports, specification sheets, regulatory filings, and interviews where appropriate) into internally consistent technoeconomic model blocks, with each \texttt{ProductionStep} object having various model block options. The objective is not necessarily to replicate any singular plant configuration, but to construct reasonable baselines for representative, transparent, and sensitivity-aware unit operations suitable for high-level comparative analysis. Using scaling factors or parallel production lines, these  blocks can then be sized to meet desired demand.

For each model block, two distinct but co-defined components are specified:

\begin{description}[noitemsep, topsep=0pt]
    \item[\textbullet\enskip Machine block options] Determines equipment sizing and associated capital and operational requirements.
    \item[\textbullet\enskip Input and output materials] Specifies the mass streams entering  and leaving the block.
\end{description}

These components are interconnected but decomposable: in the absence of detailed engineering...

Four major steps go into the creation of each block: (1) definition of bill of equipment and key inputs and outputs; (2) raw data collection; (3) aggregation into model blocks, including scenario development (low/mid/high);  and (4) validation and cross-checking.

\medskip
\noindent\textbf{Step 1: Define a Bill of Equipment and Key Inputs and Outputs.}\quad

The machine block is composed of a Bill of Equipment (BOE), which identifies:

\begin{description}[noitemsep, topsep=0pt]
    \item[\textbullet\enskip Primary process equipment] Singular or a key set of core process machines.
    \item[\textbullet\enskip Required secondary systems] e.g., handling, dosing, recirculation, discharge.
    \item[\textbullet\enskip Balance-of-plant auxiliary elements] e.g., piping, controls, structural and access.
\end{description}

There are three types of inputs:

\begin{description}[noitemsep, topsep=0pt]
    \item[\textbullet\enskip Primary inputs] (e.g., ore, brine, slurry) defined with a calculated flow rate or batch size and associated cycle time.
    \item[\textbullet\enskip Secondary inputs] i.e., reagents (e.g., lime, soda ash, sulfuric acid), defined with calculated dosages based on stoichiometric analyses with operational excess factors, inefficiencies, and possible recycle simplifications.
    \item[\textbullet\enskip Utilities] (e.g., process water, compressed air, natural gas) calculated as inputs that are fully consumed unless otherwise specified as a co-product.
\end{description}

There are likewise two types of outputs:

\begin{description}[noitemsep, topsep=0pt]
    \item[\textbullet\enskip Primary outputs] (e.g., excavated ore, pumped brine, dewatered slurry) the assumed result of the reaction of primary inputs, secondary inputs, and additional machine-specific processing.
    \item[\textbullet\enskip Secondary outputs] i.e., co-products (e.g., \ce{CO2} emissions, wastewater, tailings) calculated based on stoichiometric analyses.
\end{description}

As discussed in the production step formulation, primary and secondary inputs and outputs can have constituents defined. While calculation of reagent usage and co-product production are carried out in Step 3, the definition of key inputs and outputs is performed here, as they affect the engineering design and assumptions made in the model.

\medskip
\noindent\textbf{Step 2: Raw Data Collection.}\quad

Raw data are compiled from:

\begin{description}[noitemsep, topsep=0pt]
    \item[\textbullet\enskip Feasibility studies and technical reports] SEC Technical Report Summaries (S-K 1300) and NI 43-101 reports.
    \item[\textbullet\enskip Vendor specification sheets and package data]
    \item[\textbullet\enskip Standard engineering references] e.g., the Australasian Institute of Mining and Metallurgy's \textit{Cost Estimation Handbook}; \citet{petersPlantDesignEconomics1991}.
    \item[\textbullet\enskip Literature sources] where detailed data exist.
    \item[\textbullet\enskip Interviews and technical discussions] with operators, vendors, and other modelers.
\end{description}

The primary sources for each modeling block are listed at the top of each modeling block.

\medskip
\noindent\textbf{Step 3: Aggregate into Model Blocks.}\quad

Raw data are aggregated into standardized model blocks using a consistent structure:
\begin{equation}
    \text{Block output} = f(\text{equation of merit},\ \text{volume defining basis})
\end{equation}
Each model block provides a baseline for: the volume defining basis (i.e., the design variable, such as slurry flow rate, solids throughput, working volume, or other key sizing metric); core physical relationships as a function of the volume defining basis (e.g., mixing intensity, stoichiometric demand); installed capacity with oversizing margin; and CAPEX and OPEX components.

Given the heterogeneous and uncertain nature of raw data, we develop an envelope approach to estimate conservative (pessimistic), midpoint, and optimistic values for key parameters, anchored to reported values with uncertainty bounds introduced transparently. These envelopes help capture extremes of performance and encompass the possible space of models. Assumptions are stated explicitly to maintain traceability.

\medskip
\noindent\textit{CAPEX Aggregation.}\quad

While capital equipment costs are not a primary driver of this analysis given the focus on operational competitiveness, they are nonetheless important as a point of comparison against established studies, and operational maintenance costs are often sized as a percentage of capital costs. We use an anchor equipment approach, consistent with standard chemical and mineral processing practice:
\begin{equation}
    C_{\text{total}} = C_{\text{anchor}} + C_{\text{auxiliary}}
\end{equation}
where $C_{\text{anchor}}$ is primary equipment cost derived from vendor data, scaling relationships, and/or reported feasibility values, and $C_{\text{auxiliary}}$ represents explicit adders for integration and balance-of-plant elements.

Auxiliary CAPEX adders typically cover systems such as pumps and piping, structural steel and platforms, instrumentation and controls, electrical distribution, and reagent handling and dosing. Adders are applied as either explicit percentage ranges derived from standard cost breakdowns, or discrete subsystem packages when vendor costs are available. Standard default estimates are given in \ref{tab:c-capex-adders} drawn from Table 4 and Table 6 of Chapter 6 of \citet{petersPlantDesignEconomics1991}. Where possible, equipment-specific adders from \citet{petersPlantDesignEconomics1991} were used. 

\begin{table}[h]
\centering
\caption{Standard auxiliary CAPEX adder ranges by equipment class and scenario.}
\label{tab:c-capex-adders}
\begin{tabular}{lccc}
\hline
\textbf{Scenario} & \textbf{Mechanically simple,} & \textbf{Mechanically intensive,} & \textbf{Thermal systems,} \\
                  & \textbf{pump-dominated}       & \textbf{multiple-drive systems}  & \textbf{fans/ducts/controls} \\
\hline
Conservative & 50\% & 60\% & 55\% \\
Midpoint     & 35\% & 45\% & 40\% \\
Optimistic   & 20\% & 30\% & 25\% \\
\hline
\end{tabular}
\end{table}

\medskip
\noindent\textit{Maintenance Estimation.}\quad

Maintenance is estimated as an annual cost, as a percentage of scaled primary equipment cost, depending on the complexity of operation. We directly quote maintenance costs for chemical processes from \citet{petersPlantDesignEconomics1991}, but as not all of the processes modeled in this system are necessarily chemical in nature, we construct a similar table for physical processes (e.g., bulk ore crushing and processing, material movement, pumping and other auxiliaries), based on the SME Mining Engineering Handbook \citep{darlingSMEMiningEngineering2011}. 

\begin{table}[h]
\centering
\caption{Maintenance cost as \% of fixed-capital investment (annual basis)}
\label{tab:c-maintenace-fraction}
\begin{tabular}{p{7cm}ccc}
\hline
\textbf{Type of operation} & \textbf{Wages} & \textbf{Materials} & \textbf{Total} \\
\hline
Simple chemical process & 1--3 & 1--3 & 2--6 \\
Average chemical processes with normal operating conditions & 2--4 & 3--5 & 5--9 \\
Complicated chemical processes, severe corrosion, or extensive instrumentation & 3--5 & 4--6 & 7--11 \\
Materials movement only & 0.5--1 & 0 & 0.5--1 \\
Simple mechanical processes & 4 & 1--3 & 5--7 \\
\hline
\end{tabular}
\end{table}

\medskip
\noindent\textit{Labor Requirements Estimation.}\quad

Operating labor requirements are estimated using published operator-per-shift factors
and a staffing multiplier; the full method is described in
\ref{app:04-methods-machineblocks-labor}.

\medskip
\noindent\textit{Electricity Estimation.}\quad

Electrical demand for each block is decomposed analogously to CAPEX:
\begin{equation}
    P_{\text{total}} = P_{\text{anchor}} + P_{\text{auxiliary}}
\end{equation}
The maximum electrical demand for each machine is typically reported in specification sheets; load factors (LF) are estimated to approximate average power consumption during operation. Auxiliary electrical loads are estimated using an adder structure analogous to CAPEX, but with different ranges reflecting the distinct physical meaning of electrical auxiliaries. Standard default estimates are given in \ref{tab:c-elec-adders}.

\begin{table}[h]
\caption{Standard auxiliary electrical load adder ranges by equipment class and scenario.}
\label{tab:c-elec-adders}
\begin{tabular}{lccc}
\hline
\textbf{Scenario} & \textbf{Mechanically simple,} & \textbf{Mechanically intensive,} & \textbf{Thermal systems,} \\
                  & \textbf{pump-dominated}       & \textbf{multiple-drive systems}  & \textbf{fans/ducts/controls} \\
\hline
Conservative & 25\% & 55\% & 65\% \\
Midpoint     & 15\% & 40\% & 50\% \\
Optimistic   & 10\% & 25\% & 35\% \\
\hline
\end{tabular}
\end{table}

\medskip
\noindent\textit{Other Utilities.}\quad

The use of other utilities (water, compressed air, natural gas, steam, etc.) varies significantly across blocks depending on the process and equipment being modeled. Utility usage is derived from project disclosures, mass balance or thermodynamic calculations, and vendor specifications. Where appropriate, the same anchor-plus-adder approach is applied: identifying the dominant utility consumer, anchoring to reported usage, and applying auxiliary adders as needed.

\medskip
\noindent\textbf{Step 4: Validation and Cross-Checking.}\quad

Model outputs are validated through cross-comparison with published feasibility studies, academic literature, and grey literature; internal consistency checks across unit operations; sensitivity testing of dominant parameters. This iterative validation ensures that no block is isolated from industrial reality, while preserving transparency about sources of data and uncertainty. While interviews and discussions with industry professionals and other modelers informed initial design choices and some machine block designs, additional discussions could further validate the model inputs and results.

\paragraph{Methods for estimating capital costs}\leavevmode\par

Throughout this analysis, capital costs are estimated using a structured `anchor equipment + adders' methodology, consistent with standard mining engineering and chemical processing guidance (e.g. \citet{petersPlantDesignEconomics1991}, \citet{ausimmMonograph27Cost}, \citet{turtonAnalysisSynthesisDesign2009}, \citet{ulrichGuideChemicalEngineering1984}, \citet{darlingSMEMiningEngineering2011}). 

For each model block, a primary process equipment or set of primary process equipment are selected as the anchor, based on vendor specifications and published references. Installed capital costs are then estimated by applying installation factors, using optimistic, midpoint, and conservative values. Adders are included for auxiliary equipment and support subsystems that are required for operations, but not explicitly listed in vendor quotations or specifications. 

These adders typically include pumping and piping, structural steel and access, electrical and instrumentation systems, reagent handling subsystems, and other balance-of-plant requirements. For clarity and transparency, rather than embedding these costs within a singular multiplier, we explicitly define each as additional adders on top of the installed equipment cost. 

Percentage ranges were selected based on the aforementioned guidance as well as published cost breakdowns, and adapted to each specific block model context. Midpoint values aim to reflect a baseline plant layout; conservative and optimistic values represent extreme values that are believed to be possible.

\paragraph{Methods for estimating operating costs}\leavevmode\par

Beyond reagents and other chemistry-related production costs, certain operating costs are tied to machines, such as the amount of labor necessary to operate a machine, the average electricity draw, or direct fuel consumption. 

\paragraph{Method for estimating labor intensity}\label{app:04-methods-machineblocks-labor}\leavevmode\par

Operating labor requirements are estimated at the equipment level, consistent with standard chemical process engineering guidance (\citet{ulrichGuideChemicalEngineering1984}, \citet{turtonAnalysisSynthesisDesign2009}). For each model block, operator requirements per shift are assigned to each key piece of equipment or equipment train based on published operator-per-unit-per-shift values from Table 6-2 of~\citet{ulrichGuideChemicalEngineering1984}. Annual full-time equivalent (FTE) operators are then derived by applying a staffing factor of 4.5, accounting for the mostly continuous nature of chemical processes, as well as vacation and sick leave ($168 / 40 = 4.2$). Ulrich's method implicitly assumes labor intensity is invariant with capacity; where machines require scaling factors (i.e. continuous machines), we differentiate by process type. For comminution and materials handling blocks, we apply a scaling exponent of 0.5, consistent with AusIMM guidance~\cite{ausimmMonograph27Cost} for treatment plant labor. For downstream chemical processing blocks (e.g. leaching, precipitation, solids-liquid separation), we apply a lower exponent of 0.2, consistent with fluid processing guidance in \citet{petersPlantDesignEconomics1991}, reflecting the higher automation and lower solids-handling intensity of these unit operations. Maintenance, laboratory, support, and supervisory staff are excluded and treated as separate maintenance and overhead adders. 

\paragraph{Method for estimating chemical compositions}\label{app:04-methods-modelblocks-oreconc}\leavevmode\par

\medskip
\noindent\textbf{Brine Concentration Estimations.}\quad

For Silver Peak, \citet{garrettHandbookLithiumNatural2004} reports various analyses of Clayton Valley Brine. The Silver Peak disclosure indicates a concentration of 95 mg/L (Table 1-2)\cite{srkconsultingSECTechnicalReport2023a} but does not report other ionic concentrations; we use Garrett's 1996 measurements for non Li compounds, as well as reported density ($\rho=1.058$ \si{g/cm^3}) as a baseline, resulting in a concentration of approximately 90 ppm Li. While Garrett reports higher lithium concentrations, we use the number stated in the disclosure under the assumption that lithium grades in brine have declined over time, whereas other major compounds have remained generally constant. The disclosure notes that at about a concentration of 0.54 wt\% Li, the strong brine is then processed in the lithium carbonate plant. At an assumed density of $\rho=1.25$ \si{g/cm^3}, this results in a concentration of 6,750 ppm Li, or a brine concentration factor $F$ of $75 \times$. We generalize to assume this concentration factor applies to conservative species that stay in solution (Li, Mg, B, Rb), while other compounds largely precipitate out (e.g. about 95\% of Ca as calcite; Na, K, Cl, \ce{CO3}, \ce{SO4} buffered by saturation points \cite{flexerLithiumRecoveryBrines2018, heEnvironmentalLifeCycle2025}. 

\begin{table}[h]
\centering
\caption[Assumed Brine Compositions, Silver Peak]{Assumed Brine Compositions, Silver Peak \\
\small *Reported concentration, \cite{srkconsultingSECTechnicalReport2023a}; all others from \citet{garrettHandbookLithiumNatural2004} reporting from Garrett (1996).}
\label{tab:c-silverpeak-compositions}
\begin{tabular}{|l|l|c|c|c|c|c|c|c|c|c|c|}
\hline
 & & Li* & Na & K & Mg & Ca & Rb & Cl & \ce{SO4} & \ce{CO3} & B \\
\hline
\multirow{2}{*}{Raw Brine}
  & wt\% &       & 4.69   & 0.40  &       &      &      & 7.26   & 0.34  &      &      \\
\cline{2-12}
  & ppm  & 90*   & 49610  & 4232  & 190   & 450  & 21   & 76811  & 3597  & 74   & 67   \\
\hline
Strong Brine & ppm & 6750 & 140000 & 20000 & 14250 & 1688 & 1575 & 220000 & 10000 & 1000 & 5025 \\
\hline
\end{tabular}
\end{table}

\medskip
\noindent\textbf{Thacker Pass Concentration Estimations.}\quad

For the Thacker Pass hectorite clay deposit, only lithium concentration is tracked through the processing circuit; detailed multi-constituent leach feed and PLS compositions are not publicly disclosed in the S-K 1300 technical reports~\cite{sgscanadaincSK1300Technical2024, rothFeasibilityStudyNational2022}, and the reagent dosages in the generic clay flowsheet are calibrated directly to disclosed per-tonne consumption rates rather than to constituent-specific stoichiometry (see Appendix~\ref{app:04-methods-facilitychoices-clayleach}). This represents a modeling simplification relative to the Silver Peak and Jianxiawo pathways, where multi-ion brine and ore compositions are available from the literature.

The ROM ore lithium grade of 2,538 ppm is the proven and probable reserve-weighted (85 year) average feed grade reported in Table 14-6 of the 2024 S-K 1300 disclosure~\cite{sgscanadaincSK1300Technical2024}. No intermediate grade is reported, so the leach feed grade is closed from the same disclosure's life-of-mine mass balance. Table 13-4 reports 611.8 Mt of plant feed against 1,056.7 Mt of dry ore mined, so classification must reject $1-611.8/1,056.7=0.421$ of feed solids. Section 22.6.2 further reports an average lithium recovery of 92\%. The leach feed grade follows as $2,538 \times 0.92/0.579 = 4,033$ ppm. The same tables cross-check the acid intensity independently: 611.8 Mt of leach ore against 11.5 Mt of lithium carbonate produced is 53.20 t leach ore per tonne of product, which at the disclosed 490 kg \ce{H2SO4} per tonne of leach feed solids is 26.07 t \ce{H2SO4} per tonne LCE.

The lithium concentration in the pregnant leach solution at input to the impurity removal step is estimated at 1,290 ppm Li. Starting from the leach feed grade (4,033 ppm), converting to absolute Li mass per unit volume of the 34 wt\% leach slurry via the mixture density CCF, dividing by the physical CCD volume ratio of $~\sim$1.22 and by the PLS density of $~\sim$1.11 \si{t/m^3}, and applying the two constituent recoveries (92\% leach extraction and 99\% CCD wash recovery) gives 

$$\frac{4033 \times 4.757 \times 10^{-4}}{1.2208 \times 1.110} \times 10^{3} \times 0.92 \times 0.99 = 1,416 \times 0.9108 = 1,290 \text{ppm Li}$$

Neither recovery removes volume: the 8\% unleached lithium remains in the solids and the 1\% unwashed lithium in the cake moisture, and both leave the flowsheet once, in the CCD filter cake routed to \texttt{claycake\_stack\_nv\_62}.

\begin{table}[h]
\centering
\caption[Assumed Lithium Concentrations, Thacker Pass]{Assumed Lithium Concentrations, Thacker Pass \\
\small Lithium grades from \citet{sgscanadaincSK1300Technical2024} and its reported mass flows; only Li is tracked in the generic clay flowsheet. Other ionic species are not modeled due to the absence of publicly disclosed multi-constituent leach feed and PLS compositions.}
\label{tab:c-thackerpass-compositions}
\begin{tabular}{|l|c|}
\hline
 & \textbf{Li (ppm)} \\
\hline
\textbf{ROM Ore}                         & 2,538 \\
\hline
\textbf{Acid Leach Circuit Input}        & 4,033 \\
\hline
\textbf{PLS (Impurity Removal Input)}    & 1,290 \\
\hline
\end{tabular}
\end{table}

\medskip
\noindent\textbf{Lepidolite Concentration Estimations.}\quad

For the Jianxiawo lepidolite deposit, we use an initial ROM ore Li concentration of 1,858 ppm (0.40\% \ce{Li2O}), consistent with the upper end of CRU's reported range for kaolinite-rich Jiangxi lepidolite deposits (0.27--0.40\% \ce{Li2O}) and with the regulatory threshold below which local governments have historically issued kaolinite rather than lithium mining permits~\cite{zhangChinaHasDomestic2025}. Other major ionic species in ROM ore (K, Na, Mg, Ca, Si, Rb, F) are drawn from \citet{garrettHandbookLithiumNatural2004} and the peer-reviewed lepidolite mineralogy literature (e.g. \cite{gaoLithiumExtractionHard2023, parkEffectiveLithiumExtraction2023}).

During the key flotation step, the run-of-mine (ROM) ore is upgraded into an intermediate lepidolite concentrate. This concentrate yields a lithium grade of $7,576$ ppm (1.63\% \ce{Li2O}), representing an enrichment factor of approximately 4.08 $\times$ relative to the raw ROM ore ($1,858$ ppm). The concentrate is subsequently thickened, ground to a product size of approximately 50--75 $\mu$m by high-energy mechanical activation milling, and dewatered before entering the sulfate roasting step. The flotation step is critical for this pathway: at typical Jiangxi ore grades of less than 0.4\% \ce{Li2O} ($\sim$1,860 ppm Li), direct roasting of ROM material would require processing enormous reagent and energy quantities per unit of lithium product. Froth flotation using amine collectors and phosphate depressants selectively recovers lepidolite minerals into a concentrate, reducing the mass throughput to the energy-intensive roasting circuit by approximately 60\%~\cite{jiaoKeyTechnologiesDevelopment2024}. The implied mass yield of approximately 13-18\% (6-8 tonnes feed per tonne of classified slurry, depending on assumed recovery) is broadly consistent with CRU's industrial benchmark of 8 tonnes of 0.40\% \ce{Li2O} ore per tonne of 2.20\% \ce{Li2O} concentrate~\cite{fisherScrutinisingLithiumTechnology2022}, noting that our intermediate is a relatively lower grade (1.63\% vs. 2.20\% \ce{Li2O}) based on reporting from \citet{garrettHandbookLithiumNatural2004}.

After sulfate roasting at 800-950$^\circ$~C and water leaching, the pregnant leach solution (PLS) Li concentration is 909.9 ppm on a liquor basis, equivalent to 0.995 \si{kg~Li/m^3} at a solution density of 1.049 \si{t/m^3}. This concentration is derived from the leach mass balance: the 20 wt\% slurry, the roast reagent doses, the concentrate grade, and the leach yield together fix both the dissolved lithium and the dissolved load that sets the liquor density. Note that this figure is reported on the liquor carrier, and is therefore not directly comparable to the solids-basis concentrations quoted for ROM ore and flotation concentrate. The associated major species in the PLS (K, Na, Rb, F, Si, Cs) reflect the solubilization of alkali and fluoride constituents during roasting and leaching, consistent with the phase transformations reported in \citet{yanExtractionLithiumLepidolite2012} and \citet{dengSystematicReviewSulfate2026}.

\begin{table}[h]
\centering
\caption[Assumed Lepidolite Compositions]{Assumed Lepidolite Compositions, used for Jianxiawo \\
\small *Reported concentration, \citet{zhangChinaHasDomestic2025}; all others from~\cite{garrettHandbookLithiumNatural2004}.}
\label{tab:c-lepidolite-compositions}
\begin{tabular}{|l|c|c|c|c|c|c|c|c|}
\hline
 & \textbf{Li} & \textbf{K} & \textbf{Mg} & \textbf{Ca} & \textbf{Na} & \textbf{Si} & \textbf{Rb} & \textbf{F} \\
\hline
\textbf{ROM Ore*}        & 1,858  & ---    & ---  & ---  & ---    & ---     & ---    & ---    \\
\hline
\textbf{Flotation Concentrate} & 7,576 & 24,734 & 60   & 429  & 38,070 & 321,000 & 3,016  & 18,600 \\
\hline
\textbf{PLS}             & 909.9 & --- & 0    & 0    & ---  & --- & --- & --- \\
\hline
\end{tabular}
\end{table}




\clearpage
\subsubsection{Machine Model Block Definitions}\label{app:04-methods-blockdefs}


\paragraph{\texttt{brine\_pump}}\label{app:04-methods-blockdefs-brinepump}\leavevmode\par

\medskip
\noindent\textbf{Sources.}\quad

\citet{srkconsultingSECTechnicalReport2023a}, \citet{becerraTechnicalReportSummary2024}, \citet{srkconsultingSECTechnicalReport2022}, \citet{montgomeryassociatesconsultoreslimitadaSECTechnicalReport2023}, \citet{brookerSECTechnicalReport2023};
See Brine Evaporation Datasheet in Supplementary Data Spreadsheet. 

\medskip
\noindent\textbf{Process description.}\quad

The brine pump block captures extraction of lithium-bearing brine from subsurface aquifers and delivery to the evaporation pond system. The block is parameterized per well-pump unit; total system capacity is obtained by multiplying by the number of active wells. Capacity scaling is therefore discrete (additional parallel units) rather than continuous, and no size-scaling exponent is applied to primary equipment cost.

\medskip
\noindent\textbf{Equation of merit.}\quad
\[
  \dot{V}_\text{total}\;(\si{m^3/h}) = \dot{V}_\text{pump}\;(\si{m^3/h}) \times N_\text{wells}
\]
where $\dot{V}_\text{pump}$ is the flow rate of a single well-pump unit and $N_\text{wells}$ is the number of active extraction wells.

\medskip
\noindent\textbf{Bill of Equipment.}\quad

\begin{table}[h!]
\caption{Bill of equipment for the brine pump block.}
\label{tab:c-equip-evappond}
\centering
\begin{tabular}{>{\raggedright\arraybackslash}p{5.2cm}
                >{\raggedright\arraybackslash}p{4cm}
                >{\raggedright\arraybackslash}p{4.5cm}}
\toprule
\textbf{Equipment} & \textbf{Existing machine data used} & \textbf{Source(s)} \\
\midrule
Pumping package with piping, electrical, and other support & Brine pumping reported system analysis & Brine Evaporation Datasheet; company disclosures \\
\bottomrule
\end{tabular}
\end{table}

\medskip
\noindent\textbf{Sizing basis.}\quad

The base volume is set to \textbf{72,000 L/h} (72 \si{m^3/h}) per well-pump unit, derived from an average reported extraction rate of 20 L/s. Silver Peak reports individual well extraction rates ranging from 5.2 \si{m^3/h} to 187.5 \si{m^3/h}; the 20 L/s (72 \si{m^3/h}) value is used as a central representative rate for a single active well, calibrated to other reported average pump rates (see Brine Evaporation Datasheet in Supplementary Data Spreadsheet).

The block is assigned \texttt{process\_type = batch} with \texttt{batch\_cycle\_time = 1 h} and zero setup time. This is a modeling convenience: brine pumping is physically a continuous operation, but the hourly batch representation is consistent with the model's time-step structure and produces equivalent annual throughput. It does not imply any actual batch cycling of the pump.

\medskip
\noindent\textbf{CAPEX.}\quad

Capital cost per well-pump unit is derived from a linear regression of total reported well-field capital expenditure against number of wells across seven lithium brine project disclosures. The Salar de Atacama-SQM disclosure (438 wells, \$370.76M) was excluded as a clear outlier --- its scale is more than an order of magnitude larger than all other data points and would dominate the regression. The fitted relationship is:
\[
  C_\text{total}\;(\$\text{M}) = 0.903 \times N_\text{wells} + 9.07 \qquad R^2 = 0.784
\]
The slope (\$0.903M/well) represents the marginal installed cost of each additional well-pump unit. The intercept (\$9.07M) captures fixed site mobilization and infrastructure costs --- pump house, electrical switchgear, surface manifolding, and SCADA installation --- that are shared across the well field regardless of well count. In the per-unit model block, the slope is rounded to approximately \$1M/well as a central estimate; scenario variation (\$0.8--1.3M/well) spans the observed scatter in reported total-cost-per-well values across the regression dataset (\$0.89--1.91M/well), and reflects uncertainty in drilling depth, casing specification, and pump selection across the wide reported head range (90--550~m):

\begin{table}[h!]
\caption{CAPEX scenario values for the \texttt{brine\_pump} block (per well-pump unit).}
\centering
\begin{tabular}{lc}
\toprule
\textbf{Scenario} & \texttt{prim\_equip\_price\_base} (USD) \\
\midrule
Conservative & \$1,300,000 \\
Midpoint     & \$1,100,000 \\
Optimistic   & \$800,000    \\
\bottomrule
\end{tabular}
\end{table}

\noindent Equipment life is set to 20 years (\texttt{prim\_equip\_life = 20}), consistent with standard submersible pump replacement intervals in brine service. No throughput scaling exponent is applied; additional capacity is obtained by adding parallel well-pump units at the same unit cost. The fixed site infrastructure captured by the regression intercept is not separately modeled; it is implicitly absorbed into project-level CAPEX estimates for the well-field as a whole.

\medskip
\noindent\textbf{Electricity.}\quad

Electrical demand per well-pump unit is derived from the hydraulic power relation:
\[
  E \;=\; \frac{\rho g H}{\eta \cdot 3.6 \times 10^6} \quad (\si{kWh/m^3})
\]
For Silver Peak, with brine density $\rho = 1,058 \si{kg/m^3}$, pump efficiency $\eta \approx 70\%$, and a central head estimate of $H = 300 \si{m}$ (reported range: 90--550 m):
\[
  E = \frac{1,058 \times 9.81 \times 300}{0.70 \times 3.6 \times 10^6} \approx 1.235 \si{kWh/m^3}
\]
At a flow of 72 \si{m^3/h}, this gives a base motor demand of approximately $1.235 \times 72 \approx \mathbf{90 \si{kW}}$ per well-pump unit.

An additional load factor is applied to account for intra-playa brine transfer pumping and the elevated specific gravity of increasingly concentrated brine at later evaporation stages --- loads not captured by the single-well extraction head alone. The factor is applied as a multiplier on the base motor demand:

\begin{table}[h!]
\caption{Electricity scenario derivation for the \texttt{brine\_pump} block.}
\centering
\begin{tabular}{lcccc}
\toprule
\textbf{Scenario} & \textbf{Flow (\si{m^3/h})} & \textbf{Base motor (kW)} & \textbf{Playa adder} & \texttt{electricity\_base\_total} \\
\midrule
Conservative & 72 & 90 & $\times 1.3$ & \textbf{117 kWh/h} \\
Midpoint     & 72 & 90 & $\times 1.2$ & \textbf{108 kWh/h} \\
Optimistic   & 72 & 90 & $\times 1.1$ & \textbf{99 kWh/h}  \\
\bottomrule
\end{tabular}
\end{table}

\medskip
\noindent\textbf{Other utilities.}\quad

None. Natural gas, diesel, propane, cooling water, steam, and compressed air are all zero for this block.

\medskip
\noindent\textbf{Labor.}\quad

Labor is assigned per well-pump unit in FTE, reflecting the monitoring and intervention requirements of a single operating well. At Silver Peak's 6 active extraction wells, the scaling exponent of 0.5 yields a total staffing of approximately 0.10--0.15 FTE for well monitoring and pump intervention --- consistent with automated pump operation requiring only periodic site visits rather than continuous attendance.

\begin{table}[h!]
\caption{Labor base values for the \texttt{brine\_pump} block (per well-pump unit).}
\centering
\begin{tabular}{lcc}
\toprule
\textbf{Scenario} & \texttt{labor\_base} (FTE) & \textbf{Implied FTE at 6 wells} \\
\midrule
Conservative & 0.06 & $0.06 \times \sqrt{6} \approx 0.15$ \\
Midpoint     & 0.05 & $0.05 \times \sqrt{6} \approx 0.12$ \\
Optimistic   & 0.04 & $0.04 \times \sqrt{6} \approx 0.10$ \\
\bottomrule
\end{tabular}
\end{table}

\noindent The \texttt{labor\_scaling\_exponent} of 0.5 is the AusIMM default. For well monitoring, which scales sub-linearly with the number of wells (shared site visits, centralized monitoring), this should be physically reasonable.

\medskip
\noindent\textbf{Footprint.}\quad

Each well-pump unit adds a fixed 5~$\si{m^2}$ surface footprint (scaling exponent 1.0) for the wellhead, electrical switchgear, and access pad.

\medskip
\noindent\textbf{Availability.}\quad

\texttt{proc\_avail\_factor} = 0.90, with \texttt{unplanned\_downtime} = 0.05 and \texttt{scheduled\_maintenance} = 0.05 (10\% total downtime, split equally between planned and unplanned), held constant across scenarios.

\medskip
\noindent\textbf{Maintenance.}\quad

Maintenance costs apply to the pump and surface infrastructure only; wells are not practically maintainable once installed and are excluded from the recurring maintenance budget. Costs are specified via \texttt{maint\_override} as an annual fraction of CAPEX, reflecting the pump-only maintenance scope:

\begin{table}[h!]
\caption{Maintenance override fractions for the \texttt{brine\_pump} block.}
\centering
\begin{tabular}{lc}
\toprule
\textbf{Scenario} & \texttt{maint\_override} (fraction of CAPEX/yr) \\
\midrule
Conservative & 0.020 \\
Midpoint     & 0.015 \\
Optimistic   & 0.010 \\
\bottomrule
\end{tabular}
\end{table}

\medskip
\noindent\textbf{Notes and flags.}\quad

Several parameters carry meaningful uncertainty that is not fully captured in the conservative--optimistic range and should be noted in the methodology:

The central head estimate of 300 m is used as an average across Silver Peak's reported range of 90--550 m. This range is unusually wide and the distribution of active well depths is not disclosed; a more detailed well-field characterisation could substantially tighten the electricity estimate.

The playa adder (10--30\%) is intended to capture intra-playa transfer pumping and brine concentration effects but is not directly calibrated to a disclosed value. It is a modeling assumption and represents a source of systematic uncertainty in the electricity estimate.

The CAPEX regression is based on seven project disclosures after outlier exclusion, with $R^2=\sim 0.784$. The fixed intercept (\$9.07M) is not separately modeled as a site infrastructure line item; this is a limitation for small well-field configurations where the intercept represents a proportionally larger share of total cost.

The \texttt{process\_type = batch} designation with a 1-hour cycle is a modeling convenience only and has no physical meaning for this block. It should not be interpreted as implying that brine pumping operates intermittently.

\clearpage
\paragraph{\texttt{evaporation\_pond} (Lined and Unlined)}\label{app:04-methods-blockdefs-evappond}\leavevmode\par

\medskip
\noindent\textbf{Sources.}\quad

\citet{heEnvironmentalLifeCycle2025}, \citet{flemingLongrunAvailabilityMineral2024}, \citet{stampEnvironmentalImpactsTransition2012}, \citet{srkconsultingSECTechnicalReport2023a}, \citet{becerraTechnicalReportSummary2024}, \citet{srkconsultingSECTechnicalReport2022}, \citet{montgomeryassociatesconsultoreslimitadaSECTechnicalReport2023}, \citet{brookerSECTechnicalReport2023}

\medskip
\noindent\textbf{Process description.}\quad

Evaporation ponds concentrate lithium brine through solar evaporation in large, shallow, lined or unlined basins. The primary design variable is pond surface area, sized to remove water based on the aridity and weather conditions of the region in which it is installed --- hence its typical usage in high-desert environments. In operation, ponds require ongoing maintenance to repair the lining, harvest accumulated salts, and conduct regular monitoring.

\medskip
\noindent\textbf{Equation of merit.}\quad
\[
  A_\text{pond}\;(\si{m^2}) = \frac{\dot{V}_\text{brine}\;(\si{m^3/yr})}{\dot{E}\;(\si{m/yr})}
\]
where $\dot{V}_\text{brine}$ is the brine processing plant intake and $\dot{E}$ is the effective evaporation rate. The evaporation rate is estimated from the reported freshwater evaporation rate, multiplied by derating factors ($<$1) that account for the reduced vapour pressure of brine due to its salinity, and for practical efficiency losses at industrial-scale pond sizes. For example, the Sal de Vida disclosure applies a derating factor of 0.7 to the reported freshwater evaporation rate~\cite{montgomeryassociatesconsultoreslimitadaSECTechnicalReport2023}. Evaporation rates are typically reported in mm/day and range between 2.4 and 7.4 mm/day depending on location, for example, \citet{heEnvironmentalLifeCycle2025} report monthly raw (unadjusted) evaporation rates of 4.6--8.6 mm/day for Salar de Atacama, one of the driest and sunniest lithium-producing regions in the world.

Given the strong sensitivity of evaporation rate to geographic location, the specific rate reported in each project disclosure is used where available, rather than a generic default. Silver Peak reports a value of 89 cm/year on average, or 2.44 mm/day~\cite{srkconsultingSECTechnicalReport2023a}. 

\medskip
\noindent\textbf{Evaporative precipitation chemistry.}\quad

As brine concentrates through evaporation, increasing ion concentrations cause progressive precipitation of dissolved salts. With salar brines, precipitation typically occurs in approximate order of decreasing solubility: calcium compounds (calcite, gypsum --- \ce{CaSO4.2H2O}), sodium chloride (halite --- \ce{NaCl}), potassium chloride (sylvite --- \ce{KCl}), and potassium-magnesium compounds (carnallite --- \ce{KCl.MgCl2.6H2O}; bischofite --- \ce{MgCl2.6H2O}). Lithium, being highly soluble, generally remains in solution through this phase. The order and quantity of precipitation depend strongly on brine composition and therefore vary across projects.

\medskip
\noindent\textbf{Bill of Equipment.}\quad
 
\begin{table}[h!]
\caption{Bill of equipment for the evaporation pond block.}
\centering
\begin{tabular}{>{\raggedright\arraybackslash}p{5.2cm}
                >{\raggedright\arraybackslash}p{4cm}
                >{\raggedright\arraybackslash}p{4.5cm}}
\toprule
\textbf{Equipment} & \textbf{Existing machine data used} & \textbf{Source(s)} \\
\midrule
Installed ponds (earthworks, lining, civil) & N/A --- ponds only & Lithium brine evaporation pond analysis (project disclosures) \\
\bottomrule
\end{tabular}
\end{table}
 
\noindent No electricity is assumed in this step; pond circulation pumping is included in the \texttt{brine\_pump} block.
 
\medskip
\noindent\textbf{CAPEX estimation --- lined ponds.}\quad

Installed pond costs include land purchase, earthworks, and lining (HDPE or compacted clay liner). Based on analysis of capital cost data from brine production disclosures, costs vary between approximately \$12/$\si{m^2}$ and \$30/$\si{m^2}$, with one outlier at \$52/$\si{m^2}$.
 
\begin{table}[h!]
\caption{Lined pond unit CAPEX (\texttt{evaporation\_pond\_lined}).}
\label{tab:c-capex-lined}
\centering
\begin{tabular}{lc}
\toprule
\textbf{Scenario} & \texttt{prim\_equip\_price\_base} (\$/$\si{m^2}$) \\
\midrule
Conservative & \$30 \\
Midpoint     & \$18 \\
Optimistic   & \$12 \\
\bottomrule
\end{tabular}
\end{table}
 
\medskip
\noindent\textbf{CAPEX estimation --- unlined ponds.}\quad

Costs for unlined ponds are estimated from the Silver Peak 2023 disclosure, in which the expansion includes 8 new ``unlined earthen dike structures'' at a total estimated expenditure of \$48.9M. The Phase 3 pond (1 large pond, 200 acres, \$7.1M) yields a unit cost of approximately \$8.77/$\si{m^2}$, used as the central estimate.
 
\begin{table}[h!]
\caption{Unlined pond unit CAPEX (\texttt{evaporation\_pond\_unlined}).}
\label{tab:c-capex-unlined}
\centering
\begin{tabular}{lc}
\toprule
\textbf{Scenario} & \texttt{prim\_equip\_price\_base} (\$/$\si{m^2}$) \\
\midrule
Conservative & \$12 \\
Midpoint     & \$8  \\
Optimistic   & \$5  \\
\bottomrule
\end{tabular}
\end{table}
 
\medskip
\noindent\textbf{Electricity.}\quad

Zero; no electricity is consumed at this step. Pond circulation and inter-pond transfer pumping are included in the \texttt{brine\_pump} block (Appendix \ref{app:04-methods-blockdefs-brinepump}).

\medskip
\noindent\textbf{Other utilities.}\quad

None. All thermal and mechanical utilities are zero for this block.

\medskip
\noindent\textbf{Labor parameterization.}\quad

The \texttt{labor\_base} values are expressed in FTE per $\si{m^2}$ of pond area and are calibrated to the Silver Peak disclosure, which reports approximately 8--12 FTE managing salt removal and dike upkeep \cite{srkconsultingSECTechnicalReport2023a}. Because pond management is a fundamentally area-maintenance task --- involving periodic lining inspection, salt harvesting from pond floors, brine level monitoring, and dike upkeep --- labor requirements scale approximately linearly with pond area. The values are therefore derived by dividing target FTE directly by a reference pond area:
 
\[
  \texttt{labor\_base} = \frac{\text{Target FTE}}{A_\text{ref}} \qquad (\si{FTE/m^2})
\]
 
Using a reference area $A_\text{ref} = 26 \times 10^6\;\si{m^2}$ ($\approx$6,400 acres), consistent with the active evaporation pond footprint at Silver Peak, target FTE values of 10, 8, and 6 for the conservative, midpoint, and optimistic scenarios respectively yield:
 
\begin{table}[h!]
\caption{Labor base values for evaporation ponds, calibrated to Silver Peak disclosure.}
\label{tab:c-labor-base-ponds}
\centering
\begin{tabular}{lccc}
\toprule
\textbf{Scenario} & \textbf{Target FTE at $A_\text{ref}$} & \textbf{\texttt{labor\_base} (FTE/$\si{m^2}$)} & \textbf{Equivalent fraction} \\
\midrule
Conservative & 10 & $3.846 \times 10^{-7}$ & $10 / 26,000,000$ \\
Midpoint     &  8 & $3.077 \times 10^{-7}$ & $8 / 26,000,000$  \\
Optimistic   &  6 & $2.308 \times 10^{-7}$ & $6 / 26,000,000$  \\
\bottomrule
\end{tabular}
\end{table}
 
\noindent The midpoint value (8 FTE) corresponds to the lower bound of the disclosed range; the conservative value (10 FTE) sits near the midpoint of that range. The optimistic value (6 FTE) falls modestly below the disclosed range, representing a best-case labor efficiency assumption. The \texttt{labor\_scaling\_exponent} is set to 0.5, consistent with the standard template default from AusIMM processing plant guidance; however, for a passive area-maintenance task, linear scaling is the physically appropriate interpretation. Since \texttt{base\_volume} is set to 1 $\si{m^2}$, the formal expression \texttt{labor\_base} $\times A^{0.5}$ is not equivalent to the intended linear relationship, but the values above remain the correct calibration reference for the Silver Peak scale.

\medskip
\noindent\textbf{Footprint.}\quad

Footprint is not parameterized separately for this block. The pond surface area is the sizing basis itself (\texttt{base\_volume\_unit} = $\si{m^2}$), and \texttt{prim\_equip\_scaling\_exponent} = 1.0, reflecting that area scales linearly with throughput.

\medskip
\noindent\textbf{Availability.}\quad

Not applicable in the mechanical sense. Evaporation ponds are passive civil structures with no rotating equipment; the block is assigned \texttt{process\_type = batch} as a modeling convenience representing the seasonal pond management cycle. No availability derating is applied.

\medskip
\noindent\textbf{Maintenance cost fractions.}\quad

Annual maintenance is parameterized as a fraction of installed CAPEX (\texttt{opex\_fraction\_of\_capex}). Two considerations motivate the use of substantially lower fractions than the 1--3\% range commonly cited for chemical processing equipment.
 
First, evaporation ponds are passive civil structures with no rotating equipment, no chemical dosing systems, and no instrumentation networks requiring calibration or replacement. Maintenance activities consist of periodic liner inspection and targeted patch repairs (lined ponds only), salt harvesting from pond floors, dike inspection and minor earthworks, and brine level monitoring. These activities are labor-intensive but not capital-intensive relative to the total asset base.
 
Second, applying a standard chemical-plant maintenance fraction to a large pond area produces implausibly high implied costs. At Silver Peak's approximate pond area of $\sim$12 \si{km^2} and a midpoint lined CAPEX of \$18/$\si{m^2}$, a fraction of 1.5\% would imply \$3.2M/yr in maintenance expenditure. The Silver Peak disclosure notes approximately 8--12 FTE managing salt removal and dike upkeep \cite{srkconsultingSECTechnicalReport2023a}; at a loaded labor cost of roughly \$80,000/FTE/yr, the labor component alone is approximately \$640,000--\$960,000/yr. Total maintenance inclusive of materials and equipment is unlikely to exceed twice this figure --- placing a reasonable order-of-magnitude estimate at \$1.5--2M/yr, consistent with a fraction of approximately 0.6--0.9\% rather than 1.5\%. Lower fractions for unlined ponds reflect their reduced maintenance burden: the absence of a liner largely eliminates significant inspection and repair activity, leaving only dike, berm, and salt management.
 
The adopted fractions, calibrated to this cross-check, are:
 
\begin{table}[h!]
\caption{Annual maintenance cost fractions for lined and unlined evaporation ponds.}
\label{tab:c-maint-fractions-ponds}
\centering
\begin{tabular}{lcc}
\toprule
\textbf{Scenario} & \textbf{Lined (\texttt{opex\_fraction\_of\_capex})} & \textbf{Unlined (\texttt{opex\_fraction\_of\_capex})} \\
\midrule
Conservative & 0.010  & 0.005 \\
Midpoint     & 0.0075 & 0.004 \\
Optimistic   & 0.005  & 0.003 \\
\bottomrule
\end{tabular}
\end{table}
 
 
\medskip
\noindent\textbf{Key measurements and cross-checks.}\quad
\begin{itemize}[noitemsep]
  \item Effective evaporation rate (mm/day), including derating factors applied
  \item Lithium recovery efficiency through the evaporative phase (accounts for co-precipitation losses and seepage; 0.51 for Silver Peak~\cite{srkconsultingSECTechnicalReport2022})
  \item Salt accumulation rate and harvest frequency (primary driver of maintenance cost)
  \item Pond lining integrity for lined variants (leak rate, repair interval)
\end{itemize}
 
\medskip
\noindent\textbf{Notes and flags.}\quad

The evaporation rate is the dominant sizing driver and is highly site-specific; the model preserves sensitivity to this variable by accepting it as a project-specific input rather than applying a generic default. At Silver Peak, the comparatively low evaporation rate (2.44 mm/day versus 4--7 mm/day at South American salars), in addition to lower grade brines~\ref{app:04-methods-modelblocks-oreconc}, results in proportionally larger required pond area for a given production target, which is the primary reason evaporation pond costs dominate the Silver Peak cost structure.
 
The distinction between lined and unlined variants is material: lined ponds are appropriate where regulatory requirements mandate containment or where substrate permeability necessitates a barrier; unlined earthen structures are used where the subsurface is sufficiently impermeable and permitting conditions allow. The Silver Peak expansion uses unlined structures, reflecting both the site geology and the relatively benign chemistry of evaporation residues in a low-magnesium brine system.

\clearpage
\paragraph{\texttt{blasting\_generic}}\label{app:04-methods-blockdefs-blastgeneric}\leavevmode\par

\medskip
\noindent\textbf{Sources.}\quad

\citet{darlingSMEMiningEngineering2011} (Ch.\ 4.9 Cost Estimating for Surface Mines; Ch.\ 7.3 Explosives and Blasting); \citet{warrenGeorgiaLakeLithium2022}; \citet{bbainc.SECTechnicalReport2023}; Sandvik Leopard DI550 specification sheet (\citet{sandvikSandvikLeopardDI5502022}) and Cat C13 Diesel Engine (\citet{caterpillarCatC13ACERT2016}) for diesel derivation. 

\medskip
\noindent\textbf{Process description.}\quad

This block models the combined drilling and blasting operation that converts in-situ rock to fragmented material suitable for excavation and haulage. The block is parameterized per single production drill rig; total system capacity is obtained by multiplying by the number of active drills, analogous to the \texttt{brine\_pump} block. The dominant OPEX driver is explosive consumption, which is set at the production step level via a project-specific powder factor and is therefore not a machine block parameter; the machine block captures equipment CAPEX, diesel consumption, and labor for the drill rig and ancillary blasting equipment.

\medskip
\noindent\textbf{Equation of merit.}\quad

The equation of merit relates single-drill throughput to the drilling rate and blast pattern geometry:
\[
	\dot{m}_{\text{rock}}\;(\si{t/h}) = R_{\text{drill}}\;(\si{t/h}) \times f_{\text{rock}}\;(\text{t/m})
\]
where $R_{\text{drill}}$ is the effective drilling rate in meters of hole per operating hour (including fixed time for repositioning, rod changes, and collar setup) and $f_{\text{rock}}$ is the rock broken per meter of hole drilled, which consolidates bench height, blast pattern spacing and burden, sub-drill allowance, and in-situ rock density into a single productivity factor.

\medskip
\noindent\textbf{Sizing basis.}\quad

The step basis is \texttt{rock\_blast\_throughput} in tonnes per hour and \texttt{base\_volume} = 650, representing 650 t/h of in-situ rock broken per single production drill. This value is derived from the Whabouchi drilling disclosure as follows:

\begin{itemize}[noitemsep]
	\item Hole diameter: 115--140 mm; bench height: 12 m
	\item Blast pattern: 3.3 m spacing $\times$ 2.8 m burden
	\item In-situ rock density: 2.7 \si{t/m^3}
	\item Rock broken per hole: $12 \times 3.3 \times 2.8 = 110.9$ \si{m^3}; $110.9 \times 2.7 = 299$ t $\approx$ 300 t
	\item Pure drilling rate: 35 m/h; fixed time per hole: 7 min (0.117 h)
	\item Effective cycle time per hole: $12/35 + 0.117 = 0.460$ h/hole
	\item Single-drill throughput: $300 / 0.460 \approx 652$ t/h, rounded to 650 t/h
\end{itemize}

\medskip
\noindent\textbf{Bill of Equipment.}\quad

\begin{table}[h!]
\caption{Bill of equipment for the \texttt{blasting\_generic} block.}
\centering
\begin{tabular}{>{\raggedright\arraybackslash}p{5.2cm}
                >{\raggedright\arraybackslash}p{4cm}
                >{\raggedright\arraybackslash}p{4.5cm}}
\toprule
\textbf{Equipment} & \textbf{Machine data used} & \textbf{Source(s)} \\
\midrule
Production drill rig & Composite of Whabouchi and Georgia Lake drill parameters; Sandvik DI550 class & NI 43-101 disclosures; Sandvik specification sheet \\
Bulk explosives delivery truck (ANFO/emulsion) & Ancillary & Not explicit $\rightarrow$ inferred \\
Drill steel, bits, and down-the-hole hammers & Consumables $\rightarrow$ maintenance fraction & --- \\
\bottomrule
\end{tabular}
\end{table}

\noindent Unlike fixed chemical processing equipment, drilling and blasting equipment is mobile and arrives pre-assembled. As a result, the standard Lang-type installation factors (2.3--3.3$\times$) used for fixed-plant process equipment are not applicable. Instead, a delivery, commissioning, and initial tooling factor of 1.15$\times$ is applied to bare equipment cost, consistent with mobile mining fleet practice.

\medskip
\noindent\textbf{CAPEX.}\quad

Anchor equipment is a single production down-the-hole (DTH) drill rig in the 5-inch hammer class with a bare purchase price of \$3.5--4.5M, reflecting the market range for rigs of this capacity (e.g.\ Sandvik Leopard DI550/DI650i, Epiroc FlexiROC D65). The bulk explosives truck is estimated at \$0.8--1.2M bare. Both are multiplied by 1.15$\times$ for delivery and commissioning.

Auxiliary and indirect costs (engineering, site infrastructure for explosive storage, maintenance facilities, and contingency) are applied as explicit percentage adders on the installed equipment subtotal, following the standard three-scenario structure:

\begin{table}[h!]
\caption{CAPEX derivation for the \texttt{blasting\_generic} block (single drill unit).}
\centering
\begin{tabular}{lccccc}
\toprule
\textbf{Scenario} & \textbf{Drill inst.} & \textbf{Truck inst.} & \textbf{Equip.\ subtotal} & \textbf{Aux.\ \%} & \texttt{prim\_equip\_price\_base} \\
\midrule
Conservative & \$5.18M & \$1.38M & \$6.56M & 30\% & \textbf{\$8,500,000} \\
Midpoint     & \$4.60M & \$1.15M & \$5.75M & 20\% & \textbf{\$6,900,000} \\
Optimistic   & \$4.03M & \$0.92M & \$4.95M & 15\% & \textbf{\$5,700,000} \\
\bottomrule
\end{tabular}
\end{table}

\noindent Note that drill rig and truck bare purchase prices are drawn from the Georgia Lake disclosure. The 30\%/20\%/15\% auxiliary adder range is consistent with the ``mechanically intensive, multiple-drive systems'' column in \ref{tab:c-capex-adders}, though at the lower end, reflecting the simpler integration requirements of mobile equipment relative to fixed-plant installations.

Equipment life is set to 15 years, reflecting heavy mobile equipment replacement cycles in hard rock mining service. The CAPEX scaling exponent is 0.60 (six-tenths rule).

\medskip
\noindent\textbf{Electricity.}\quad

\texttt{electricity\_base\_total} = 0. The production drill rig is entirely diesel-powered. No grid electricity is consumed at the blasting step.

\medskip
\noindent\textbf{Diesel.}\quad

Diesel is the sole thermal utility for this block; electricity, natural gas, and propane are zero (the drill rig is entirely diesel-powered). Diesel consumption is derived from the Sandvik DI550 engine specification: a Caterpillar C13B rated at 328 kW, with a brake-specific fuel consumption (BSFC) of 0.235 kg/kWh. At a diesel density of 0.84 kg/L, full-load consumption is:
\[
	\dot{V}_{\text{full}} = \frac{328 \times 0.235}{0.84} = 91.8 \text{ L/h}
\]

Effective fuel consumption depends on the operating load factor, which varies with rock hardness, compressor demand, and the fraction of time spent on non-drilling activities (tramming, repositioning). A load factor range of 0.65--0.85 is applied across scenarios. The bulk explosives truck operates intermittently (charging loaded holes while the drill repositions) and is pro-rated to drill operating hours at 8--15 L/h:

\begin{table}[h!]
\caption{Diesel derivation for the \texttt{blasting\_generic} block at 650 t/h base volume.}
\centering
\begin{tabular}{lcccc}
\toprule
\textbf{Scenario} & \textbf{Load factor} & \textbf{Drill (L/h)} & \textbf{Truck (L/h)} & \texttt{diesel\_base\_total} \\
\midrule
Conservative & 0.85 & 78 & 15 & \textbf{95 L/h} \\
Midpoint     & 0.75 & 69 & 12 & \textbf{80 L/h} \\
Optimistic   & 0.65 & 60 &  8 & \textbf{65 L/h} \\
\bottomrule
\end{tabular}
\end{table}

We note that the drill fuel consumption is derived from rated engine power and a standard BSFC; actual fuel consumption depends on rock competence, drilling depth, and site-specific operating practices. The explosives truck pro-rated values (8--15 L/h) are estimates based on intermittent operation. The derivation does not apply the $\div 0.80$ combustion efficiency adjustment used for thermal process equipment, because diesel here is a direct mechanical fuel input (powering the compressor and hydraulics), not a thermal heating source with stack losses.

\medskip
\noindent\textbf{Other utilities.}\quad

Natural gas, propane, steam, cooling water, and compressed air base totals are all zero. Compressed air for the DTH hammer is generated by the rig's onboard diesel-driven screw compressor and is captured within the diesel consumption above. No external compressed air supply is expected to be required.

\medskip
\noindent\textbf{Consumables.}\quad

The primary consumable is explosives (emulsion or ANFO), which is dosed at the production step level via the project-specific powder factor and is therefore not a machine block parameter. Drill bits, DTH hammers, and drill steel are significant wear items; their cost is captured implicitly within the \texttt{opex\_fraction\_of\_capex} maintenance fraction rather than as a separately tracked consumable at this level of model granularity.

\medskip
\noindent\textbf{Labor.}\quad

Labor is 4.5 FTE at base scale with a scaling exponent of 0.5 and \texttt{dedicated\_labor} = FALSE. The drill rig requires a single operator; the explosives truck is operated by the same crew during non-drilling intervals. The 0.5 exponent reflects the discrete nature of drill fleet staffing but shared supervisory and support staff:
\[
	\text{FTE} = 1.0 \times 4.5 = 4.5
\]

\medskip
\noindent\textbf{Footprint.}\quad

No fixed plant footprint is assigned. The drill rig and explosives truck operate within the active pit area, which is determined by mine planning rather than process layout.

\medskip
\noindent\textbf{Availability.}\quad

Process availability (\texttt{proc\_avail\_factor}) is set to 0.90, comprising approximately 5\% scheduled maintenance (bit changes, hammer rebuilds, engine service) and 5\% unplanned downtime (mechanical failures, weather delays, blast clearance). This is held constant across scenarios; throughput variation is driven by the production step level, not the machine block.

\medskip
\noindent\textbf{Maintenance.}\quad

Maintenance is parameterized as \texttt{opex\_fraction\_of\_capex}. Drilling and blasting equipment falls between the ``materials movement only'' (0.5--1\%) and ``simple mechanical processes'' (5--7\%) categories in \ref{tab:c-maintenace-fraction}. The drill rig has significant wear items (DTH hammers, drill bits, drill steel) that drive replacement costs well above the materials-movement baseline. A maintenance fraction of 5\%/4\%/3\% (conservative/midpoint/optimistic) is adopted, corresponding to the lower end of the simple mechanical processes range.

\medskip
\noindent\textbf{Key measurements and cross-checks.}\quad

\begin{itemize}[noitemsep]
	\item \textit{Throughput:} The base volume of 650 t/h (Whabouchi) is consistent with the Georgia Lake cross-check of 522 t/h for a shallower bench geometry.
	\item \textit{Powder factor:} The Georgia Lake disclosure reports 50.1 kg explosives per hole and 202 t rock broken per hole, giving a powder factor of $50.1/202 = 0.248$ kg/t. This is consistent with the SME hard rock range of 0.11--0.38 kg/t (converted from 0.3--1.0 \si{kg/m^3} at 2.6--2.8 \si{t/m^3} density). These powder factors are applied at the production step level, not the machine block.
	\item \textit{Life of Mine unit cost:} The Georgia Lake disclosure reports life-of-mine drilling costs of \$40.9M and blasting costs of \$22.9M for 29.6 Mt of material, giving a combined drilling and blasting unit cost of \$2.15/t. This aggregate value encompasses CAPEX amortization, diesel, labor, maintenance, and consumables, and serves as a high-level validation target for the assembled block once throughput and project parameters are applied.
	\item \textit{Diesel:} The Sandvik DI550 specification (328 kW, Cat C13B engine) yields 92 L/h at full load; at 75\% load factor, 69 L/h for the drill alone. Combined with the explosives truck, the midpoint of 80 L/h is consistent with the operating range for this equipment class.
\end{itemize}

\clearpage
\paragraph{\texttt{excavation\_loading\_generic}}\label{app:04-methods-blockdefs-excavload}\leavevmode\par

\medskip
\noindent\textbf{Sources.}\quad

\citet{CatPerformanceHandbookfromVSTfuelconsumption20221209T212009}; \citet{warrenGeorgiaLakeLithium2022}; \citet{mccrackenNI43101Technical2023}; \citet{miningplusptyltd.SECTechnicalReport2023}; \citet{petersPlantDesignEconomics1991}; \citet{ulrichGuideChemicalEngineering1984}

\medskip
\noindent\textbf{Process description.}\quad

This block models the combined excavation and loading operation at an open-pit mine face: mechanical fragmentation of blasted or unconsolidated material, loading into haul trucks or conveyors, and associated mine maintenance. The block encompasses one primary hydraulic excavator, one wheel loader (for stockpile management, truck loading from temporary stockpiles, and cleanup), and one track dozer (general other material movement and maintenance). All three machines are diesel-powered; the block consumes no grid electricity.

\medskip
\noindent\textbf{Equation of merit.}\quad

\[
	\dot{m}_{\text{loose}}\;(\si{t/h}) = \frac{V_{\text{bucket}}\;(\si{m^3}) \times \rho_{\text{loose}}\;(\si{t/m^3}) \times f_{\text{fill}}}{t_{\text{cycle}}\;(\text{min}) / 60}
\]
where $V_{\text{bucket}}$ is the excavator bucket volume, $\rho_{\text{loose}}$ is the loose (bank-swelled) rock density, $f_{\text{fill}}$ is the bucket fill factor, and $t_{\text{cycle}}$ is the dig--swing--dump cycle time. For the modeled base case, the cycle time and bucket payload are combined into a single \texttt{base\_volume} of 10 t/cycle.

\medskip
\noindent\textbf{Sizing basis.}\quad

The step basis is \texttt{loose\_rock\_throughput} in tonnes of loose rock per hour. The block is assigned \texttt{process\_type = batch} with a base volume of \textbf{10 t/cycle}, representing the payload of a single excavator bucket pass. Batch cycle times are 1.5 / 1.0 / 0.8 minutes (conservative / midpoint / optimistic), giving throughputs of 400 / 600 / 750 t/h per block unit; setup time is incorporated into these cycle times (set to 0). The 10 t/cycle payload is consistent with a Cat 395-class excavator equipped with a 6.5 \si{m^3} bucket at $\sim$1.8 $\si{t/m^3}$ loose rock density and $\sim$85\% fill factor ($6.5 \times 1.8 \times 0.85 \approx 9.9$ t); these are generic equipment-sizing values for the block and are distinct from the project-specific $\rho_{\text{loose}}$ used in the mining mass balance, which is lower for both modeled deposits. 

The scaling exponent is 1.0, reflecting that mobile mining equipment scales as discrete integer units: each additional block unit represents one additional excavator--loader--dozer set. \texttt{dedicated\_line = FALSE}, as multiple block units can share haul roads and ancillary infrastructure.

\medskip
\noindent\textbf{Bill of Equipment.}\quad

\begin{table}[h!]
\caption{Bill of equipment for the \texttt{excavation\_loading\_generic} block.}
\centering
\begin{tabular}{>{\raggedright\arraybackslash}p{5.2cm}
				>{\raggedright\arraybackslash}p{4cm}
				>{\raggedright\arraybackslash}p{4.5cm}}
\toprule
\textbf{Equipment} & \textbf{Machine data used} & \textbf{Source(s)} \\
\midrule
Hydraulic excavator (Cat 390/395 class, $\sim$90 t) & Excavators & Cat Performance Handbook Ed. 46 \\
Wheel loader (Cat 988 class, $\sim$50 t) & Loaders & Cat Performance Handbook Ed. 46 \\
Track dozer (Cat D9 class, $\sim$49 t) & Dozers & Cat Performance Handbook Ed. 46 \\
Auxiliary (service truck, fuel/lube trailer) & Not explicit $\rightarrow$ inferred & --- \\
\bottomrule
\end{tabular}
\end{table}

\medskip
\noindent\textbf{CAPEX.}\quad

Equipment costs are based on delivered mine-specification pricing for new machines. One excavator at \$2.0M purchase $\rightarrow$ \$2.5M delivered, one wheel loader at \$2.0M purchase $\rightarrow$ \$2.5M delivered, and one track dozer at \$1.5M delivered. The ``delivered'' price includes mine-specification outfitting and freight/commissioning, representing a 25\% markup over base purchase price for the excavator and loader.

Primary equipment subtotal: \$2.5M + \$2.5M + \$1.5M = \$6.5M per block unit. Auxiliary equipment (service vehicles, fuel/lube trailer, initial spare parts inventory) is added as a percentage of the primary subtotal:

\begin{table}[h!]
\caption{CAPEX by scenario for the \texttt{excavation\_loading\_generic} block.}
\centering
\begin{tabular}{>{\raggedright\arraybackslash}p{5.2cm}
				>{\raggedright\arraybackslash}p{4cm}
				>{\raggedright\arraybackslash}p{4.5cm}}
\toprule
\textbf{Scenario} & \textbf{Auxiliary factor} & \textbf{Total installed} \\
\midrule
Conservative & 25\% & $\approx$ \$8.1M \\
Midpoint & 15\% &  $\approx$ \$7.5M \\
Optimistic & 10\% &  $\approx$ \$7.2M \\
\bottomrule
\end{tabular}
\end{table}

\noindent Equipment life is 20 years.

\medskip
\noindent\textbf{Electricity.}\quad

Zero; all equipment is diesel-powered and no grid electricity is consumed at the excavation face.

\medskip
\noindent\textbf{Diesel.}\quad

Diesel consumption is derived from the Caterpillar Performance Handbook, Edition 46, which reports hourly fuel consumption by duty category based on VisionLink telematics data from customer machines worldwide (top and bottom 5\% excluded). The duty categories are defined by idle time fraction: Low ($\sim$50\% idle), Medium ($\sim$30\% idle), and High ($\sim$10\% idle).

Equipment fuel consumption data (L/h):

\begin{table}[h!]
\caption[Hourly fuel consumption by duty category.]{Hourly fuel consumption by duty category.\cite{CatPerformanceHandbookfromVSTfuelconsumption20221209T212009}}
\centering
\begin{tabular}{>{\raggedright\arraybackslash}p{4.5cm}
				>{\raggedright\arraybackslash}p{2.8cm}
				>{\raggedright\arraybackslash}p{2.8cm}
				>{\raggedright\arraybackslash}p{2.8cm}}
\toprule
\textbf{Machine} & \textbf{Low (L/h)} & \textbf{Medium (L/h)} & \textbf{High (L/h)} \\
\midrule
Cat 390D (Tier 2, C18) & 29.5--49.4 & 49.4--63.2 & 63.2--73.2 \\
Cat 988K & 23.8--34.1 & 34.1--44.7 & 44.7--55.3 \\
Cat D9R & 32.0--46.0 & 46.0--59.5 & 59.5--73.0 \\
\bottomrule
\end{tabular}
\end{table}

\noindent Duty assignments are based on the operational role of each machine. The excavator is the primary production machine, operating in continuous truck-loading mode with minimal idle time: Medium duty at minimum, High duty in hard rock or high-demand scenarios. The wheel loader and track dozer are support machines operating intermittently (stockpile management, bench cleanup, push-up); their duty is assigned as Low when intermittent and Medium when operating under sustained demand. Representative values are taken as the midpoint of each duty band.

\begin{table}[h!]
\caption{Diesel consumption derivation by scenario for the \texttt{excavation\_loading\_generic} block.}
\centering
\begin{tabular}{>{\raggedright\arraybackslash}p{4.5cm}
				>{\raggedright\arraybackslash}p{3.2cm}
				>{\raggedright\arraybackslash}p{3.2cm}
				>{\raggedright\arraybackslash}p{3.2cm}}
\toprule
\textbf{Component} & \textbf{Optimistic} & \textbf{Midpoint}  & \textbf{Conservative} \\
\midrule
Excavator duty & Medium--low & Medium & High \\
Excavator (L/h) & 49 & 56 & 68 \\
Loader duty & Low & Low & Medium \\
Loader (L/h) & 24 & 29 & 39 \\
Dozer duty & Low & Low & Medium \\
Dozer (L/h) & 32 & 39 & 53 \\
\midrule
\textbf{Fleet total (L/h)} & \textbf{105} & \textbf{124} & \textbf{150} \\
\bottomrule
\end{tabular}
\end{table}

\medskip
\noindent\textbf{Other utilities.}\quad

None. No steam, cooling water, compressed air, or natural gas is consumed.

\medskip
\noindent\textbf{Consumables --- Tires.}\quad

The machines all require tires, which we model as consumable tools. We estimate the tire cost below, assuming four tires per wheel loader:

\begin{table}[h!]
\caption{Annual tire cost estimation for the wheel loader.}
\centering
\begin{tabular}{>{\raggedright\arraybackslash}p{4.7cm}
				>{\raggedright\arraybackslash}p{2.9cm}
				>{\raggedright\arraybackslash}p{2.7cm}
				>{\raggedright\arraybackslash}p{2.9cm}}
\toprule
\textbf{Parameter} & \textbf{Optimistic} & \textbf{Midpoint} & \textbf{Conservative} \\
\midrule
Tire unit cost (\$/tire) & 8,000 & 12,000 & 15,000 \\
Tire life (operating hours) & 5,000 & 4,000 & 3,500 \\
Loader annual hours (h/yr) & 4,000 & 4,400 & 4,700 \\
Tires replaced per year & 3.2 & 4.4 & 5.4 \\
Annual tire cost (\$/yr) & $\sim$\$26,000 & $\sim$\$53,000 & $\sim$\$81,000 \\
\bottomrule
\end{tabular}
\end{table}

\noindent Loader annual operating hours are estimated as block availability (90\%) $\times$ 8,760 h/yr $\times$ loader utilization factor (50--60\%), reflecting the loader's intermittent support role. Tire life of 3,500--5,000 hours spans the range for mining applications on maintained haul roads~\citep{CatPerformanceHandbookfromVSTfuelconsumption20221209T212009}.

\medskip
\noindent\textbf{Labor.}\quad

Labor is 4.5 FTE at base scale with a scaling exponent of 0.5 and \texttt{dedicated\_labor} = FALSE. The base represents the excavator operator only; loader and dozer operators are assumed to be covered at the facility level through labor pooling, keeping total facility labor within reasonable bounds when multiple mining blocks are instantiated in parallel:
\[
	\text{FTE} = 1.0 \times 4.5 = 4.5
\]

\medskip
\noindent\textbf{Footprint.}\quad

No fixed plant footprint is assigned. Mobile mining equipment operates within the active pit area, which does not contribute to building cost calculations.
 
\medskip
\noindent\textbf{Availability.}\quad

Availability for this block is set to \texttt{proc\_avail\_factor} = 0.90, with \texttt{unplanned\_downtime} = 0.05 and \texttt{scheduled\_maintenance} = 0.05, held constant across all three scenarios. The excavation block does not require a dedicated production line (\texttt{dedicated\_line = FALSE}); multiple block units can share ancillary infrastructure (fuel bays, maintenance workshops, haul road access).

\medskip
\noindent\textbf{Maintenance.}\quad

Annual maintenance cost is expressed as a fraction of scaled equipment cost. For mobile mining equipment, maintenance includes engine and hydraulic system servicing, replacement, structural repairs, and scheduled component rebuilds. Tires are excluded (captured in tooling above). 

\begin{table}[h!]
\caption{Maintenance cost fractions for the \texttt{excavation\_loading\_generic} block.}
\centering
\begin{tabular}{>{\raggedright\arraybackslash}p{2.8cm}
				>{\raggedright\arraybackslash}p{3.8cm}
				>{\raggedright\arraybackslash}p{6.5cm}}
\toprule
\textbf{Scenario} & \textbf{\texttt{maint\_override}} & \textbf{Basis} \\
\midrule
Conservative & 0.08 & Hard rock, heavy undercarriage wear \\
Midpoint & 0.06 & Typical mining conditions \\
Optimistic & 0.04 & Favorable ground, newer equipment \\
\bottomrule
\end{tabular}
\end{table}

\noindent Published mining equipment maintenance benchmarks (excluding tires) typically fall in the range of 4--10\% of replacement cost annually, depending on operating conditions and rebuild cycle timing \citep[see Ch.~4.9]{darlingSMEMiningEngineering2011}. The midpoint of 6\% is a central estimate for medium-scale open pit operations.

\clearpage
\paragraph{\texttt{crushing\_cycle}}\label{app:04-methods-blockdefs-crushing}\leavevmode\par

\medskip
\noindent\textbf{Sources.}\quad

\citet{gaoLithiumExtractionHard2023}; \citet{khakmardanLifeCycleAssessment2024}; \citet{liuCarbonFootprintBatteryGrade2025}; \citet{sgscanadaincSK1300Technical2024}; \citet{samariMineralResourceEstimate2024}; \citet{riordanTonopahLithiumClaims2023}; \citet{darlingSMEMiningEngineering2011}, Chapters 3.3--3.4 (Comminution --- Jaw, Impact, Gyratory, and Cone Crushers); \citet{ausimmMonograph27Cost}, Chapter 11 (Beneficiation --- Comminution); \citet{metso:outotecBasicsMineralsProcessing2024}; \citet{michaudCrushingMineralProcessing2015}; \citet{sandvikSandvikCH890iCH895i2019}; \citet{2025JawCrusher2025}; other vendor data.

\medskip
\noindent\textbf{Process description.}\quad

This block models a three-stage crushing circuit reducing run-of-mine ore to a product size of approximately 5--25 mm suitable for downstream classification, attrition scrubbing, or flotation feed preparation. The circuit comprises a primary jaw crusher, two secondary/tertiary cone crushers operating in closed circuit with sizing screens, and associated material handling (ROM receiving, conveyors, transfer chutes, dust collection). Both hard rock (spodumene, lepidolite) and clay pathways use this block, though target product sizes differ: hard rock circuits target $\sim$15 mm for flotation feed, while clay circuits may accept coarser product for slurry preparation.

\medskip
\noindent\textbf{Sizing basis.}\quad

The step basis is \texttt{crushed\_ore\_throughput} in t/h, with \texttt{process\_type = continuous} and a base volume of \textbf{240 t/h} as a baseline.

The block's internal cone-crusher/sizing-screen recirculation stays internal, captured by the crusher-recirculation equipment adder below without entering the step basis. A circuit-level circulating load (classifier oversize returned for re-crushing) can be added at the step level, with the crusher scaling to that combined \texttt{crushed\_ore\_throughput}. No oversizing factor is assumed (1.0). The scaling exponent is 0.65, consistent with the six-tenths rule for comminution equipment.

The crushing stages are sized as follows (after 911 Metallurgist):

\begin{table}[h!]
\caption{Crushing stage size reduction targets.}
\centering
\begin{tabular}{>{\raggedright\arraybackslash}p{3.0cm}
				>{\raggedright\arraybackslash}p{4.0cm}
				>{\raggedright\arraybackslash}p{3.0cm}
				>{\raggedright\arraybackslash}p{3.0cm}}
\toprule
\textbf{Stage} & \textbf{Crusher type} & \textbf{Feed size} & \textbf{Product size} \\
\midrule
Primary & Jaw (C6X class) & up to 1.5 m & 150--300 mm \\
Secondary & Cone (CH890 class) & 150--450 mm & 40--130 mm \\
Tertiary & Cone (CH890 class) & 20--150 mm & 5--25 mm \\
\bottomrule
\end{tabular}
\end{table}


\medskip
\noindent\textbf{Bill of Equipment.}\quad

\begin{table}[h!]
\caption{Bill of equipment for the \texttt{crushing\_cycle} block.}
\centering
\begin{tabular}{>{\raggedright\arraybackslash}p{5.2cm}
				>{\raggedright\arraybackslash}p{4cm}
				>{\raggedright\arraybackslash}p{4.5cm}}
\toprule
\textbf{Equipment} & \textbf{Machine data used} & \textbf{Source(s)} \\
\midrule
Primary jaw crusher (C6X class, 160 kW) & Jaw crushers & 911 Metallurgist; Metso:Outotec; vendor data \\
Secondary cone crusher (Sandvik CH890 class, 750 kW) & Cone crushers & 911 Metallurgist; AusIMM Ch.\ 11; vendor data \\
Tertiary cone crusher (Sandvik CH890 class, 750 kW) & Cone crushers & 911 Metallurgist; AusIMM Ch.\ 11; vendor data \\
ROM receiving (dump pocket, grizzly, feeder) & Not explicit $\rightarrow$ inferred & --- \\
Sizing screens (scalping, product, recirculation) & Not explicit $\rightarrow$ inferred & --- \\
Crusher recirculation (return conveyors, splitters, bins) & Not explicit $\rightarrow$ inferred & --- \\
Conveyors and transfer chutes & Not explicit $\rightarrow$ inferred & --- \\
Dust collection (baghouses, fans, ducting) & Not explicit $\rightarrow$ inferred & --- \\
\bottomrule
\end{tabular}
\end{table}

\clearpage
\medskip
\noindent\textbf{CAPEX.}\quad
 
Crusher capital costs are estimated using a power-based scaling approach anchored to published pricing for a mid-scale primary jaw crusher (C6X125 class, 240 t/h, 160 kW installed motor). The bare equipment cost for this anchor unit is USD 0.45--0.68M, with costs scaling sub-linearly by installed power at exponent $n = 0.65$. Installation multipliers of 1.5--2.0$\times$ bare cost are applied for the jaw crusher, increasing to 1.6--2.1$\times$ for the cone crushers to reflect additional structural steel, hydraulic systems, and controls scope.
 
The secondary and tertiary crushers are both Sandvik CH890-class cone crushers at 750 kW rated power. Because the same machine class is used for both stages, they are costed identically at \$3.5M / \$4.5M / \$5.5M installed (optimistic / midpoint / conservative). Auxiliary equipment is added as a percentage of the installed crusher subtotal, itemized in \ref{tab:crushing-aux-capex}. 
 
\begin{table}[h!]
\caption{Auxiliary CAPEX adders for the \texttt{crushing\_cycle} block (\% of crusher subtotal).}
\label{tab:crushing-aux-capex}
\centering
\begin{tabular}{>{\raggedright\arraybackslash}p{5.2cm}
				>{\raggedright\arraybackslash}p{2.5cm}
				>{\raggedright\arraybackslash}p{2.5cm}
				>{\raggedright\arraybackslash}p{2.5cm}}
\toprule
\textbf{Component} & \textbf{Optimistic} & \textbf{Midpoint} & \textbf{Conservative} \\
\midrule
ROM receiving & 20\% & 30\% & 45\% \\
Sizing screens & 15\% & 25\% & 35\% \\
Crusher recirculation & 10\% & 15\% & 25\% \\
Conveyors and transfer chutes & 35\% & 60\% & 80\% \\
Dust collection & 8\% & 12\% & 20\% \\
\midrule
\textbf{Auxiliary total} & \textbf{88\%} & \textbf{142\%} & \textbf{205\%} \\
\bottomrule
\end{tabular}
\end{table}
 
\begin{table}[h!]
\caption{CAPEX summary for the \texttt{crushing\_cycle} block.}
\centering
\begin{tabular}{>{\raggedright\arraybackslash}p{6.3cm}
				>{\raggedright\arraybackslash}p{2.5cm}
				>{\raggedright\arraybackslash}p{2.5cm}
				>{\raggedright\arraybackslash}p{2.5cm}}
\toprule
\textbf{Item} & \textbf{Optimistic} & \textbf{Midpoint} & \textbf{Conservative} \\
\midrule
Primary jaw (160 kW), installed & \$0.7M & \$1.0M & \$1.4M \\
Secondary cone (750 kW), installed & \$3.5M & \$4.5M & \$5.5M \\
Tertiary cone (750 kW), installed & \$3.5M & \$4.5M & \$5.5M \\
\midrule
Crusher subtotal & \$7.7M & \$10.0M & \$12.4M \\
Auxiliary adder & $\times$1.88 & $\times$2.42 & $\times$3.05 \\
\midrule
\textbf{Total installed CAPEX} & \textbf{\$14.5M} & \textbf{\$24.2M} & \textbf{\$37.8M} \\
\bottomrule
\end{tabular}
\end{table}
\noindent Equipment life is assumed to be 15 years, given the heavy wear and tear these machines undergo. 
 
\medskip
\noindent\textbf{Electricity.}\quad
 
Total rated crusher motor power is $160 + 750 + 750 = 1,660$ kW. Actual operating power is less than rated; load factors of 65\% / 70\% / 75\% (optimistic / midpoint / conservative) are applied to reflect typical draw during continuous operation. Auxiliary electrical loads (screens, feeders, conveyors, dust collection fans, lubrication systems) are added as a percentage of crusher power.
 
\begin{table}[h!]
\caption{Electricity derivation for the \texttt{crushing\_cycle} block.}
\centering
\begin{tabular}{>{\raggedright\arraybackslash}p{6.0cm}
				>{\raggedright\arraybackslash}p{2.5cm}
				>{\raggedright\arraybackslash}p{2.5cm}
				>{\raggedright\arraybackslash}p{2.5cm}}
\toprule
\textbf{Parameter} & \textbf{Optimistic} & \textbf{Midpoint} & \textbf{Conservative} \\
\midrule
Rated motor power (kW) & 1,660 & 1,660 & 1,660 \\
Load factor & 65\% & 70\% & 75\% \\
Crusher draw (kW) & 1,079 & 1,162 & 1,245 \\
Auxiliary electrical adder & 30\% & 40\% & 50\% \\
\midrule
\textbf{Total electricity (kWh/h)} & \textbf{1,400} & \textbf{1,630} & \textbf{1,870} \\
Specific energy (kWh/t) & 5.8 & 6.8 & 7.8 \\
\bottomrule
\end{tabular}
\end{table}

\clearpage
\medskip
\noindent\textbf{Other utilities.}\quad

None; no thermal utilities as all equipment is assumed to be electrically driven. Dust suppression water is assumed to be used but limited in quantity and not tracked in this model. Run-of-mine moisture is not tracked at the crusher for either project (assumed 0\%, or completely dry at the input to the crusher). For Thacker Pass this follows from the arid northern Nevada setting, though the actual disclosure indicates that an in-situ wet slurrying process is used to extract the clay ore. For Jianxiawo it is an assumed conservatism; adhering moisture entering with the ore would displace part of the water added at classification, so setting it to zero sizes the beneficiation make-up at its upper bound. As such, the water usage at the mining/extraction step of both models is a general conservativism that is accounted for by the water added at classification. 

\medskip
\noindent\textbf{Consumables.}\quad

Crusher wear parts (manganese jaw plates, cone mantles and bowl liners) are significant consumables for this block. These are currently not populated (\texttt{tooling\_cost\_base} = 0) and are flagged as a gap requiring project-specific liner consumption data or vendor estimates.


\medskip
\noindent\textbf{Maintenance.}\quad

\texttt{maint\_override} is not currently populated for this block, and thus inherits a standard 5\% value across all scenarios.

\medskip
\noindent\textbf{Labor.}\quad

Labor is 9.0 FTE at base scale with a scaling exponent of 0.5 and \texttt{dedicated\_labor} = FALSE, summing operator-per-shift factors across the crushing circuit --- primary crusher (0.5), secondary/tertiary crushers (0.5), screens and conveyors (0.5), and helper (0.5):
\[
	\text{FTE} = (0.5 + 0.5 + 0.5 + 0.5) \times 4.5 = 9.0
\]

\medskip
\noindent\textbf{Footprint.}\quad

The base plot area is 500~$\si{m^2}$ with a scaling exponent of 0.7, covering the crusher structures, ROM handling, screening, inter-stage conveyors, dust collection, and control room.

\medskip
\noindent\textbf{Availability.}\quad

\texttt{proc\_avail\_factor} = 0.92, with \texttt{unplanned\_downtime} = 0.03 and \texttt{scheduled\_maintenance} = 0.05, held constant across scenarios. The 92\% availability is representative of continuous crushing circuits with scheduled liner change-outs and preventive maintenance.

\clearpage
\paragraph{\texttt{attrition\_scrubber}}\label{app:04-methods-blockdefs-attritionscrubber}\leavevmode\par

\medskip
\noindent\textbf{Sources.}\quad

\citet{laneNI43101Technical2024}; \citet{sgscanadaincSK1300Technical2024}; \citet{riordanTonopahLithiumClaims2023}; \citet{flsmidthWEMCOAttritionScrubber2023} WEMCO Attrition Scrubber product datasheet (No.\ 120 cell); \citet{ulrichGuideChemicalEngineering1984}

\medskip
\noindent\textbf{Process description.}\quad

This block models an attrition scrubbing circuit that uses high-intensity mechanical abrasion to liberate clay coatings and deagglomerate particles prior to downstream classification or leaching. Attrition scrubbers operate at high pulp densities (70--80 wt\% solids) to maximize inter-particle contact and scrubbing intensity. The block is configured as a single bank of four cells in series, consistent with vendor guidance of 1--2 minutes residence time per cell and 4--6 cells for effective scrubbing. Additional capacity is added as parallel banks (lines).

\medskip
\noindent\textbf{Sizing basis.}\quad

The step basis is \texttt{scrubbed\_slurry\_throughput} in $\si{m^3/h}$ of slurry at operating density. The base volume of \textbf{158 $\si{m^3/h}$} is derived from the WEMCO No.\ 120 cell maximum capacity of 225 t/h dry solids at 75 wt\% solids:
\[
	\rho_{\text{slurry}} = \frac{1}{w_s/\rho_s + (1-w_s)/\rho_w} = \frac{1}{0.75/2.7 + 0.25/1.0} = 1.895\;\si{t/m^3}
\]
\[
	\dot{m}_{\text{slurry}} = 225/0.75 = 300\;\si{t/h} \implies \dot{V} = 300/1.895 \approx 158\;\si{m^3/h}
\]
where $\rho_s = 2.7\;\si{t/m^3}$ is a representative hard-rock solids density. The bank residence time is $4 \times 1.2 = 4.8$ minutes at the vendor maximum-capacity point. As such we set the following other process design parameters: \texttt{process\_type = continuous}, \texttt{scaling\_exponent} = 1.0 (discrete parallel banks), \texttt{oversizing\_factor} = 1.15, \texttt{dedicated\_line = TRUE}.

\medskip
\noindent\textbf{Bill of Equipment.}\quad

\begin{table}[h!]
\caption{Bill of equipment for the \texttt{attrition\_scrubber} block.}
\centering
\begin{tabular}{>{\raggedright\arraybackslash}p{5.2cm}
				>{\raggedright\arraybackslash}p{4cm}
				>{\raggedright\arraybackslash}p{4.5cm}}
\toprule
\textbf{Equipment} & \textbf{Machine data used} & \textbf{Source(s)} \\
\midrule
Attrition scrubber package (4$\times$ WEMCO No.\ 120 cells) & Attrition scrubbers & WEMCO datasheet \\
Feed preparation and slurry make-up & Not explicit $\rightarrow$ inferred & --- \\
Slurry pumping (feed, discharge, transfer) & Not explicit $\rightarrow$ inferred & --- \\
\bottomrule
\end{tabular}
\end{table}


\medskip
\noindent\textbf{CAPEX.}\quad

The bare equipment cost for a four-cell WEMCO No.\ 120 package is estimated at approximately \$5M. An installation multiplier of 1.6$\times$ (structural steel, erection, alignment, electrical terminations, commissioning) yields a base installed cost of \$8M. A $\pm$25\% uncertainty band defines the scenario range: \$10M conservative, \$8M midpoint, \$6M optimistic.

Auxiliary equipment is added as a percentage of the installed package cost:

\begin{table}[h!]
\caption{Auxiliary CAPEX adders for the \texttt{attrition\_scrubber} block (\% of installed package).}
\label{tab:attrition-aux-capex}
\centering
\begin{tabular}{>{\raggedright\arraybackslash}p{5.2cm}
				>{\raggedright\arraybackslash}p{2.5cm}
				>{\raggedright\arraybackslash}p{2.5cm}
				>{\raggedright\arraybackslash}p{2.5cm}}
\toprule
\textbf{Component} & \textbf{Optimistic} & \textbf{Midpoint} & \textbf{Conservative} \\
\midrule
Slurry pumps and piping & 15\% & 22.5\% & 30\% \\
Structural and access steel & 10\% & 15\% & 20\% \\
Instrumentation and controls & 5\% & 10\% & 15\% \\
Feed conditioning and density control & 10\% & 17.5\% & 25\% \\
\midrule
\textbf{Auxiliary total} & \textbf{40\%} & \textbf{65\%} & \textbf{90\%} \\
\bottomrule
\end{tabular}
\end{table}

\begin{table}[h!]
\caption{CAPEX summary for the \texttt{attrition\_scrubber} block.}
\centering
\begin{tabular}{>{\raggedright\arraybackslash}p{5.2cm}
				>{\raggedright\arraybackslash}p{2.5cm}
				>{\raggedright\arraybackslash}p{2.5cm}
				>{\raggedright\arraybackslash}p{2.5cm}}
\toprule
\textbf{Item} & \textbf{Optimistic} & \textbf{Midpoint} & \textbf{Conservative} \\
\midrule
Installed package & \$6.0M & \$8.0M & \$10.0M \\
Auxiliary adder & $\times$1.40 & $\times$1.65 & $\times$1.90 \\
\midrule
\textbf{Total installed CAPEX} & \textbf{\$8.4M} & \textbf{\$13.2M} & \textbf{\$19.0M} \\
\bottomrule
\end{tabular}
\end{table}

\noindent Equipment life is 15 years. \texttt{prim\_equip\_scaling\_exponent} = 0.6.


\medskip
\noindent\textbf{Electricity.}\quad

Per-cell motor ratings from the WEMCO datasheet range from 75--150 HP (56--112 kW). Load factors of 65\% / 70\% / 75\% are applied to the four-cell bank to obtain cell demand. Auxiliary electrical loads are itemized:

\begin{table}[h!]
\caption{Auxiliary electrical adders for the \texttt{attrition\_scrubber} block (\% of cell demand).}
\centering
\begin{tabular}{>{\raggedright\arraybackslash}p{5.2cm}
				>{\raggedright\arraybackslash}p{2.5cm}
				>{\raggedright\arraybackslash}p{2.5cm}
				>{\raggedright\arraybackslash}p{2.5cm}}
\toprule
\textbf{Component} & \textbf{Optimistic} & \textbf{Midpoint} & \textbf{Conservative} \\
\midrule
Feed pump(s) & 10\% & 12\% & 15\% \\
Discharge pump(s) & 10\% & 12\% & 15\% \\
Discharge screen & 5\% & 5\% & 5\% \\
Density control & 2\% & 3\% & 5\% \\
Slurry transfer to downstream & 20\% & 35\% & 50\% \\
Instrumentation and controls & 3\% & 3\% & 5\% \\
\midrule
\textbf{Auxiliary total} & \textbf{50\%} & \textbf{70\%} & \textbf{95\%} \\
\bottomrule
\end{tabular}
\end{table}

\noindent Pump power is estimated from the slurry hydraulic equation $P = \rho g Q H / \eta$ at $\rho = 1,895\;\si{kg/m^3}$, $Q = 158\;\si{m^3/h}$, and pump efficiency $\eta \approx 0.60$, giving approximately 0.82 kW per meter of head. The slurry transfer adder (20--50\%) is the dominant variable, driven by inter-area pumping distance and elevation change; the conservative end is consistent with the Thacker Pass disclosure, which reports total ``Attrition Scrubbing Area (w/ Slurry Transfer)'' demand of 1,204 kW against 600 kW installed cell power.

\begin{table}[h!]
\caption{Electricity summary for the \texttt{attrition\_scrubber} block.}
\centering
\begin{tabular}{>{\raggedright\arraybackslash}p{5.2cm}
				>{\raggedright\arraybackslash}p{2.5cm}
				>{\raggedright\arraybackslash}p{2.5cm}
				>{\raggedright\arraybackslash}p{2.5cm}}
\toprule
\textbf{Parameter} & \textbf{Optimistic} & \textbf{Midpoint} & \textbf{Conservative} \\
\midrule
Motor rating per cell (HP) & 75 & 125 & 150 \\
Motor rating per cell (kW) & 56 & 93 & 112 \\
Rated bank power, 4 cells (kW) & 224 & 373 & 448 \\
Load factor & 0.65 & 0.70 & 0.75 \\
Cell demand (kW) & 146 & 261 & 336 \\
Auxiliary electrical adder & +50\% & +70\% & +95\% \\
\midrule
\textbf{Total electricity (kWh/h)} & \textbf{219} & \textbf{443} & \textbf{655} \\
\bottomrule
\end{tabular}
\end{table}

\medskip
\noindent\textbf{Other utilities.}\quad

None. Process water for slurry make-up is dosed and costed at the production step rather than as a block utility, since the requirement is set by the target slurry density and the circuit's recycle configuration rather than by the equipment.


\medskip
\noindent\textbf{Consumables.}\quad

None populated; impeller wear parts are a potential consumable but are not currently modeled, which we flag as a gap requiring project-specific liner consumption data or vendor estimates.


\medskip
\noindent\textbf{Labor.}\quad

Labor is 4.5 FTE at base scale with a scaling exponent of 0.5 and \texttt{dedicated\_labor} = FALSE, assuming one operator per shift for the four-cell bank (cells operate in series as a single train):
\[
	\text{FTE} = 1.0 \times 4.5 = 4.5
\]

\medskip
\noindent\textbf{Footprint.}\quad

The base plot area is 300~$\si{m^2}$ with a scaling exponent of 0.7, covering the four-cell scrubber bank, feed and discharge sumps, slurry pumps, density control equipment, and local motor control.


\medskip
\noindent\textbf{Availability.}\quad

\texttt{proc\_avail\_factor} = 0.92, with \texttt{unplanned\_downtime} = 0.03 and \texttt{scheduled\_maintenance} = 0.05, held constant across scenarios. The 92\% availability is typical of continuous scrubbing circuits with scheduled inspections and preventive maintenance.

\medskip
\noindent\textbf{Maintenance.}\quad

\texttt{maint\_override} is not currently populated for this block, and thus inherits a standard 5\% value across all scenarios.

\clearpage
\paragraph{\texttt{classifier\_generic}}\label{app:04-methods-blockdefs-classifier}\leavevmode\par

\medskip
\noindent\textbf{Sources.}\quad

\citet{flsmidthKREBSGMAXMolded2024}; \citet{ausimmMonograph27Cost}; \citet{metso:outotecBasicsMineralsProcessing2024}; \citet{sgscanadaincSK1300Technical2024}; \citet{riordanTonopahLithiumClaims2023}; \citet{laneNI43101Technical2024}; \citet{nationalrenewableenergylaboratorynrelgoldencoImprovingPumpingSystem2006}; \citet{zhangInnovativeRecoveryLepidolite2025}

\medskip
\noindent\textbf{Process description.}\quad

Classification separates particles by size using hydraulic or mechanical means, typically following crushing, attrition, or grinding. In lithium mineral processing flowsheets, classification is most commonly achieved using hydrocyclone clusters, often followed by screening to reject oversize particles. This block models a hydrocyclone cluster with associated feed pumping, piping, and product handling. For consistency with other slurry-based unit operations (attrition, leaching, thickening), the block is parameterized on slurry volumetric flow rate ($\si{m^3/h}$) rather than solids throughput.


\medskip
\noindent\textbf{Sizing basis.}\quad

The step basis is \texttt{classified\_slurry\_throughput} in $\si{m^3/h}$, with \texttt{process\_type = continuous} and a base volume of \textbf{160 $\si{m^3/h}$}. Capital scales at \texttt{prim\_equip\_scaling\_exponent} = 0.60, consistent with the six-tenths rule; utilities scale at \texttt{scaling\_exponent\_throughput} = 1.0. There is no oversizing (\texttt{oversizing\_factor} = 1.0), and lines must be dedicated \texttt{dedicated\_line = TRUE}.

\medskip
\noindent\textbf{Bill of Equipment.}\quad

\begin{table}[h!]
\caption{Bill of equipment for the \texttt{classifier\_generic} block.}
\centering
\begin{tabular}{>{\raggedright\arraybackslash}p{5.2cm}
				>{\raggedright\arraybackslash}p{4cm}
				>{\raggedright\arraybackslash}p{4.5cm}}
\toprule
\textbf{Equipment} & \textbf{Machine data used} & \textbf{Source(s)} \\
\midrule
Hydrocyclone cluster (FLSmidth Krebs gMAX, 2 operating + 1 standby) & Hydrocyclone classifiers & Krebs gMAX datasheet \\
Auxiliary (slurry pumps, piping, product handling, instrumentation) & Not explicit $\rightarrow$ inferred & --- \\
\bottomrule
\end{tabular}
\end{table}


\medskip
\noindent\textbf{CAPEX.}\quad

Installed equipment cost for the hydrocyclone cluster (2 operating + 1 standby units, manifold, distributor, and structural support) is estimated at \$1.5M / \$2.0M / \$2.8M (optimistic / midpoint / conservative). Auxiliary equipment is added as a percentage of the installed cluster cost:

\begin{table}[h!]
\caption{Auxiliary CAPEX adders for the \texttt{classifier\_generic} block (\% of installed cluster).}
\label{tab:classifier-aux-capex}
\centering
\begin{tabular}{>{\raggedright\arraybackslash}p{6.2cm}
				>{\raggedright\arraybackslash}p{2.5cm}
				>{\raggedright\arraybackslash}p{2.5cm}
				>{\raggedright\arraybackslash}p{2.5cm}}
\toprule
\textbf{Component} & \textbf{Optimistic} & \textbf{Midpoint} & \textbf{Conservative} \\
\midrule
Slurry pumps and piping & 20\% & 25\% & 35\% \\
Structural and access steel & 10\% & 15\% & 20\% \\
Instrumentation and controls & 8\% & 12\% & 18\% \\
Screen package (oversize rejection) & 15\% & 20\% & 30\% \\
\midrule
\textbf{Auxiliary total} & \textbf{53\%} & \textbf{72\%} & \textbf{103\%} \\
\bottomrule
\end{tabular}
\end{table}

\begin{table}[h!]
\caption{CAPEX summary for the \texttt{classifier\_generic} block.}
\centering
\begin{tabular}{>{\raggedright\arraybackslash}p{5.2cm}
				>{\raggedright\arraybackslash}p{2.5cm}
				>{\raggedright\arraybackslash}p{2.5cm}
				>{\raggedright\arraybackslash}p{2.5cm}}
\toprule
\textbf{Item} & \textbf{Optimistic} & \textbf{Midpoint} & \textbf{Conservative} \\
\midrule
Installed cluster & \$1.5M & \$2.0M & \$2.8M \\
Auxiliary adder & $\times$1.53 & $\times$1.72 & $\times$2.03 \\
\midrule
\textbf{Total installed CAPEX} & \textbf{\$2.3M} & \textbf{\$3.4M} & \textbf{\$5.7M} \\
\bottomrule
\end{tabular}
\end{table}

\noindent The equipment life is estimated at 15 years, based on the heavy expected wear. The primary scaling exponent is \texttt{prim\_equip\_scaling\_exponent} = 0.6, consistent with the six-tenths rule. 

\medskip
\noindent\textbf{Electricity.}\quad

Hydrocyclones have no moving parts; energy demand is dominated by slurry feed pumping. Pumping power is estimated from the hydraulic equation $P = \rho g Q H / \eta$, expressed as a specific energy intensity:
\[
	E_{\text{pump}} = \frac{\rho\, g\, H}{\eta \times 3.6 \times 10^6}\;\si{kWh/m^3}
\]
The range $E_{\text{pump}} = 0.15$--$0.40$ \si{kWh/m^3} spans slurry densities of 1,100--1,800 $\si{kg/m^3}$, pump heads of 30--100 m, and pump efficiencies of 60--80\%. Auxiliary loads (instrumentation, valve actuators, lubrication) add 10--16 kW.

\begin{table}[h!]
\caption{Electricity summary for the \texttt{classifier\_generic} block.}
\centering
\begin{tabular}{>{\raggedright\arraybackslash}p{6.0cm}
				>{\raggedright\arraybackslash}p{2.5cm}
				>{\raggedright\arraybackslash}p{2.5cm}
				>{\raggedright\arraybackslash}p{2.5cm}}
\toprule
\textbf{Parameter} & \textbf{Optimistic} & \textbf{Midpoint} & \textbf{Conservative} \\
\midrule
Specific energy (\si{kWh/m^3}) & 0.15 & 0.25 & 0.40 \\
Pump power at 160 $\si{m^3/h}$ (kW) & 24 & 40 & 64 \\
Auxiliary (kW) & 10 & 12 & 16 \\
\midrule
\textbf{Total electricity (kWh/h)} & \textbf{34} & \textbf{52} & \textbf{80} \\
\bottomrule
\end{tabular}
\end{table}


\medskip
\noindent\textbf{Other utilities.}\quad

None. Process water for slurry density adjustment is dosed and costed at the production step rather than as a block utility, since the requirement is set by the inlet and outlet slurry densities rather than by the equipment.


\medskip
\noindent\textbf{Consumables.}\quad

Hydrocyclone wear liners (urethane or ceramic) are the principal consumable for this block, but are not estimated for this block (\texttt{tooling\_cost\_base} = 0); liner replacement data from vendor quotations or site operating records are required to populate this field.


\medskip
\noindent\textbf{Maintenance.}\quad

\texttt{maint\_override} is not currently populated for this block, and thus inherits the facility-level default maintenance fraction.


\medskip
\noindent\textbf{Labor.}\quad
Labor is 4.5 FTE at base scale with a scaling exponent of 0.5 and \texttt{dedicated\_labor} = FALSE, assuming one operator per shift for the hydrocyclone cluster and associated pumping:
\[
	\text{FTE} = 1.0 \times 4.5 = 4.5
\]

\medskip
\noindent\textbf{Footprint.}\quad

The base plot area is 50~$\si{m^2}$ with a scaling exponent of 0.7, covering the cyclone cluster structure, feed sump and pump, and discharge launders.

\medskip
\noindent\textbf{Availability.}\quad

\texttt{proc\_avail\_factor} = 0.92, with \texttt{unplanned\_downtime} = 0.03 and \texttt{scheduled\_maintenance} = 0.05, held constant across scenarios.

\clearpage
\paragraph{\texttt{flotation\_generic}}\label{app:04-methods-blockdefs-flotation}\leavevmode\par
 
\medskip
\noindent\textbf{Sources.}\quad
\citet{jiaoKeyTechnologiesDevelopment2024}; \citet{garrettHandbookLithiumNatural2004}; \citet{darlingSMEMiningEngineering2011} Ch.~12 (Flotation); \citet{metso:outotecBasicsMineralsProcessing2024}; \citet{ausimmMonograph27Cost}; \citet{petersPlantDesignEconomics1991}; \citet{zhangInnovativeRecoveryLepidolite2025}
 
\medskip
\noindent\textbf{Process description.}\quad
This block models a mechanical froth-flotation circuit that upgrades deslimed lepidolite fines to a lithium concentrate. Feed is conditioned with a dodecylamine collector and a sodium hexametaphosphate depressant, then floated across rougher, scavenger, and cleaner stages to produce a concentrate at approximately 1.6\% \ce{Li2O}; flotation tailings report to the tailings sink. For consistency with the other slurry-based unit operations, the block is parameterized on flotation feed slurry volumetric flow rate (\si{m^3/h}) at approximately 30 wt\% solids.
 
\medskip
\noindent\textbf{Sizing basis.}\quad
The step basis is \texttt{flotation\_feed\_throughput} in \si{m^3/h}, with \texttt{process\_type = continuous} and a base volume of \textbf{160 \si{m^3/h}}. The \texttt{scaling\_exponent} = 0.67, consistent with the six-tenths rule; \texttt{oversizing\_factor} = 1.2 (design throughput 192 \si{m^3/h}); lines are not dedicated (\texttt{dedicated\_line = FALSE}). At 30 wt\% solids and grain density 2.7 \si{t/m^3} (slurry density 1.233 \si{t/m^3}), the base volume corresponds to 197 t/h slurry and 59 t/h dry solids.
 
\medskip
\noindent\textbf{Equation of merit.}\quad
Installed cell volume is set by feed rate and total residence time:
\[
	V_{\text{cells}}\;(\si{m^3}) = \frac{\dot{V}_{\text{slurry}}\,(\si{m^3/h}) \times \tau_{\text{res}}\,(\text{min})}{60}
\]
across the rougher, scavenger, and cleaner stages, giving 59 / 73 / 103 \si{m^3} of installed cell volume (optimistic / midpoint / conservative) in 8 / 10 / 14 cells.
 
\medskip
\noindent\textbf{Bill of Equipment.}\quad
\begin{table}[h!]
\caption{Bill of equipment for the \texttt{flotation\_generic} block.}
\centering
\begin{tabular}{>{\raggedright\arraybackslash}p{5.2cm}
				>{\raggedright\arraybackslash}p{4cm}
				>{\raggedright\arraybackslash}p{4.5cm}}
\toprule
\textbf{Equipment} & \textbf{Machine data used} & \textbf{Source(s)} \\
\midrule
Conditioning tank and high-torque agitator & Inferred & --- \\
Rougher / scavenger / cleaner cells (mechanical) & Mechanical flotation cells & Jiao 2024; SME Ch.~12 \\
Flotation air blower & Inferred & SME Ch.~12 \\
Froth and tailings pumps & Inferred & --- \\
Reagent dosing (collector + depressant), instrumentation & Dosages from literature & Jiao 2024 \\
\bottomrule
\end{tabular}
\end{table}
 
\medskip
\noindent\textbf{CAPEX.}\quad
Built bottom-up from the flotation cells (sized from the required cell volume above), the conditioning tank, and flotation auxiliaries (blower, dosing, froth/tails pumps, instrumentation).
 
\begin{table}[h!]
\caption{CAPEX summary for the \texttt{flotation\_generic} block.}
\centering
\begin{tabular}{>{\raggedright\arraybackslash}p{5.2cm}
				>{\raggedright\arraybackslash}p{2.5cm}
				>{\raggedright\arraybackslash}p{2.5cm}
				>{\raggedright\arraybackslash}p{2.5cm}}
\toprule
\textbf{Item} & \textbf{Optimistic} & \textbf{Midpoint} & \textbf{Conservative} \\
\midrule
Flotation cells, installed & \$1.5M & \$2.8M & \$5.8M \\
Conditioning tank, installed & \$0.3M & \$0.4M & \$0.5M \\
Flotation auxiliaries & \$0.4M & \$0.7M & \$1.0M \\
\midrule
\textbf{Total installed CAPEX} & \textbf{\$2.2M} & \textbf{\$3.9M} & \textbf{\$7.3M} \\
\bottomrule
\end{tabular}
\end{table}
 
\noindent Equipment life is estimated at 15 years (rubber-lined, reagent-wetted service). The primary scaling exponent is \texttt{prim\_equip\_scaling\_exponent} = 0.6, consistent with the six-tenths rule.
 
\medskip
\noindent\textbf{Electricity.}\quad
Bottom-up from cell drives and blower, the conditioning agitator, froth/tails pumping, and reagent dosing.
 
\begin{table}[h!]
\caption{Electricity summary for the \texttt{flotation\_generic} block.}
\centering
\begin{tabular}{>{\raggedright\arraybackslash}p{6.0cm}
				>{\raggedright\arraybackslash}p{2.5cm}
				>{\raggedright\arraybackslash}p{2.5cm}
				>{\raggedright\arraybackslash}p{2.5cm}}
\toprule
\textbf{Component} & \textbf{Optimistic} & \textbf{Midpoint} & \textbf{Conservative} \\
\midrule
Flotation cells + blower (1--2 \si{kW/m^3}) & 59 & 110 & 206 \\
Conditioning agitator & 30 & 45 & 75 \\
Pumps (froth + tailings) & 18 & 35 & 66 \\
Reagent dosing and instrumentation & 3 & 5 & 8 \\
\midrule
\textbf{Total electricity (kWh/h)} & \textbf{110} & \textbf{195} & \textbf{355} \\
\bottomrule
\end{tabular}
\end{table}
 
\noindent This corresponds to 1.9 / 3.3 / 6.0 kWh per tonne dry solids, within the 2--8 kWh/t range reported for fine-particle flotation (SME Ch.~12).
 
\medskip
\noindent\textbf{Compressed air.}\quad
Aeration is set from the industry-standard superficial gas velocity $J_g = Q_{\text{air}}/A_{\text{cell}}$, with $J_g = 1.0$--$2.0$ cm/s across scenarios (industrial mechanical cells; 1 cm/s a conservative design value) at an assumed 3 m cell depth.
 
\begin{table}[h!]
\caption{Compressed air for the \texttt{flotation\_generic} block.}
\centering
\begin{tabular}{>{\raggedright\arraybackslash}p{6.0cm}
				>{\raggedright\arraybackslash}p{2.5cm}
				>{\raggedright\arraybackslash}p{2.5cm}
				>{\raggedright\arraybackslash}p{2.5cm}}
\toprule
\textbf{Parameter} & \textbf{Optimistic} & \textbf{Midpoint} & \textbf{Conservative} \\
\midrule
Superficial gas velocity $J_g$ (cm/s) & 1.0 & 1.5 & 2.0 \\
Compressed air (\si{Nm^3/h}) & 710 & 1,310 & 2,470 \\
\bottomrule
\end{tabular}
\end{table}
 
\noindent The blower electrical load is carried in the cells-plus-blower electricity line above; this field records air volume for mass-balance and emissions accounting only.
 
\medskip
\noindent\textbf{Other utilities.}\quad
None.
 
\medskip
\noindent\textbf{Labor.}\quad
Labor is 2.25 FTE at base scale with a scaling exponent of 0.5 and \texttt{dedicated\_labor} = FALSE, assuming half an operator per shift for the flotation train under shared coverage with the classification operator:
\[
	\text{FTE} = 0.5 \times 4.5 = 2.25
\]
 
\medskip
\noindent\textbf{Footprint.}\quad
The base plot area is 700~\si{m^2} (320--950~\si{m^2} across scenarios) with a scaling exponent of 0.7, covering the conditioning tank, cell bank, blower, reagent dosing, and froth/tails pumps.
 
\medskip
\noindent\textbf{Availability.}\quad
\texttt{proc\_avail\_factor} = 0.90, with  \texttt{unplanned\_downtime} = 0.05 and \texttt{scheduled\_maintenance} = 0.05, held constant across scenarios; the lower availability relative to purely mechanical steps reflects flotation's sensitivity to feed quality.
 
\medskip
\noindent\textbf{Maintenance.}\quad
\texttt{maint\_override} is not currently populated for this block, and thus inherits the facility-level default maintenance fraction.

\clearpage
\paragraph{\texttt{high\_energy\_grinding\_mill}}\label{app:04-methods-blockdefs-hegm}\leavevmode\par

\medskip
\noindent\textbf{Sources.}\quad

\citet{debakkerEnergyUseFine2014}; \citet{brissetteEnergySavingsTechnology2010}; Glencore Technology IsaMill product literature and compendium of papers (\citet{glencoretechnologyIsaMill2020Compendium2020}); \citet{kumarEnergyEfficientAdvancedUltrafine2023}; \citet{nanMethodsExtractingLithium2019}; \citet{nanMethodsExtractingLithium2022}; \citet{zouMethodEfficientLithium2024}; 
\citet{ulrichGuideChemicalEngineering1984}.

\medskip
\noindent\textbf{Process description.}\quad

High-energy grinding (also termed mechanical activation) uses a stirred media mill to reduce particle size and disrupt crystalline structure prior to roasting or leaching. In the lepidolite pathway, this step mechanically activates the flotation concentrate to increase reactive surface area --- a fundamentally more energy-intensive duty than conventional regrinding. The dominant equipment class is the horizontal high-speed stirred mill (IsaMill / Netzsch platform) or vertical stirred mill (Metso HIGMill), selected based on target product size and circuit configuration. Stirred mills operate at power intensities of 100--350 \si{kW/m^3} (vs. $\sim$20 \si{kW/m^3} for conventional ball mills), enabling fine and ultrafine grinding in a compact footprint.


\medskip
\noindent\textbf{Equation of merit.}\quad

\[
	P_{\text{motor}}\;(\text{kW}) = E_{\text{sp}}\;(\text{kWh/t}) \times \dot{m}_{\text{DS}}\;(\si{t/h})
\]
where $E_{\text{sp}}$ is the specific energy consumption and $\dot{m}_{\text{DS}}$ is the dry solids throughput. Specific energy is the primary sizing variable for stirred mills; it is determined by laboratory or pilot-scale test work and depends on feed size, product size, ore hardness, and media selection.


\medskip
\noindent\textbf{Sizing basis.}\quad

The step basis is \texttt{ground\_slurry\_throughput} in $\si{m^3/h}$ of slurry at approximately 50 wt\% solids. The \texttt{base\_volume} is 21 $\si{m^3/h}$, which at 50 wt\% solids and a slurry density of 1.46 $\si{t/m^3}$ (grain density 2.7 $\si{t/m^3}$) corresponds to a dry solids throughput of $21 \times 1.46 \times 0.50 = 15.3$ t/h. The \texttt{scaling\_exponent} = 1.0, reflecting that stirred mills scale as discrete integer units (each additional unit adds one mill); an \texttt{oversizing\_factor} of 1.2 is applied to accommodate feed variability. The line is not dedicated \texttt{dedicated\_line = FALSE}.

\medskip
\noindent\textbf{Bill of Equipment.}\quad

\begin{table}[h!]
\caption{Bill of equipment for the \texttt{high\_energy\_grinding\_mill} block.}
\centering
\begin{tabular}{>{\raggedright\arraybackslash}p{5.2cm}
				>{\raggedright\arraybackslash}p{4cm}
				>{\raggedright\arraybackslash}p{4.5cm}}
\toprule
\textbf{Equipment} & \textbf{Machine data used} & \textbf{Source(s)} \\
\midrule
Horizontal high-speed stirred mill (IsaMill M3000 class, 1.12 MW) & IsaMill M3000 as sizing anchor & Glencore Technology \\
Ceramic grinding media supply system & Not explicit $\rightarrow$ inferred & --- \\
Slurry feed/discharge handling & Not explicit $\rightarrow$ inferred & --- \\
Auxiliary (electrical, structural, controls) & Not explicit $\rightarrow$ inferred & --- \\
\bottomrule
\end{tabular}
\end{table}


\medskip
\noindent\textbf{CAPEX.}\quad

The CAPEX anchor is an IsaMill M3000-class unit (3,000 L grinding chamber, 1,120 kW motor) at an estimated installed cost of \$4.5M (2025 USD), based on vendor literature and AusIMM cost relationships for fine-grinding equipment. Required motor power is derived from the specific energy and throughput; CAPEX is then scaled from the anchor using the six-tenths rule ($n = 0.65$):
\[
	\text{CAPEX} = \$4.5\text{M} \times \left(\frac{P_{\text{motor}}}{1,120\;\text{kW}}\right)^{0.65}
\]

\begin{table}[h!]
\caption{CAPEX summary for the \texttt{high\_energy\_grinding\_mill} block.}
\centering
\begin{tabular}{>{\raggedright\arraybackslash}p{5.2cm}
				>{\raggedright\arraybackslash}p{2.5cm}
				>{\raggedright\arraybackslash}p{2.5cm}
				>{\raggedright\arraybackslash}p{2.5cm}}
\toprule
\textbf{Parameter} & \textbf{Optimistic} & \textbf{Midpoint} & \textbf{Conservative} \\
\midrule
Specific energy (kWh/t) & 40 & 70 & 120 \\
Motor power (kW) & 610 & 1,070 & 1,840 \\
Scaling ratio $(P/1,120)^{0.65}$ & 0.67 & 0.97 & 1.38 \\
\midrule
\textbf{Total installed CAPEX} & \textbf{\$3.0M} & \textbf{\$4.4M} & \textbf{\$6.2M} \\
\bottomrule
\end{tabular}
\end{table}

\noindent We assume the equipment life is 20 years, with a \texttt{prim\_equip\_scaling\_exponent} = 0.65 consistent with the six-tenths rule. 

\medskip
\noindent\textbf{Electricity.}\quad

Electricity is typically the dominant operating cost for stirred mills, and is derived directly from the specific energy requirement. The scenario variation reflects the difference between conventional regrind duty (40 kWh/t at P80 $\sim$40--60 $\mu$m) and full mechanical activation ($\geq$120 kWh/t at P80 $<$15 $\mu$m with crystallographic disruption):

\begin{table}[h!]
\caption{Electricity derivation for the \texttt{high\_energy\_grinding\_mill} block.}
\centering
\begin{tabular}{>{\raggedright\arraybackslash}p{5.2cm}
				>{\raggedright\arraybackslash}p{2.5cm}
				>{\raggedright\arraybackslash}p{2.5cm}
				>{\raggedright\arraybackslash}p{2.5cm}}
\toprule
\textbf{Parameter} & \textbf{Optimistic} & \textbf{Midpoint} & \textbf{Conservative} \\
\midrule
Target duty & Coarse regrind & Fine regrind & Mech.\ activation \\
Specific energy (kWh/t DS) & 40 & 70 & 120 \\
Dry solids throughput (t/h) & 15.3 & 15.3 & 15.3 \\
\midrule
\textbf{Total electricity (kWh/h)} & \textbf{610} & \textbf{1,070} & \textbf{1,840} \\
\bottomrule
\end{tabular}
\end{table}

\noindent The specific energy ranges are anchored to \citet{debakkerEnergyUseFine2014} and \citet{brissetteEnergySavingsTechnology2010} for conventional regrind (20--60 kWh/t), and to the CN patent literature (CN110983071B, CN118835103A, CN107475537B) for mechanical activation duties (40--120 kWh/t). No auxiliary electrical adder is applied; the motor draw is the total block demand, as feed/discharge handling is gravity-fed or included in the mill package.

\medskip
\noindent\textbf{Other utilities.}\quad

None.

\medskip
\noindent\textbf{Consumables.}\quad

Ceramic grinding media (typically 2--4 mm beads) is the principal consumable and a significant OPEX driver, commonly accounting for 30--50\% of total circuit operating cost \citep{debakkerEnergyUseFine2014}. Consumption rates of 0.3--1.0 kg/t solids at media costs of \$1--3/kg give an estimated media cost of \$0.30--2.50/t dry solids. \texttt{tooling\_cost\_base} is not currently populated; this is an area for future work.

\medskip
\noindent\textbf{Labor.}\quad

Labor is 2.25 FTE at base scale with a scaling exponent of 0.5 and \texttt{dedicated\_labor} = FALSE. The stirred mill can be modeled under high automation, requiring only 0.5 operators per shift (with operators shared with adjacent circuit duties):
\[
	\text{FTE} = 0.5 \times 4.5 = 2.25
\]

\medskip
\noindent\textbf{Footprint.}\quad

The base plot area is 150~$\si{m^2}$ with a scaling exponent of 0.6, covering the mill, media handling system, feed/discharge sumps, and controls. Stirred mills have a substantially smaller footprint per installed kW than conventional ball mills due to their high power intensity.

\medskip
\noindent\textbf{Availability.}\quad

\texttt{proc\_avail\_factor} = 0.92, with \texttt{unplanned\_downtime} = 0.03 and \texttt{scheduled\_maintenance} = 0.05, held constant across scenarios. This is consistent with IsaMill vendor-reported availability of 90--93\%.

\medskip
\noindent\textbf{Maintenance.}\quad

\texttt{maint\_override} is not currently populated for this block, and thus inherits the facility-level default maintenance fraction.


\clearpage
\paragraph{\texttt{sulfate\_roasting\_train\_lepidolite}}\label{app:04-methods-blockdefs-sulfateroasting}\leavevmode\par

\medskip
\noindent\textbf{Sources.}\quad

\citet{dengSystematicReviewSulfate2026}, \citet{yanExtractionLithiumLepidolite2012}, \citet{gaoLithiumExtractionHard2023}, \citet{parkEffectiveLithiumExtraction2023}, \citet{nanMethodsExtractingLithium2019}, \citet{nanMethodsExtractingLithium2022}, \citet{zouMethodEfficientLithium2024}, \citet{jiaoKeyTechnologiesDevelopment2024}, \citet{garrettHandbookLithiumNatural2004}; Chinese equipment vendor quotes for mineral sulfation roasting at 50--80 t/h scale.

This block models the complete roasting train for sulfate roasting of dewatered lepidolite concentrate. The train encompasses feed conditioning (blending of lepidolite cake, \ce{Na2SO4}, \ce{K2SO4}, and lime), a rotary kiln operating at 800--950$^\circ$C, an afterburner or gas cleaning stage for fluorine management, a heat recovery exchanger, and discharge conveying to the water leach circuit. The roasting step converts lithium from the tightly bound mica lattice into water-soluble lithium sulfate via the central reaction:
\begin{equation}
    \ce{(Li2O)(Al2O3)(3SiO2)(KF)(s) + Na2SO4(s) ->
    Li2SO4(s) + Al2O3(s) + 3SiO2(s) + KF(g/s)}
\end{equation}
Fluorine species are captured by lime addition as \ce{CaF2} and \ce{AlF3} precipitates, rather than volatilizing as toxic \ce{HF} gas. This is the characteristic feature distinguishing sulfate-based from acid- or chlorine-based lepidolite processing routes.

\medskip
\noindent\textbf{Equation of merit.}\quad
\[
    \dot{Q}_\text{thermal}\;(\text{MMBTU/h}) = q_\text{specific}\;(\text{MMBTU/t}) \times \dot{m}_\text{mix}\;(\si{t/h})
\]
where $q_\text{specific}$ is the specific fuel consumption per tonne of roasting mix and $\dot{m}_\text{mix}$ is the roasting mix throughput. This relationship also determines the required kiln size via the thermal load per unit kiln volume.

\medskip
\noindent\textbf{Sizing basis.}\quad

The step basis is \texttt{roasting\_mix} in tonnes per hour at 87.5 wt\% solids entering the kiln. The base volume is \textbf{60 t/h roasting mix}, consistent with the capacity range of industrial rotary kilns used for mineral sulfation (typically 50--80 t/h for analogous applications). The roasting mix includes dry ore solids (0.8 t per tonne wet cake), sodium sulfate addition ($0.5 \times 0.8 = 0.4$ t), lime addition ($0.15 \times 0.8 = 0.12$ t), potassium sulfate addition ($0.1 \times 0.8 = 0.08$ t), and residual moisture (0.2 t), giving a total input-to-process conversion factor (IPCF) of 1.6 per tonne of wet cake feed. At Jianxiawo's approximate annual roasting mix throughput of $\sim$3.2 Mt/yr, the required number of parallel kiln trains at 60 t/h and 92\% availability is approximately 7 trains.

The block is assigned \texttt{process\_type = batch} with \texttt{batch\_cycle\_time = 1 h} and zero setup time, as a modeling convenience for a continuous thermal operation. The oversizing factor is 1.20, reflecting the standard 20\% design margin applied to continuous thermal equipment to accommodate feed variability and maintain campaign availability.

\medskip
\noindent\textbf{Bill of Equipment.}\quad
\begin{table}[h!]
\caption{Bill of equipment for the \texttt{sulfate\_roasting\_train\_lepidolite} block.}
\centering
\begin{tabular}{>{\raggedright\arraybackslash}p{5.2cm}
                >{\raggedright\arraybackslash}p{4cm}
                >{\raggedright\arraybackslash}p{4.5cm}}
\toprule
\textbf{Equipment} & \textbf{Machine data used} & \textbf{Source(s)} \\
\midrule
Feed conditioning mixer & Inferred from analogous mineral processing trains & \citet{dengSystematicReviewSulfate2026} \\
Rotary kiln (50--80 t/h class) & Rotary Kiln data & Vendor quotes; \citet{jiaoKeyTechnologiesDevelopment2024} \\
Afterburner / fluorine gas scrubber & Not explicit $\rightarrow$ inferred & \citet{dengSystematicReviewSulfate2026}; \citet{yanExtractionLithiumLepidolite2012} \\
Heat recovery exchanger & Not explicit $\rightarrow$ inferred & --- \\
Feed and discharge conveyors & Not explicit $\rightarrow$ inferred & --- \\
Exhaust gas fans and ducting & Not explicit $\rightarrow$ inferred & --- \\
Instrumentation and controls & Not explicit $\rightarrow$ inferred & --- \\
\bottomrule
\end{tabular}
\end{table}

\medskip
\noindent\textbf{CAPEX.}\quad

Capital cost per kiln train is estimated using a bottom-up approach anchored to Chinese equipment vendor pricing, cross-checked against the Jianxiawo project-level capital intensity reported by CRU~\cite{zhangChinaHasDomestic2025}. The previous values in this block (\$67M/\$42M/\$25M) reflected Western-specification installed plant costs and are not applicable to a Chinese domestic operation.
 
\textit{Bottom-up estimate.} Chinese-fabricated metallurgical rotary kilns at 50--80~t/h capacity are quoted at approximately \textyen5--15M bare equipment (\$0.7--2.1M at 7.1~RMB/USD). A full installed train requires the conditioning mixer, afterburner or gas scrubber, heat recovery exchanger, feed and discharge conveyors, civil foundations, electrical distribution, and instrumentation, together with a fluorine-resistant refractory upcharge for the HF-bearing operating environment. Applying a full-train installation multiplier of 3.0--4.5$\times$ the bare kiln cost yields an installed range of approximately \$2--10M per train for domestic Chinese specification.
 
\textit{Top-down cross-check.} CRU reports Jianxiawo's all-in capital intensity at \$5,500/t~LCE/yr~\cite{zhangChinaHasDomestic2025}, giving a total project CAPEX of approximately \$253M for 46,000~t~LCE/yr capacity. With approximately 7 parallel roasting trains and the thermal conversion step representing a plausible 25--35\% of total processing capital (consistent with \citet{petersPlantDesignEconomics1991} guidance for calcination-dominated flowsheets), the implied per-train cost is approximately \$7--11M at midpoint --- consistent with the bottom-up range.

Based on these cross-defined estimates, we estimate the conservative, midpoint, and optimistic CAPEX costs in \ref{tab:c-sulfateroastingtrainlepidolite-capex}.
 
\begin{table}[h!]
\caption[CAPEX scenario values for the \texttt{sulfate\_roasting\_train\_lepidolite} block.]{CAPEX scenario values for the \texttt{sulfate\_roasting\_train\_lepidolite} block (per 60~t/h train, fully installed).}
\label{tab:c-sulfateroastingtrainlepidolite-capex}

\begin{tabular}{>{\raggedright\arraybackslash}p{2.5cm}
                >{\raggedright\arraybackslash}p{4.3cm}
                >{\raggedright\arraybackslash}p{8.5cm}}
\toprule
\textbf{Scenario} & \texttt{prim\_equip\_price\_base} & \textbf{Basis} \\
\midrule
Conservative & \$18,000,000  & Chinese spec, corrosion-resistant refractory upgrade, full balance of plant (BOP), 20\% contingency \\
Midpoint     & \$10,000,000 & Standard Chinese mineral roasting train, full installed BOP; consistent with CRU top-down anchor \\
Optimistic   & \$5,000,000  & Streamlined domestic fabrication, minimal F-control allowance, low contingency \\
\bottomrule
\end{tabular}
\end{table}

\noindent The CAPEX scaling exponent is 0.60, following the six-tenths rule for process equipment cost. Equipment life is 20~years across all scenarios, with major refractory relining campaigns expected every 5--7~years within that period, though this is currently not modeled. The fixed site infrastructure shared across multiple trains (central control room, switchgear building, shared reagent storage) is not separately modeled at this block level; it is implicitly distributed across the per-train estimates.

\noindent\textbf{Natural gas.}\quad

Natural gas is the dominant utility and the primary operating cost driver for this block. The three scenarios span 151--226 MMBTU/h at the 60 t/h base scale, equivalent to 2.52--3.77 MMBTU/t roasting mix (2.65--3.97 GJ/t).

The process heat demand has three components. First, sensible heat is needed to raise the dry solids from ambient temperatures to 850$^\circ$C ($\sim$0.74 GJ/t roasting mix). Second, the estimated reaction enthalpy for decomposition is ($\sim$0.50 GJ/t). Third, drying load is required to evaporate the residual moisture in the 80 wt\% solids input cake. At 0.125 t \ce{H2O} per tonne roasting mix, the drying duty totals $0.125 \times 2.77 \approx 0.35$ GJ/t, where 2.77 GJ per tonne \ce{H2O} is composed of sensible heating of liquid water from 25$^\circ$C to 100$^\circ$C ($4.18 \times 75 = 313.5$), latent heat of evaporation (2257), and superheating of the vapour to a $\sim$200$^\circ$C stack ($2.0 \times 100 = 200$). The total demand is therefore $\sim$1.59 GJ/t. With the outlet temperature assumed at 200$^\circ$C, the effect of the heat recovery exchanger on the drying load is contained in this figure. 
 
The demand excludes flue-gas loss, shell loss, and the heat leaving in solids discharged at kiln temperature, none of which is recovered. Fired duty is therefore $q_\text{demand}/\eta$, and the three scenarios are set by the assumed thermal efficiency. 

\begin{table}[h!]
\caption{Natural gas scenario derivation for the \texttt{sulfate\_roasting\_train\_lepidolite} block.}
\centering
\begin{tabular}{>{\raggedright\arraybackslash}p{2.5cm}
                >{\raggedright\arraybackslash}p{2.4cm}
                >{\raggedright\arraybackslash}p{2.4cm}
                >{\raggedright\arraybackslash}p{2.7cm}
                >{\raggedright\arraybackslash}p{5cm}}
\toprule
\textbf{Scenario} & \textbf{MMBTU/h} & \textbf{MMBTU/t mix} & \textbf{Implied $\eta$} & \textbf{Basis} \\
\midrule
Conservative & 226 & 3.77 & 0.40 & Older kiln, feed preheat only, higher excess air and shell radiation \\
Midpoint     & 181 & 3.02 & 0.50 & Modern kiln with feed preheat \\
Optimistic   & 151 & 2.52 & 0.60 & Adds some level of calcine discharge heat recovery \\
\bottomrule
\end{tabular}
\end{table}
 
\noindent All three scenarios sit above the 1.59 GJ/t process demand. The implied efficiency band of 0.40--0.60 is a rule of thumb for a direct-fired rotary kiln with feed preheat; the $\sim$0.50 GJ/t reaction enthalpy inside the demand, 31\% of the total, is likewise an estimate. Natural gas dominates operating costs for this step and is the primary driver of its contribution to the modeled utility costs per tonne of lithium carbonate, as well as a large direct contributor to greenhouse gas emissions.

\medskip
\noindent\textbf{Electricity.}\quad

Electrical demand covers kiln rotation drive, feed and discharge conveying, exhaust gas fans, and instrumentation. It is secondary to natural gas across all scenarios:

\begin{table}[h!]
\caption{Electricity scenario values for the \texttt{sulfate\_roasting\_train\_lepidolite} block.}
\centering
\begin{tabular}{lcc}
\toprule
\textbf{Scenario} & \texttt{electricity\_base\_total} (kWh/h) & \textbf{Intensity (kWh/t mix)} \\
\midrule
Conservative & 270 & 4.50 \\
Midpoint     & 200 & 3.33 \\
Optimistic   & 160 & 2.67 \\
\bottomrule
\end{tabular}
\end{table}

\noindent While electricity plays only a minor role in total operating cost, it contributes meaningfully to Scope 2 emissions given the carbon intensity of the Chinese grid.

\medskip
\noindent\textbf{Compressed air.}\quad

Compressed air consumption is set at 150 \si{Nm^3/h} as a round number across all scenarios, held constant across scenarios as it is driven by the control system specification, rather than process throughput quantities. This value assumes wet scrubbing for fluorine management rather than dry bag filtration, which would use substantially more compressed air.

\medskip
\noindent\textbf{Labor.}\quad

Labor is 5.85 FTE per train at base scale with a scaling exponent of 0.2 and \texttt{dedicated\_labor} = FALSE. The low exponent reflects the reality that one operator shift crew covers multiple kilns once above a minimum staffing threshold; labor grows only weakly with kiln size in this equipment class:
\[
	\text{FTE} = (0.5 + 0.5 + 0.2 + 0.1) \times 4.5 = 5.85
\]

\medskip
\noindent\textbf{Footprint.}\quad

Each train occupies a base plot area of 3,800~$\si{m^2}$ with a scaling exponent of 0.7, covering the kiln shell and drive, feed and discharge buildings, gas handling equipment, heat recovery system, and surrounding access ways.

\medskip
\noindent\textbf{Availability.}\quad

Process availability is 92\% with 3\% unplanned downtime and 5\% scheduled maintenance (totaling 8\% total downtime). The availability figure is slightly lower than for purely mechanical steps, reflecting refractory maintenance requirements and the corrosive operating environment characteristic of fluorine-bearing feedstocks.

\medskip
\noindent\textbf{Other utilities.}\quad

None. Diesel, propane, cooling water, and steam are all zero for this block. The dominant energy input is natural gas for kiln firing, described above; electricity covers all remaining mechanical and control loads.


\medskip
\noindent\textbf{Notes and flags.}\quad

\texttt{tooling\_cost} is zero and \texttt{maint\_override} is unset, so the block inherits the facility-level default maintenance rate of 5\% of scaled primary equipment cost. This is a notable gap for a refractory-lined rotary kiln: annual maintenance costs for refractory repair, refractory relining campaigns, and corrosion-related component replacement are material and should be estimated in more detail.

\clearpage
\paragraph{\texttt{sulfur\_burner}}\label{app:04-methods-blockdefs-sulfurburner}\leavevmode\par
 
\medskip
\noindent\textbf{Sources.}\quad
\citet{davenportSulfuricAcidManufacture2006}, Sec.~2.2 (plant cost), Table~3.1 plant~S1 (equipment list and operating data), Sec.~8.6--8.8 (converter and catalyst geometry); \citet{mousavinezhadSustainableLithiumProduction2025}, \citet{ulrichChemicalEngineeringProcess2004} Table~6.2 (operator staffing); Feasibility studies \citet{sgscanadaincSK1300Technical2024} and \citet{rothFeasibilityStudyNational2022}.

\medskip
\noindent\textbf{Process description.}\quad
 
This block covers the first half of the acid plant: molten sulfur filtration and firing, catalytic \ce{SO2} conversion over four beds in a 3--1 double-contact arrangement, waste-heat recovery, and the power island that converts that heat to electricity. The air drying tower is assigned here as air preparation, although the acid it circulates comes from the absorption circuit. The power island is assigned here because $\sim$75\% of the plant's reaction enthalpy, and all of its high-grade fraction, is released in this block.
 
The block is anchored throughout on a single physical plant: the 4,400~t/d \ce{H2SO4} sulfur-burning train of \citet{davenportSulfuricAcidManufacture2006}, which is both the capital cost anchor (Sec.~2.2) and the equipment and operating anchor (Table~3.1, plant~S1). The disclosed sulfur feed rate of 60.1~t~S/h is reproduced by the block's \texttt{base\_volume} of 60,000~kg~S/h to within 0.17\%, so the sizing basis and the cost basis are not independent estimates of different plants.

\medskip
\noindent\textbf{Bill of Equipment.}\quad
 
\begin{table}[h!]
\caption{Bill of equipment for the \texttt{sulfur\_burner} block.}
\label{tab:c-equip-sulfurburner}
\centering
\begin{tabular}{>{\raggedright\arraybackslash}p{5.2cm}
                >{\raggedright\arraybackslash}p{4cm}
                >{\raggedright\arraybackslash}p{4.5cm}}
\toprule
\textbf{Equipment} & \textbf{Machine data used} & \textbf{Source(s)} \\
\midrule
Sulfur pressure leaf filters (3, 2 online) & Unit count (staffing) & Table 3.1 \\
Sulfur furnace & Unit count (staffing) & Table 3.1 \\
Air drying tower & Unit count (staffing) & Table 3.1 \\
Main air blower & Unit count (staffing) & Table 3.1 \\
Fire-tube waste-heat boilers (2) & Unit count (staffing) & Table 3.1 \\
Superheaters (2) & Unit count (staffing) & Table 3.1 \\
Four-bed converter & Unit count (staffing); catalyst bed volume & Table 3.1; Sec.~8.6 \\
Interbed gas-to-gas exchangers (3) & Unit count (staffing) & Table 3.1 \\
Economizers (3) & Unit count (staffing) & Table 3.1 \\
Hot gas filter & Unit count (staffing) & Table 3.1 \\
Turbogenerator and condenser & Unit count (staffing) & Not explicit $\rightarrow$ inferred \\
\bottomrule
\end{tabular}
\end{table}

 
\medskip
\noindent\textbf{Sizing basis.}\quad
 
At the Thacker Pass sulfur demand of 633,130~t~S/yr, required line time is $633,130,493 / 60,000 = 10,552$ h against an available line time of $7,350 \times (1 - 0.03 - 0.05) = 6,762$~h, giving 1.561 trains. The block is configured on the \texttt{batch} path with a one-hour cycle, zero setup, and \texttt{dedicated\_line} FALSE, mirroring \texttt{sulfate\_roasting\_train\_lepidolite}. This is a modeling device rather than a claim that the plant operates in batches; this forces capacity to be added by adding additional trains rather than by scaling one train continuously, which is more practically realistic. 

\medskip
\noindent\textbf{CAPEX.}\quad
 
\citet{davenportSulfuricAcidManufacture2006} Sec.~2.2 gives approximately US\$75~million (2005) for a new 4,400~t/d plant. A $\pm$30\% envelope is applied as a Class~4--5 estimate, escalated to 2025 dollars at $\text{CEPCI}_{2025}/\text{CEPCI}_{2005} = 810/468.2 = 1.7300$, and split as an assumption as 0.6 to this block and 0.4 to absorption:
\[
    \texttt{prim\_equip\_price\_base} = 75 \times 10^6 \times 1.7300 \times 0.6 \times \{1.3,\, 1.0,\, 0.7\}
\]
giving $\sim$\$101,200,000 / $\sim$\$77,850,000 / \$54,500,000 conservative / midpoint / optimistic, at a 0.6 scaling exponent and an assumed 20-year equipment life. Since the capital expenses are estimated for the whole plant, this split only changes the reported division of capital costs and maintenance costs between the two steps in the step-level plots, but does not affect overall reported costs. 

\medskip
\noindent\textbf{Electricity.}\quad
 
This block is a net electricity \textit{producer}, and it is the only parameter in either acid plant block taken from the project disclosure rather than derived from the equipment and energy balances of \citet{davenportSulfuricAcidManufacture2006}.

Sulfur combustion and \ce{SO2} conversion release approximately 4,035~MJ per tonne of acid, corresponding to roughly 75\% of the plant's total reaction enthalpy, and its only high-grade fraction. That heat is captured through steam, which is expanded through a condensing turbogenerator to produce electricity. The plant then consumes part of that output itself, dominated by the main air blower.

Rather than deriving the steam cycle, the block adopts the net figure disclosed for the project. \citet{sgscanadaincSK1300Technical2024} report acid plant generation of 1,016.9 GWh/y against acid plant consumption of 427.9 GWh/y at full build, on a disclosed sulfur consumption of 9.15 t S per tonne of lithium carbonate, or 1.24 Mt S/yr. Expressed on the basis of sulfur burned, that is 823 kWh generated and 346 kWh consumed per tonne of sulfur, for a net export of \textbf{476 kWh/t S}, which is then applied to this model's sulfur consumption of 9.35 t S per tonne LCE. 

\begin{table}[h!]
\caption{Electricity for the \texttt{sulfur\_burner} block, per tonne of S.}
\centering
\begin{tabular}{>{\raggedright\arraybackslash}p{6.5cm}
                >{\raggedright\arraybackslash}p{2.5cm}
                >{\raggedright\arraybackslash}p{2.5cm}
                >{\raggedright\arraybackslash}p{2.5cm}}
\toprule
\textbf{Component} & \textbf{Cons.} & \textbf{Mid.} & \textbf{Opt.} \\
\midrule
Disclosed generation (kWh/t S)      & \multicolumn{3}{c}{823} \\
Disclosed plant consumption (kWh/t S) & \multicolumn{3}{c}{346} \\
\textbf{Disclosed net export (kWh/t S)} & \multicolumn{3}{c}{\textbf{476}} \\
\midrule
Parametric band ($\pm$30\%)             & 333.6 & 476.5  & 619.5 \\
\midrule
\texttt{electricity\_base\_total} (kWh/h at 60,000 kg S/h) & \textbf{$-$20,000} & \textbf{$-$28,600} & \textbf{$-$37,200} \\
\bottomrule
\end{tabular}
\end{table}

\noindent As this relies on a single estimate, a $\pm$30\% band is applied parametrically as plausible estimate confidence variation, the same envelope applied to the capital cost anchor. The whole net electricity usage is reported in this block, with the subsequent \texttt{acid\_absorption} step carrying zero, because the disclosure reports acid plant power at plant level with no split between the burning and absorption sections. Because the credit is booked at the local electricity price, the acid plant's net position propagates directly into both project cost and scope~2 emissions: at the midpoint, $-$4,518~kWh per tonne of lithium carbonate, worth \$407/t and $-1.44$~t \ce{CO2eq}/t.

\medskip
\noindent\textbf{Other utilities.}\quad

Natural gas, cooling water and steam are all reported as zero. Sulfur melting and steam tracing are served by low-pressure steam generated by the plant. The condenser and acid coolers are air-cooled, following the \citet{davenportSulfuricAcidManufacture2006}'s description of the acid circuits, so no cooling water circuit is drawn. No steam leaves the block; the turbine condenses rather than passing steam to process, which keeps the plant's recovered heat as electricity only.

\medskip
\noindent\textbf{Consumables.}\quad
 
Vanadium pentoxide catalyst is tracked at the production step level at $9.354\times10^{-5}$~kg per kilogram of sulfur burned. The converter holds four beds totalling roughly 690 \si{m^3}, sized from the bed thickness and converter diameter correlations of \citet{davenportSulfuricAcidManufacture2006} Sec.~8.6, at an assumed bulk density of 550 \si{kg/m^3} and 10\% of the charge is replaced each year:
\[
    \frac{690 \times 550 \times 0.10}{60,000 \times 6,762} = 9.354\times10^{-5}
\]
The replacement fraction is an assumption; Sec.~8.8 describes screening the beds every one to two years with a small makeup addition but gives no quantity. \texttt{tooling\_cost\_base} is zero, as the catalyst is the block's only recurring replacement item and is tracked as a reagent instead.
 
\medskip
\noindent\textbf{Labor.}\quad
 
Labor is 29.7 FTE at base scale with a scaling exponent of 0.5 and \texttt{dedicated\_labor} = FALSE, built from \citet{ulrichChemicalEngineeringProcess2004} Table~6.2 operator-per-shift factors applied to the equipment list:
\[
    \text{FTE} = 4.5 \times (3 \times 0.2 + 1.0 + 0.2 + 0.5 + 3.0 + 0.5 + 3 \times 0.1 + 3 \times 0.1 + 0.2) = 4.5 \times 6.6 = 29.7
\]
taking three sulfur filters at 0.2 each, the air blower at 1.0, the drying tower at 0.2, the furnace at 0.5, the converter at 0.5 as a reactor, three interbed exchangers and three economizers at 0.1 each, the hot gas filter at 0.2, and the power island at 3.0. The power island is counted once, as a generating plant covering both boilers, both superheaters, the turbogenerator and the condenser; counting those units individually would staff them twice.
 
\medskip
\noindent\textbf{Footprint.}\quad
 
The base plot area is assumed at 20,000~$\si{m^2}$ with a scaling exponent of 0.7, covering the sulfur melting and filtration area, the furnace and boiler house, the converter and its exchanger train, the blower and drying tower, and the power island.
 
\medskip
\noindent\textbf{Availability.}\quad
 
Process availability is 92\%, with 3\% unplanned downtime and 5\% scheduled maintenance, matched to \texttt{acid\_absorption} because the two blocks are halves of one plant that starts and stops together. The scheduled fraction is dominated by the catalyst screening turnaround rather than by routine outages.
 
\medskip
\noindent\textbf{Maintenance.}\quad
 
\texttt{opex\_fraction\_of\_capex} is 0.0 for this block; maintenance is carried at the facility level rather than as a block-specific adder.

\clearpage
\paragraph{\texttt{acid\_absorption}}\label{app:04-methods-blockdefs-acidabsorption}\leavevmode\par
 
\medskip
\noindent\textbf{Sources.}\quad
 
\citet{davenportSulfuricAcidManufacture2006}, Sec.~2.2 (plant cost), Table~3.1 plant~S1 (equipment list and operating data), and the absorption chapters; \citet{ulrichChemicalEngineeringProcess2004} Table~6.2 (operator staffing); Feasibility studies \citet{sgscanadaincSK1300Technical2024} and \citet{rothFeasibilityStudyNational2022}. 

\medskip
\noindent\textbf{Process description.}\quad
 
This block covers the second half of the acid plant, where the \ce{SO3} gas made upstream is turned into liquid sulfuric acid. Direct addition of water is not practical; adding it to the gas results in a highly exothermic reaction, creating a vapor that is slow and expensive to condense. Instead, the gas is absorbed into a large circulating stream of $\sim$ 98.5~wt \% sulfuric acid, which is large enough to absorb the thermal load. Absorbing \ce{SO3} makes the acid stronger, so water is added to hold the acid strength, and product acid is drawn off the circuit at the rate it is being made. The heat absorbed up by the circulating acid is still significant, at roughly 1,350~MJ per tonne of product, about a quarter of all the heat the plant releases, but it is rejected to atmosphere through air-cooled exchangers; at around 100\textdegree C it is too low-grade to raise steam in this plant configuration.

\medskip
\noindent\textbf{Bill of Equipment.}\quad
 
\begin{table}[h!]
\caption{Bill of equipment for the \texttt{acid\_absorption} block.}
\label{tab:c-equip-acidabsorption}
\centering
\begin{tabular}{>{\raggedright\arraybackslash}p{5.2cm}
                >{\raggedright\arraybackslash}p{4cm}
                >{\raggedright\arraybackslash}p{4.5cm}}
\toprule
\textbf{Equipment} & \textbf{Machine data used} & \textbf{Source(s)} \\
\midrule
Absorption towers (2, interpass and final) & Unit count (staffing) & Table 3.1 \\
Acid coolers (3) & Unit count (staffing) & Table 3.1 \\
Acid circulation pumps and pump tanks (3) & Unit count (staffing) & Table 3.1 \\
Candle mist eliminators (2) & Unit count (staffing) & Table 3.1 \\
Product acid cooler & Unit count (staffing) & Table 3.1 \\
\bottomrule
\end{tabular}
\end{table}

\medskip
\noindent\textbf{Sizing basis.}\quad
 
The block basis is \texttt{acid\_plant\_throughput} in kilograms of \ce{H2SO4} at 100~wt\%, with a \texttt{base\_volume} of 183,500~kg/h, based on size of the 4,400~t/d anchor train from \citet{davenportSulfuricAcidManufacture2006}. The block uses the \texttt{batch} path with a one-hour cycle and no setup time, for the same reason as the burner. 

\medskip
\noindent\textbf{CAPEX.}\quad
 
The acid plant capital anchor of approximately US\$75~million (2005) for a 4,400~t/d plant \cite{davenportSulfuricAcidManufacture2006} is escalated at $\text{CEPCI}_{2025}/\text{CEPCI}_{2005} = 810/468.2 = 1.7300$ and split 0.4 to this block, giving $\sim$\$67,470,000 / $\sim$\$51,900,000 / $\sim$\$36,330,000 conservative / midpoint / optimistic under the same $\pm$30\% envelope, at a 0.6 scaling exponent and an assumed 20-year life. The source prices the plant as a whole and gives no breakdown by section, so the 0.6/0.4 split between the two blocks is assumed to balance the estimated capital cost of this step. 

\medskip
\noindent\textbf{Electricity.}\quad
 
This block carries no electricity. The acid plant's power production is reported in the disclosure for the plant as a whole, with no split between the burning and absorption sections, so we record the net production in its entirety in \texttt{sulfur\_burner} (Appendix \ref{app:04-methods-blockdefs-sulfurburner}). The circulation pumps and cooler fans that physically sit in this step are included inside that disclosed figure.

\medskip
\noindent\textbf{Other utilities.}\quad

Natural gas, cooling water and steam are all reported as zero. The heat of absorption is assumed to be rejected to atmosphere through air-cooled exchangers rather than recovered, due to the low quality of the heat. 

\medskip
\noindent\textbf{Consumables.}\quad
 
Make-up water is tracked at the production step level at $1.837\times10^{-4}$ \si{m^3} per kilogram of product acid, which is the water chemically bound in the acid, $M_{\ce{H2O}}/M_{\ce{H2SO4}}$. Because that water ends up inside the 100~wt\% acid the step basis counts, adding it does not change the step volume. No tooling or replacement consumables, alongside cooling tower make-up or boiler blowdown water, are tracked.

\medskip
\noindent\textbf{Labor.}\quad
 
Labor is 8.1 FTE at base scale with a scaling exponent of 0.5 and \texttt{dedicated\_labor} = FALSE, built from \citet{ulrichChemicalEngineeringProcess2004} Table~6.2 operator-per-shift factors applied to the equipment list:
\[
    \text{FTE} = 4.5 \times (2 \times 0.2 + 3 \times 0.1 + 3 \times 0.2 + 2 \times 0.2 + 0.1) = 4.5 \times 1.8 = 8.1
\]
taking the two towers at 0.2 each, three acid coolers at 0.1 each as heat exchangers, three circulation pumps and tanks at 0.2 each, two mist eliminators at 0.2 each as filters, and the product acid cooler at 0.1. The air drying tower circulates acid drawn from this circuit but is staffed under \texttt{sulfur\_burner} as air preparation, and is not counted twice.

\medskip
\noindent\textbf{Footprint.}\quad
 
The base plot area is 8,000~$\si{m^2}$ with a scaling exponent of 0.7, covering the towers, acid coolers and pump tanks, mist eliminators, and product acid storage and loading.
 
\medskip
\noindent\textbf{Availability.}\quad
 
Process availability is 92\%, with 3\% unplanned downtime and 5\% scheduled maintenance, matched to \texttt{sulfur\_burner} because the two blocks are halves of one plant that are assumed to start and stop together. 
 
\medskip
\noindent\textbf{Maintenance.}\quad
 
\texttt{opex\_fraction\_of\_capex} is 0.0 for this block; maintenance is carried at the facility level rather than as a block-specific adder.


\clearpage
\paragraph{\texttt{acid\_leach\_circuit}}\label{app:04-methods-blockdefs-acidleach}\leavevmode\par

\medskip
\noindent\textbf{Sources.}\quad

\citet{iyerLifecycleAnalysisLithium2024, royLithiumClayAssessing2024, khakmardanLifeCycleAssessment2024, mousavinezhadNextLithiumBoom2025, mousavinezhadSustainableLithiumProduction2025, wesselkaemperOptimizationProcessingCosts2025, sgscanadaincSK1300Technical2024, riordanTonopahLithiumClaims2023, fluorenterprisesRhyoliteRidgeLithiumBoron2022}

The acid leach circuit dissolves lithium and associated metals from classified clay slurry using sulfuric acid, producing a pregnant leach solution (PLS) for downstream processing. Unlike hard-rock pathways, no thermal activation is required; acid attack on the clay mineral lattice proceeds directly at near-ambient temperatures. High-torque agitation is required throughout because the elevated fines content of clay slurries produces viscous, non-Newtonian behavior that standard low-shear impellers cannot adequately handle.

\medskip
\noindent\textbf{Equation of merit.}\quad
\[
    \tau_\text{res}\;(\text{h}) = \frac{V_\text{tank}\;(\si{m^3})}{\dot{V}_\text{slurry}\;(\si{m^3/h})} \times N_\text{tanks}
\]
The circuit is sized to deliver a total residence time sufficient for lithium extraction. At 1 h residence time per tank and 3 tanks in cascade, the total design residence time is 3 h, consistent with disclosed operating conditions at Tonopah Flats and Thacker Pass.

\medskip
\noindent\textbf{Bill of Equipment.}\quad

\begin{table}[h!]
\caption{Bill of equipment for the \texttt{acid\_leach\_circuit} block.}
\centering
\begin{tabular}{>{\raggedright\arraybackslash}p{5.2cm}
                >{\raggedright\arraybackslash}p{4cm}
                >{\raggedright\arraybackslash}p{4.5cm}}
\toprule
\textbf{Equipment} & \textbf{Machine data used} & \textbf{Source(s)} \\
\midrule
Acid leach tanks (3-tank cascade) & AusIMM leach duty tank estimates & AusIMM Monograph 27 \\
High-torque agitator systems & Mixtec high-torque agitator series & Agitator tank series data \\
Feed and slurry make-up systems & Not explicit $\rightarrow$ inferred & --- \\
Acid receiving and storage & Not explicit $\rightarrow$ inferred & --- \\
Slurry and recycle pumping & Not explicit $\rightarrow$ inferred & --- \\
Balance-of-plant (structural, piping, electrical) & Not explicit $\rightarrow$ inferred & --- \\
\bottomrule
\end{tabular}
\end{table}

\medskip
\noindent\textbf{Sizing basis.}\quad

The base volume is 1,100 \si{m^3/h} slurry throughput, corresponding to the operating flow rate of a single 3-tank leach cascade. The design tank volume is approximately 1,300 \si{m^3} per tank, giving an operating fill of approximately 1,100 \si{m^3} at a 85\% fill level (oversizing factor 1.15, i.e.\ $1100 \times 1.15 \approx 1265~ \si{m^3} \approx 1300~\si{m^3}$ design). Three tanks operate in series with a 1 h residence time per tank, for a total circuit residence time of 3 h, consistent with Thacker Pass and Tonopah Flats disclosures.

The block is assigned \texttt{process\_type = continuous} and \texttt{scaling\_exponent\_throughput = 1.0}, as additional throughput capacity is added by operating multiple parallel cascades rather than by scaling individual tank size.

\medskip
\noindent\textbf{CAPEX.}\quad

Installed cost per tank-agitator vessel is estimated from AusIMM leach duty tank benchmarks for heavy-duty, corrosion-resistant steel construction, combined with agitator costs from Mixtec product data. Installed agitator cost is derived by applying an installation factor of 2.0--3.0$\times$ to the raw equipment price of approximately \$83,000 per agitator unit, yielding \$166,000--\$249,000 per agitator. Combined with the tank vessel, the fully installed cost per tank-agitator unit spans \$1.67M--\$3.75M across scenarios.

Per Association for the Advancement of Cost Engineering (AACE) guidance for Class~4--5 estimates~\cite{bredehoeft47R11CostEstimate2020}, primary vessels typically represent 30--50\% of total installed circuit cost, with the remainder covering interconnecting piping, structural steel, instrumentation, electrical distribution, civil foundations, and engineering. The midpoint fraction (40\%) is used for the midpoint scenario, with the conservative (30\%) and optimistic (50\%) fractions applied to the respective bounds:

\begin{table}[h!]
\caption{CAPEX derivation for the \texttt{acid\_leach\_circuit} block (3-tank cascade).}
\centering
\begin{tabular}{lcccc}
\toprule
\textbf{Scenario} & \textbf{Per-tank, inst.} & \textbf{3-tank subtotal} & \textbf{Vessel frac.} & \texttt{prim\_equip\_price\_base} \\
\midrule
Conservative & \$3.75M & \$11.25M & 30\% & \$37.7M \\
Midpoint     & \$2.71M & \$8.13M  & 40\% & \$20.3M \\
Optimistic   & \$1.67M & \$5.01M  & 50\% & \$10.0M \\
\bottomrule
\end{tabular}
\end{table}

\noindent The CAPEX scaling exponent is 0.60 (six-tenths rule). Equipment life is set to 15 years, shorter than the 20-year default used for most mechanical blocks, reflecting the anticipated replacement cycle for rubber-lined steel tanks exposed to continuous sulfuric acid service.

\medskip
\noindent\textbf{Electricity.}\quad

Electrical demand has two components: agitation and pumping. Agitation is fixed at 35 kW per tank (3 tanks = 105 kW total), derived from the Mixtec high-torque agitator specification at operating volume. Pumping demand is estimated using a volumetric energy-intensity approach rather than detailed hydraulic modeling, collapsing pump head, slurry density, and pump-motor efficiency into a single specific energy coefficient.

For slurry transport, a specific energy intensity of 0.15--0.40 \si{kWh/m^3} is assumed, corresponding to moderate transfer heads of approximately 30--60m, a slurry specific gravity of 1.4--1.5, and overall pump-motor efficiencies of approximately 65\%. For wash and recycle water, a lower intensity of 0.10--0.25 \si{kWh/m^3} is applied, with water flow estimated at 10--50\% of the slurry volumetric flow rate.

\begin{table}[h!]
\caption{Electricity scenario derivation for the \texttt{acid\_leach\_circuit} block.}
\centering
\begin{tabular}{lp{2cm}p{2cm}p{2cm}p{2cm}p{3cm}}
\toprule
\textbf{Scenario} & \textbf{Agitation (kW)} & \textbf{Slurry pumping (kW)} & \textbf{Water pumping (kW)} & \textbf{Total pumping (kW)} & \texttt{electricity\_base\_total} \textbf{kWh/h} \\
\midrule
Conservative & 105 & 440 & 110 & 550 & 655 \\
Midpoint     & 105 & 302 &  48 & 350 & 455 \\
Optimistic   & 105 & 165 &  11 & 200 & 305 \\
\bottomrule
\end{tabular}
\end{table}

\noindent Slurry pumping is the dominant electrical consumer, accounting for 65--84\% of total demand across scenarios. The agitation contribution (105 kW) is fixed regardless of scenario, as it is determined by tank volume and impeller specification rather than by head or flow uncertainty.

\medskip
\noindent\textbf{Steam.}\quad

Steam is used to maintain the leach circuit at the target operating temperature of 75--90$^\circ$C. The heating duty has two components: the gross sensible load required to raise the slurry from ambient feed temperature ($\sim$22$^\circ$C) to the leach temperature, and a substantial offset from the exothermic heat of dilution of concentrated sulfuric acid.

\textit{Acid dilution credit.} At the disclosed acid addition rate of 0.49~t \ce{H2SO4} per tonne~\cite{sgscanadaincSK1300Technical2024} and an ore mass flow of approximately 686,000 \si{kg/h} at base scale, the total acid addition is roughly 336,000 \si{kg/h}. The integral heat of dilution from 98\% to process concentration (~15\%) is approximately 750 \si{kJ/kg} \ce{H2SO4}, yielding approximately 70,000 kW of exothermic heat --- sufficient to raise 1,595 \si{t/h} of slurry (effective $c_p \approx 2.71~\text{kJ/(kg$\cdot^\circ$C)}$ at 43\% solids) by approximately 58$^\circ$C from ambient. Acid dilution alone therefore brings the slurry to approximately 80$^\circ$C under ambient Nevada desert conditions.

\textit{Net steam demand.} The residual heating duty after the acid credit is small and strongly dependent on the target leach temperature:

\begin{table}[h!]
\caption{Steam demand derivation for the \texttt{acid\_leach\_circuit} block.}
\centering
\begin{tabular}{lp{2cm}p{2.3cm}p{3.5cm}p{3cm}}
\toprule
\textbf{Scenario} & \textbf{Target T ($^\circ$C)} & \textbf{Gross duty (kW)} & \textbf{Net duty after acid credit (kW)} & \texttt{steam\_base\_total} \textbf{(kg/h)} \\
\midrule
Conservative & 90 & 81,700 & 11,700 & $\sim$21,100 \\
Midpoint     & 84 & 74,500 &  4,500 & $\sim$8,000  \\
Optimistic   & 75 & 63,700 &      0 & $\sim$0      \\
\bottomrule
\end{tabular}
\end{table}

\noindent Steam demand is calculated assuming saturated steam at approximately 130$^\circ$C (3~bar~g) with condensate return at 80$^\circ$C, giving approximately 2,000~kJ/kg useful heat per kilogram of steam. The midpoint value of $\sim$8,000~kg/h implies a target leach temperature of approximately 84$^\circ$C, consistent with the 75--90$^\circ$C operating range.

We use a constant \texttt{steam\_base\_total = 8,000}~kg/h across all scenarios, which corresponds to the midpoint derivation above. The conservative and optimistic scenarios differ substantially (21,100 vs.\ 0~kg/h), and the constant value should be treated as a central estimate only. At the optimistic end, acid dilution heat alone is expected to be sufficient to reach the target temperature and no external steam input is required. This result is sensitive to the assumed acid concentration, ambient temperature, and heat losses; it should be revisited if detailed heat balance data become available.

\medskip
\noindent\textbf{Labor.}\quad

Labor is 4.5 FTE at base scale with a scaling exponent of 0.2 and \texttt{dedicated\_labor} = FALSE, assuming one operator per shift for the continuous three-tank leach cascade:
\[
	\text{FTE} = 1.0 \times 4.5 = 4.5
\]

\medskip
\noindent\textbf{Footprint.}\quad

The base plot area is 800~$\si{m^2}$ with a scaling exponent of 0.7, covering the three leach tanks, electrical equipment, and additional civil infrastructure.

\medskip
\noindent\textbf{Availability.}\quad

Process availability is 92\% with 3\% unplanned downtime and 5\% scheduled maintenance, giving 8\% total downtime. These values are shared with the sulfate roasting block and reflect the corrosive operating environment. The continuous process type means any unplanned downtime directly affects circuit throughput without buffer.

\medskip
\noindent\textbf{Key measurements and cross-checks.}\quad
\begin{itemize}[noitemsep]
    \item Lithium extraction efficiency (fraction of feed Li reporting to PLS); validated against Thacker Pass disclosed recoveries
    \item Acid consumption per tonne ore: the 2024 S-K 1300 disclosure gives a design dose of 490 kg \ce{H2SO4} per tonne of leach feed solids, the same figure reported by Roy et al.\ 2024 from the 2022 feasibility study (0.49 t/t) for a deposit then carrying a ROM grade of 3,270 ppm Li; Mousavinezhad et al.\ 2025 report 15.97 t \ce{H2SO4} per tonne LCE produced; the same disclosure's Table 13-4 gives 611.8 Mt of leach ore against 11.5 Mt of lithium carbonate produced, i.e. 26.07 t \ce{H2SO4} per tonne LCE at that dose, against this model's 28.58. 
    \item Total circuit residence time (3 h target) consistent with Thacker Pass and Tonopah Flats disclosures
    \item CAPEX cross-check: midpoint of \$20.3M for a 1,100~\si{m^3/h} cascade implies approximately \$18.5/(\si{m^3/h}), consistent with the Class~4--5 estimate range for stirred-tank leach circuits in the AusIMM cost database
\end{itemize}

\medskip
\noindent\textbf{Notes and flags.}\quad

The electricity derivation uses a volumetric energy-intensity approach rather than detailed hydraulic modeling. The assumed specific energy range (0.15--0.40 \si{kWh/m^3} for slurry) encompasses significant uncertainty in actual pump head, which depends on circuit elevation differences and pipe run lengths not specified in the available disclosures.

\clearpage
\paragraph{\texttt{water\_leach}}\label{app:04-methods-blockdefs-waterleach}\leavevmode\par

\medskip
\noindent\textbf{Sources.}\quad

\citet{dengSystematicReviewSulfate2026, gaoLithiumExtractionHard2023, parkEffectiveLithiumExtraction2023, yanExtractionLithiumLepidolite2012}; Chinese patents (\citet{zouMethodEfficientLithium2024, nanMethodsExtractingLithium2019, nanMethodsExtractingLithium2022}). Tank costs from AusIMM benchmarks (\citet{ausimmMonograph27Cost}); agitator costs from Mixtec product data (\citet{baguleyMixtechAustraliaFluid2009}).


\medskip
\noindent\textbf{Process description.}\quad

The water leach circuit dissolves water-soluble lithium sulfate and other sulfate species from sulfate-roasted calcine into aqueous solution, producing a pregnant leach solution (PLS) for downstream impurity removal. Unlike the acid leach circuit used in the clay pathway, no acid is added --- the lithium was converted to water-soluble \ce{Li2SO4} during upstream sulfate roasting, and leaching proceeds by simple dissolution in hot water at approximately 80$^\circ$C. The circuit also separates the insoluble leach residue (primarily amorphous silica, alumina, and \ce{CaF2} from fluoride capture during roasting) from the PLS for disposal to tailings. Separation is by pressure filtration to a 65 wt\% solids cake, followed by a single-stage displacement wash returning the entrained lithium-bearing liquor to the PLS.

The feed is roasted calcine slurried with process water to approximately 20 wt\% solids. At this low solids fraction and with crystalline rather than clay-derived particles, the slurry is Newtonian and substantially less viscous than the 43 wt\% clay feed in the acid leach circuit. The leach medium is water with dissolved sulfate salts rather than sulfuric acid, but the hot sulfate solution ($\sim$80$^\circ$C) may require lined or stainless steel construction depending on dissolved sulfate concentrations.

\medskip
\noindent\textbf{Equation of merit.}\quad

\[
	\tau_\text{res}\;(\text{h}) = \frac{V_\text{tank}\;(\si{m^3})}{\dot{V}_\text{slurry}\;(\si{m^3/h})} \times N_\text{tanks}
\]
The circuit is sized to deliver a total residence time of 2 hours (1 hour per tank, 2 tanks in cascade), consistent with reported leaching times of 1--2 hours at 60--90$^\circ$C for sulfate-roasted lepidolite~\citep{yanExtractionLithiumLepidolite2012, dengSystematicReviewSulfate2026}.

\medskip
\noindent\textbf{Bill of Equipment.}\quad

\begin{table}[h!]
\caption{Bill of equipment for the \texttt{water\_leach} block.}
\centering
\begin{tabular}{>{\raggedright\arraybackslash}p{5.2cm}
				>{\raggedright\arraybackslash}p{4cm}
				>{\raggedright\arraybackslash}p{4.5cm}}
\toprule
\textbf{Equipment} & \textbf{Machine data used} & \textbf{Source(s)} \\
\midrule
Leach tanks (2-tank cascade, $\sim$72 \si{m^3} each) & AusIMM tank benchmarks & AusIMM Monograph 27 \\
Standard agitator systems & Mixtec agitator series & Agitator tank series data \\
Filter press and cake wash & \texttt{mechanical\_dewatering} block & \citet{porvooIndustrialFilterPress2025} \\
Feed and slurry make-up systems & Not explicit $\rightarrow$ inferred & --- \\
Slurry and recycle pumping & Not explicit $\rightarrow$ inferred & --- \\
Steam heating system (jacket or coil) & Not explicit $\rightarrow$ inferred & --- \\
Balance-of-plant (structural, piping, electrical) & Not explicit $\rightarrow$ inferred & --- \\
\bottomrule
\end{tabular}
\end{table}

\medskip
\noindent\textbf{Sizing basis.}\quad

The base volume is 60 $\si{m^3/h}$ of leaching slurry at 20 wt\% solids with calcine grain density 2.7 \si{t/m^3}, giving a slurry density of 1.144 \si{t/m^3} and a solids throughput of 13.7 t/h. Each tank has a working volume of 60 \si{m^3} (equal to 1 hour of feed at the base flow rate) and a design volume of 72 \si{m^3} after the \texttt{oversizing\_factor} = 1.20. Two tanks operate in series for a total circuit residence time of 2 hours.

The \texttt{scaling\_exponent\_throughput} = 0.70, \texttt{process\_type} = continuous, \texttt{dedicated\_line} = FALSE. Equipment life is assumed to be 15 years, given the caustic nature of the outputs of the sulfuric acid-based roasting process that feed into this water leach. 

\medskip
\noindent\textbf{CAPEX.}\quad

CAPEX is built bottom-up from AusIMM tank volumetric benchmarks and Mixtec agitator costs, with lower installation factors than the acid leach circuit to reflect the smaller tank sizes and absence of acid storage and dosing equipment.

\textit{Tanks.} Raw tank cost = AusIMM volumetric cost $\times$ total design volume (144 \si{m^3}). At \$4,500--\$9,000/\si{m^3} with a midpoint of \$6,500, the raw tank cost is \$648K--\$936K--\$1,296K.

\textit{Agitators.} Raw agitator cost is anchored to the Mixtec product series~\citep{baguleyMixtechAustraliaFluid2009}, the same source used for the acid leach circuit's high-torque units ($\sim$\$83K raw each). The water leach duty is substantially lighter: at 0.25--0.65 \si{kW/m^3} on a 60 \si{m^3} working volume, each tank requires a 15--40 kW drive for suspension of fine, non-viscous solids, versus the larger high-torque units specified for the 43 wt\% clay slurry. Scaling the Mixtec anchor down to this motor class gives an estimated \$40K--\$60K--\$80K per installed-ready unit (bare equipment), or \$80K--\$120K--\$160K raw for the two-tank cascade. The per-unit prices are an interpolation within the Mixtec series rather than quoted list prices for a specific model, and are flagged as such.

\textit{Installation.} A factor of 2.0--3.0$\times$ bare equipment is applied, consistent with the effective range for other agitated tank leach circuits in the model. The scope includes erection, foundations, piping, electrical distribution, and the steam heating circuit.

\textit{Auxiliaries.} Slurry pumps, piping, and instrumentation are applied as 20--40\% of the installed cost.

\textit{Residue filtration.} The circuit discharges 20\% solids output to a filter press producing the 65 wt\% cake routed to \texttt{slag\_stack\_cn\_65}, with a single-stage displacement wash returning the entrained liquor to the PLS. Rather than deriving this equipment separately, the block adapts the \texttt{mechanical\_dewatering} parameterization (Appendix \ref{app:04-methods-blockdefs-dewater}), which is already defined at a base volume of 60 \si{m^3/h} a heavy duty dewatering process. As such, the press quantities can be added directly to the base cost and adders derived in this block without rescaling: \$30.8M--\$13.3M--\$6.75M of primary equipment, 115--74--45 kWh/h, 95--55--30 \si{Nm^3/h} compressed air, 3.0--1.5--0.7 \si{m^3/h} cooling water, and \$960K--\$300K--\$133K per year of filter cloth tooling. Of note - the \texttt{mechanical\_dewatering} block uses a primary equipment scaling exponent of 0.7 and a throughput exponent of 1.0, against 0.6 and 0.7 for the leach circuit, but this is regarded as an engineering simplification given the slightly different operating conditions of this filter press. 

\begin{table}[h!]
\caption{CAPEX derivation for the \texttt{water\_leach} block (2-tank cascade).}
\centering
\begin{tabular}{>{\raggedright\arraybackslash}p{2.0cm}
				>{\raggedright\arraybackslash}p{1.8cm}
				>{\raggedright\arraybackslash}p{1.8cm}
				>{\raggedright\arraybackslash}p{1.8cm}
				>{\raggedright\arraybackslash}p{1.5cm}>{\raggedright\arraybackslash}p{1.7cm}
				>{\raggedright\arraybackslash}p{3.0cm}}
\toprule
\textbf{Scenario} & \textbf{Raw equip.} & \textbf{Inst.\ sub.} & \textbf{Aux.} & \textbf{Cascade tot.} & \textbf{Press} & \texttt{prim\_equip\_price\_base} \\
Cons. & \$1,456K & \$4,368K & \$1,747K & \$6,115K & \$30,800K & \textbf{\$36,900,000} \\
Mid.  & \$1,056K & \$2,640K & \$792K   & \$3,432K & \$13,300K & \textbf{\$16,730,000} \\
Opt.  & \$728K   & \$1,456K & \$291K   & \$1,747K & \$6,750K  & \textbf{\$8,500,000} \\
\bottomrule
\end{tabular}
\end{table}
 
\noindent \textit{Derivation detail (midpoint row).} Raw tanks: \$6,500/\si{m^3} $\times$ 144 \si{m^3} = \$936K. Raw agitators: 2 $\times$ \$60K = \$120K. Total raw: \$1,056K $\times$ 2.5 installation = \$2,640K. Auxiliaries at 30\%: \$792K. Subtotal: \$3,432K, rounded to \$3,430K. Including the \$13,300K \texttt{mechanical\_dewatering} press gives \$16,730K.


\medskip
\noindent\textbf{Electricity.}\quad

Electricity demand is dominated by agitation, which accounts for approximately 70\% of the midpoint total. The agitation power intensity (0.25--0.65 \si{kW/m^3}) is substantially lower than the 0.5--1.5 \si{kW/m^3} used in the acid leach circuit because the 20 wt\% calcine slurry is Newtonian and non-viscous. Pumping loads are modest given the low volumetric throughput.

\begin{table}[h!]
\caption{Electricity derivation for the \texttt{water\_leach} block.}
\centering
\begin{tabular}{>{\raggedright\arraybackslash}p{6.0cm}
				>{\raggedright\arraybackslash}p{2.8cm}
				>{\raggedright\arraybackslash}p{2.8cm}
				>{\raggedright\arraybackslash}p{2.8cm}}
\toprule
\textbf{Component} & \textbf{Opt.\ (kW)} & \textbf{Mid.\ (kW)} & \textbf{Cons.\ (kW)} \\
\midrule
Agitation (2 $\times$ 60 \si{m^3})     & 30  & 48  & 78  \\
Slurry and transfer pumps           & 6   & 12  & 20  \\
Heating circuit pump                & 2   & 4   & 6   \\
Instrumentation and auxiliaries     & 2   & 3   & 4   \\
Press feed, hydraulics and cake discharge & 45  & 74  & 115 \\
\midrule
\texttt{electricity\_base\_total}   & \textbf{85 kWh/h}  & \textbf{141 kWh/h} & \textbf{223 kWh/h} \\
\bottomrule
\end{tabular}
\end{table}

\medskip
\noindent\textbf{Steam.}\quad

We assume no additional steam \texttt{steam\_base\_total} = 0 across all scenarios. The water leach circuit operates at approximately 80$^\circ$C. Any heating duty would result from the raising of make-up leach water from ambient to leach temperature:
\[
	\dot{Q}_\text{water} = 54.9\;\text{t/h} \times 4.18\;\text{kJ/(kg$\cdot^\circ$C)} \times 55\;^\circ\text{C} \approx 3,500\;\text{kW}
\]
We assign a zero steam value, reflecting the assumption that residual sensible heat from the upstream sulfate roasting step (which operates at 800--950$^\circ$C) and preheating of leach water using roaster off-gas or recycle as a warm filtrate from downstream separation is sufficient to bring the calcine--water slurry to the target leach temperature without external steam input. Direct quenching of hot calcine into the leach tanks is not assumed. Additional auxiliary heating units are retained in the bill of equipment for startup and upset conditions only.

\medskip
\noindent\textbf{Other utilities.}\quad

Compressed air is consumed at 125--85--60 \si{Nm^3/h} and cooling water at 3.0--1.5--0.7 \si{m^3/h}. The increment over the 30 \si{Nm^3/h} required for tank cleaning and instrumentation due to the filter press demand, as derived per \texttt{mechanical\_dewatering} (Appendix \ref{app:04-methods-blockdefs-dewater}).

\medskip
\noindent\textbf{Consumables.}\quad

\texttt{tooling\_cost\_base} = \$960K--\$300K--\$133K per year, entirely filter cloth replacement on the press, per \texttt{mechanical\_dewatering} (Appendix \ref{app:04-methods-blockdefs-dewater}). The leach itself consumes no tooling; process water dosed at the production step level as a reagent.

\medskip
\noindent\textbf{Labor.}\quad

Labor is 4.05 FTE at base scale with a scaling exponent of 0.2 and \texttt{dedicated\_labor} = FALSE. The press adds 0.9 FTE per \texttt{mechanical\_dewatering} (Appendix \ref{app:04-methods-blockdefs-dewater}):
\[
	\text{FTE} = 0.7 \times 4.5 + 0.9 = 4.05
\]

\medskip
\noindent\textbf{Footprint.}\quad

The base plot area is 600~$\si{m^2}$ with a scaling exponent of 0.6, covering the two leach tanks, slurry pumps, piping, and access (350 $\si{m^2}$) plus 250$\si{m^2}$ for the filter press per \texttt{mechanical\_dewatering} (Appendix \ref{app:04-methods-blockdefs-dewater}).

\medskip
\noindent\textbf{Availability.}\quad

Process availability is 90\% with 5\% unplanned downtime and 5\% scheduled maintenance - though the water leaching circuit itself likely has higher uptime, the availability is dependent on the filter operation time, and thus inherits the assumptions for \texttt{mechanical\_dewatering} (Appendix \ref{app:04-methods-blockdefs-dewater}).

\medskip
\noindent\textbf{Maintenance.}\quad

Neither \texttt{maint\_override} nor a nonzero \texttt{opex\_fraction\_of\_capex} is set; the block inherits the facility-level default maintenance rate of 5\% of scaled primary equipment cost, charged against the total undiscounted equipment price. Maintenance requirements are lower than the acid leach circuit because the circuit handles a near-neutral sulfate solution rather than concentrated sulfuric acid, with no rubber-lined vessels or acid-rated dosing systems to maintain. The corrosion exposure that remains (hot sulfate solution with potential residual fluoride) is reflected in the shortened 15-year equipment life rather than in an elevated maintenance rate.

\medskip
\noindent\textbf{Key measurements and cross-checks.}\quad

\begin{itemize}[noitemsep]
	\item \textit{CAPEX}: the midpoint of \$3.98M is built from AusIMM tanks at \$6,500/\si{m^3} plus Mixtec agitators, with a 2.5$\times$ installation factor and 30\% auxiliaries. The installation factor is consistent with other agitated tank leach circuits in the model. 
	\item \textit{Electricity}: Agitation dominates ($\sim$70\% of midpoint) at a power intensity of 0.40 \si{kW/m^3} --- approximately half the acid leach circuit's intensity, reflecting the less viscous 20 wt\% calcine slurry.
	\item \textit{Steam}: the zero steam demand assumes that residual heat from the upstream sulfate roasting circuit (800--950$^\circ$C) is sufficient to reach the 80$^\circ$C leach temperature. The gross heating duty is approximately 3,700 kW, which is small relative to the thermal energy in the hot calcine.
	\item \textit{Residence time}: 2 hours total (1 h $\times$ 2 tanks) is consistent with the 1--2 h range reported in the literature for sulfate-roasted lepidolite water leaching at 60--90$^\circ$C~\citep{yanExtractionLithiumLepidolite2012, dengSystematicReviewSulfate2026}.
\end{itemize}

\clearpage
\paragraph{\texttt{solid\_liquid\_separation}}\label{app:04-methods-blockdefs-sls}\leavevmode\par

\medskip
\noindent\textbf{Sources.}\quad

Filter press costs from published benchmarks (\citet{porvooIndustrialFilterPress2025}). 

\medskip
\noindent\textbf{Process description.}\quad

The solid--liquid separation block performs polishing clarification of a low-solids process solution by direct pressure filtration. It is used in the Jianxiawo (lepidolite) pathway to remove residual precipitates ($<$1 wt\% solids) from neutralized pregnant leach solution after staged neutralization, yielding a clarified lithium sulfate solution suitable for downstream carbonate precipitation.

At the low solids loadings in this application ($\sim$0.5 wt\%, or approximately 13 t/day of precipitate fines at base scale), the block assumes a pressure leaf filter or a small plate-and-frame filter press as the primary equipment.

\medskip
\noindent\textbf{Bill of Equipment.}\quad

\begin{table}[h!]
\caption{Bill of equipment for the \texttt{solid\_liquid\_separation} block.}
\centering
\begin{tabular}{>{\raggedright\arraybackslash}p{5.2cm}
				>{\raggedright\arraybackslash}p{4cm}
				>{\raggedright\arraybackslash}p{4.5cm}}
\toprule
\textbf{Equipment} & \textbf{Machine data used} & \textbf{Source(s)} \\
\midrule
Pressure leaf filter or plate-and-frame press (470--630 mm plates) & Porvoo 2025 benchmarks & \citet{porvooIndustrialFilterPress2025} \\
Feed pump (centrifugal, 100 \si{m^3/h}) & Not explicit $\rightarrow$ inferred & --- \\
Backwash system & Not explicit $\rightarrow$ inferred & --- \\
Coagulant/flocculant dosing & Not explicit $\rightarrow$ inferred & --- \\
Instrumentation \& controls & Not explicit $\rightarrow$ inferred & --- \\
\bottomrule
\end{tabular}
\end{table}

\noindent No gravity thickener is included. The entire liquid flow passes through the filter; solids are collected on the filter media and periodically discharged. This distinguishes the block from \texttt{thickener\_generic} and \texttt{multi\_stage\_ccd\_thickening}, both of which use gravity settling as the primary separation mechanism.


\medskip
\noindent\textbf{Sizing basis.}\quad

The base volume is 100 $\si{m^3/h}$ of feed solution to be clarified. At $\sim$0.5 wt\% solids and a feed density of $\sim$1.07 \si{t/m^3}, the solids throughput is approximately 0.5 t/h (13 t/day). The filter is sized by liquid throughput and filtration rate rather than by solids settling flux. At a specific filtration rate of 3--8 $\si{m^3/m^2/h}$ (typical for polishing-grade pressure filtration), the required filtration area is 12--33 $\si{m^2}$, which is within the capacity of a single 470--630~mm plate-and-frame press or a bank of pressure leaf filter elements.

The \texttt{oversizing\_factor} = 1.15, \texttt{scaling\_exponent} = 0.70 (throughput) / 0.60 (equipment), \texttt{process\_type} = continuous, \texttt{dedicated\_line} = FALSE. Equipment life is 20 years.

\medskip
\noindent\textbf{CAPEX.}\quad

CAPEX is built bottom-up from the Porvoo filter press benchmark plus auxiliaries. No gravity thickener is included, so the cost is substantially lower than for the \texttt{thickener\_generic} block.

\textit{Polishing filter.} Bare equipment costs for 470--630~mm plate-and-frame presses handling 60--150 $\si{m^3/h}$ are \$45,000--\$160,000 from published benchmarks~\citep{porvooIndustrialFilterPress2025}. At 100~$\si{m^3/h}$, a single 630~mm press is adequate for the midpoint and optimistic scenarios; the conservative scenario assumes a larger press or parallel units to provide redundancy.

\textit{Auxiliaries.} The feed pump, backwash system, coagulant dosing, and instrumentation are estimated at \$15,000--\$50,000 depending on the complexity of the backwash and chemical dosing arrangement.

\textit{Total plant cost factor.} A factor of 2.0--3.0$\times$ bare equipment is applied, reflecting the simple installation (no large foundations, no concrete tank, no rake mechanism).

\begin{table}[h!]
\caption{CAPEX derivation for the \texttt{solid\_liquid\_separation} block.}
\centering
\begin{tabular}{>{\raggedright\arraybackslash}p{2.2cm}
				>{\raggedright\arraybackslash}p{2.2cm}
				>{\raggedright\arraybackslash}p{2.0cm}
				>{\raggedright\arraybackslash}p{2.0cm}
				>{\raggedright\arraybackslash}p{1.5cm}
				>{\raggedright\arraybackslash}p{3.0cm}}
\toprule
\textbf{Scenario} & \textbf{Filter bare} & \textbf{Aux bare} & \textbf{Tot bare} & \textbf{PF} & \texttt{prim\_equip\_price\_base} \\
\midrule
Cons. & \$160,000 & \$50,000 & \$210,000 & 3.0$\times$ & \textbf{\$630,000} \\
Mid.  & \$100,000 & \$30,000 & \$130,000 & 2.5$\times$ & \textbf{\$325,000} \\
Opt.  & \$55,000  & \$15,000 & \$70,000  & 2.0$\times$ & \textbf{\$140,000} \\
\bottomrule
\end{tabular}
\end{table}

\noindent \textit{Derivation detail (midpoint row).} A single 630~mm plate-and-frame press at \$100,000 bare handles the full 100 $\si{m^3/h}$ liquid throughput at 5 $\si{m^3/m^2/h}$ filtration rate (20 $\si{m^2}$ required area). Auxiliaries (feed pump, backwash, coagulant dosing) add \$30,000. The total bare equipment cost of \$130,000 at a 2.5$\times$ plant factor gives \$325,000 installed.

\medskip
\noindent\textbf{Electricity.}\quad

Electricity demand is dominated by the feed pump, which must push the full 100 $\si{m^3/h}$ liquid flow through the filter media at 2--5 bar pressure drop. 

\begin{table}[h!]
\caption{Electricity derivation for the \texttt{solid\_liquid\_separation} block.}
\centering
\begin{tabular}{>{\raggedright\arraybackslash}p{5.0cm}
				>{\raggedright\arraybackslash}p{2.5cm}
				>{\raggedright\arraybackslash}p{2.5cm}
				>{\raggedright\arraybackslash}p{2.5cm}}
\toprule
\textbf{Component} & \textbf{Opt.\ (kW)} & \textbf{Mid.\ (kW)} & \textbf{Cons.\ (kW)} \\
\midrule
Feed pump (100 $\si{m^3/h}$ through filter)  & 5   & 8   & 15  \\
Backwash pump (intermittent)                  & 2   & 3   & 6   \\
Compressed air (cake discharge)               & 1   & 2   & 5   \\
Filter mechanism (plate shifter if press)     & 2   & 2   & 4   \\
Coagulant/flocculant dosing                   & 1   & 1   & 2   \\
Instrumentation \& controls                   & 2   & 2   & 3   \\
\midrule
\texttt{electricity\_base\_total}                                & \textbf{13} & \textbf{18} & \textbf{35} \\
\bottomrule
\end{tabular}
\end{table}

\noindent The feed pump accounts for approximately 45\% of the midpoint total, consistent with a configuration where the dominant energy cost is pushing liquid through filter media rather than settling solids by gravity. The scenario spread (13--35 kWh/h) is driven primarily by the assumed filter pressure drop (2 bar optimistic, 3 bar midpoint, 5 bar conservative), which affects the feed pump power via $P = \Delta p \cdot Q / \eta$.

\textit{Feed pump cross-check (midpoint).} At 100 $\si{m^3/h}$ = 0.028 $\si{m^3/s}$, 3 bar = 300 kPa, and $\eta = 0.60$:
\[
	P = \frac{\Delta p \times Q}{\eta} = \frac{300,000 \times 0.028}{0.60} = 14,000 \;\text{W} = 14 \;\text{kW (installed)}
\]
The tabulated 8 kW reflects the average draw at partial load (filter media resistance increases during the filtration cycle; the pump operates below design pressure for most of the cycle). An average-to-installed ratio of 0.57 is consistent with centrifugal pump part-load operation.

\medskip
\noindent\textbf{Consumables.}\quad

Filter cloths for the press are a recurring consumable that would be tracked through \texttt{tooling\_cost\_base}, which are currently unmodeled.

\medskip
\noindent\textbf{Labor.}\quad

Labor is 0.9 FTE at base scale with a scaling exponent of 0.2 and \texttt{dedicated\_labor} = FALSE. The highly automated polishing filter requires only periodic monitoring, backwash initiation, and cloth inspection:
\[
	\text{FTE} = 0.2 \times 4.5 = 0.9
\]

\medskip
\noindent\textbf{Footprint.}\quad

The base plot area is 400~$\si{m^2}$ with a scaling exponent of 0.6. Although the filter press and feed pump alone have a modest footprint, the total includes the coagulant storage and dosing system, backwash water tank, and access.

\medskip
\noindent\textbf{Availability.}\quad

Process availability is 90\% with 5\% unplanned downtime and 5\% scheduled maintenance.

\medskip
\noindent\textbf{Maintenance.}\quad

Neither \texttt{maint\_override} nor a nonzero \texttt{opex\_fraction\_of\_capex} is set; the block inherits the facility-level default maintenance rate of 5\% of scaled primary equipment cost, charged against the total undiscounted equipment price. The polishing filter has modest maintenance requirements (cloth/element replacement, pump seals).

\medskip
\noindent\textbf{Key measurements and cross-checks.}\quad

\begin{itemize}[noitemsep]
	\item \textit{CAPEX}: the midpoint of \$325K at 100 $\si{m^3/h}$ is 19\% of the \texttt{thickener\_generic} midpoint (\$1.70M), consistent with the absence of a gravity thickener and the minimal solids handling equipment. The CAPEX is driven by the filter and feed pump, not by settling infrastructure.
\end{itemize}

\clearpage
\paragraph{\texttt{thickener\_generic}}\label{app:04-methods-blockdefs-thickener}\leavevmode\par

\medskip
\noindent\textbf{Sources.}\quad

\citet{sgscanadaincSK1300Technical2024}, \citet{petersPlantDesignEconomics1991}, \citet{mularCAPCOSTSHandbookEstimating1998}; vendor data (gravity thickener and filter press cost curves). Thickener bare equipment costs from the CAPCOSTS cost correlation~\citep{mularCAPCOSTSHandbookEstimating1998}; filter press costs from published filter press price benchmarks~\citep{porvooIndustrialFilterPress2025}. Cross-checked against a published municipal thickening and dewatering project cost estimate of \$9.4M--\$10.9M installed (\citet{cleargovCostEstimateDraft}); rake drive motor power confirmed against manufacturer specification of 5.5 kW for a 24 m center-drive thickener~\cite{jinhaominingHotItemLarge}. 


\medskip
\noindent\textbf{Process description.}\quad

The thickener block combines gravity settling with mechanical dewatering to concentrate a dilute slurry into a dense cake or thickened underflow. A high-rate gravity thickener (with flocculant) settles the feed to an intermediate underflow density, followed by a filter press that mechanically dewaters the underflow to the target output solids concentration. The block is used in three contexts across the model pathways, listed below in decreasing order of solids throughput at the block's base scale. 

\begin{enumerate}[noitemsep]
	\item \textit{Thacker Pass clay classification thickening} (15 wt\% $\rightarrow$ 55 wt\% solids). Classified smectite clay slurry is thickened to prepare a dense leach feed for acid leaching, with the overflow recycled as process water. This is the heaviest duty and the primary basis for the block's sizing and CAPEX derivation.
	\item \textit{Silver Peak brine separation} ($\sim$9 wt\% $\rightarrow$ clarified brine). Precipitated \ce{Mg(OH)2} and \ce{CaCO3} solids from batch chemical purification are removed by settling and filtration, yielding a clarified lithium brine. The solids throughput is substantially lower than the clay application at comparable volumetric feed rates.
	\item \textit{\ce{Li2CO3} thickening and dewatering} (20 wt\% $\rightarrow$ 70 wt\% solids). Precipitated \ce{Li2CO3} slurry from the soda ash reactor is concentrated for downstream washing and drying. \ce{Li2CO3} crystals ($\rho = 2.11$ \si{t/m^3}, 20--50 $\mu$m) settle rapidly compared to clay, so the thickener is small-diameter at this duty; the filter press handles the final dewatering from $\sim$35 to 70 wt\%.
\end{enumerate}
 
All three applications share the same two-stage equipment structure (gravity thickener followed by filter press) and differ only in the size of equipment required. The block is parameterized for the heaviest duty (clay thickening at 100 \si{m^3/h}), and the model's six-tenths scaling rule reduces the cost for the lighter applications, which involve faster-settling feed materials and lower solids throughputs. Modern high-rate gravity thickeners use polymer flocculant to promote rapid particle aggregation, which can achieve the same throughput at a fraction of the diameter of a conventional unit. Of note - the Thacker Pass disclosure uses decanter centrifuges rather than filter presses for the dewatering stage; this block uses filter presses as the generic dewatering equipment for broader applicability.

\medskip
\noindent\textbf{Bill of Equipment.}\quad

\begin{table}[h!]
\caption{Bill of equipment for the \texttt{thickener\_generic} block.}
\centering
\begin{tabular}{>{\raggedright\arraybackslash}p{5.2cm}
				>{\raggedright\arraybackslash}p{4cm}
				>{\raggedright\arraybackslash}p{4.5cm}}
\toprule
\textbf{Equipment} & \textbf{Machine data used} & \textbf{Source(s)} \\
\midrule
High-rate gravity thickener (17--22 m diameter at base scale) & CAPCOSTS Fig.~190 & \citet{mularCAPCOSTSHandbookEstimating1998} \\
Filter press package (630--800 mm plates) & Porvoo 2025 benchmarks & \citet{porvooIndustrialFilterPress2025} \\
Liquids handling (feed, underflow, overflow, recycle pumps) & Not explicit $\rightarrow$ inferred & --- \\
Output discharge equipment & Not explicit $\rightarrow$ inferred & --- \\
Auxiliaries (utilities, civil, structural) & Not explicit $\rightarrow$ inferred & --- \\
\bottomrule
\end{tabular}
\end{table}

\medskip
\noindent\textbf{Sizing basis.}\quad

The base volume is 100 $\si{m^3/h}$ of thickener feed slurry at 15 wt\% solids with grain density 2.6 \si{t/m^3}, giving a slurry density of 1.102 \si{t/m^3} and a solids throughput of 16.5 t/h (396 t/day). Design capacity is 115 $\si{m^3/h}$ after the \texttt{oversizing\_factor} = 1.15. The \texttt{scaling\_exponent} = 0.60 (six-tenths rule). Lines are not dedicated (\texttt{dedicated\_line} = FALSE). Equipment life is 15 years.

The gravity thickener is sized from the solids settling flux. The Thacker Pass equipment list specifies a single 55~m diameter classification thickener handling approximately 3,600 tpd of solids, implying a unit settling area of 7.1 ft$^2$/tpd~\citep{sgscanadaincSK1300Technical2024}. This falls within the copper-tailings range of 4--10 ft$^2$/tpd tabulated in CAPCOSTS~\citep[p.~270]{mularCAPCOSTSHandbookEstimating1998}, confirming that lithium clay with flocculant (dosed at 260 g/t per the disclosure) settles in the base-metal range rather than the potash or uranium-precipitate range.

At the base volume of 100 $\si{m^3/h}$ (396 tpd solids), the required settling area at 6--10 ft$^2$/tpd is 2,400--3,960 ft$^2$, corresponding to a thickener diameter of 55--71 ft (17--22 m). The filter press handles the thickener underflow ($\sim$45 $\si{m^3/h}$ at 30 wt\% solids), pressing it to the 55 wt\% target.

\medskip
\noindent\textbf{CAPEX.}\quad

CAPEX is built bottom-up from the CAPCOSTS high-capacity thickener correlation and the Porvoo filter press benchmark.

\textit{Thickener.} The CAPCOSTS correlation for high-capacity (flocculant-assisted) thickeners is~\citep[Figure~190, p.~271]{mularCAPCOSTSHandbookEstimating1998}:
\[
	\text{Price}_{\text{1998}} = 5,809 \times D^{0.8786} \quad (\text{USD, } D \text{ in feet, range 16--100 ft, M\&S = 1400})
\]
The price includes the drive, motor, lift, controls, rake mechanism (mild steel), tank shell, and anchor channel. It excludes erection, foundations, freight, and remote controls. Prices are inflated from 1998 to 2025 using the CEPCI ratio of $810/389.5 \approx 2.08$.

\textit{Filter press.} Bare equipment costs for 630--800~mm plate filter presses are taken from the Porvoo benchmark at \$75,000--\$200,000~\citep{porvooIndustrialFilterPress2025}.

\textit{Total plant cost factor.} The bare equipment total (thickener + filter press) is multiplied by a total plant cost factor of 2.5--3.5$\times$ to account for erection, foundations, piping, electrical, and instrumentation.

\begin{table}[h!]
\caption{CAPEX derivation for the \texttt{thickener\_generic} block.}
\centering
\begin{tabular}{>{\raggedright\arraybackslash}p{2.2cm}
				>{\raggedright\arraybackslash}p{1.5cm}
				>{\raggedright\arraybackslash}p{1.5cm}
				>{\raggedright\arraybackslash}p{2.0cm}
				>{\raggedright\arraybackslash}p{1.5cm}
				>{\raggedright\arraybackslash}p{1.5cm}
				>{\raggedright\arraybackslash}p{1.2cm}
				>{\raggedright\arraybackslash}p{2.0cm}}
\toprule
\textbf{Scenario} & \textbf{SA (ft$^2$/tpd)} & \textbf{D (ft)} & \textbf{Thk bare} & \textbf{FP bare} & \textbf{Tot bare} & \textbf{PF} & \texttt{prim\_equip\_price\_base} \\
\midrule
Cons. & 10 & 71 & \$504K & \$200K & \$704K & 3.5$\times$ & \textbf{\$2,450,000} \\
Mid.  & 7  & 59 & \$431K & \$130K & \$561K & 3.0$\times$ & \textbf{\$1,700,000} \\
Opt.  & 6  & 55 & \$403K & \$75K  & \$478K & 2.5$\times$ & \textbf{\$1,200,000} \\
\bottomrule
\end{tabular}
\end{table}

\noindent \textit{Derivation detail (midpoint row).} The settling area of 7 ft$^2$/tpd (rounded from the Thacker Pass value of 7.1) at 396 tpd gives a required area of 2,772 ft$^2$ and a thickener diameter of 59 ft (18 m). The Mular correlation gives a bare thickener cost of \$210,200 in 1998 USD, which inflates to \$437,100 in 2025 USD via CEPCI. Adding a 630--800 mm filter press at \$130,000 bare gives a total bare equipment cost of \$567,100. Applying a total plant cost factor of 3.0$\times$ yields the installed cost of \$1,701,300, rounded to \$1,700,000.

\medskip
\noindent\textbf{Electricity.}\quad

Electricity demand is the sum of continuous loads from the gravity thickener (rake drive, underflow pump, flocculant dosing) and cycle-averaged loads from the filter press (hydraulic power unit, feed pump, plate shifter, cloth wash, and compressed air). The filter press operates in batch cycles while the thickener runs continuously; the cycle-averaged loads for the press reflect intermittent operation at 50--70\% load factor.

The rake drive power is anchored to two independent data points: the Thacker Pass classification thickener at 30 kW for a 55 m unit (Table~14-4 of~\citet{sgscanadaincSK1300Technical2024}), and a manufacturer specification of 5.5 kW for a 24 m unit~\cite{jinhaominingHotItemLarge}. Both data points are consistent with power scaling as $D^2$:
\[
	P_{\text{rake}} = 30 \times (18/55)^2 = 3.2 \;\text{kW}; \quad P_{\text{rake}} = 5.5 \times (18/24)^2 = 3.1 \;\text{kW}
\]
The midpoint thickener at 18 m diameter therefore draws approximately 3 kW for the rake mechanism.

\begin{table}[h!]
\caption{Electricity derivation for the \texttt{thickener\_generic} block.}
\centering
\begin{tabular}{>{\raggedright\arraybackslash}p{5.0cm}
				>{\raggedright\arraybackslash}p{2.0cm}
				>{\raggedright\arraybackslash}p{2.0cm}
				>{\raggedright\arraybackslash}p{2.0cm}}
\toprule
\textbf{Component} & \textbf{Opt.\ (kW)} & \textbf{Mid.\ (kW)} & \textbf{Cons.\ (kW)} \\
\midrule
Thickener rake drive              & 3   & 3    & 3   \\
Underflow pump                    & 5   & 8    & 12  \\
Overflow pump + flocculant dosing & 4   & 5    & 5   \\
Filter press hydraulic unit (cycle-avg.)     & 8   & 11   & 15  \\
Filter press feed pump (cycle-avg.) & 8 & 10   & 14  \\
Press auxiliaries (shifter, air, cloth wash) & 3 & 5 & 7 \\
Instrumentation \& controls       & 2   & 3    & 3   \\
\midrule
\textbf{Total}                    & \textbf{33}  & \textbf{45} & \textbf{59} \\
\texttt{electricity\_base\_total} & \textbf{35 kWh/h} & \textbf{45 kWh/h} & \textbf{60 kWh/h} \\
\bottomrule
\end{tabular}
\end{table}

\noindent The thickener's continuous loads (rake, pumps, flocculant) account for approximately 36\% of the midpoint total, with the filter press cycle-averaged loads accounting for the remainder. The Thacker Pass disclosure uses decanter centrifuges for the dewatering stage, which are substantially more power-intensive (355 kW main drive per unit); the filter press assumption in this block gives a lower electricity draw and represents a generic rather than Thacker Pass-specific configuration.

\medskip
\noindent\textbf{Consumables.}\quad

Polymer flocculant is tracked at the production step level as a secondary input with a dosage of $2.6 \times 10^{-4}$ \si{kg/m^3} of slurry feed, calibrated to the Thacker Pass disclosed consumption of approximately 260 g/t of solids~\citep{sgscanadaincSK1300Technical2024}. Filter cloths for the press are a recurring consumable tracked through \texttt{tooling\_cost\_base}, which are currently unmodeled.

\medskip
\noindent\textbf{Labor.}\quad

Labor is 1.8 FTE at base scale with a scaling exponent of 0.2 and \texttt{dedicated\_labor} = FALSE, consistent with a highly automated thickener (e.g. automated rake controls and flocculant dosing) combined with a semi-automated filter press requiring periodic supervisory attention for cake discharge and cloth condition:
\[
	\text{FTE} = 0.4 \times 4.5 = 1.8
\]

\medskip
\noindent\textbf{Footprint.}\quad

The base plot area is 600~$\si{m^2}$ with a scaling exponent of 0.7. At the midpoint thickener diameter of 18~m, the plan area of the thickener tank alone is approximately 250~$\si{m^2}$, leaving 350~$\si{m^2}$ for the filter press, piping, and access.

\medskip
\noindent\textbf{Availability.}\quad

Process availability is 92\% with 3\% unplanned downtime and 5\% scheduled maintenance. The high availability reflects the robust, low-complexity nature of gravity settling operations. The filter press component has a somewhat lower individual availability, but the combined block maintains 92\% overall because the thickener provides buffer capacity in its underflow inventory.

\medskip
\noindent\textbf{Maintenance.}\quad

Neither \texttt{maint\_override} nor a nonzero \texttt{opex\_fraction\_of\_capex} is set; the block inherits the facility-level default maintenance rate of 5\% of scaled primary equipment cost, charged against the total undiscounted equipment price. The gravity thickener has modest maintenance requirements (rake bearings, overflow weir adjustment), while the filter press requires regular cloth replacement and hydraulic seal maintenance.

\clearpage
\paragraph{\texttt{multi\_stage\_ccd\_thickening}}\label{app:04-methods-blockdefs-ccd}\leavevmode\par

\medskip
\noindent\textbf{Sources.}\quad

\citet{sgscanadaincSK1300Technical2024}, \citet{mularCAPCOSTSHandbookEstimating1998}; filter press costs from published benchmarks (\citet{porvooIndustrialFilterPress2025}). 

\medskip
\noindent\textbf{Process description.}\quad

The multi-stage CCD block separates the pregnant leach solution (PLS) from the solid leach residue through eight stages of counter-current decantation (including one filter press step), recovering dissolved lithium from entrained solution while minimizing lithium losses to the filter cake. The input is post-neutralization slurry at approximately 34 wt\% solids. In each stage, the slurry is fed to a thickener where the solids settle under gravity (aided by flocculant), producing a clarified overflow and a dense underflow. Fresh wash water enters the last stage and flows counter-current to the solids, progressively diluting the entrained mother liquor. The first-stage overflow is the clarified PLS containing the dissolved lithium; the eighth-stage underflow is the washed leach residue routed to final filtration.

The Thacker Pass disclosure describes the CCD thickeners as ``Hi-Density'' units, achieving a denser underflow ($\sim$33 wt\% solids per stage) than a conventional high-rate thickener. The block also includes a final leach-residue filter press to dewater the eighth-stage underflow for disposal or tailings stacking.

\medskip
\noindent\textbf{Bill of Equipment.}\quad

\begin{table}[h!]
\caption{Bill of equipment for the \texttt{multi\_stage\_ccd\_thickening} block.}
\centering
\begin{tabular}{>{\raggedright\arraybackslash}p{5.2cm}
				>{\raggedright\arraybackslash}p{4cm}
				>{\raggedright\arraybackslash}p{4.5cm}}
\toprule
\textbf{Equipment} & \textbf{Machine data used} & \textbf{Source(s)} \\
\midrule
8 Hi-Density CCD thickeners (17--25 m diameter at base scale) & CAPCOSTS Fig.~190 & \citet{mularCAPCOSTSHandbookEstimating1998} \\
Leach residue filter press (630--800 mm plates) & Porvoo 2025 & \citet{porvooIndustrialFilterPress2025} \\
Underflow and wash-water pumps (per stage) & Not explicit $\rightarrow$ inferred & --- \\
Flocculant dosing system & Not explicit $\rightarrow$ inferred & --- \\
Counter-current piping manifolds & Not explicit $\rightarrow$ inferred & --- \\
Instrumentation \& controls & Not explicit $\rightarrow$ inferred & --- \\
\bottomrule
\end{tabular}
\end{table}

\medskip
\noindent\textbf{Sizing basis.}\quad

The base volume is 100 $\si{m^3/h}$ of CCD feed slurry at 34 wt\% solids. The stream is acid-inclusive, so at grain density 2.6 \si{t/m^3} its density is 1.399 \si{t/m^3} and the cake solids loading, leach residue plus neutralization gypsum, is 36.5 t/h (877 t/day). The \texttt{oversizing\_factor} = 1.15, \texttt{scaling\_exponent} = 0.60, \texttt{process\_type} = continuous, \texttt{dedicated\_line} = FALSE. Equipment life is assumed to be 15 years.

The CCD stages are sized by their area, which fixes each stage's diameter. The feed is flocculated leach residue, which should settle quickly, so each unit's area belongs at or under the lower edge of the 4--10 ft$^2$/tpd copper-tailings range~\citep[p.~270]{mularCAPCOSTSHandbookEstimating1998}. The Thacker Pass disclosure specifies an eight-stage 36~m CCD train at 3.9 ft$^2$/tpd~\citep{sgscanadaincSK1300Technical2024}. The scenarios thus sweep 2.6 / 3.9 / 5.2 ft$^2$/tpd; bare cost goes as $D^{0.8786}$ with $D \propto \mathrm{SA}^{1/2}$. 

\medskip
\noindent\textbf{CAPEX.}\quad

CAPEX is built bottom-up from the CAPCOSTS high-capacity thickener correlation applied to each of the 8 CCD stages, plus auxiliary equipment.

\textit{Thickener per stage.} The Mular correlation for high-capacity thickeners (16--100~ft range) is $\text{Price}_{1998} = 5,809 \times D^{0.8786}$, inflated to 2025 via CEPCI ($\times$2.08). Each CCD stage requires settling area = SA $\times$ 877 tpd, giving diameter $D = \sqrt{4A/\pi}$. The eight-stage bare thickener cost is 8$\times$ the per-stage cost.

\textit{Auxiliaries.} The CCD circuit requires underflow pumps (one per stage), wash-water distribution pumps, a shared flocculant dosing system, counter-current piping manifolds, and a leach-residue filter press. Auxiliary costs are estimated at \$350K--\$900K (scenario-dependent) based on the Porvoo filter press benchmark (\$75K--\$200K) plus inferred pump and piping costs.

\textit{Plant factor.} The total plant cost factor for an 8-stage CCD circuit is lower than the 2.5--3.5$\times$ used for a standalone thickener because the eight thickeners share common foundations, flocculant systems, piping manifolds, control systems, and EPCM. A factor of 1.3--1.7$\times$ is applied.

\begin{table}[h!]
\caption{CAPEX derivation for the \texttt{multi\_stage\_ccd\_thickening} block.}
\centering
\begin{tabular}{>{\raggedright\arraybackslash}p{2.0cm}
				>{\raggedright\arraybackslash}p{1.3cm}
				>{\raggedright\arraybackslash}p{1.5cm}
				>{\raggedright\arraybackslash}p{1.8cm}
				>{\raggedright\arraybackslash}p{1.5cm}
				>{\raggedright\arraybackslash}p{1.5cm}
				>{\raggedright\arraybackslash}p{1.0cm}
				>{\raggedright\arraybackslash}p{2.5cm}}
\toprule
\textbf{Scenario} & \textbf{SA} & \textbf{D/stg} & \textbf{8$\times$stg bare} & \textbf{Aux} & \textbf{Tot bare} & \textbf{PF} & \texttt{prim\_equip\_price\_base} \\
\midrule
Cons. & 5.2 & 76.2 ft & \$4,289K & \$900K & \$5,189K & 1.5$\times$ & \textbf{\$7,800,000} \\
Mid.  & 3.9 & 66.0 ft & \$3,780K & \$600K & \$4,380K & 1.5$\times$ & \textbf{\$6,600,000} \\
Opt.  & 2.6 & 53.9 ft & \$3,163K & \$350K & \$3,513K & 1.5$\times$ & \textbf{\$5,300,000} \\
\bottomrule
\end{tabular}
\end{table}

\noindent \textit{Derivation detail (midpoint row).} At 3.9 ft$^2$/tpd and 877 tpd, the required area per stage is 3,419 ft$^2$, giving $D = 66.0$ ft (20.11 m). The Mular correlation gives \$230K per stage in 1998 USD, inflating to \$473K in 2025. Eight stages total \$3,780K bare. Adding \$600K in auxiliaries gives \$4,380K total bare, which at 1.5$\times$ plant factor yields \$6,600K installed.

\medskip
\noindent\textbf{Electricity.}\quad

Electricity demand is derived from two independently anchored components per CCD stage: a rake drive and an underflow pump, plus shared loads.
 
\textit{Rake drive.} The rake motor power is anchored to the Jinhao Mining NZ-24 specification of 5.5~kW for a 24~m center-drive thickener (a pure rake motor rating), and scaled to the per-stage diameter via $P \propto D^2$~\cite{jinhaominingHotItemLarge}. At the midpoint diameter of 20.11~m, the rake power per stage is $5.5 \times (20.11/24)^2 = 3.9$~kW. 
 
\textit{Underflow pump.} Each stage's underflow at 33 wt\% solids (density 1.255 \si{t/m^3}) is pumped to the next stage. At 100 \si{m^3/h} CCD feed with 36.5 t/h cake solids, the underflow volumetric flow per stage is $36.5/0.33/1.255 = 88.2$ \si{m^3/h} $= 0.0245$ \si{m^3/s}. Pump power is $P = \rho g Q H / \eta = 1255 \times 9.81 \times 0.0245 \times 10 / 0.60 = 5.0$ kW per stage at 10 m head and 60\% efficiency.

\textit{Wash-water and overflow pumping.} Two duties named in the auxiliary equipment scope carry loads of the same order as the interstage pumps. Fresh wash water enters the last stage at 0.615 \si{m^3} per \si{m^3} of CCD feed, drawing 2.8 kW at the same head and efficiency. The first-stage overflow is the clarified PLS at 1.206 \si{m^3} per \si{m^3} of feed, 120.6 \si{m^3/h} and 1.115 \si{t/m^3}, drawing 6.1 kW to move it forward.
 
\textit{Cross-check at Thacker Pass scale.} At Phase 1 (36 m thickeners, 2,900 tpd of residue), the same derivation gives 12.4 kW rake $+ 16.6$ kW pump $= 29.0$ kW per stage, against the disclosed 30 kW per thickener (Table~14-4 of~\citet{sgscanadaincSK1300Technical2024}). 
 
\begin{table}[h!]
\caption{Electricity derivation for the \texttt{multi\_stage\_ccd\_thickening} block.}
\centering
\begin{tabular}{>{\raggedright\arraybackslash}p{7.0cm}
				>{\raggedright\arraybackslash}p{2.5cm}
				>{\raggedright\arraybackslash}p{2.5cm}
				>{\raggedright\arraybackslash}p{2.5cm}}
\toprule
\textbf{Component} & \textbf{Opt.\ (kW)} & \textbf{Mid.\ (kW)} & \textbf{Cons.\ (kW)} \\
\midrule
Rake drive per stage ($D^2$ from Jinhao)            & 2.6  & 3.9  & 5.2 \\
Underflow pump per stage ($P = \rho g Q H / \eta$)  & 4.0  & 5.0  & 6.0 \\
$\times$ 8 stages                                   & 53   & 71   & 89  \\
Shared (flocculant, I\&C)                           & 8    & 8    & 8   \\
Wash-water distribution pump                        & 2.2  & 2.8  & 3.4 \\
Stage-1 overflow pump                               & 4.9  & 6.1  & 7.3 \\
\midrule
\texttt{electricity\_base\_total} (Rounded)     & \textbf{70 kWh/h} & \textbf{90 kWh/h} & \textbf{110 kWh/h} \\
\bottomrule
\end{tabular}
\end{table}
 
\noindent The scenario variation in the underflow pump (4.0--6.0~kW) reflects different assumed pumping heads: 8~m for optimistic (compact installation), 10~m for midpoint, and 12~m for conservative (taller thickeners with longer transfer lines); the same head applies to the wash and overflow pumps. The rake and pump components contribute roughly equally to the per-stage total at the midpoint (3.9 vs 5.0~kW), with the pump dominating slightly because of the dense underflow. The leach-residue filter press remains in this block's bill of equipment and its capital cost sits in the auxiliary allowance, but its operating duty is not priced here. A press taking last stage of underflow from 33 to 62~wt\% solids would draw 140~kW at the block's base scale, against the 8~kW allowance above, so pricing it in a shared rounding term would understate it by more than an order of magnitude. It is charged instead at \texttt{claycake\_stack\_nv\_62} through that sink's $s_{\text{feed}}$ declaration, where the duty scales on the calculated amount of water removed rather than on the volume fed.

\medskip
\noindent\textbf{Consumables.}\quad

Flocculant is tracked at the production step level at 0.04 \si{kg/m^3} of CCD feed, calibrated to the Thacker Pass disclosure's reported polymer dosage for the CCD circuit. Filter cloths for the press are a recurring consumable tracked through \texttt{tooling\_cost\_base}, which are currently unmodeled.

\medskip
\noindent\textbf{Labor.}\quad

Labor is 1.8 FTE at base scale with a scaling exponent of 0.2 and \texttt{dedicated\_labor} = FALSE. The low operator count reflects the highly automated nature of CCD circuits, where counter-current flow rates, flocculant dosing, and underflow densities could be managed automatically with minimal manual intervention:
\[
	\text{FTE} = 0.4 \times 4.5 = 1.8
\]

\medskip
\noindent\textbf{Footprint.}\quad

At the midpoint thickener diameter of 20.11~m per stage, the plan area of eight thickeners arranged in a double row is approximately $8 \times \pi \times 10.06^2 = 2,541$~$\si{m^2}$

\medskip
\noindent\textbf{Availability.}\quad

Process availability is 92\% with 3\% unplanned downtime and 5\% scheduled maintenance. The CCD circuit has inherent redundancy in that the loss of a single stage degrades washing efficiency but does not halt production; however, pump failures in the counter-current circuit could cascade rapidly.

\medskip
\noindent\textbf{Maintenance.}\quad

Neither \texttt{maint\_override} nor a nonzero \texttt{opex\_fraction\_of\_capex} is set; the block inherits the facility-level default maintenance rate of 5\% of scaled primary equipment cost, charged against the total undiscounted equipment price. The CCD thickeners themselves have low maintenance requirements (rake bearings, flocculant injection nozzles), but the inter-stage pumps and final filter press have higher wear rates due to the abrasive post-leach slurry.

\medskip
\noindent\textbf{Key measurements and cross-checks.}\quad

\begin{itemize}[noitemsep]
	\item \textit{CAPEX}: the midpoint of \$6.6M represents eight 20.11~m Hi-Density thickeners plus a filter press and auxiliaries, at a 1.5$\times$ total plant cost factor for the series installation. The implied settling area of 3.9 ft$^2$/tpd per stage is anchored to the Thacker Pass 36~m CCD thickeners and falls below the Cu-tail range (4--10), consistent with faster-settling post-leach amorphous silica residue.
	\item \textit{Thacker Pass note}: the actual Thacker Pass Phase~1 CCD uses 8 thickeners at 36~m diameter for an estimated throughput of $\sim$280 \si{m^3/h}. The model would represent this as $\sim$2.8 machines at base\_volume = 100 \si{m^3/h}, scaled via the six-tenths rule: CAPEX = \$6.6M $\times$ $2.8^{0.6}$ = \$12.2M. This is broadly consistent with the scale of the Thacker Pass process plant capital budget.
\end{itemize}

\clearpage
\paragraph{\texttt{mechanical\_dewatering}}\label{app:04-methods-blockdefs-dewater}\leavevmode\par

\medskip
\noindent\textbf{Sources.}\quad

\citet{mularCAPCOSTSHandbookEstimating1998}, \citet{porvooIndustrialFilterPress2025}. Bare equipment costs are cross-checked against both the CAPCOSTS vertical automated plate-and-frame filter press correlation (Figure~182, p.~263)~\citep{mularCAPCOSTSHandbookEstimating1998} and the Porvoo published filter press price benchmarks~\citep{porvooIndustrialFilterPress2025}.

\medskip
\noindent\textbf{Process description.}\quad

The mechanical dewatering block converts thickened or classified slurry into a handleable filter cake by pressure filtration, removing interstitial water through a combination of hydraulic pressure and membrane squeeze. The block is sized for the lepidolite dewatering duty in the Jianxiawo pathway, where ground lepidolite concentrate at 60 wt\% solids is dewatered to 80 wt\% solids prior to sulfate roasting. The roasting train requires a low-moisture feed; excess water would consume thermal energy in evaporation and could cause handling problems.

The primary equipment is a bank of large automated plate-and-frame filter presses operating in parallel batch cycles. At 60 wt\% feed solids (already well within the plate-and-frame operating range of 25--60\%), no upstream belt filter stage is assumed --- the presses receive the dense slurry directly from the grinding circuit.

\medskip
\noindent\textbf{Bill of Equipment.}\quad

\begin{table}[h!]
\caption{Bill of equipment for the \texttt{mechanical\_dewatering} block.}
\centering
\begin{tabular}{>{\raggedright\arraybackslash}p{5.2cm}
				>{\raggedright\arraybackslash}p{4cm}
				>{\raggedright\arraybackslash}p{4.5cm}}
\toprule
\textbf{Equipment} & \textbf{Machine data used} & \textbf{Source(s)} \\
\midrule
Automated membrane filter presses (5--8 units, 1000--1500 mm plates) & CAPCOSTS Fig.~182; Porvoo 2025 & \citet{mularCAPCOSTSHandbookEstimating1998, porvooIndustrialFilterPress2025} \\
High-pressure feed slurry pumps and piping & Not explicit $\rightarrow$ inferred & --- \\
Cake discharge conveyors & Not explicit $\rightarrow$ inferred & --- \\
Filtrate collection and pumping & Not explicit $\rightarrow$ inferred & --- \\
Cloth washing system & Included in press packages & Vendor standard \\
Compressed air system & Not explicit $\rightarrow$ inferred & --- \\
Structural platform and access & Not explicit $\rightarrow$ inferred & --- \\
Instrumentation and controls & Not explicit $\rightarrow$ inferred & --- \\
\bottomrule
\end{tabular}
\end{table}

\medskip
\noindent\textbf{Sizing basis.}\quad

The base volume is 60 \si{m^3} per batch of slurry being pressed. At 60 wt\% solids with grain density 2.7 \si{t/m^3}, the slurry density is 1.607 \si{t/m^3}. Each batch cycle is 1.2 hours (1.0 hour pressing + 0.2 hours setup and cake discharge), giving an effective throughput of 50 $\si{m^3/h}$. The solids throughput is $50 \times 1.607 \times 0.60 = 48.2$ t/h. The filtrate production is approximately 20 \si{m^3/h} (the difference between the water content at 60 wt\% and 80 wt\% solids).

The product cake at 80 wt\% solids has a density of $1/(0.8/2.7 + 0.2/1.0) = 2.015$ \si{t/m^3} and a volume of approximately 36 \si{m^3} per batch. Each press completes one full cycle per 1.2 h batch, consistent with the \texttt{batch\_cycle\_time} = 1.0 h and \texttt{batch\_setup\_time} = 0.2 h carried in the model. A 1500 mm press with 100 chambers and 40 mm cake thickness holds $100 \times 1.95\;\si{m^2} \times 0.040\;\si{m} = 7.8$ \si{m^3} of cake per cycle, so the base duty requires $36 / 7.8 = 4.6 \rightarrow 5$ presses. Thus the optimistic scenario carries 5 presses (exactly sized), the midpoint carries 6 (with one on maintenance/standby rotation), and the conservative scenario carries 8 (to account for derating). Sizing here is to the base duty; the model's \texttt{oversizing\_factor} = 1.25 is applied on top of this and is not double-counted in the integer press counts.

The \texttt{scaling\_exponent\_throughput} = 1.0 (linear volumetric batch-scheduling basis) and \texttt{prim\_equip\_scaling\_exponent} = 0.70. The 0.70 CAPEX exponent reflects that capacity growth in press plants is met first by moving to larger press frames rather than by replicating units: industrial presses span 470 mm to 2500 mm plates ($>$20$\times$ filtration area per machine), so within the modeled range equipment cost follows the area-based economy of scale of the CAPCOSTS correlation ($\sim$0.6 on filtration area) plus dilution of shared auxiliaries, blending to $\sim$0.7 at the block level. Unit replication, which would push the exponent toward 1.0, becomes binding only beyond the largest available frames. Equipment life is 20 years.

\medskip
\noindent\textbf{CAPEX.}\quad

CAPEX is built bottom-up from the per-press bare equipment cost multiplied by the number of presses, plus auxiliaries and a total plant cost factor.

\textit{Filter presses.} The CAPCOSTS automated vertical plate-and-frame correlation gives \$146,000 (1998 USD) at 800 ft$^2$ (74 \si{m^2}) filtration area, inflating to \$303,600 in 2025 via CEPCI (index ratio 810/389.5; the 2025 CEPCI value of $\sim$810 is an assumption pending the published annual index). A 1500 mm, 100-chamber press has $100 \times 2 \times 1.95 = 390$ \si{m^2} of filtration area; scaling the anchor with an area exponent of 0.55--0.65 (the correlation slope; rule-of-thumb range, as the Fig.~182 slope was not re-digitized) gives \$757K--\$894K per press before membrane equipment. Porvoo reports a 35--50\% cost premium for membrane configurations over standard presses~\citep{porvooIndustrialFilterPress2025}. Two membrane-inclusive cross-checks from Porvoo bracket the result: scaling the 1000 mm top-end (\$350K, $\sim$97 \si{m^2}) to 390 \si{m^2} at exponent 0.6 gives \$807K, and scaling the 800 mm membrane benchmark (\$180--280K at $\sim$34 \si{m^2}) gives $\sim$\$998K. Thus, adopted per-press bare costs are \$750K (optimistic; competitive procurement at the low end of the anchors), \$850K (midpoint; within the membrane-inclusive cross-check band), and \$1,100K (conservative; CAPCOSTS mid-scaling plus the full 35\% membrane premium).

\textit{Installation.} A multiplicative factor of 1.5--2.5$\times$ bare equipment is applied for auxiliaries, given that presses arrive as largely pre-assembled packaged units; the scope covers erection on the elevated platform, structural steel, piping, and electrical distribution.

\textit{Auxiliaries.} High-pressure feed pumps, cake conveyors, filtrate collection, the squeeze-air compressor, cloth-wash system, and instrumentation are applied as 20--40\% of the installed subtotal.

\begin{table}[h!]
\caption{CAPEX derivation for the \texttt{mechanical\_dewatering} block.}
\centering
\begin{tabular}{>{\raggedright\arraybackslash}p{1.8cm}
				>{\raggedright\arraybackslash}p{2.3cm}
				>{\raggedright\arraybackslash}p{1.8cm}
				>{\raggedright\arraybackslash}p{1.8cm}
				>{\raggedright\arraybackslash}p{1.3cm}
				>{\raggedright\arraybackslash}p{2.7cm}}
\toprule
\textbf{Scenario} & \textbf{Presses (raw)} & \textbf{Presses (tot.)} & \textbf{Inst.\ subtotal} & \textbf{Aux} & \texttt{prim\_equip\_price\_base} \\
\midrule
Cons. & 8 $\times$ \$1,100K & \$8,800K & \$22,000K & \$8,800K &  \textbf{\$30,800,000} \\
Mid.  & 6 $\times$ \$850K & \$5,100K & \$10,200K & \$3,060K  & \textbf{\$13,300,000} \\
Opt.  & 5 $\times$ \$750K & \$3,750K & \$5,625K & \$1,125K  & \textbf{\$6,750,000} \\
\bottomrule
\end{tabular}
\end{table}

\noindent \textit{Derivation detail (midpoint row).} Six 1500 mm membrane presses at \$850K each give \$5,100K raw. Installed subtotal: \$5,100K $\times$ 2.0 = \$10,200K. Auxiliaries at 30\% of installed: \$3,060K. Total: \$13,260K, rounded to \$13,300K. The effective total factor on bare equipment is 1.8--2.6--3.5$\times$ across scenarios, with the midpoint consistent with Mular total-plant ratios.

\medskip
\noindent\textbf{Electricity.}\quad

\textit{Feed pumps.} $P = \dot{V} \Delta P / \eta$ with $\dot{V}$ = 50 \si{m^3/h} continuous-equivalent, a cycle-averaged delivery pressure of 6--8--10 bar (assumption: feed pressure ramps from $\sim$2 bar during fill to 15 bar at the end of pressing), and pump-motor efficiency 0.65--0.60--0.55:
\[
	P_\text{mid} = \frac{(50/3600)\;\si{m^3/s} \times 8\times10^5\;\text{Pa}}{0.60} = 18.5\;\text{kW}
\]
 
\textit{HPUs.} Installed HPU power of 45--55 kW per 1500 mm press (assumption), at full load $\sim$10\% of the cycle plus 1--2 kW clamp-holding, gives a cycle-averaged 4--6--8 kW per press (assumption).
 
\textit{Squeeze/blow air compressor.} Derived from the compressed-air demand below at $\sim$0.11 kW per \si{Nm^3/h} delivered at 8 bar (screw-compressor rule-of-thumb).
 
\begin{table}[h!]
\caption{Electricity derivation for the \texttt{mechanical\_dewatering} block.}
\centering
\begin{tabular}{>{\raggedright\arraybackslash}p{6.2cm}
				>{\raggedright\arraybackslash}p{2.0cm}
				>{\raggedright\arraybackslash}p{2.0cm}
				>{\raggedright\arraybackslash}p{2.0cm}}
\toprule
\textbf{Component} & \textbf{Opt.\ (kW)} & \textbf{Mid.\ (kW)} & \textbf{Cons.\ (kW)} \\
\midrule
Feed pumps (50 \si{m^3/h} at 6--10 bar avg.) & 13 & 19 & 25 \\
Press HPUs (cycle-avg., $N \times$ 4--8 kW) & 20 & 36 & 64 \\
Squeeze/blow air compressor & 4 & 7 & 10 \\
Cake discharge conveyors ($\sim$60 t/h cake) & 4 & 6 & 8 \\
Cloth wash pump (periodic) & 2 & 3 & 5 \\
Instrumentation and controls & 2 & 3 & 3 \\
\midrule
\texttt{electricity\_base\_total} & \textbf{45 kWh/h} & \textbf{74 kWh/h} & \textbf{115 kWh/h} \\
\bottomrule
\end{tabular}
\end{table}

\noindent \textit{Cross-check.} Specific energy is 0.93--1.54--2.39 kWh/t solids across scenarios, within or adjacent to the 1--3 kWh/t rule-of-thumb band for pressure filtration; the optimistic value sits just below the band, reflecting large frames and long cycles.

\medskip
\noindent\textbf{Other utilities.}\quad

\textit{Compressed air.} Membrane squeeze displaces the filtrate expressed during the squeeze phase, taken as 10--20\% of the 24 \si{m^3} of filtrate per batch (assumption), delivered at 8--10 bar; cake blow-down consumes 0.8--1.2$\times$ the squeeze volume (assumption). Midpoint: $24 \times 0.15 \times 9\;\text{(pressure ratio)} / 1.2\;\text{h} = 27$ \si{Nm^3/h} squeeze $+$ 27 \si{Nm^3/h} blow-down $\approx$ 55 \si{Nm^3/h}. \texttt{compressed\_air\_base\_total} = 30--55--95 \si{Nm^3/h}. 

\textit{Cooling water.} HPU oil-cooler duty is taken as 40--60\% of cycle-averaged HPU electrical load (assumption). Midpoint: $36 \times 0.5 = 18$ kW rejected; at $\Delta T$ = 10$^\circ$C,
\[
	\dot{m} = \frac{18\;\text{kW}}{4.18\;\text{kJ/(kg$\cdot^\circ$C)} \times 10\;^\circ\text{C}} = 0.43\;\text{kg/s} \approx 1.5\;\si{m^3/h}
\]

Thus  \texttt{cooling\_water\_base\_total} = 0.7--1.5--3.0 \si{m^3/h}.

\medskip
\noindent\textbf{Consumables.}\quad

Filter cloths are the block's recurring consumable and are carried as a tooling cost. Porvoo reports cloth costs of \$200--\$800 per piece with replacement every 6--18 months depending on application severity~\citep{porvooIndustrialFilterPress2025}. The abrasive, blinding-prone lepidolite duty would likely sit mid-to-high in that severity range, giving $c_\text{cloth}$ = \$400--\$500--\$600 (assumption: selection within the cited range).
 
The number of tools consumed per year is the installed cloth population divided by cloth life. Each press carries $\sim$100 cloths (assumption: one drape-over cloth per plate, plates $\approx$ chambers), so the installed population is $N_\text{press} \times 100$, and cloth life is 18--12--6 months (assumption: within the cited 6--18 month range, worsening from optimistic to conservative):
\[
	\dot{n}_\text{cloth} = \frac{N_\text{press} \times 100\;\text{cloths}}{L_\text{cloth}\;(\text{yr})}
\]
 
\begin{itemize}[noitemsep]
	\item Optimistic: $\dot{n} = (5 \times 100)/1.5 = 333$ cloths/yr $\times$ \$400 = \$133K/yr
	\item Midpoint: $\dot{n} = (6 \times 100)/1.0 = 600$ cloths/yr $\times$ \$500 = \$300K/yr
	\item Conservative: $\dot{n} = (8 \times 100)/0.5 = 1,600$ cloths/yr $\times$ \$600 = \$960K/yr
\end{itemize}
 
\noindent \texttt{tooling\_cost\_base} = \$133K--\$300K--\$960K per year at base scale, with \texttt{tooling\_scaling\_exponent} = 1.0: at fixed cloth life, $\dot{n}_\text{cloth}$ scales linearly with the press count and hence with throughput, so the tool count per year --- not the cost per tool --- carries the scaling. This follows the unit convention of the \texttt{excavation\_loading\_generic} tire cost (tools/machine block/year $\times$ \$/tool). The equivalent intensity is \$0.35--\$0.79--\$2.53 per tonne of solids. Cloth costs are excluded from the maintenance scope to avoid double-counting.

\medskip
\noindent\textbf{Labor.}\quad

Labor is 0.9 FTE at base scale with a scaling exponent of 0.2 and \texttt{dedicated\_labor} = FALSE, consistent with a bank of automated presses requiring periodic supervisory attention for cake discharge quality, cloth condition monitoring, and cycle parameter adjustment:
\[
	\text{FTE} = 0.2 \times 4.5 = 0.9
\]

\medskip
\noindent\textbf{Footprint.}\quad

The base plot area is 250~$\si{m^2}$ with a scaling exponent of 0.7, covering a linear bank of 5--8 plate-and-frame presses with auxiliaries.

\medskip
\noindent\textbf{Availability.}\quad

Process availability is 90\% with 5\% unplanned downtime and 5\% scheduled maintenance. The parallel press configuration provides partial redundancy: the loss of a single press reduces throughput but does not halt production.

\medskip
\noindent\textbf{Maintenance.}\quad

Neither \texttt{maint\_override} nor a nonzero \texttt{opex\_fraction\_of\_capex} is set; the block inherits the facility-level default maintenance rate of 5\% of scaled primary equipment cost, charged against the total undiscounted equipment price. Filter presses are mechanically intensive equipment with regular cloth replacement, hydraulic seal maintenance, and plate inspection requirements; cloth replacement is carried separately in the tooling cost.

\medskip
\noindent\textbf{Key measurements and cross-checks.}\quad

\begin{itemize}[noitemsep]
	\item \textit{CAPEX}: the midpoint of \$13.3M is 6 presses at \$850K raw, a 2.0$\times$ installation factor, and 30\% auxiliaries, giving an effective 2.6$\times$ total factor on bare equipment. The per-press price is bracketed by two independent membrane-inclusive scalings of the Porvoo benchmarks (\$807K, \$998K) and by the CAPCOSTS correlation with and without the membrane premium. We can further bracket the overall installed costs through two other sources: \citet{hahnWhitepaper_Tailings2024} reports a 15-press, 150,000 t/d installation at \texteuro 195M equipment-plus-auxiliaries (\texteuro 13M per press for $\sim$2000--2500 \si{m^2} machines; area-normalizing to this block's 390 \si{m^2} frame gives $\sim$\texteuro 5,800/m$^2 \rightarrow \sim$\$2.45M equipment per press, above the adopted \$0.75--1.1M raw. The same source cites a turnkey multiplier of 2.5--3.5$\times$ on equipment for building, piping, electrical, engineering, and commissioning, which brackets this block's effective 1.8--2.6--3.5$\times$ total factor. Simultaneously, a technical report for the New Mexico Copper Concentrator (\citep{wspusaincCakeStackFiltered2025}) reports \$40M equipment $+$ \$14M installed labor for a full tailings pressure-filtration circuit, the right order of magnitude for a multi-press plant but with no per-press count in the excerpt to derive a unit cost. Both benchmarks are tailings duties --- larger-tonnage and less value-dense than the assumed lepidolite-concentrate duty --- and no lepidolite-specific press-cost benchmark was located; on these references the midpoint \$13.3M is supported and the conservative \$30.8M sits within the industry turnkey-factor range.
	\item \textit{Electricity}: 0.93--2.39 kWh/t solids across scenarios against a 1--3 kWh/t rule-of-thumb band for pressure filtration, with demand shared between feed pumping and cycle-averaged hydraulics rather than HPU-dominated.
\end{itemize}

\clearpage
\paragraph{\texttt{staged\_neutralization}}\label{app:04-methods-blockdefs-neutralization}\leavevmode\par

\medskip
\noindent\textbf{Sources.}\quad

\citet{sgscanadaincSK1300Technical2024}, \citet{battagliaRecoveryLithiumCarbonate2022}, \citet{liuProcessAnalysisStudy2021}, \citet{dengSystematicReviewSulfate2026}, \citet{yanExtractionLithiumLepidolite2012}, \citet{ausimmMonograph27Cost}, \citet{petersPlantDesignEconomics1991}; vendor data. Cost methodology scaled from the \texttt{liming\_system} block (Appendix \ref{app:04-methods-blockdefs-limingsystem}).

\medskip
\noindent\textbf{Process description.}\quad

The staged neutralization block removes dissolved metal impurities from pregnant leach solution (PLS) through multi-stage pH-controlled precipitation. Target species include \ce{Al^{3+}}, \ce{Fe^{3+}}, \ce{Mg^{2+}}, and \ce{Ca^{2+}}, with the specific ionic loading depending on the upstream pathway. Some key reactions could include:
\begin{equation}
	\ce{Fe^{3+}(aq) + 3OH^{-} -> Fe(OH)3(s) v} \quad\text{(pH $\sim$3.5--4)}
\end{equation}
\begin{equation}
	\ce{Al^{3+}(aq) + 3OH^{-} -> Al(OH)3(s) v} \quad\text{(pH $\sim$4.5--5.5)}
\end{equation}
\begin{equation}
	\ce{Mg^{2+}(aq) + Ca(OH)2 -> Mg(OH)2(s) v + Ca^{2+}(aq)} \quad\text{(pH $\sim$9--10)}
\end{equation}

Each species precipitates at a distinct pH range, necessitating staged addition of neutralizing agents (limestone for bulk acid neutralization, lime for \ce{Mg} removal, soda ash for residual \ce{Ca} polishing) with intermediate settling or clarification between stages. This block is modeled as a more complex variant of the \texttt{liming\_system} block: it handles a more diverse ionic loading, requires multiple pH-controlled stages, and serves a different functional role (impurity removal from PLS while keeping lithium in solution, rather than simple brine softening). For example, in the sulfate-roast water-leach case, aluminum and iron are largely rejected to the leach residue during roasting and leaching and enter this block only in trace quantities; reagent-derived calcium, supplied by the lime dosed at roasting, is instead the principal dissolved impurity, and thallium the principal regulated one.

This block is used in both the Thacker Pass and Jianxiawo pathways: for bulk acid neutralization and Fe/Al precipitation following sulfuric acid leaching of clay (Thacker Pass), for Mg/Ca removal from washed PLS (Thacker Pass impurity removal), and for staged neutralization of water leach PLS from lepidolite calcine (Jianxiawo). Reagent dosages, elimination fractions, and ionic loading are defined at the production step level and vary by pathway; the block provides the shared equipment-level cost and utility basis.

\medskip
\noindent\textbf{Equation of merit.}\quad
\[
	V_\text{working}\;(\si{m^3}) = \dot{V}_\text{PLS}\;(\si{m^3/h}) \times \tau_\text{total}\;(\text{h})
\]
The working volume of the multi-stage circuit $V_\text{working}$ is sized to deliver a PLS throughput $\dot{V}_\text{PLS}$ with total residence time $\tau_\text{total}$ across multiple neutralization and conditioning stages. At a base flow of 100 $\si{m^3/h}$ and total residence time of $\sim$1.2 h (consistent with the \texttt{liming\_system} and \texttt{precipitator} blocks), the working volume is approximately 120 \si{m^3}, distributed across 3--4 reaction stages plus conditioning. The additional stages relative to the \texttt{liming\_system} (which uses a single-stage cascade) are the primary driver of the complexity premium applied in the CAPEX derivation.

\medskip
\noindent\textbf{Bill of Equipment.}\quad

\begin{table}[h!]
\caption{Bill of equipment for the \texttt{staged\_neutralization} block.}
\centering
\begin{tabular}{>{\raggedright\arraybackslash}p{6.0cm}
				>{\raggedright\arraybackslash}p{4.0cm}
				>{\raggedright\arraybackslash}p{4.5cm}}
\toprule
\textbf{Equipment} & \textbf{Machine data used} & \textbf{Source(s)} \\
\midrule
Multi-stage neutralization/ conditioning tanks (3--4 stages, $\sim$120 \si{m^3} total) & Based on \texttt{liming\_system} tank train & Methodology writeup \\
Lime slaking system (5 \si{m^3} slaker + 25 \si{m^3} storage) & See \texttt{liming\_system} & Methodology writeup \\
Reagent handling (lime, limestone, soda ash) & Not explicit $\rightarrow$ inferred & --- \\
pH control instrumentation (per-stage probes, dosing loops) & Not explicit $\rightarrow$ inferred & --- \\
Slurry transfer pumps \& piping & Not explicit $\rightarrow$ inferred & --- \\
Instrumentation \& controls & Not explicit $\rightarrow$ inferred & --- \\
\bottomrule
\end{tabular}
\end{table}

The equipment configuration is more complex than the \texttt{liming\_system}: multiple reaction stages operate at different pH setpoints (typically pH $\sim$3.5 for Fe, pH $\sim$5 for Al, pH $\sim$9--10 for Mg), each requiring independent reagent dosing and pH control. Intermediate clarification or settling may occur between stages to remove precipitated solids before the next pH adjustment.


\medskip
\noindent\textbf{Sizing basis.}\quad

The base volume is 100 $\si{m^3/h}$ of pregnant leach solution. Design capacity is 115 $\si{m^3/h}$ after the \texttt{oversizing\_factor} = 1.15. The \texttt{scaling\_exponent} = 0.60 (six-tenths rule). Lines are not dedicated (\texttt{dedicated\_line} = FALSE). Equipment life is 15 years, matching the \texttt{liming\_system} to reflect lime-service wear.


\medskip
\noindent\textbf{CAPEX.}\quad

CAPEX is derived by scaling the \texttt{liming\_system} installed cost from 60 $\si{m^3/h}$ to 100 $\si{m^3/h}$ using the six-tenths rule, then applying a complexity multiplier to account for the additional stages, instrumentation, and multi-reagent handling:

\textit{Step 1: Scale \texttt{liming\_system} to 100 $\si{m^3/h}$.} The scaling factor is $(100/60)^{0.60} = 1.667^{0.60} = 1.359$.

\textit{Step 2: Apply complexity multiplier.} The multiplier increases from 1.50$\times$ (conservative) to 1.75$\times$ (optimistic). This inverted pattern reflects a range-compression effect: the staged neutralization requires a minimum equipment configuration (multi-stage pH control, per-stage instrumentation, multi-reagent handling) that establishes a cost floor. This floor prevents the optimistic cost from scaling down as aggressively as the underlying \texttt{liming\_system} base, compressing the scenario range from 2.07$\times$ (\texttt{liming\_system}) to 1.77$\times$:

\begin{table}[h!]
\caption{CAPEX derivation for the \texttt{staged\_neutralization} block.}
\centering
\begin{tabular}{>{\raggedright\arraybackslash}p{2.2cm}
				>{\raggedright\arraybackslash}p{2.5cm}
				>{\raggedright\arraybackslash}p{2.0cm}
				>{\raggedright\arraybackslash}p{2.2cm}
				>{\raggedright\arraybackslash}p{3.5cm}}
\toprule
\textbf{Scenario} & \textbf{Lime\_soft scaled} & \textbf{$\times$ 1.359} & \textbf{Complexity} & \texttt{prim\_equip\_price\_base} \\
\midrule
Conservative & \$1,971K & \$2,679K & $\times$1.50 & \textbf{\$4,010,000} \\
Midpoint     & \$1,346K & \$1,829K & $\times$1.60 & \textbf{\$2,930,000} \\
Optimistic   & \$950K   & \$1,291K & $\times$1.75 & \textbf{\$2,260,000} \\
\bottomrule
\end{tabular}
\end{table}

\medskip
\noindent\textbf{Electricity.}\quad

Electricity demand is 170 / 120 / 90 kWh/h (conservative / midpoint / optimistic). These values are cross-checked against the \texttt{liming\_system} block scaled linearly from 50 $\si{m^3/h}$ to 100 $\si{m^3/h}$:

\begin{table}[h!]
\caption{Electricity cross-check for the \texttt{staged\_neutralization} block.}
\centering
\begin{tabular}{>{\raggedright\arraybackslash}p{2.5cm}
				>{\raggedright\arraybackslash}p{3.5cm}
				>{\raggedright\arraybackslash}p{3.5cm}
				>{\raggedright\arraybackslash}p{4cm}}
\toprule
\textbf{Scenario} & \textbf{\texttt{liming\_system} at 100 $\si{m^3/h}$} & \texttt{electricity\_base\_total} & \textbf{Premium over \texttt{liming\_system}} \\
\midrule
Conservative & 150 kWh/h & \textbf{170 kWh/h} & $\sim$15\% \\
Midpoint     & 100 kWh/h & \textbf{120 kWh/h} & 20\% \\
Optimistic   & 70 kWh/h  & \textbf{90 kWh/h}  & $\sim$30\% \\
\bottomrule
\end{tabular}
\end{table}

\noindent The $\sim$15--30\% premium is consistent with the additional electrical loads of a multi-stage circuit: per-stage pH probes with dosing loops, additional inter-stage transfer pumps, and more intensive mixing at the Fe/Al precipitation stages (which require rapid reagent dispersion to avoid local pH overshoot). The midpoint specific electricity of 1.2 \si{kWh/m^3} is comparable to the \texttt{liming\_system} (1.0 \si{kWh/m^3} at 50 $\si{m^3/h}$).

\medskip
\noindent\textbf{Other utilities.}\quad

No thermal utilities are consumed. Lime is assumed to be purchased (rather than quicklime), given the lack of water resources at Thacker Pass that would be needed to produce lime from quicklime. 

\medskip
\noindent\textbf{Consumables.}\quad

Limestone (\ce{CaCO3}), quicklime (\ce{CaO}), soda ash (\ce{Na2CO3}), and flocculant are tracked at the production step level as secondary inputs with pathway-specific dosages; no tooling is used in this step - \texttt{tooling\_cost\_base} = 0.

\medskip
\noindent\textbf{Labor.}\quad

Labor is 4.5 FTE at base scale with a scaling exponent of 0.2 and \texttt{dedicated\_labor} = FALSE. The same base FTE as the simpler \texttt{liming\_system} is used despite the greater operational complexity because the additional stages can be fully automated with controlled pH dosing loops, requiring supervisory rather than manual intervention:
\[
	\text{FTE} = 1.0 \times 4.5 = 4.5
\]

\medskip
\noindent\textbf{Footprint.}\quad

The base plot area is 500~$\si{m^2}$ with a scaling exponent of 0.7, matching the \texttt{liming\_system}. The additional neutralization stages add modestly to the tank footprint, but the lime slaking and reagent storage systems (which dominate the plot area) are comparable in size.

\medskip
\noindent\textbf{Availability.}\quad

Process availability is 92\% with 3\% unplanned downtime and 5\% scheduled maintenance. The slightly higher availability (vs.\ 90\% for the \texttt{precipitator}) reflects the robust operating envelope of neutralization chemistry, which is less sensitive to temperature and stoichiometric control than crystallization circuits.

\medskip
\noindent\textbf{Maintenance.}\quad

Neither \texttt{maint\_override} nor a nonzero \texttt{opex\_fraction\_of\_capex} is set; the block inherits the facility-level default maintenance rate of 5\% of scaled primary equipment cost, charged against the total undiscounted equipment price.

\medskip
\noindent\textbf{Key measurements and cross-checks.}\quad

\begin{itemize}[noitemsep]
	\item \textit{CAPEX}: midpoint of \$2.93M at 100 $\si{m^3/h}$ implies \$29.3/($\si{m^3/h}$), approximately 2.2$\times$ the \texttt{precipitator}'s \$13/($\si{m^3/h}$) and 1.6$\times$ the \texttt{liming\_system}'s scaled cost. This premium is consistent with the additional stages and instrumentation.
	\item \textit{Electricity}: 120 kWh/h midpoint at 100 $\si{m^3/h}$ = 1.2 \si{kWh/m^3}, comparable to the \texttt{liming\_system} (1.0) and \texttt{batch\_treatment} (1.5). The scaling from \texttt{liming\_system} passes within 20\%.
	\item \textit{Pathway usage}: this block is shared across two pathways (Thacker Pass clay leach and Jianxiawo lepidolite water leach), each with different ionic loadings and reagent combinations. The 100 $\si{m^3/h}$ base volume and multi-stage design accommodate both use cases without pathway-specific equipment.
\end{itemize}

\clearpage
\paragraph{\texttt{precipitator}}\label{app:04-methods-blockdefs-precipitator}\leavevmode\par

\noindent\textbf{Sources.}\quad

\citet{battagliaRecoveryLithiumCarbonate2022}, \citet{liuProcessAnalysisStudy2021}, \citet{khakmardanLifeCycleAssessment2024}, \citet{iyerLifecycleAnalysisLithium2024}, \citet{royLithiumClayAssessing2024}, \citet{sgscanadaincSK1300Technical2024}, \citet{samariMineralResourceEstimate2024}, \citet{srkconsultingSECTechnicalReport2023a}, \citet{ausimmMonograph27Cost}, \citet{petersPlantDesignEconomics1991}; vendor data (e.g. ~\citet{baguleyMixtechAustraliaFluid2009}).

The precipitator block converts dissolved lithium --- either as a chloride or sulfate species depending on the upstream pathway, as noted in Appendix \ref{app:04-methods-facilitychoices-pls} --- into lithium carbonate via reaction with soda ash (\ce{Na2CO3}), yielding \ce{Li2CO3} as a solid product:
\begin{equation}
    \ce{2LiCl(aq) + Na2CO3(aq) -> Li2CO3(s) + 2NaCl(aq)}
\end{equation}
\begin{equation}
    \ce{Li2SO4(aq) + Na2CO3(aq) -> Li2CO3(s) + Na2SO4(aq)}
\end{equation}
Precipitation is conducted at 80--95$^\circ$C to exploit the solubility of \ce{Li2CO3} and maximize recovery yield. The block encompasses pre-heating of the polished liquor feed, reagent addition, reaction and crystallization, and conditioning/aging for crystal agglomeration prior to downstream solid-liquid separation.

\medskip
\noindent\textbf{Equation of merit.}\quad
\[
    V_\text{working}\;(\si{m^3}) = \dot{V}_\text{liquor}\;(\si{m^3/h}) \times \tau_\text{res}\;(\text{h})
\]
The working volume of the circuit $V_\text{working}$ is sized to deliver a liquor volume $\dot{V}_\text{liquor}$ with total residence time $\tau_\text{res}$, across a three-tank cascade.~\cite{srkconsultingSECTechnicalReport2023a} For example, with a base liquor flow of 50 \si{m^3/h} and total residence time of 1.2 h, this gives a working volume of 60 \si{m^3}. 

\medskip
\noindent\textbf{Bill of Equipment.}\quad
 
\begin{table}[h!]
\caption{Bill of equipment for the \texttt{precipitator} block.}
\centering
\begin{tabular}{>{\raggedright\arraybackslash}p{5.2cm}
                >{\raggedright\arraybackslash}p{4cm}
                >{\raggedright\arraybackslash}p{4.5cm}}
\toprule
\textbf{Equipment} & \textbf{Machine data used} & \textbf{Source(s)} \\
\midrule
Steam preheater & Not explicit $\rightarrow$ inferred & --- \\
Rapid mix tank (reagent addition, $\sim$5~\si{m^3}) & AusIMM Rapid Mix Tank & Rapid Mix Tanks \\
Reaction / crystallization tank ($\sim$35~\si{m^3}) & Mixtec high-torque agitator series & Agitator Tanks \\
Conditioning / aging tank ($\sim$10~\si{m^3}) & Mixtec high-torque agitator series & Agitator Tanks \\
Slurry transfer pumps \& piping & Not explicit $\rightarrow$ inferred & --- \\
Reagent handling (\ce{Na2CO3}) & Not explicit $\rightarrow$ inferred & --- \\
Instrumentation \& controls & Not explicit $\rightarrow$ inferred & --- \\
\bottomrule
\end{tabular}
\end{table}

\medskip
\noindent\textbf{Sizing basis.}\quad

The base volume is 50,000 L/h (50 \si{m^3/h}) of polished brine or purified leach liquor. This corresponds to a reference production rate of approximately 10~t~\ce{Li2CO3}/day. The three-tank configuration splits working volume between rapid mixing ($\sim$5 \si{m^3}), reaction and crystallization ($\sim$35 \si{m^3}), and conditioning/aging ($\sim$10 \si{m^3}). Throughput scaling uses an exponent of 0.7, reflecting the sub-linear cost scaling of vessel-dominated circuits with instrumentation.

\medskip
\noindent\textbf{CAPEX.}\quad

Tank costs are scaled from AusIMM rapid mix tank and Mixtec agitator tank benchmarks using the six-tenths rule. The volumetric cost (USD per \si{m^3} of working volume) is applied to the 60 \si{m^3} working volume to obtain the raw tank cost. Subsystem adders cover slurry pumps and piping (20\%), soda ash handling and dosing (15\%), and instrumentation and controls (10\%) applied to the tank cost. An implicit installation factor of $\sim$1.15$\times$ is embedded in the volumetric unit costs, giving an effective total multiplier of approximately 1.67$\times$:

\begin{table}[h!]
\caption{CAPEX derivation for the \texttt{precipitator} block.}
\centering
\begin{tabular}{>{\raggedright\arraybackslash}p{2.5cm}
                >{\raggedright\arraybackslash}p{2.8cm}
                >{\raggedright\arraybackslash}p{2.8cm}
                >{\raggedright\arraybackslash}p{2cm}
                >{\raggedright\arraybackslash}p{4.5cm}}
\toprule
\textbf{Scenario} & \textbf{Volumetric cost} & \textbf{Raw tank CAPEX} & \textbf{Total mult.} & \texttt{prim\_equip\_price\_base} \\
\midrule
Conservative & \$9,000/\si{m^3} & \$540,000 & $\times$1.67 & \textbf{\$900,000} \\
Midpoint     & \$6,500/\si{m^3} & \$390,000 & $\times$1.67 & \textbf{\$650,000} \\
Optimistic   & \$4,800/\si{m^3} & \$288,000 & $\times$1.67 & \textbf{\$480,000} \\
\bottomrule
\end{tabular}
\end{table}

\noindent Equipment life is 20~years. The CAPEX scaling exponent is 0.60 (six-tenths rule).

\medskip
\noindent\textbf{Pre-heating and thermal demand.}\quad

Pre-heating the feed liquor to the precipitation temperature is required to exploit the retrograde solubility of \ce{Li2CO3} and is noted across multiple project disclosures (Thacker Pass, Bonnie Claire, Silver Peak). Precipitation is conducted at 80--95$^\circ$C. The thermal duty depends on the feed temperature entering the precipitator, which varies significantly between pathways: chloride-pathway feeds (Silver Peak) arrive near ambient ($\sim$22$^\circ$C) after batch chemical purification, while sulfate-pathway feeds (Thacker Pass) may retain some heat from upstream processing ($\sim$40--50$^\circ$C). Given the relatively small amount of material and temperature increase, this heating demand is described here but not yet implemented into the model. 

\medskip
\noindent\textbf{Electricity.}\quad
Electricity demand for precipitation was estimated as the sum of (i) tank mixing power, (ii) liquor transfer pumping power, and (iii) auxiliary motor loads for reagent make-up and dosing (lime slaking and soda ash dissolution).

\textit{Agitation.} Mixing power is estimated from the SME Mining Engineering Handbook envelope of 0.2--0.6 \si{kW/m^3}, with a midpoint estimate of 0.4 \si{kW/m^3}, applied to the 50 \si{m^3/h} base flow volume:
 
\textit{Transfer pumping.} Estimated from the hydraulic power relation $P = \rho g Q H / \eta$ at $Q = 50/3600 = 0.0139$ \si{m^3/s}, $\rho = 1300$ \si{kg/m^3}, $\eta = 0.65$, and transfer heads of 20/30/40 m (optimistic/midpoint/conservative), giving 5.5/8.2/10.9kW respectively.
 
\textit{Auxiliary loads} (soda ash dissolution, reagent dosing, instrumentation) account for the remaining demand to reach the total:
 
\begin{table}[h!]
\caption{Electricity scenario derivation for the \texttt{precipitator} block.}
\centering
\begin{tabular}{>{\raggedright\arraybackslash}p{2.5cm}
                >{\raggedright\arraybackslash}p{2.8cm}
                >{\raggedright\arraybackslash}p{2.8cm}
                >{\raggedright\arraybackslash}p{2cm}
                >{\raggedright\arraybackslash}p{4.5cm}}
\toprule
\textbf{Scenario} & \textbf{Mixing (kW)} & \textbf{Transfer pump (kW)} & \textbf{Auxiliary (kW)} & \texttt{electricity\_base\_total} \\
\midrule
Conservative & 30 & 10.9 & 35 & \textbf{75 kWh/h} \\
Midpoint     & 20 &  8.2 & 20 & \textbf{50 kWh/h} \\
Optimistic   & 10 &  5.5 &  10 & \textbf{25 kWh/h} \\
\bottomrule
\end{tabular}
\end{table}
 
\noindent The auxiliary load (35/20/10 kW) is approximately equal in magnitude to the mixing power across all scenarios, consistent with the additional drive loads for soda ash conveying, dissolution tank agitation, and metering pumps. Total electricity load is then rounded. 

\medskip
\noindent\textbf{Labor.}\quad

Labor is 4.5 FTE at base scale with a scaling exponent of 0.2 and \texttt{dedicated\_labor} = FALSE, assuming one operator per shift for the continuous three-tank precipitation cascade:
\[
	\text{FTE} = 1.0 \times 4.5 = 4.5
\]

\medskip
\noindent\textbf{Footprint.}\quad

The base plot area is 150~$\si{m^2}$ with a scaling exponent of 0.7, covering the three tanks, preheater, reagent handling, and surrounding access.

\medskip
\noindent\textbf{Availability.}\quad

Process availability is 90\% with 5\% unplanned downtime and 5\% scheduled maintenance. The slightly lower availability (relative to the 92\% used for leach circuit blocks) reflects the sensitivity of crystallization circuits to upsets in temperature and reagent dosing.

\medskip
\noindent\textbf{Key measurements and cross-checks.}\quad
\begin{itemize}[noitemsep]
    \item \ce{Li2CO3} precipitation yield (fraction of dissolved Li reporting to product); literature values 85--95\% for optimized circuits \cite{liuProcessAnalysisStudy2021, battagliaRecoveryLithiumCarbonate2022}; Thacker Pass is taken at 0.97 as a recycle-inclusive plant recovery consistent with the disclosed circuit losses
    \item Crystal size distribution and filterability (conditioning/aging tank residence time is the primary lever)
    \item Soda ash stoichiometry: approximately 3.26 kg \ce{Na2CO3} per kg \ce{Li2CO3} produced at Thacker Pass and 3.51 at Jianxiawo, with the difference due to different recovery rates, against a stoichiometric floor of $\text{M}_{\ce{Na2CO3}} / (2\,\text{M}_{\ce{Li}}) = 7.63$ kg per kg Li, or $2.17\times$ that floor
    \item Midpoint CAPEX of \$650,000 for a 50 \si{m^3/h}/h circuit implies \$13 (\si{m^3/h}), consistent with Class~4--5 benchmarks for small stirred-tank precipitation at this scale
\end{itemize}

\clearpage
\paragraph{\texttt{liming\_system}}\label{app:04-methods-blockdefs-limingsystem}\leavevmode\par

\medskip
\noindent\textbf{Sources.}\quad
\citet{battagliaRecoveryLithiumCarbonate2022}, \citet{liuProcessAnalysisStudy2021}, \citet{khakmardanLifeCycleAssessment2024}, \citet{iyerLifecycleAnalysisLithium2024}, \citet{royLithiumClayAssessing2024}, \citet{sgscanadaincSK1300Technical2024}, \citet{ausimmMonograph27Cost}, \citet{petersPlantDesignEconomics1991}; vendor data (BHT, Westpro, Sodimate lime slaking systems). 

\medskip
\noindent\textbf{Process description.}\quad

The liming system block removes divalent cation impurities---primarily \ce{Mg^{2+}}, \ce{Ca^{2+}}, \ce{Fe^{2+}}, and \ce{Mn^{2+}}---from brine or pregnant leach solution by reaction with slaked lime (\ce{Ca(OH)2}). The primary reaction for magnesium removal is:
\begin{equation}
	\ce{Mg^{2+}(aq) + Ca(OH)2(s) -> Mg(OH)2(s) v + Ca^{2+}(aq)}
\end{equation}
Precipitated \ce{Mg(OH)2} reports to the solids phase, while the calcium introduced by lime addition remains in solution and must be subsequently removed, typically by soda ash addition in a downstream polishing step. Unlike the \texttt{precipitator} block, which converts dissolved lithium into a solid product, liming system removes unwanted impurities while keeping lithium in solution. The resulting slurry has a substantially higher solids loading than a lithium carbonate precipitation circuit, requiring more robust agitation and solids-handling equipment. The block encompasses quicklime receiving and slaking, milk-of-lime storage and dosing, rapid reagent mixing, reaction and crystallization in a continuous stirred-tank reactor (CSTR), and conditioning/aging for floc agglomeration prior to downstream solid-liquid separation.

This block is not actively used in the current model but serves as a reference baseline for other tank-based circuits, e.g. \texttt{precipitator}~\ref{app:04-methods-blockdefs-precipitator} and \texttt{batch\_treatment\_generic}~\ref{app:04-methods-blockdefs-batchtreat}.

\medskip
\noindent\textbf{Equation of merit.}\quad
\[
	V_\text{circuit}\;(\si{m^3}) = \dot{V}_\text{liquor}\;(\si{m^3/h}) \times \tau_\text{circuit}\;(\text{h})
\]
The working volume of the softening circuit $V_\text{circuit}$ is sized to deliver a liquor flow $\dot{V}_\text{liquor}$ with total circuit residence time $\tau_\text{circuit}$ across a three-tank cascade: rapid mix ($\sim$5 \si{m^3}, $\tau \approx 6$ min), reaction CSTR ($\sim$35 \si{m^3}, $\tau \approx 42$ min), and conditioning ($\sim$20 \si{m^3}, $\tau \approx 24$ min), giving $\tau_\text{circuit} \approx 1.2$ h and $V_\text{circuit} = 60$ \si{m^3} at the 50 $\si{m^3/h}$ base flow. An additional 30 \si{m^3} of agitated volume is required for the lime slaking subsystem (5 \si{m^3} detention slaker + 25 \si{m^3} milk-of-lime storage), sized by the required lime production rate rather than the brine throughput. Total agitated volume across the block is 90 \si{m^3}.

\medskip
\noindent\textbf{Bill of Equipment.}\quad
\begin{table}[h!]
\caption{Bill of equipment for the \texttt{liming\_system} block.}
\centering
\begin{tabular}{>{\raggedright\arraybackslash}p{5.2cm}
				>{\raggedright\arraybackslash}p{4cm}
				>{\raggedright\arraybackslash}p{4.5cm}}
\toprule
\textbf{Equipment} & \textbf{Machine data used} & \textbf{Source(s)} \\
\midrule
Detention lime slaker ($\sim$5 \si{m^3}) & AusIMM Rapid Mix Tank & Rapid Mix Tanks \\
Milk-of-lime storage tank ($\sim$25 \si{m^3}) & Mixtec Agitator Tank Series & Agitator Tanks \\
Rapid mix tank --- reagent addition ($\sim$5 \si{m^3}) & AusIMM Rapid Mix Tank & Rapid Mix Tanks \\
Reaction / crystallization CSTR ($\sim$35 \si{m^3}) & Continuous Stirred Tank Reactor & CSTRs \\
Conditioning / aging tank ($\sim$20 \si{m^3}) & Mixtec Agitator Tank Series & Agitator Tanks \\
Quicklime feed system (silo, screw feeder) & Not explicit $\rightarrow$ inferred & --- \\
Grit removal (classifying screen/rake) & Not explicit $\rightarrow$ inferred & --- \\
Slurry transfer pumps \& piping & Not explicit $\rightarrow$ inferred & --- \\
Lime dosing (metering pumps) & Not explicit $\rightarrow$ inferred & --- \\
Instrumentation \& controls & Not explicit $\rightarrow$ inferred & --- \\
\bottomrule
\end{tabular}
\end{table}

The lime slaking subsystem is substantially more equipment-intensive than the soda ash handling in the \texttt{precipitator} block. A quicklime conditioning system encompasses receiving, slaking, grit removal, slurry storage, and dosing. The solids-handling and mechanical components are expected to dominate the cost of the slaking subsystem.

\medskip
\noindent\textbf{Sizing basis.}\quad

The base volume is 50,000 L/h (50 $\si{m^3/h}$) of brine or leach solution, matching the \texttt{precipitator} block throughput for cross-comparability. Design capacity is 60 $\si{m^3/h}$ after applying the \texttt{oversizing\_factor} = 1.2. The \texttt{scaling\_exponent} = 0.60 (six-tenths rule), consistent with the cost scaling of vessel-dominated circuits. Lines are not dedicated (\texttt{dedicated\_line} = FALSE). Equipment life is 15 years, shorter than the precipitator's 20-year life to reflect the higher abrasive and scaling demands of lime slurry service.

\medskip
\noindent\textbf{CAPEX.}\quad

CAPEX is built bottom-up from two subsystems: the lime slaking system and the softening circuit proper. Tank costs are scaled from AusIMM rapid mix tank and Mixtec agitator tank volumetric benchmarks, with the same unit costs (\$/\si{m^3} of working volume) used in the \texttt{precipitator} block. An implicit installation factor of $\sim$1.15$\times$ is embedded in the volumetric unit costs.

\textit{Lime slaking system} (30 \si{m^3} total tank volume). The raw tank cost covers the 5 \si{m^3} slaker and 25 \si{m^3} storage vessel. A +250\% subsystem adder (total 3.5$\times$) is applied to account for the quicklime silo, screw feeder, grit removal, dust scrubber, dosing pumps, and associated piping and controls:

\begin{table}[h!]
\caption{Lime slaking system CAPEX derivation for the \texttt{liming\_system} block.}
\centering
\begin{tabular}{>{\raggedright\arraybackslash}p{2.5cm}
				>{\raggedright\arraybackslash}p{2.8cm}
				>{\raggedright\arraybackslash}p{2.5cm}
				>{\raggedright\arraybackslash}p{2cm}
				>{\raggedright\arraybackslash}p{3cm}}
\toprule
\textbf{Scenario} & \textbf{Vol.\ cost} & \textbf{Raw tanks} & \textbf{$\times$3.5} & \textbf{Slaking total} \\
\midrule
Conservative & \$9,000/\si{m^3} & \$270,000 & \$945,000 & \$945,000 \\
Midpoint     & \$6,500/\si{m^3} & \$195,000 & \$682,500 & \$683,000 \\
Optimistic   & \$4,800/\si{m^3} & \$144,000 & \$504,000 & \$504,000 \\
\bottomrule
\end{tabular}
\end{table}

\noindent\textit{Flag}: the +250\% subsystem adder is an engineering estimate with no direct vendor or published source. The multiplier is plausible given the number of mechanical components (silo, feeder, grit classifier, dust scrubber) but should be validated against a vendor quote.

\textit{Softening circuit} (60 \si{m^3} working volume). Subsystem adders are applied to the raw circuit tank cost: slurry pumps and piping (20\%), lime handling and dosing (25/40/60\% by scenario), and instrumentation and controls (10\%). The total is then combined with the lime slaking system:

\begin{table}[h!]
\caption{Softening circuit and total CAPEX derivation for the \texttt{liming\_system} block.}
\centering
\begin{tabular}{>{\raggedright\arraybackslash}p{2.3cm}
				>{\raggedright\arraybackslash}p{2.3cm}
				>{\raggedright\arraybackslash}p{2.2cm}
				>{\raggedright\arraybackslash}p{2.0cm}
				>{\raggedright\arraybackslash}p{2.0cm}
				>{\raggedright\arraybackslash}p{2.8cm}}
\toprule
\textbf{Scenario} & \textbf{Raw circuit tanks} & \textbf{Adder \%} & \textbf{Circuit sub.} & \textbf{+ Slaking} & \texttt{prim\_equip\_price\_base} \\
\midrule
Conservative & \$540,000 & +90\% & \$1,026,000 & \$945,000 & \textbf{\$1,971,000} \\
Midpoint     & \$390,000 & +70\% & \$663,000  & \$683,000 & \textbf{\$1,346,000} \\
Optimistic   & \$288,000 & +55\% & \$446,000  & \$504,000 & \textbf{\$950,000} \\
\bottomrule
\end{tabular}
\end{table}

\noindent The adder percentages differ by scenario: the conservative case assumes a larger lime handling and dosing fraction (60\%) reflecting corrosion-resistant materials and redundant dosing pumps, while the optimistic case (25\%) reflects a simplified system with standard carbon steel construction.

\medskip
\noindent\textbf{Electricity.}\quad

Electricity demand is rederived from first principles as the sum of (i) tank agitation power across all five agitated vessels, (ii) lime handling system drives, (iii) liquor and slurry transfer pumping, and (iv) auxiliary loads.

\textit{Agitation.} Mixing power is estimated from the power intensity (\si{kW/m^3}) applied to the \textit{working volume} of each agitated vessel, not the throughput flow rate. Intensity ranges are assigned by vessel duty: the lime slaker requires vigorous agitation (1.0--2.0 \si{kW/m^3}) to maintain slaking temperature uniformity and prevent drowning; the reaction CSTR uses moderate agitation (0.15--0.40 \si{kW/m^3}) for solids suspension, consistent with the 0.2--0.6 \si{kW/m^3} SME envelope applied in the \texttt{precipitator} block; and the conditioning tank uses gentle mixing (0.08--0.20 \si{kW/m^3}) for floc growth:

\begin{table}[h!]
\caption{Agitation power derivation by vessel for the \texttt{liming\_system} block.}
\centering
\begin{tabular}{>{\raggedright\arraybackslash}p{4.5cm}
				>{\raggedright\arraybackslash}p{1.3cm}
				>{\raggedright\arraybackslash}p{2.5cm}
				>{\raggedright\arraybackslash}p{2.5cm}
				>{\raggedright\arraybackslash}p{2.5cm}}
\toprule
\textbf{Vessel} & \textbf{Vol.\ (\si{m^3})} & \textbf{Optimistic (kW)} & \textbf{Midpoint (kW)} & \textbf{Conservative (kW)} \\
\midrule
Lime slaker         & 5  & 5.0 \scriptsize{(1.0 \si{kW/m^3})} & 7.5 \scriptsize{(1.5 \si{kW/m^3})} & 10.0 \scriptsize{(2.0 \si{kW/m^3})} \\
Lime storage        & 25 & 2.5 \scriptsize{(0.10)}          & 3.8 \scriptsize{(0.15)}          & 6.3 \scriptsize{(0.25)}          \\
Rapid mix (circuit) & 5  & 2.5 \scriptsize{(0.50)}          & 5.0 \scriptsize{(1.0)}           & 7.5 \scriptsize{(1.5)}           \\
Reaction CSTR       & 35 & 5.3 \scriptsize{(0.15)}          & 8.8 \scriptsize{(0.25)}          & 14.0 \scriptsize{(0.40)}         \\
Conditioning        & 20 & 1.6 \scriptsize{(0.08)}          & 2.4 \scriptsize{(0.12)}          & 4.0 \scriptsize{(0.20)}          \\
\midrule
\textbf{Agitation subtotal} & \textbf{90} & \textbf{16.9} & \textbf{27.4} & \textbf{41.8} \\
\bottomrule
\end{tabular}
\end{table}

\textit{Lime handling.} The slaking subsystem drives include a screw feeder (3--5 kW), grit removal rake or classifying screen (2--4 kW), and lime dosing metering pump (1--2 kW), totaling 6/8.5/11 kW (optimistic/midpoint/conservative).

\textit{Transfer pumping.} Estimated from the hydraulic power relation $P = \rho g Q H / \eta$ at $Q = 50/3600 = 0.0139$ $\si{m^3/s}$, $\rho = 1050$ $\si{kg/m^3}$ (dilute brine), $\eta = 0.65$, and transfer heads of 20/30/40 m (optimistic/midpoint/conservative), giving 4.4/6.6/8.8 kW for the main brine pump. An additional 1--2 kW for the lime slurry dosing pump gives a pumping subtotal of 5.4/7.6/10.8 kW.

\textit{Auxiliary loads} (instrumentation, lighting, sampling systems) account for the remaining 4/6/9 kW.

\begin{table}[h!]
\caption{Electricity scenario derivation for the \texttt{liming\_system} block.}
\centering
\begin{tabular}{>{\raggedright\arraybackslash}p{3.0cm}
				>{\raggedright\arraybackslash}p{2.5cm}
				>{\raggedright\arraybackslash}p{2.2cm}
				>{\raggedright\arraybackslash}p{2.0cm}
				>{\raggedright\arraybackslash}p{1.8cm}
				>{\raggedright\arraybackslash}p{2.8cm}}
\toprule
\textbf{Scenario} & \textbf{Agitation (kW)} & \textbf{Lime hdlg.\ (kW)} & \textbf{Pumping (kW)} & \textbf{Aux.\ (kW)} & \texttt{elec\_base\_total} \\
\midrule
Conservative & 41.8 & 11 & 10.8 & 9 & \textbf{75 kWh/h} \\
Midpoint     & 27.4 & 8.5 & 7.6 & 6 & \textbf{50 kWh/h} \\
Optimistic   & 16.9 & 6 & 5.4 & 4 & \textbf{35 kWh/h} \\
\bottomrule
\end{tabular}
\end{table}

\medskip
\noindent\textbf{Other utilities.}\quad

No thermal utilities are consumed. Lime is assumed to be purchased (rather than quicklime), given the lack of water resources at Thacker Pass that would be needed to produce lime from quicklime. 

\medskip
\noindent\textbf{Consumables.}\quad

Equipment-tied consumables (e.g., pump seals, agitator wear parts) are not separately tracked in this model; \texttt{tooling\_cost\_base} = 0. 

\medskip
\noindent\textbf{Labor.}\quad

Labor is 4.5 FTE at base scale with a scaling exponent of 0.2 and \texttt{dedicated\_labor} = FALSE, assuming one operator per shift for the continuous softening circuit. The lime slaking subsystem is assumed to not add incremental operators; modern slaking systems can use automated controls for e.g. lime-to-water ratios, and require only periodic supervisory attention:
\[
	\text{FTE} = 1.0 \times 4.5 = 4.5
\]

\medskip
\noindent\textbf{Footprint.}\quad

The base plot area is 500~$\si{m^2}$ with a scaling exponent of 0.7, covering the softening tanks, lime slaking system (including silo and feed equipment), dosing and pumping skids, and surrounding access. The substantially larger footprint compared to the \texttt{precipitator} (150~$\si{m^2}$) reflects the lime receiving and storage infrastructure.

\medskip
\noindent\textbf{Availability.}\quad

Process availability is 92\% with 3\% unplanned downtime and 5\% scheduled maintenance. The slightly higher availability relative to the \texttt{precipitator} (90\%) reflects the simpler chemistry (no retrograde solubility sensitivity) and more robust operating envelope of lime softening.

\medskip
\noindent\textbf{Maintenance.}\quad

Neither \texttt{maint\_override} nor a nonzero \texttt{opex\_fraction\_of\_capex} is set; the block inherits the facility-level default maintenance rate of 5\% of scaled primary equipment cost. Lime slaking systems require above-average maintenance due to grit accumulation, calcium carbonate scaling on piping and valves, and abrasive wear on slurry pumps, so a \texttt{maint\_override} above the 5\% default should be considered if this block is activated for production use.

\clearpage
\paragraph{\texttt{batch\_treatment\_generic}}\label{app:04-methods-blockdefs-batchtreat}\leavevmode\par

\medskip
\noindent\textbf{Sources.}\quad

\citet{srkconsultingSECTechnicalReport2023a}; \citet{ausimmMonograph27Cost}; \citet{petersPlantDesignEconomics1991}; vendor data. Cost methodology adapted from the \texttt{precipitator} (Appendix \ref{app:04-methods-blockdefs-precipitator}) and \texttt{liming\_system} (Appendix \ref{app:04-methods-blockdefs-limingsystem}) blocks.

\medskip
\noindent\textbf{Process description.}\quad

The batch treatment block precipitates dissolved magnesium and calcium impurities from concentrated lithium brine, leaving lithium in solution. It is modeled on the Silver Peak (Nevada) batch chemical purification process, in which discrete batches of concentrated brine are sequentially treated with slaked lime (\ce{Ca(OH)2}) and soda ash (\ce{Na2CO3}). The primary reactions are:
\begin{equation}
	\ce{Mg^{2+}(aq) + Ca(OH)2(s) -> Mg(OH)2(s) v + Ca^{2+}(aq)}
\end{equation}
\begin{equation}
	\ce{Ca^{2+}(aq) + Na2CO3(aq) -> CaCO3(s) v + 2Na^{+}(aq)}
\end{equation}
Lime targets magnesium removal; the calcium introduced by the lime is subsequently removed by soda ash addition. Native calcium in the brine is also removed by soda ash. Reagent dosages, stoichiometry, and elimination fractions are described in the production step writeup (Appendix \ref{app:04-methods-facilitychoices-batchchem}). The block encompasses quicklime slaking and storage, soda ash dissolution and dosing, batch reactor agitation, and conditioning prior to downstream solid-liquid separation. This block is intentionally structured as a composite of the \texttt{liming\_system} and \texttt{precipitator} blocks, with the addition of a soda ash dissolution and dosing subsystem, reflecting the dual-reagent treatment performed in a single batch vessel.

\medskip
\noindent\textbf{Equation of merit.}\quad
\[
	V_\text{batch}\;(\si{m^3}) = \text{batch size}\;(\text{L}) \times 10^{-3}
\]
The batch reactor must accommodate the full batch volume $V_\text{batch}$ with headspace for agitation and reagent addition. At a base volume of 60,000 L/batch = 60 \si{m^3/batch}, the design vessel volume is 69 \si{m^3} after the \texttt{oversizing\_factor} = 1.15 is applied. Each batch cycle takes 7 hours total: 6 hours reaction time (\texttt{batch\_cycle\_time}) plus 1 hour setup (\texttt{batch\_setup\_time}) for filling, draining, and changeover, giving an effective throughput of $60/7 \approx 8.6$ $\si{m^3/h}$. The 60 \si{m^3} batch volume is consistent with the 60 \si{m^3} tank cascade working volume priced in the CAPEX derivation (rapid mix 5 + CSTR 35 + conditioning 20 \si{m^3}).

\medskip
\noindent\textbf{Bill of Equipment.}\quad

\begin{table}[h!]
\caption{Bill of equipment for the \texttt{batch\_treatment\_generic} block.}
\centering
\begin{tabular}{>{\raggedright\arraybackslash}p{5.2cm}
				>{\raggedright\arraybackslash}p{4cm}
				>{\raggedright\arraybackslash}p{4.5cm}}
\toprule
\textbf{Equipment} & \textbf{Machine data used} & \textbf{Source(s)} \\
\midrule
Detention lime slaker ($\sim$5 \si{m^3}) & AusIMM Rapid Mix Tank & Rapid Mix Tanks \\
Milk-of-lime storage tank ($\sim$25 \si{m^3}) & Mixtec Agitator Tank Series & Agitator Tanks \\
Rapid mix tank --- reagent addition ($\sim$5 \si{m^3}) & AusIMM Rapid Mix Tank & Rapid Mix Tanks \\
Reaction / crystallization tank ($\sim$35 \si{m^3}) & Continuous Stirred Tank Reactor & CSTRs \\
Conditioning / aging tank ($\sim$20 \si{m^3}) & Mixtec Agitator Tank Series & Agitator Tanks \\
Soda ash dissolution tank ($\sim$5 \si{m^3}) & Not explicit $\rightarrow$ inferred & --- \\
Slurry transfer pumps \& piping & Not explicit $\rightarrow$ inferred & --- \\
Reagent handling (\ce{CaO} + \ce{Na2CO3}) & Not explicit $\rightarrow$ inferred & --- \\
Instrumentation \& controls & Not explicit $\rightarrow$ inferred & --- \\
\bottomrule
\end{tabular}
\end{table}

The equipment list mirrors the \texttt{liming\_system} tank cascade (30 \si{m^3} lime slaking + 60 \si{m^3} circuit = 90 \si{m^3} total), with the addition of a soda ash dissolution and dosing subsystem. In batch operation, the three circuit vessels (rapid mix, CSTR, conditioning) collectively hold the 60 \si{m^3} batch; reagent addition, reaction, and conditioning occur sequentially across the cascade rather than in a single vessel.

\medskip
\noindent\textbf{Sizing basis.}\quad

The base volume is 60,000 L/batch (60 \si{m^3}) of concentrated lithium brine, with a total cycle time of 7 hours. \texttt{process\_type} = batch, \texttt{dedicated\_line} = FALSE. The \texttt{scaling\_exponent} = 0.60 (six-tenths rule), though this exponent is not applied for batch processes in the model: batch CAPEX scales as \texttt{machines\_required} $\times$ \texttt{prim\_equip\_price\_base}, where machines required is determined by line time demand.

\medskip
\noindent\textbf{CAPEX.}\quad

CAPEX is built bottom-up from three subsystems, following the methodology established in the \texttt{liming\_system} (Appendix \ref{app:04-methods-blockdefs-limingsystem}) and \texttt{precipitator} (Appendix \ref{app:04-methods-blockdefs-precipitator}) blocks:

\begin{itemize}[noitemsep]
	\item \textit{Tank train} (60 \si{m^3} working volume): volumetric cost $\times$ volume, using the same AusIMM/Mixtec benchmarks as the precipitator. This gives raw tank costs of \$540,000 / \$390,000 / \$288,000 (conservative / midpoint / optimistic).
	\item \textit{Soda ash dissolution and dosing system}: 15\% of raw tank train cost, reflecting a small dissolution tank ($\sim$5 \si{m^3}), metering pumps, and dry storage/handling. This gives \$81,000 / \$59,000 / \$43,000.
	\item \textit{Lime slaking package}: identical to the \texttt{liming\_system} lime slaking subsystem (30 \si{m^3} at 3.5$\times$ multiplier): \$945,000 / \$683,000 / \$504,000.
\end{itemize}

The pre-auxiliary subtotal is then adjusted with subsystem adders applied to the full package: slurry pumps and piping (20\%), additional reagent handling and dosing (15\%), and instrumentation and controls (10\%), totaling a 45\% adder:

\begin{table}[h!]
\caption{CAPEX derivation for the \texttt{batch\_treatment\_generic} block.}
\centering
\begin{tabular}{>{\raggedright\arraybackslash}p{2.2cm}
				>{\raggedright\arraybackslash}p{2.0cm}
				>{\raggedright\arraybackslash}p{1.6cm}
				>{\raggedright\arraybackslash}p{1.6cm}
				>{\raggedright\arraybackslash}p{2.0cm}
				>{\raggedright\arraybackslash}p{1.6cm}
				>{\raggedright\arraybackslash}p{3cm}}
\toprule
\textbf{Scenario} & \textbf{Tank train} & \textbf{Soda ash} & \textbf{Lime slk.} & \textbf{Pre-aux sub.} & \textbf{$\times$1.45} & \texttt{prim\_equip\_price\_base} \\
\midrule
Conservative & \$540K & \$81K & \$945K & \$1,566K & \$2,271K & \textbf{\$2,270,000} \\
Midpoint     & \$390K & \$59K & \$683K & \$1,132K & \$1,641K & \textbf{\$1,640,000} \\
Optimistic   & \$288K & \$43K & \$504K & \$835K  & \$1,211K & \textbf{\$1,210,000} \\
\bottomrule
\end{tabular}
\end{table}

\noindent Equipment life is 15 years, similar to the operating conditions of the liming system. 

\medskip
\noindent\textbf{Electricity.}\quad

Electricity demand is derived for batch operation, where the power draw varies across the 7-hour cycle. The model field \texttt{electricity\_base\_total} represents the cycle-averaged power draw (kWh/h), since the code computes total consumption as $\texttt{electricity\_base\_total} \times \text{line time required}$, where line time includes both reaction and setup phases.

\textit{Batch tank agitation.} The 60 \si{m^3} batch reactor requires time-varying agitation: vigorous mixing during reagent addition ($\sim$2 hr at 0.15--0.45 \si{kW/m^3}), gentle agitation during reaction and settling ($\sim$3 hr at 0.03--0.08 \si{kW/m^3}), conditioning ($\sim$1 hr at 0.04--0.12 \si{kW/m^3}), and no agitation during setup ($\sim$1 hr). The time-weighted average over the 7-hour cycle:

\begin{table}[h!]
\caption{Batch tank agitation power derivation for the \texttt{batch\_treatment\_generic} block.}
\centering
\begin{tabular}{>{\raggedright\arraybackslash}p{4.5cm}
				>{\raggedright\arraybackslash}p{1.3cm}
				>{\raggedright\arraybackslash}p{2.5cm}
				>{\raggedright\arraybackslash}p{2.5cm}
				>{\raggedright\arraybackslash}p{2.5cm}}
\toprule
\textbf{Phase} & \textbf{Dur.\ (hr)} & \textbf{Opt.\ (kW)} & \textbf{Mid.\ (kW)} & \textbf{Cons.\ (kW)} \\
\midrule
Lime addition + mixing           & 1.0 & 12 \scriptsize{(0.20 \si{kW/m^3})} & 18 \scriptsize{(0.30)} & 27 \scriptsize{(0.45)} \\
Soda ash addition + mixing       & 1.0 & 9 \scriptsize{(0.15)}           & 12 \scriptsize{(0.20)} & 18 \scriptsize{(0.30)} \\
Reaction / settling              & 3.0 & 1.8 \scriptsize{(0.03)}         & 3 \scriptsize{(0.05)}  & 4.8 \scriptsize{(0.08)} \\
Conditioning                     & 1.0 & 2.4 \scriptsize{(0.04)}         & 4.2 \scriptsize{(0.07)} & 7.2 \scriptsize{(0.12)} \\
Setup (fill / drain)             & 1.0 & 0                               & 0                      & 0                      \\
\midrule
\textbf{Time-wtd.\ avg.\ ($\div$7 hr)} & & \textbf{4.1} & \textbf{6.2} & \textbf{9.5} \\
\bottomrule
\end{tabular}
\end{table}

\textit{Ancillary loads} (cycle-averaged). The lime slaking system operates intermittently ($\sim$1.5 hr per 7-hour cycle, proportional to the 60 \si{m^3} batch reagent demand); the 25 \si{m^3} storage tank runs intermittent keep-alive agitation during the remainder, totaling 2.5--5.0 kW cycle-averaged across slaker, storage, and lime handling drives. Soda ash dissolution ($\sim$0.8 hr at 3--5 kW) contributes $\sim$0.4 kW average. Pumping (fill and drain, $\sim$1 hr at 4--7 kW) contributes 0.5--1.0 kW average. Instrumentation and controls add 1.5--2.5 kW:

\begin{table}[h!]
\caption{Total electricity derivation for the \texttt{batch\_treatment\_generic} block.}
\centering
\begin{tabular}{>{\raggedright\arraybackslash}p{4.0cm}
				>{\raggedright\arraybackslash}p{2.5cm}
				>{\raggedright\arraybackslash}p{2.5cm}
				>{\raggedright\arraybackslash}p{4.5cm}}
\toprule
\textbf{Component} & \textbf{Opt.\ (kW)} & \textbf{Mid.\ (kW)} & \textbf{Cons.\ (kW)} \\
\midrule
Batch agitation (time-wtd.)      & 4.1  & 6.2  & 9.5  \\
Lime slaking + storage (cycle-avg.) & 2.5  & 3.5  & 5.0  \\
Soda ash dissolution             & 0.3  & 0.4  & 0.5  \\
Pumping (cycle-avg.)             & 0.5  & 0.7  & 1.0  \\
Auxiliaries (I\&C)               & 1.5  & 2.0  & 2.5  \\
\midrule
\textbf{Total (cycle-avg.)}      & \textbf{8.9} & \textbf{12.8} & \textbf{18.5} \\
\texttt{electricity\_base\_total}    & \textbf{9 kWh/h} & \textbf{13 kWh/h} & \textbf{18 kWh/h} \\
\bottomrule
\end{tabular}
\end{table}

\noindent These cycle-averaged values are much lower than a continuous-mode estimate would suggest, because batch operation entails intermittent equipment usage with substantial settling periods where agitation is minimal.

\medskip
\noindent\textbf{Other utilities.}\quad

No thermal utilities are consumed. Process water is required for lime slaking, if quicklime were used as the input, at an approximately 2:1 water-to-quicklime ratio by weight.

\medskip
\noindent\textbf{Consumables.}\quad

Quicklime (\ce{CaO}) and soda ash (\ce{Na2CO3}) are the primary reagents and are tracked at the production step level as secondary inputs with stoichiometrically derived dosages. \texttt{tooling\_cost\_base} = 0.

\medskip
\noindent\textbf{Labor.}\quad

Labor is 5.4 FTE at base scale with a scaling exponent of 0.2 and \texttt{dedicated\_labor} = TRUE. The higher base (5.4 vs.\ 4.5 for the continuous blocks) and dedicated-labor flag reflect the nature of batch operation; operators must be present throughout each 7-hour cycle for reagent charging, pH monitoring, settling observation, and valve sequencing, with additional labor associated with auxiliary equipment and control:
\[
	\text{FTE} = (1.0 + 0.2) \times 4.5 = 5.4
\]

\medskip
\noindent\textbf{Footprint.}\quad

The base plot area is 300~$\si{m^2}$ with a scaling exponent of 0.6, covering the batch reactor, lime slaking system, soda ash storage, and surrounding access.

\medskip
\noindent\textbf{Availability.}\quad

Process availability is 90\% with 5\% unplanned downtime and 5\% scheduled maintenance, matching the \texttt{precipitator} block.

\medskip
\noindent\textbf{Maintenance.}\quad

Neither \texttt{maint\_override} nor a nonzero \texttt{opex\_fraction\_of\_capex} is set; the block inherits the facility-level default maintenance rate of 5\% of scaled primary equipment cost. As noted in the \texttt{liming\_system} block, lime slaking systems require above-average maintenance, so a \texttt{maint\_override} above the 5\% default should be considered if this block is refined beyond its current baseline role.

\medskip
\noindent\textbf{Key measurements and cross-checks.}\quad

\begin{itemize}[noitemsep]
	\item \textit{Specific electricity}: midpoint of 13 kWh/h at 8.6 $\si{m^3/h}$ effective throughput implies 1.5 \si{kWh/m^3} of brine, compared to $\sim$1.0 \si{kWh/m^3} for the continuous \texttt{precipitator} and \texttt{liming\_system} blocks. The higher intensity is consistent with the combined lime + soda ash treatment and the overhead of the lime storage keep-alive agitation spread over a smaller batch throughput.
	\item \textit{CAPEX}: midpoint of \$1.64M at $\sim$8.6 $\si{m^3/h}$ effective throughput implies \$191/($\si{m^3/h}$), approximately 15$\times$ the precipitator's \$13/($\si{m^3/h}$). The premium is driven by the lime slaking package (\$683K, 42\% of midpoint total), the dual-reagent system, and the lower effective throughput of batch operation.
\end{itemize}

\clearpage
\paragraph{\texttt{wash\_purify\_combo}}\label{app:04-methods-blockdefs-washpurify}\leavevmode\par

\medskip
\noindent\textbf{Sources.}\quad

\citet{srkconsultingSECTechnicalReport2023a}, \citet{sgscanadaincSK1300Technical2024}, \citet{petersPlantDesignEconomics1991}; vendor data. Tank costs from AusIMM benchmarks (Rapid Mix Tanks, Agitator Tanks); re-dewatering filter costs from published filter press price benchmarks~\citet{porvooIndustrialFilterPress2025}. 


\medskip
\noindent\textbf{Process description.}\quad

The wash\_purify\_combo block removes residual mother liquor and soluble impurities (principally \ce{Na^+}, \ce{K^+}, \ce{SO4^{2-}} or \ce{Cl^-}) from \ce{Li2CO3} filter cake by multi-stage counter-current repulp washing. The input is dewatered \ce{Li2CO3} cake at approximately 70 wt\% solids from the upstream \texttt{mechanical\_dewatering} step. In each wash stage, the cake is re-slurried (``repulped'') with clean or partially clean wash water in an agitated tank, dissolving and displacing entrained impurities from the crystal surfaces and interstices. The dilute slurry is then re-dewatered on a modest filter to recover the cake at 70 wt\% solids before the next stage. Two to three counter-current stages are required to achieve battery-grade purity ($>$99.5\% \ce{Li2CO3}), with the cleanest water contacting the cleanest cake. The ``purification'' in the block name refers to this physical displacement of impurities, not to chemical redissolution or recrystallization.

The primary equipment is therefore agitated repulp tanks with a secondary re-dewatering filter --- structurally analogous to the tank-based circuits in the \texttt{precipitator} and \texttt{liming\_system} blocks, not to the high-duty hydraulic press in the \texttt{mechanical\_dewatering} block. The re-dewatering filter handles a lighter duty than the upstream dewatering step: it re-presses dilute slurry from $\sim$35 wt\% back to 70 wt\% solids, not from 20 wt\% to 70 wt\%.

This block is used in both the chloride (Silver Peak) and sulfate (Thacker Pass, Jianxiawo) pathways. Wash liquor handling differs by pathway and, within the sulfate pathway, by project: in the chloride route, wash water is recycled to the precipitation reactor or evaporation pond; in the sulfate route it is retained on site at Thacker Pass under the zero-discharge permit and discharged to wastewater treatment at Jianxiawo.

\medskip
\noindent\textbf{Equation of merit.}\quad

\[
	V_\text{repulp}\;(\si{m^3}) = V_\text{cake}\;(\si{m^3}) + V_\text{wash}\;(\si{m^3})
\]
Each repulp tank must hold the cake plus wash water. At a 1:1 wash ratio by volume, 60 \si{m^3} of cake produces approximately 120 \si{m^3} of repulp slurry at $\sim$35 wt\% solids. With three tanks at $\sim$35 \si{m^3} each (105 \si{m^3} total working volume), the cake is processed in sub-batches of $\sim$20 \si{m^3} through the counter-current wash circuit within the 1-hour cycle.

\medskip
\noindent\textbf{Bill of Equipment.}\quad

\begin{table}[h!]
\caption{Bill of equipment for the \texttt{wash\_purify\_combo} block.}
\centering
\begin{tabular}{>{\raggedright\arraybackslash}p{5.2cm}
				>{\raggedright\arraybackslash}p{4cm}
				>{\raggedright\arraybackslash}p{4.5cm}}
\toprule
\textbf{Equipment} & \textbf{Machine data used} & \textbf{Source(s)} \\
\midrule
Repulp tanks, SS316 ($3 \times \sim$35 \si{m^3}) & Mixtec Agitator Tank Series & Agitator Tanks \\
Re-dewatering filter (630--800 mm plate press or vacuum belt) & Published benchmarks & Porvoo 2025 \\
Wash water pumps \& piping & Not explicit $\rightarrow$ inferred & --- \\
Slurry transfer pumps & Not explicit $\rightarrow$ inferred & --- \\
Compressed air system (cake blow-down) & Not explicit $\rightarrow$ inferred & --- \\
Cake discharge conveyor & Not explicit $\rightarrow$ inferred & --- \\
Instrumentation \& controls & Not explicit $\rightarrow$ inferred & --- \\
\bottomrule
\end{tabular}
\end{table}

\medskip
\noindent\textbf{Sizing basis.}\quad

The base volume is 60 \si{m^3/batch} of \ce{Li2CO3} cake at 70 wt\% solids (approximately 95 tonnes wet cake containing 66.5 tonnes \ce{Li2CO3}). While additional water is added and thus the equipment physically is sized to contain more than 120 \si{m^3} of the process basis slurry, the batched nature (\texttt{process\_type} = batch) means that sizing to the contained cake volume is appropriate as long as units are consistent. Each batch cycle is 1.2 hours (1 hour wash + 0.2 hours setup), giving an effective throughput of 50 $\si{m^3/h}$. Equipment life is assumed to be 20 years.

\medskip
\noindent\textbf{CAPEX.}\quad

CAPEX is built bottom-up from two components: repulp tanks (costed from AusIMM volumetric benchmarks, as in the \texttt{precipitator} block) and a re-dewatering filter (costed from published filter press price benchmarks). A stainless steel premium of 1.3--1.5$\times$ is applied to all wetted equipment to meet battery-grade \ce{Li2CO3} purity requirements.

\textit{Repulp tanks} (105 \si{m^3} total working volume, 3 tanks $\times$ 35 \si{m^3}): raw tank cost = volumetric cost $\times$ volume, then $\times$ SS premium.

\textit{Re-dewatering filter}: a 630--800 mm plate filter press sized for intermittent re-pressing of repulped slurry from $\sim$35 wt\% to 70 wt\% solids. Bare cost anchored to the published benchmark of \$75,000--\$220,000 for 630--800 mm presses, with a stainless-steel premium and installation applied.

\textit{Subsystem adders}: wash water pumps and piping (15\%), instrumentation and controls (10\%), totaling 25\% on the combined tank + filter cost.

\begin{table}[h!]
\caption{CAPEX derivation for the \texttt{wash\_purify\_combo} block.}
\centering
\begin{tabular}{>{\raggedright\arraybackslash}p{2.2cm}
				>{\raggedright\arraybackslash}p{2.5cm}
				>{\raggedright\arraybackslash}p{2.3cm}
				>{\raggedright\arraybackslash}p{2.4cm}
				>{\raggedright\arraybackslash}p{3.5cm}}

\toprule
\textbf{Scenario} & \textbf{Tanks (stainless)} & \textbf{Filter (installed)} & \textbf{Wash water sys.} &  \texttt{prim\_equip\_price\_base} \\
\midrule
Conservative & \$1,418K & \$330K & \$150K & \textbf{\$2,370,000} \\
Midpoint     & \$956K   & \$200K & \$100K & \textbf{\$1,570,000} \\
Optimistic   & \$655K   & \$130K & \$60K  & \textbf{\$1,060,000} \\
\bottomrule
\end{tabular}
\end{table}

\noindent \textit{Derivation details (midpoint row).} The raw tank cost is \$6,500/\si{m^3} $\times$ 105 \si{m^3} = \$683K, which becomes \$956K after applying a 1.4$\times$ stainless steel premium. The filter cost starts from a 630 mm press at \$108K bare, multiplied by 1.4$\times$ for installation and 1.3$\times$ for stainless steel construction, giving \$197K $\approx$ \$200K installed. The wash water system (holding tank, heater stub, metering pumps) is estimated at \$100K. The subtotal of \$1,256K is then multiplied by 1.25$\times$ for piping and instrumentation adders, giving a total of \$1,570K.

\medskip
\noindent\textbf{Electricity.}\quad

Electricity demand is derived from the major loads during the 1.2-hour batch cycle. Each of the 2--3 wash stages involves repulping (agitation), pumping, and re-dewatering on the filter. Equipment operates intermittently; \texttt{electricity\_base\_total} represents the cycle-averaged power draw.

\textit{Agitation.} Repulp tanks at 0.20--0.40 \si{kW/m^3} on the working slurry volume ($\sim$35 \si{m^3} per tank), with 1--2 tanks active, running $\sim$30--40 min of the 72-min cycle.

\textit{Re-dewatering filter.} The filter re-presses repulped slurry 2--3 times per batch, running $\sim$30--45 min total. Installed motor power for a 630--800 mm filter press with hydraulic closure is 15--55 kW; cycle-averaged draw is lower due to intermittent HPU operation.

\textit{Pumping.} Wash water and slurry transfer pumps from $P = \rho g Q H / \eta$, running during fill and drain phases.

\begin{table}[h!]
\caption{Electricity derivation for the \texttt{wash\_purify\_combo} block.}
\centering
\begin{tabular}{>{\raggedright\arraybackslash}p{5.2cm}
				>{\raggedright\arraybackslash}p{3.0cm}
				>{\raggedright\arraybackslash}p{3.0cm}
				>{\raggedright\arraybackslash}p{3.0cm}}
\toprule
\textbf{Component} & \textbf{Opt.\ (kW)} & \textbf{Mid.\ (kW)} & \textbf{Cons.\ (kW)} \\
\midrule
Repulp agitation (cycle-avg.)    & 8   & 15  & 22  \\
Re-dewatering filter (cycle-avg.) & 18  & 38  & 55  \\
Wash water pumps (cycle-avg.)    & 8   & 13  & 18  \\
Slurry transfer pumps            & 7   & 12  & 16  \\
Compressed air (blow-down)       & 3   & 5   & 7   \\
Auxiliaries (I\&C, conveyor)     & 5   & 7   & 9   \\
\midrule
\texttt{electricity\_base\_total} & $\sim$\textbf{50 kWh/h} & $\sim$\textbf{90 kWh/h} & $\sim$\textbf{130 kWh/h} \\
\bottomrule
\end{tabular}
\end{table}

\noindent The re-dewatering filter remains the dominant load ($\sim$40\% of midpoint total), consistent with the repeated pressing of repulped slurry.

\medskip
\noindent\textbf{Other utilities.}\quad

Cooling water is consumed at 0.2 $\si{m^3/h}$ for hydraulic oil cooling on the re-dewatering filter. Compressed air is consumed at 1 \si{Nm^3/h} for cake blow-down between wash stages. No thermal utilities are tracked; wash water temperature conditioning (if used to exploit retrograde solubility of \ce{Li2CO3}) would require steam or process heat integration, but is not currently modeled.

\medskip
\noindent\textbf{Consumables.}\quad

\texttt{tooling\_cost\_base} = 0. Filter cloths for the re-dewatering filter are a recurring consumable (replacement every 500--2,000 press cycles) but are not yet populated. This gap is less significant than for the \texttt{mechanical\_dewatering} block given the lighter filter duty.

\medskip
\noindent\textbf{Labor.}\quad

Labor is 3.15 FTE at base scale with a scaling exponent of 0.5 and \texttt{dedicated\_labor} = FALSE, consistent with a semi-automated wash operation sharing operators with adjacent dewatering and drying equipment. The higher scaling exponent (0.5 vs.\ 0.2 for tank-only circuits) reflects the additional operator involvement in managing the re-dewatering filter:
\[
	\text{FTE} = (0.5 + 0.2) \times 4.5 = 3.15
\]

\medskip
\noindent\textbf{Footprint.}\quad

The base plot area is 100~$\si{m^2}$ with a scaling exponent of 0.7, covering the repulp tanks, re-dewatering filter, wash water system, and cake discharge area.

\medskip
\noindent\textbf{Availability.}\quad

Process availability is 90\% with 5\% unplanned downtime and 5\% scheduled maintenance.

\medskip
\noindent\textbf{Maintenance.}\quad

Neither \texttt{maint\_override} nor a nonzero \texttt{opex\_fraction\_of\_capex} is set; the block inherits the facility-level default maintenance rate of 5\% of scaled primary equipment cost. The re-dewatering filter's cloth and seal replacement requirements are the maintenance driver here, and are not separately parameterized.

\medskip
\noindent\textbf{Key measurements and cross-checks.}\quad

\begin{itemize}[noitemsep]
	\item \textit{CAPEX}: bottom-up midpoint of \$1.57M at 50 $\si{m^3/h}$ effective throughput implies \$31/($\si{m^3/h}$), approximately 2.4$\times$ the \texttt{precipitator}'s \$13/($\si{m^3/h}$). The premium is consistent with a stainless steel construction requirement and the re-dewatering filter.
	\item \textit{Electricity}: midpoint of 90 kWh/h (1.5 \si{kWh/m^3} of cake) is dominated by the re-dewatering filter ($\sim$42\%).
	\item \textit{Wash efficiency}: 2--3 counter-current stages at $\sim$1:1 wash ratio total 120--180 \si{m^3} of wash water per 60 \si{m^3} batch. This demand is tracked at the production step level.
\end{itemize}

\clearpage
\paragraph{\texttt{drying\_packaging} (Generic and Propane options)}\label{app:04-methods-blockdefs-drypack}\leavevmode\par

\medskip
\noindent\textbf{Sources.}\quad

\citet{laneNI43101Technical2024, sgscanadaincSK1300Technical2024, riordanFalchaniLithiumProject2024, samariMineralResourceEstimate2024, srkconsultingSECTechnicalReport2023a, daiEverBattClosedloopBattery2019}

\medskip
\noindent\textbf{Process description.}\quad

Following precipitation and solid-liquid separation, lithium carbonate is produced as a wet cake at approximately 70 wt\% solids and 30 wt\% moisture. This model block dries the cake to a product specification of $\leq$0.4 wt\% moisture and packages the dry product for shipment. The block is modeled as a compact finishing circuit centered on an indirect gas-fired dryer with baghouse dust capture and a bulk bagging line. Two fuel variants share identical equipment configurations and all non-fuel parameters: \texttt{drying\_packaging\_generic} uses natural gas (\texttt{natural\_gas\_base\_total}) and is used for the Thacker Pass and Jianxiawo pathways; \texttt{drying\_packaging\_propane} substitutes propane (\texttt{propane\_base\_total}) and is used for Silver Peak, where no natural gas supply is available at the remote Nevada site \cite{srkconsultingSECTechnicalReport2023a}.

\medskip
\noindent\textbf{Equation of merit.}\quad
\[
	\dot{m}_{\ce{H2O}}\;(t/h) \approx \dot{m}_\text{cake}\;(t/h) \times 0.297
\]
where the factor 0.297 is the water evaporated per tonne of 70~wt\% solids wet cake, derived from:
\[
	w_\text{evap} = w_\text{water} - \frac{w_\text{solids} \times w_\text{product}}{1 - w_\text{product}}
	= 0.30 - \frac{0.70 \times 0.004}{0.996} = 0.2972~\text{t/t cake}
\]
Thermal duty scales with the water evaporation rate; electricity and CAPEX scale with wet cake throughput.

\medskip
\noindent\textbf{Sizing basis.}\quad

The step basis is \texttt{cake\_dry\_throughput} and \texttt{base\_volume} = 2.5, representing 2.5~t/h of wet cake feed to the dryer. This value originates from the EverBatt battery materials cost model \cite{daiEverBattClosedloopBattery2019}, where a capacity adjustment of $1/0.4 = 2.5$ is applied to a baseline dryer parameterized at $x = 0.4$. At the 2.5~t/h wet cake base:
\begin{itemize}[noitemsep]
	\item Water evaporated: $2.5 \times 0.297 = \mathbf{0.743}$~t/h
	\item Dry \ce{Li2CO3} product: $2.5 \times 0.7028 = 1.757$ t/h at $\leq$0.4 wt\% moisture
\end{itemize}

\medskip
\noindent\textbf{Bill of Equipment.}\quad
 
\begin{table}[h!]
\caption{Bill of equipment for the \texttt{drying\_packaging} blocks.}
\centering
\begin{tabular}{>{\raggedright\arraybackslash}p{5.2cm}
                >{\raggedright\arraybackslash}p{4cm}
                >{\raggedright\arraybackslash}p{4.5cm}}
\toprule
\textbf{Equipment} & \textbf{Machine data used} & \textbf{Source(s)} \\
\midrule
Indirect gas-fired dryer & Indirect dryer & Feasibility disclosures, \cite{daiEverBattClosedloopBattery2019} \\
Baghouse + exhaust fan & Inferred from dryer package & --- \\
Bulk bagging line & Standard \ce{Li2CO3} packaging & Feasibility disclosures \\
Feed handling \& conveying & Not explicit $\rightarrow$ inferred & --- \\
Product bins, silo \& loadout & Not explicit $\rightarrow$ inferred & --- \\
Utilities, controls \& instrumentation & Not explicit $\rightarrow$ inferred & --- \\
\bottomrule
\end{tabular}
\end{table}

\vspace{1cm}
\medskip
\noindent\textbf{CAPEX.}\quad

The EverBatt battery materials cost model (Argonne National Laboratory) parameterizes the bare dryer purchase cost as $C = (591,236 \times x^{0.6}) \times 1.483$ (2019\$), giving \$506,014 at $x = 0.4$. Inflated to 2025\$ yields \$666,000. A $\pm$30\% range on this anchor reflects Class~4--5 uncertainty in equipment specification, giving bare dryer costs of \$466,000 / \$666,000 / \$866,000 (optimistic / midpoint / conservative).
 
Auxiliary bare equipment is added as fractions of the bare dryer cost following Peters \& Timmerhaus Table 6-11: baghouse and exhaust fan at 25\% / 32\% / 40\% and bulk bagging line at 18\% / 22\% / 30\% (optimistic / midpoint / conservative). Balance-of-plant items (feed handling and conveying, product silo and loadout, utilities and controls) are added at 30\% / 40\% / 52\% of the primary equipment subtotal per Peters \& Timmerhaus Table 6-17 for fluid-solid processing. A Lang-type installation factor of 2.3$\times$ / 2.8$\times$ / 3.3$\times$ (optimistic / midpoint / conservative) is then applied to total bare equipment, covering direct installation labor, piping, instrumentation, electrical, structural, civil, engineering, and contingency:
 
\begin{table}[h!]
\caption{Total installed circuit cost for the \texttt{drying\_packaging} blocks.}
\centering
\begin{tabular}{lcccc}
\toprule
\textbf{Scenario} & \textbf{Bare dryer} & \textbf{Total bare} & \textbf{Install factor} & \texttt{prim\_equip\_price\_base} \\
\midrule
Conservative & \$866,000 & \$2,238,000 & $\times$3.3 & \textbf{\$7,400,000} \\
Midpoint     & \$666,000 & \$1,437,000 & $\times$2.8 & \textbf{\$4,000,000} \\
Optimistic   & \$466,000 & \$867,000   & $\times$2.3 & \textbf{\$2,000,000} \\
\bottomrule
\end{tabular}
\end{table}

\noindent These values apply to both variants. The CAPEX scaling exponent is 0.60, following the six-tenths rule for process equipment cost (which is also used in EverBatt), and equipment life is 20 years. 

\noindent\textbf{Electricity.}\quad

Electrical demand covers the dryer drive and exhaust fan / baghouse system, with all minor motor loads (conveyor drives, bagging line, instrumentation) treated as auxiliary adders.

\noindent Electricity is identical across the natural gas and propane variants; only the fuel supply differs.

\medskip
\noindent\textbf{Natural gas (\texttt{drying\_packaging\_generic}).}\quad

Thermal duty is derived from first principles. The drying load per tonne of water evaporated comprises sensible heat of liquid water from 25$^\circ$C to 100$^\circ$C ($\approx$419~kJ/kg), latent heat of evaporation at 100$^\circ$C ($\approx$2,257~kJ/kg), and a 10\% solids heating penalty. An additional air-overheating penalty of 10\% / 20\% / 30\% (optimistic / midpoint / conservative) accounts for energy required to superheat the exhaust stream. Dividing by 80\% combustion efficiency gives the fuel input per tonne of water evaporated:
 
\begin{table}[h!]
\caption{Natural gas derivation for \texttt{drying\_packaging\_generic} at 2.5~t/h wet cake.}
\centering
\begin{tabular}{>{\raggedright\arraybackslash}p{2.3cm}
				>{\raggedright\arraybackslash}p{3.2cm}
				>{\raggedright\arraybackslash}p{3.2cm}
				>{\raggedright\arraybackslash}p{4.8cm}}
\toprule
\textbf{Scenario} & \textbf{Net} (MMBTU/t \ce{H2O}) & \textbf{Fuel} (MMBTU/t \ce{H2O}) & \texttt{natural\_gas\_base\_total} (At 0.743~t/h \ce{H2O}) \\
\midrule
Conservative & 3.63 & 4.54 & \textbf{3.37} \\
Midpoint	 & 3.35 & 4.18 & \textbf{3.11} \\
Optimistic	 & 3.07 & 3.84 & \textbf{2.85} \\
\bottomrule
\end{tabular}
\end{table}

\medskip
\noindent\textbf{Propane (\texttt{drying\_packaging\_propane}).}\quad
Propane consumption is converted directly from the equivalent natural gas thermal demand using the propane higher heating value of 24,174~Btu/L (liquid propane at standard conditions), giving 41.37~L propane per MMBTU. \texttt{natural\_gas\_base\_total} is set to \texttt{False} for this variant.
 
\begin{table}[h!]
\caption{Propane derivation for \texttt{drying\_packaging\_propane} at 2.5~t/h wet cake.}
\centering
\begin{tabular}{>{\raggedright\arraybackslash}p{2.5cm}
				>{\raggedright\arraybackslash}p{3.5cm}
				>{\raggedright\arraybackslash}p{3.5cm}
				>{\raggedright\arraybackslash}p{3.5cm}}
\toprule
\textbf{Scenario} & \textbf{NatGas equiv.\ (MMBTU/h)} & \textbf{Conversion (L/MMBTU)} & \texttt{propane\_base\_total} \\
\midrule
Conservative & 3.37 & 41.37 & \textbf{139 L/h} \\
Midpoint	 & 3.11 & 41.37 & \textbf{129 L/h} \\
Optimistic	 & 2.85 & 41.37 & \textbf{118 L/h} \\
\bottomrule
\end{tabular}
\end{table}

\medskip
\noindent\textbf{Other utilities.}\quad
\begin{itemize}[noitemsep]
	\item \texttt{cooling\_water\_base\_total}: 0.1 \si{m^3/h} (generic variant only) --- minor seal and bearing cooling for the dryer drive
	\item \texttt{compressed\_air\_base\_total}: 0.5 \si{Nm^3/h} (both variants) --- instrumentation air only; no bag filter pulse-jet system is modeled at this scale
	\item Steam, diesel: zero for both variants
\end{itemize}

\medskip
\noindent\textbf{Labor.}\quad
Labor is 4.05 FTE at base scale with a scaling exponent of 0.5 and \texttt{dedicated\_labor} = FALSE, summing operator-per-shift factors across the three principal equipment groups in the circuit --- indirect dryer (0.2), separator/baghouse/conveying (0.2), and packaging line (0.5):
\[
	\text{FTE} = (0.2 + 0.2 + 0.5) \times 4.5 = 4.05
\]

\medskip
\noindent\textbf{Footprint.}\quad
The base plot area is 400~$\si{m^2}$ with a scaling exponent of 0.7, covering the dryer shell, baghouse, bagging line, product silo, and surrounding access, and is assumed to be identical across both fuel variants.
 
\medskip
\noindent\textbf{Availability.}\quad
Process availability is 90\% with 5\% unplanned and 5\% scheduled maintenance, consistent with a straightforward mechanical finishing circuit. Identical across both variants.
 
\medskip
\noindent\textbf{Key measurements and cross-checks.}\quad
\begin{itemize}[noitemsep]
	\item Product moisture at $\leq$0.4~wt\% --- the primary dryer performance specification
	\item Total energy: 401~kWh/t wet cake (midpoint, gas + electric combined) is consistent with the EverBatt electric-only benchmark once normalised to the same throughput basis and energy carrier
	\item CAPEX: midpoint \$4.0M implies the bare dryer (\$666k from EverBatt) represents 17\% of total installed circuit cost, with installation and auxiliaries accounting for the remainder --- consistent with Peters \& Timmerhaus guidance for small solid-fluid processing circuits
\end{itemize}

\end{document}

